\section{结论}\label{sec:conclusion}

本文的主线是把有限窗口读出所诱导的微态空间，通过可审计的构造链条组织为稳定类型、可实现的算术结构与可证书化的统计/谱指纹接口。相应的主要结论可按“输入口径—构造机制—桥接与验证”三段收束如下。

\subsection{主要结论：按机制分组}\label{subsec:core-conclusions}
\paragraph{输入口径的固定}
数论侧以反共振规范代表与有限样本差异度证书固定稳定性基线（第 \ref{sec:principles} 节），动力学侧以扫描—投影生成器（群作用与二元读出）固定可观测对象的最小输入（第 \ref{sec:spg} 节）。上述口径在全文保持不变，使后续结论具有可追溯的依赖链。
在该接口上，“扫描深度 \(\Rightarrow\) 清晰度”可被压缩为最优判错剖面与边界指数律，并在地址语义下进一步分解为“未同步（统计型 \(\mathsf{NULL}\)）”与“地址碰撞”两类失效机制的双预算容量律（定理 \ref{thm:spg-clarity-boundary-dimension}、\ref{thm:spg-double-budget-address-capacity}）；在在线制度下，该清晰度剖面又与同步扣除的 Chebotarev 容量不等式耦合，从而把定义性同步门槛与统计均匀化预算分离为可审计的两段账本（定理 \ref{thm:pom-sync-subtracted-chebotarev-capacity}）。

\paragraph{分辨率折叠与递归地址化}
本文构造有限窗口折叠 \(\Fold_m\) 并证明其跨分辨率相容与统一的 sofic 因子接口结论（定理 \ref{thm:fold-suite}；并见定义 \ref{def:fold-word} 与命题 \ref{prop:fold-basic}）。在同一可测框架内，递归地址化被严格化为“可读出前缀”的内生化机制：地址对象未给定时，评价函数 \(\ind_{C_a}\) 在指称层不成立，读出偏函数不定义并以 \(\mathsf{NULL}\) 表示不可寻址；一旦地址给定，派生柱事件族不扩张既有可见 \(\sigma\)-代数，而仅在旧层可见 \(\mathcal{G}^{(L)}\) 内重选并组织可见事件，从而 \(\mathcal{G}^{(L+1)}\subseteq \mathcal{G}^{(L)}\)（命题 \ref{prop:recursive-addressing-refinement}；推论 \ref{cor:recursive-addressing-visible-sigma-nonexpansion}）。折叠降阶的非局域性可由规范差 \(G_m\) 定量审计并给出最坏与典型的双重锚点（定义 \ref{def:fold-gauge-anomaly}；定理 \ref{thm:fold-gauge-anomaly-max}、\ref{thm:fold-gauge-anomaly-density}）。
并且滑动块因子 \(\Phi_m\) 的可逆性在 \(m=2\) 处出现唯一的两点分支：除两条常值轨道外皆一一，而常值轨道发生 \(2\)-对-\(1\) 坍缩；该阈值可被解释为一个 \(\mathbb{Z}/2\) 隐变量在 \(m=2\) 窗口下不可同步消去、在 \(m\ge 3\) 时被制度性消去的最小同调障碍（定理 \ref{thm:Phi_m-conjugacy-threshold}；注 \ref{rem:Phi_m-conjugacy-threshold-z2-homology}）。

\paragraph{稳定算术与投影本体数学}
在值保持重写与规范横截面机制下，分辨率 \(m\) 的稳定类型算术与 \(\ZZ/(F_{m+2}\ZZ)\) 同构（定理 \ref{thm:finite-resolution-mod}），从而把统一/分裂的第一性来源归结为模数的环论性质及 CRT 结构（推论 \ref{cor:field-phase-fib-prime}、\ref{cor:crt-factorization}）；并可沿 Fibonacci 模数族的整除偏序定义 Fibonacci-adic 逆极限环并得到紧致单位群（定义 \ref{def:fib-adic-ring}；命题 \ref{prop:fib-adic-basic}）。在此基础上，本文建立投影本体数学（POM）的投影词演算，并给出信息账本恒等式，将粗粒化的信息损失严格对象化为可计算残差（定理 \ref{thm:pom-kl-ledger}；推论 \ref{cor:pom-ent-increase-equals-kl}）。同时在局部读出塔片段上证明重写的终止性与合流性并给出正规形，进而得到可审计闭环与判定性接口（定理 \ref{thm:pom-rewrite-termination-confluence}；命题 \ref{prop:pom-projword-decidable}、\ref{prop:pom-rewrite-audit-loop}）。高阶碰撞矩的 moment-kernel 编译与共振判别把“共振”落实为可检验机制（定理 \ref{thm:pom-collision-kernel-ak}、\ref{thm:pom-moment-resonance}）。

\paragraph{纤维的拓扑—立方几何闭包与因子链混合下界}
在局部核 \((011\leftrightarrow 100)\) 的唯一结构输入下，本文把“纤维”从纯计数对象提升为一组可审计的拓扑/几何/概率不变量族：
首先，在定理 \ref{thm:pom-fiber-independence-complex-classification} 中把候选位点诱导图的独立集族提升为独立集复形，
证明每条纤维的同伦型要么可缩、要么同伦等价于单一球面 \(S^{\tau(x)-1}\)，并给出 \(\tau(x)\) 的闭式与唯一维度同调刚性（推论 \ref{cor:pom-fiber-independence-complex-homology}）。
其次，把纤维内逆代换次数视为随机变量 \(R_x\)，得到严格的分量独立分解与 \(\varphi\)-代数的均值/方差密度（命题 \ref{prop:pom-fiber-rewrite-independent-sum}、\ref{prop:pom-path-indset-thermo-constants}），并由实根性导出强单峰性、Poisson--binomial 因子化与局部 CLT，从而给出统一的 Bernstein/Chernoff 型浓缩（定理 \ref{thm:pom-fiber-rewrite-realroot-ulc}；推论 \ref{cor:pom-fiber-rewrite-poisson-binomial}、\ref{cor:pom-fiber-rewrite-local-clt}）。
再次，在定义 \ref{def:pom-fiber-reconstruction-graph} 的“单次开关”相邻结构下，纤维重构图 \(\mathcal F_x\) 严格分解为 Fibonacci cube 的 Cartesian 乘积（定理 \ref{thm:pom-fiber-reconstruction-cartesian-product}），从而获得 CAT(0) 立方几何与完全闭式的立方面计数（定理 \ref{thm:pom-fibcube-fvector-closed}、推论 \ref{cor:pom-fiber-fpoly-factorization}），并存在可常延迟枚举全预像的 Gray--Hamilton 结构（定理 \ref{thm:pom-fibcube-gray-hamilton}、推论 \ref{cor:pom-fiber-constant-delay-traversal}）。
并且路径分量乘法多重度在细化偏序上严格单调，融合下降量具有闭式，从而给出一条可审计的严格 Schur-凸序（命题 \ref{prop:pom-path-component-multiplicity-refinement-monotone-extrema}）。
此外，在锯齿理想格的最大链/线性扩张口径下，双不变量可达域在尺度极限上等于原子点集的闭凸包，因而任一 Pareto 前沿暴露点至多由两种原子长度混合实现（定理 \ref{thm:pom-fiber-ab-convex-hull}；推论 \ref{cor:pom-fiber-ab-two-phase-mixing}）。
最后，以 \(\Omega_m\) 上超立方体随机游走为基线，本文定义折叠因子链 \(\overline P^{\square}_m\)，证明其可逆、无自环且谱隙/对数 Sobolev 常数不劣于超立方体本身（命题 \ref{prop:pom-fold-factor-chain-reversible-noloop}、定理 \ref{thm:pom-fold-factor-chain-gap-lsi}），从而把稳定层 \(X_m\) 的混合与浓缩能力下界化为与折叠细节无关的统一常数。
该结论簇与 rotation polytope 的凸几何骨架、碰撞矩递推的 \(\mathbb F_2/(p)\)-进刚性签名共同组成“6D 读出—折叠核—算术/几何/概率证书”的闭合接口。

\paragraph{桥接接口与验证}
在谱侧，动力学 \(\zeta\) 与有限部接口把周期计数、迹与主谱半径组织为可审计指纹，并为后续的完成化与扭曲通道提供统一输入（第 \ref{sec:zeta-finite-part} 节）。在统计侧，折叠直方图的偏差指标、跨分辨率一致性残差与到 Parry(PF) 基线的误差分解式将有限样本误差与模型差异项显式分离（第 \ref{sec:statistical-stability}、\ref{sec:experiments} 节；推论 \ref{cor:tv-decomposition-parry}），并以脚本化协议给出可复现实证书输出（第 \ref{sec:experiments} 节）。
其中在金属平均族驱动 \(\alpha_A=[0;A,A,\dots]\) 下，确定性 TV 证书的 \((\log N)/N\) 主系数由 \(\kappa(A)=A/\log\lambda_A\) 控制并由黄金端点唯一最小化（推论 \ref{cor:metallic-mean-golden-unique-optimal-tv-constant}）；同时在最小原子下界的支撑护栏下，KL 误差获得按差异度平方衰减的二次提升（引理 \ref{lem:kl-from-tv-qmin}；推论 \ref{cor:rotation-folded-kl-certificate}）。
进一步，将输入不确定性编码为 \(W_1\) 传输差异后，Lipschitz 推前与读出误差在制度上保持协变，并在黄金分支的 Fibonacci 分母子序列上出现二周期极限常数与 \(\sqrt5\)-分裂（定理 \ref{thm:kronecker-w1-lipschitz-readout}；推论 \ref{cor:kronecker-w1-fibonacci-limits}；命题 \ref{prop:w1-lipschitz-pushforward}）；在 \(\zeta/\Xi\) 的解析读出链条中，Stieltjes 读出的最强尺度 Lipschitz 常数为 \((t+1)^{-2}\)，从而稳定—病态分界可由单一因子量化（命题 \ref{prop:xi-stieltjes-w1-lipschitz}）。
其中在 \(\zeta/\Xi\) 的完成化审计中，离开 unitary slice 的信息被端点二次径向压缩并导出严格的可见化预算门槛与单半径带状禁入证书，同时最小 \((2\times2)\) 自对偶模板给出显式轴外零点塔作为可证伪原型（定理 \ref{thm:xi-jensen-single-radius-band-exclusion}、\ref{thm:xi-2x2-selfdual-offcritical-template}）。
在同步核的加权压力接口中，primitive 计数的 Abel--Mertens 有限部常数、解析半径的判别式刻画与单位根扭曲阈值，以及 Edgeworth/中偏差喷射的可审计截断余项界，均可由 Perron 分支的解析延拓与谱隙一致性统一封口（命题 \ref{prop:sync-kernel-abel-mertens-analytic-family}，推论 \ref{cor:sync-kernel-analytic-radius-discriminant-formula}--\ref{cor:sync-kernel-unit-root-twist-threshold}，定理 \ref{thm:sync-kernel-even-edgeworth-c4}，命题 \ref{prop:sync-kernel-moderate-deviation-cauchy-certified}--\ref{prop:sync-kernel-curvature-bilipschitz}）。
此外，真实输入 \(40\) 状态核的零谱幂零指数把在线延迟与复位词长提升为线性代数不变量，并给出瞬态记忆深度的严格上界（命题 \ref{prop:real-input-40-nilpotent-index}，推论 \ref{cor:real-input-40-nilpotent-memory-depth}）；而折叠碰撞矩谱的递推—谱半径闭包把全部 Rényi 指纹率压缩为一族 Perron 根的极限结构（命题 \ref{prop:pom-collision-renyi-perron-closure}）。
并且最大纤维指数与 Rényi 端点在指数层面严格同一，纤维指数的大偏差/CLT 与容量饱和判别亦由碰撞谱半径族给出可审计封口；其中 Rényi 切线公式、凸锥单调性、独立折叠乘积律与“典型＝最坏”的饱和相刻画进一步把典型指数、最坏指数与多源堆叠统一为一条凸对偶链条（命题 \ref{prop:pom-maxfiber-exponent-equals-renyi-endpoint}，\ref{prop:pom-fiber-index-ldp-from-collision-spectrum}，推论 \ref{cor:pom-fiber-index-clt-variance-certificate}，命题 \ref{prop:pom-typical-fiber-exponent-renyi-tangent}--\ref{prop:pom-saturation-phase-characterization}，命题 \ref{prop:pom-capacity-saturation-discriminant}）。
并且最大纤维顶端谱层的缺口满足 \(12/13\) 步延迟自相似，从而次谱比值在奇偶上呈现不同的指数微分尺度（注 \ref{rem:pom-max-fiber-gap-delay-12-13}）。
在算子代数口径中，交叉视界的指数次乘性与条件互信息下界把“两源残余相关”压缩为对数指数缺口，并给出指数极值态的严格实现与唯一性模板；而在深度谱的粗粒化制度下，尺度增长强制谱间距指数收缩并导致反演条件数指数爆炸（命题 \ref{prop:op-algebra-index-submultiplicativity-commuting-square}，推论 \ref{cor:op-algebra-cmi-lower-bound-by-index-gap}，命题 \ref{prop:op-algebra-index-extremal-entropy-loss-maxfiber}，命题 \ref{prop:xi-coarsegraining-illposedness-exponential-law}）。
同一算子代数模板下，“视界零知识—幺正切片”的可交换极分解、联合谱梳参数化与 Li/对数 Sobolev 证书把临界线集中、零点简单性与尺度间距约束统一为可审计的算子条件，并给出近零知识—健全性—纤维规模的不可突破折衷界（定理 \ref{thm:op-algebra-hzom-commuting-polar-forces-critical-line}，推论 \ref{cor:op-algebra-hzom-critical-zeros-joint-spectrum-comb}、\ref{cor:op-algebra-hzom-simple-zeros-joint-spectrum}、\ref{cor:op-algebra-hzom-zero-spacing-scale-lower-bound}，命题 \ref{prop:op-algebra-li-criterion-operatorized}、\ref{prop:op-algebra-common-unitary-dilation-criterion}、\ref{prop:op-algebra-logsobolev-zero-band}，推论 \ref{cor:op-algebra-nearzk-soundness-fiber-tradeoff}）。
在 Hardy/Smirnov 因子分解口径中，临界线纯相位性把完成化行列式压入内函数类并消去外因子（幅度泄漏），从而泄漏积分严格上界离线零点的加权总偏移；相位导数的 Poisson 叠加律、Clark 测度与 Carleson/BMO 模板进一步把“挑战相位种子—最坏健全性—可恢复速率”统一提升为零点几何的可测度判据（定理 \ref{thm:xi-hardy-inner-characterization}，命题 \ref{prop:xi-leakage-zero-offset-sum}、\ref{prop:xi-phase-derivative-poisson-superposition}，推论 \ref{cor:xi-clark-measure-horizon-semantics}，定理 \ref{thm:xi-carleson-bmo-recovery-template}--\ref{thm:xi-rh-sufficient-template-inner-clark-conservative}）。
在 de Branges/canonical system 口径中，内函数可被唯一提升为 trace-规一化 Hamiltonian 的逆谱模型，并把“极端零知识”、谱移–相对熵恒等式、SU(1,1)–Riemann--Hilbert 正定跳跃、adelic 拼接与 Dirichlet 族零自由带统一为同一条可审计的谱实现链条；在该链条中，离线零点与非正规性障碍及信息复杂度下界互相等价制约（定理 \ref{thm:xi-debranges-canonical-extreme-zk}，命题 \ref{prop:xi-krein-spectral-shift-equals-kl}、\ref{prop:xi-debranges-dpp-sine-kernel-limit}，定理 \ref{thm:xi-su11-rhp-equivalence}，命题 \ref{prop:xi-maslov-spectralflow-entropy-quantization}、\ref{prop:xi-adelic-zk-gluing-critical-line}、\ref{prop:xi-dirichlet-zero-free-from-unified-zk}，推论 \ref{cor:xi-ae-simple-zeros-under-ac-seeds}，命题 \ref{prop:xi-nonnormality-obstruction-offcritical}、\ref{prop:xi-bm-density-forces-linear-transcript-length}）。

\paragraph{window-\(6\) 锚点下的离散统一收束：黄金时间—6D 微态—Zeckendorf--GUT 三分支}
除按机制分组的全局结论外，本文亦给出一组可独立成文、且完全以有限对象封闭表述的终端命题组：其以 window-\(6\) 作为最小分辨率锚点，把“时间=单生成扫描迭代计数”的协议化刻度、\(\Omega_6\to X_6\) 的 \(64\to 21\) 折叠压缩、uplift 尾部三分支 \(\{21,34,55\}\)、以及 Zeckendorf--Fibonacci 边界塔所诱导的 GUT 顶层项对齐与家族数锁定，统一压缩为一条可审计的离散分支结构（见 \S\ref{subsubsec:window6-terminal-propositions}；主结论为定理 \ref{thm:terminal-6d-microstate-golden-time-gut-branch}）。

\paragraph{四项结论性命题：从地址语义到证书化预算与秩稳定性}
为便于把正文中分散出现的结论直接复用到协议设计与审计流程中，可将若干核心断言压缩为四条互补的“证书接口”（均严格依赖既有编号结论而非解释性表述）：
\begin{enumerate}
  \item \textbf{地址语义的对象化：\(\mathsf{NULL}\) 作为局部截面障碍。}
  在前缀站点上把协议可接受且单值可核验的读出/证书组织为（预）层 \(\mathscr F\) 后，“地址先于取值”的语义需区分为两级：实际全局读出由 \(\mathscr F(a)\neq\varnothing\) 判定，而兼容局部家族则由 sheafification 纤维 \(\mathscr F^\sharp(a)\neq\varnothing\) 记录；只有在 \(a\) 处已经成层时，兼容局部存在才等价于原始对象层中的全局读出（\S\ref{subsec:prefix-site-null}，命题 \ref{prop:null-as-local-section-obstruction}；并见命题 \ref{prop:recursive-addressing-refinement} 与推论 \ref{cor:recursive-addressing-visible-sigma-nonexpansion}）。
  由此，语义/协议/碰撞三类空值在同一站点语言下分别对应于局部对象缺失、兼容局部仅停留在 sheafification、以及单值唯一化失败三种互不混淆的机制（推论 \ref{cor:null-trichotomy-local-section}）。

  \item \textbf{折叠“清晰度”的二阶证书：碰撞膨胀 \(\Rightarrow\) 熵缺口上界。}
  在均匀微态基线下，折叠推前分布 \(w_m(x)=d_m(x)/2^m\) 相对类型均匀分布的熵缺口 \(\mathrm{gap}\) 可被二阶碰撞指纹
  \(C_{\mathrm{col}}(m)=|X_m|\sum_x w_m(x)^2\)
  的对数上界证书化（定理 \ref{thm:pom-spectral-information-sandwich} 与定义 \ref{def:pom-collision-inflation}），并可由表 \ref{tab:fold_multiplicity_histogram} 在有限分辨率上直接审计。
  该结论将“缩略图的清晰度”从解释性说法压缩为一个仅依赖二阶统计量的可计算证书。

  \item \textbf{单射化预算的最优对偶包络：分位配额律的 \(\inf_{q\ge 2}\) 规范化。}
  prime-register 的配额阈值可由纤维指数尾部分位 \(B_m(\varepsilon)=\exp(m\gamma_m(\varepsilon))\) 给出，并在区间读出/重建口径下具有锐性（\S\ref{subsec:pom-quantile-budget-law}，推论 \ref{cor:pom-prime-register-quantile-budget}）。
  进一步，整数矩阶 \(q\) 所导出的全部 Chernoff/配额不等式族可被压缩为对 \(q\ge 2\) 的单一最优下包络，从而得到一条参数自由的寄存器指数预算律；并可用有限 \(q\)-窗口的“最强可计算尾界”实现可执行证书（推论 \ref{cor:pom-prime-register-q-quota-family} 及 \S\ref{subsec:pom-quantile-budget-law} 的对偶设计式）。

  \item \textbf{模式轴完备性的充要判据：Hankel 秩一致有界 \(\Longleftrightarrow\) ACC-Rank。}
  在统一最小性与协议闭包假设下，一旦残差轴预算已排除碰撞空值，则“不会因需要新增可见自由度而触发 \(\mathsf{NULL}\)”的模式轴完备性与 Hankel 秩的一致有界性严格等价（定理 \ref{thm:xi-completeness-iff-hankel-bounded}），从而把 completeness 的关键缺口转换为可回归测试的秩稳定性扫描（注 \ref{rem:xi-acc-rank}）。
  与分位配额证书合并后，completeness 的审计任务可被压缩为“秩稳定证书 \(\times\) 配额证书”的二元后端核验（命题 \ref{prop:xi-completeness-auditable-decomposition}）。
\end{enumerate}

\paragraph{从射影压力算子到无回头素数影子的低阶闭式}
为将“整数点证书—离散锚点—低阶行列式闭式”三类接口直接对接至连续压力、最小 Markov 完备化与 prime shadow 审计，下列断言给出一组互补的可核验命题（证明与可复核证书分布于 POM、window-\(6\) 与真实输入核附录）：
\begin{enumerate}
  \item \textbf{射影压力算子的实参数提升与整数点严格退化。}
  对确定性在线核/转导器 \(\mathcal{T}\)，由输出分桶矩阵 \(B_b\)、单纯形 \(\Delta\) 与射影动力学 \((\Phi_b,g_b)\) 定义射影转移算子
  \[
  (\mathcal{L}_q f)(w)=\sum_{b\in A_{\mathrm{out}}} g_b(w)^{\,q}\,f(\Phi_b(w))\qquad (q\ge 1)
  \]
  （定义 \ref{def:pom-projective-transfer-operator}）。
  对整数 \(q\ge 2\)，其在齐次多项式子空间上的限制严格退化为对称幂商 moment-kernel，从而谱半径在整数点满足
  \(
  \rho(\mathcal{L}_q|_{\mathcal{H}_q})=r_q
  \)
  （定理 \ref{thm:pom-projective-operator-degenerates-to-moment-kernel}）。
  因此可把离散压力 \(\log r_q\) 升级为一条可执行的连续压力曲线候选：令 \(q=t+1\ge 1\)，取 \(\mathcal{L}_{t+1}\) 的 Perron 根 \(\lambda(t+1)\) 并置 \(\Lambda(t)=\log\lambda(t+1)-\log|A_{\mathrm{out}}|\)，在整数点严格插值 \(\Lambda(q-1)=\log r_q-\log|A_{\mathrm{out}}|\)（注 \ref{rem:pom-projective-pressure-route}）。

  \item \textbf{window-\(6\) 纤维的仿射几何分类与“线性核”不可实现性。}
  将 \(\Omega_6=\{0,1\}^6\) 视为 \(\mathbb{F}_2^6\) 并记 dyadic 折叠 \(F_6=\mathrm{Fold}^{\mathrm{bin}}_6\)。
  其 \(21\) 条纤维 \(F_w=F_6^{-1}(w)\subset\mathbb{F}_2^6\) 出现完全分类：恰有 \(8\) 条为仿射 \(1\)-平坦（大小 \(2\)），恰有 \(3\) 条为仿射 \(2\)-平坦（大小 \(4\)，且稳定类型恰为 \(\{\texttt{000000},\texttt{000010},\texttt{101000}\}\)），其余 \(10\) 条均非仿射子空间（定理 \ref{thm:terminal-foldbin6-fiber-affine-geometry}；表 \ref{tab:foldbin6_fiber_affine_geometry}）。
  由此推出线性核障碍：不存在 \(\mathbb{F}_2\)-线性映射 \(L:\mathbb{F}_2^6\to\mathbb{F}_2^k\) 使纤维划分等同于陪集划分（推论 \ref{cor:terminal-foldbin6-no-linear-kernel}）。

  \item \textbf{类型邻接图 \texorpdfstring{$\Gamma_6$}{Gamma6} 的完全刚性与连续对称禁戒。}
  令 \(\Gamma_6\) 为由立方体边推送诱导的类型邻接图（命题 \ref{prop:terminal-gamma6-rigidity}）。
  则其自同构群平凡 \(\mathrm{Aut}(\Gamma_6)=\{1\}\)，并且图连通且直径为 \(3\)（推论 \ref{cor:terminal-gamma6-diameter-simple-spectrum}）。
  因此任何紧致群若在 \(X_6\) 上作用并保持邻接关系，则该作用必为平凡作用（推论 \ref{cor:terminal-gamma6-no-continuous-nonabelian-symmetry}）：类型重标号自由度被邻接结构内生消除。

  \item \textbf{edge--flux skeleton 的 Smith 判别式与 Levi 可约性强迫。}
  在四块 edge--flux skeleton 的相互作用矩阵 \(E\in\mathbb{Z}^{4\times4}\) 上，穷举审计给出完全可核验的 Smith 正规形
  \(
  \mathrm{SNF}(E)=\mathrm{diag}(1,1,3,3450)
  \)，从而 \(\mathrm{coker}(E)\cong \mathbb{Z}/3\mathbb{Z}\oplus \mathbb{Z}/3450\mathbb{Z}\) 且 \(\det(E)=10350\)（推论 \ref{cor:window6-edge-flux-skeleton-smith-nonsimple}）。
  该整数判别式排除其与秩 \(4\) 的有限型不可约 Cartan 数据同值，从而在“连续李对称的离散痕迹”解释下强迫最小 Levi 型可约性。

  \item \textbf{推送核 \texorpdfstring{$P_6$}{P6} 的最小热化闭合：可逆性、平衡测度与显式谱隙常数。}
  在“纤维内瞬时热化”的最小无记忆闭包机制下，稳定类型过程闭合为 \(X_6\) 上的可逆马尔可夫链，其转移核为推送核 \(P_6\)，平稳分布为 \(\pi_6(w)\propto d^{\mathrm{bin}}_6(w)\)（定理 \ref{thm:terminal-foldbin6-pushforward-markov}）。
  并且谱隙审计给出显式常数：\(\lambda_2(P_6)\approx 0.4841208\)、\(\lambda_\star(P_6)\approx 0.6030940\)（推论 \ref{cor:terminal-foldbin6-minimal-markov-closure-gap} 与式 \eqref{eq:window6_pushforward_markov_spectral_gap}），从而对任意初始分布 \(\mu\) 有
  \(
  \|\mu P_6^t-\pi_6\|_{\mathrm{TV}}\le C(\mu)\lambda_\star(P_6)^t
  \)
  的指数混合上界。

  \item \textbf{压力--有限部的局部双几何：恒压漂移容量与曲率错配谱。}
  在真实输入 \(40\) 状态核的三维观测下，以压力 \(P(\theta)\) 与 Abel/Mertens 有限部常数 \(\log\mathfrak{M}(\theta)\) 作为两条独立热力学量，定义恒压漂移容量向量
  \[
  v_\perp=\nabla\log\mathfrak{M}(0)-\mu\,\nabla P(0),\qquad
  \mu=\frac{\langle\nabla P(0),\nabla\log\mathfrak{M}(0)\rangle}{\|\nabla P(0)\|^2}
  \]
  （定义 \ref{def:real-input-40-isobaric-drift-vector}）。
  数值上 \(\mu\approx 0.7113818\)、\(v_\perp\neq 0\)（命题 \ref{prop:real-input-40-isobaric-drift-capacity}），从而 \(\log\mathfrak{M}\) 不可能仅为 \(P\) 的一元复合（推论 \ref{cor:real-input-40-logM-not-just-P}）。
  更进一步，曲率错配谱由广义本征问题 \(Hv=\mu\,gv\) 组织（命题 \ref{prop:real-input-40-logM-curvature-mismatch-spectrum}），并给出有限部漂移与 near-\(1\) 主极点缺口之间的二次换算律（命题 \ref{prop:real-input-40-space-time-conversion-law}）。
  此外，由 \(s\)-平面坐标 \(s=\log\rho/\log\lambda\) 与谱比混合时间 \(\tau_{\mathrm{mix}}=1/(-\log(\rho/\lambda))\) 的定义，缺口 \(1-\Re(s)\) 与混合时间满足严格互易恒等式 \((1-\Re(s))\,\tau_{\mathrm{mix}}=1/\log\lambda\)（推论 \ref{cor:near1-pole-taumix-m2}）。

  \item \textbf{二维加权 \texorpdfstring{$\zeta(z;u,v)$}{zeta(z;u,v)} 的低阶行列式闭式与 primitive 谱的有限阶递推。}
  在真实输入核上对碰撞与输出引入二维联合势并定义 \(\Delta(z;u,v)=\det(I-zM_{u,v})\)（见 \S\ref{app:real-input-40-zeta-uv}）。
  行列式分母具有显式 \((2+8)\) 次闭式分解 \(\Delta(z;u,v)=(vz^2-1)P_8(z;u,v)\)（定理 \ref{thm:real-input-40-det-uv-2plus8}），故对每个固定 \((u,v)\) 的迹序列 \(a_n(u,v)=\Tr(M_{u,v}^n)\) 满足阶数至多为 \(10\) 的常系数线性递推，并可进一步以 Möbius 反演近线性时间恢复 primitive 轨道谱（命题 \ref{prop:real-input-40-zeta-uv-recurrence}）。

  \item \textbf{Bartholdi/Ihara 型无回头（测地）prime shadow 的低阶行列式闭式。}
  在真实输入 \(40\) 状态核的 essential 核心边空间上构造无回头边移位 \(H_{\mathrm{nb}}\) 后，尽管 \(H_{\mathrm{nb}}\) 为 \(48\times48\) 矩阵，其行列式分母却压缩为一个 \(15\) 次整系数多项式（定理 \ref{thm:real-input-40-geodesic-det}）：
  \[
  \det(I-zH_{\mathrm{nb}})
  =1-4z^3-9z^4-12z^5-16z^6-11z^7-11z^8-4z^9-z^{10}
  +8z^{11}+12z^{12}+10z^{13}+8z^{14}+2z^{15}.
  \]
  由根模排序得到最大非 Perron 模 \(\Lambda_{\mathrm{nb}}\) 并给出严格 Ramanujan 裕度 \(\Lambda_{\mathrm{nb}}/\sqrt{\rho_{\mathrm{nb}}}<1\)，从而在素数长度上获得指数底严格小于 \(\rho_{\mathrm{nb}}^{1/2}\) 的次指数误差（推论 \ref{cor:real-input-40-geodesic-ramanujan-margin}）。
\end{enumerate}

\paragraph{谱—算术封口：从共振闭包到端点簇刚性}
为将关键的“不可消去结构”直接对接至能量、谱与算术的审计接口，在上文四条证书接口之外，下列断言给出一组可独立引用的封口命题（证明与证书分布于附录及 POM/同步核章节）：
\begin{enumerate}
  \item \textbf{共振闭包对碰撞能量的不可消去下界。}
  以折叠碰撞概率 \(\mathsf{Col}_m\) 与共振梯常数 \(C_\varphi(u)>0\) 为记号（定义 \ref{def:fold-normalized-amplitude} 与定义 \ref{def:fold-critical-resonance-ladder}），对任意固定 \(U\in\NN\) 有严格下界（推论 \ref{cor:fold-collision-gap-resonance-ladder-U}）
  \[
  \liminf_{m\to\infty}F_{m+2}\,\mathsf{Col}_m\ \ge\ 1+2\sum_{u=1}^{U}C_\varphi(u)^2.
  \]
  因此“碰撞趋均匀/谱收缩”在能量层面被可枚举的共振闭包结构性钉住；任何仅基于 \(\|\widehat d_m\|_\infty\) 控制来推出 \(\ell^2\) 均匀化的路线，在原始 Fold 口径下均会遭遇该下界机制的阻断。

  \item \textbf{Fold 的群代数严格分解诱导确定型 Riesz 乘积；碰撞谱等价于其能量平均。}
  群环生成函数满足严格乘积分解 \(D_m(x)\equiv\prod_{j=1}^m(1+x^{F_{j+1}})\ (\mathrm{mod}\ x^{F_{m+2}}-1)\)（命题 \ref{prop:fold-multiplicity-group-algebra}），等价于 \(d_m\) 在循环群代数上的卷积分解；其 Fourier 变换完全因子化（定理 \ref{thm:fold-spectrum-factorization}）。
  因而
  \(
  F_{m+2}\mathsf{Col}_m=\sum_k a_m(k)^2
  \)
  可视作由 Fibonacci 频率驱动的有限 Riesz 乘积能量平均（定义 \ref{def:fold-normalized-amplitude}），从而把碰撞极限问题严格转译为“振幅分布”问题而非“零点几何”问题。

  \item \textbf{实退化点给出的可审计实参数阈值：输出位势的局部展开窗口。}
  在真实输入 \(40\) 状态核的输出位势一参数加权中，正实退化点把 \(u>0\) 分段为谱支排序不变的相图区间（注 \ref{rem:real-input-40-output-potential-degeneracy-phase}；表 \ref{tab:real_input_40_output_potential_spectral_collision_roots}）。
  距离 \(u=1\) 最近的退化点给出实参数阈值，并在 \(\theta=\log u\) 坐标中读出一个统一保守半径
  \[
  |\theta|<\theta_0,\qquad \theta_0=0.0533477827\ldots,
  \]
  可作为围绕 \(\theta=0\) 的局部拟合/展开在实参数方向上的可审计有效域封口。

  \item \textbf{自对偶完成化的压力偶性定律：奇阶累积量全灭的结构根源。}
  在自对偶完成化下，压力可写为 \(P(\theta)=\theta/2+\log y(2\cosh(\theta/2))\)（命题 \ref{prop:weighted-completion-Q}），从而 \(\Lambda(\theta):=P(\theta)-\theta/2\) 为偶函数（定理 \ref{thm:sync-kernel-gallavotti-cohen}）。
  等价地 \(P(\theta)-P(-\theta)=\theta\)，故除线性项外 \(P\) 在 \(\theta=0\) 的所有奇阶导数严格为零；因此任何由 \(P\) 导出的中心化涨落展开，其首个非高斯修正必落在偶 Hermite 扇区，而所有奇阶偏斜项在结构层面被排除。

  \item \textbf{最小三态“统一$\to$破缺”有效核的 Newman 阈值：闭式相变原型。}
  对三态有效核族 \(A_{\mathrm{eff}}(\alpha)=U+\alpha R\)（命题 \ref{prop:pom-a2-gut-spectral-splitting}），核版 Newman 偏移 \(\delta(\alpha)=\log(\Lambda(\alpha)^2/\rho(\alpha))\) 存在唯一临界点 \(\alpha_\star>1\) 并发生“负实慢模态贴边”，且该阈值具有完全代数证书（命题 \ref{prop:pom-a2-gut-newman-threshold}）：
  \[
  \alpha_\star=\frac{x_\star(x_\star^3-x_\star^2+1)}{x_\star^2-x_\star+2}=2.3609261183\ldots,
  \]
  其中 \(x_\star>1\) 为 \(x^6-x^5-4x^4+8x^3-6x^2-4x+2=0\) 的唯一正实根。

  \item \textbf{Hankel 见证的 \texorpdfstring{$p$}{p}-进刚性：素数在最小实现维数证书中的激活阈值。}
  非零 Hankel 主子式不仅给出秩下界，还携带不可约的 \(p\)-进信号：在最小两例中
  \(
  \det(H_{k=4})=2^{17}\cdot 3^2
  \)
  而
  \(
  \det(H_{k=5})=2^{13}\cdot 3^5\cdot 5^7
  \)，
  因而 \(p=5\) 在 \(k=5\) 处以指数 \(7\) 首次显性进入（注 \ref{rem:pom-hankel-stiffness-p5}）。

  \item \textbf{模 \(2\) 影子的自动序列结构：差分湮灭与周期上界。}
  在共振窗口 \(q=9,\dots,17\) 上，特征多项式的模 \(2\) 影子坍缩为
  \(
  P_q(\lambda)\equiv \lambda^{a_q}(\lambda+1)^{b_q}\ (\mathrm{mod}\ 2)
  \)
  并诱导湮灭式 \(E^{a_q}(E+1)^{b_q}S_q\equiv 0\ (\mathrm{mod}\ 2)\)。
  因而 \(S_q(m)\bmod 2\) 为 \(2\)-自动序列且周期满足 \(T_q\mid 2^{\lceil\log_2 b_q\rceil}\)，并在表 \ref{tab:fold_collision_charpoly_mod2_shadow_q9_17} 给出具体上界（注 \ref{rem:pom-charpoly-mod2-shadow}）。

  \item \textbf{消元结果式的端点簇刚性：\texorpdfstring{$u$}{u}-adic 退化分解钉死端点为离散有理点。}
  对同步核速率曲线的原始消元结果式 \(R_{\mathrm{raw}}(\alpha,u)\)，其 \(u\)-adic 赋值满足 \(v_u(R_{\mathrm{raw}})=7\) 且具有刚性分解 \(7=3+4\)（命题 \ref{prop:rate-resultant-uadic-split-7}；表 \ref{tab:sync_kernel_rate_curve_resultant_degree}）。
  首项结构强制
  \[
  R_{\mathrm{raw}}(\alpha,u)=\pm\,u^7\,\alpha(\alpha-1)(4\alpha-1)^4+\text{更高阶项},
  \]
  从而把 \(\alpha=0,1\) 与 \(\alpha=\tfrac14\) 的端点簇（四重）提升为由退化几何不可避免地产生的离散有理点；并在对偶闭合 \((u,\alpha)\mapsto(1/u,1-\alpha)\) 下对称强制出现 \(\alpha=\tfrac34\) 的对应端点簇。
\end{enumerate}

\paragraph{共振窗与常数层：从 CM 相位锁定到 Abel-归一化可积核}
\begin{enumerate}
  \item \textbf{共振复谱的 CM 相位锁定：虚二次域 \texorpdfstring{$\QQ(\sqrt{-15})$}{Q(sqrt(-15))} 的自发显影。}
  在共振窗 \(q=9,\dots,17\) 中，令 \(\mu_q\) 为最大模非实根（取 \(\Im(\mu_q)\le 0\)），并置单位相位
  \(\zeta_q:=\mu_q/|\mu_q|\in\mathbb{S}^1\)。
  令
  \(
  z_\ast:=(-1-\mathrm{i}\sqrt{15})/4
  \)
  为 \(2z^2+z+2=0\) 的单位圆根，则表 \ref{tab:fold_collision_shadow_spectral_packet_algebraic_phase_residual_q9_17} 给出可复算的相位锁定证书：\(|\zeta_q-z_\ast|\) 在全窗保持 \(10^{-2}\) 量级并在 \(q=12\) 达到 \(10^{-4}\) 量级，同时代数残差 \(\rho_q=|2\zeta_q^2+\zeta_q+2|\) 保持小量级（注 \ref{rem:pom-resonance-trimodal-phase-lock}）。
  并且由式 \eqref{eq:pom-phase-lock-residual-trace-equivalence}，\(\rho_q\) 等价于实迹偏差 \(|(\zeta_q+\zeta_q^{-1})+\tfrac12|\) 的幅度，从而锁定强度可被压缩为单一实标量。
  因此共振复模态的单位圆相位可被压缩到虚二次域 \(\QQ(\sqrt{-15})\) 的单一代数轨道上，形成一个完全内禀且可审计的 CM 型谱指纹。

  \item \textbf{复模态固定频率的“通用二阶振子”：由 \texorpdfstring{$2z^2+z+2$}{2z^2+z+2} 给出的归一化湮灭子。}
  置归一化移位 \(T:=r_{q-2}^{-1}E\)（\(E\) 为 \(m\)-方向前移算子），则 \(\mathcal L_{\mathrm{CM}}(T):=2T^2+T+2\) 对 \(\nu=z_\ast,\overline{z_\ast}\) 的指数模态严格湮灭；在共振窗内该算子可作为“相位感知”的三步消振接口，与二周期 Aitken 消去器互补（注 \ref{rem:pom-cm-annihilator}）。
  等价地，对归一化实振子 \(u(m)=\zeta^{m}+\zeta^{-m}\)（\(\zeta\in\mathbb{S}^1\)）有
  \(2u(m+2)+u(m+1)+2u(m)=\bigl(2(\zeta+\zeta^{-1})+1\bigr)u(m+1)\)，
  因而其二阶综合征幅度逐点正比于 \(\rho(\zeta)=|2\zeta^2+\zeta+2|\)，从而给出无需复求根的时域审计接口（注 \ref{rem:pom-cm-annihilator}）。

  \item \textbf{两步重标度点的 Poisson 峰 \texorpdfstring{$\Rightarrow$}{=>} Lorentz 极限：Breit--Wigner 轮廓的离散预极限与“高—宽互易”。}
  令 \(\eta_q:=|\mu_q|/r_{q-2}\uparrow 1\) 且 \(a_q:=\mu_q/r_{q-2}\in\mathbb{D}\)。
  则对应 Blaschke 因子 \(B_q(z)=(z-a_q)/(1-\overline{a_q}z)\) 的边界相位导数为 Poisson 核 \(P_{\eta_q}(t-\theta_q)\)，并在缩放窗口 \(t=\theta_q+\varepsilon_q u\)（\(\varepsilon_q=1-\eta_q\)）上收敛到 Lorentz 核 \(2/(1+u^2)\)；
  同时 Poisson 核面积恒为 \(2\pi\)，峰高量级为 \(\varepsilon_q^{-1}\)（注 \ref{rem:pom-resonance-poisson-lorentz-prelimit}）。
  并且由表 \ref{tab:fold_collision_shadow_spectral_packet_q9_17} 的审计读数，若定义逐阶漂移
  \(
  \Delta\theta_q:=|\arg(\mu_{q+1})-\arg(\mu_q)|
  \)
  （\(q=9,\dots,16\)），则有一致支配
  \(
  \max_{9\le q\le 16}\Delta\theta_q/\varepsilon_q
  =0.0029/0.0204<0.143
  \)，从而 \(\Delta\theta_q\le 0.143\,\varepsilon_q\)（注 \ref{rem:pom-resonance-poisson-lorentz-prelimit}）。

  \item \textbf{\texorpdfstring{$q=12$}{q=12} 的四重交汇点：四条审计不变量的同步。}
  在共振窗内，枢轴点 \(q=12\) 同时满足：（i）\(|\zeta_q-z_\ast|\) 取全窗最小（表 \ref{tab:fold_collision_shadow_spectral_packet_algebraic_phase_residual_q9_17}）；（ii）负实外侧根通道切换 \(|\lambda_q^-|/r_{q-2}\) 在该点穿越阈值 \(1\)（表 \ref{tab:fold_collision_shadow_spectral_packet_q9_17}）；（iii）共振缺口 \(\Delta(q)\) 在已核验序列中唯一取 \(0\)（注 \ref{rem:pom-delta-cq-threshold}）；（iv）在模 \(11\) 的嵌入切片中出现 \(\chi_{A_2}\mid P_{12}\ (\mathrm{mod}\ 11)\) 的嵌入旗标（命题 \ref{prop:pom-charpoly-modp-a2-embedding}；表 \ref{tab:fold_collision_charpoly_mod11_fingerprint_q12_15}，同表 \(q=13,15\) 的嵌入列为 \(0\)）。
  该四重同刻把 \(q=12\) 识别为一个结构转折点：其同步性来自两步重标度、代数相位残差、最小递推阶数与有限域嵌入四条互不依赖的审计管线（注 \ref{rem:pom-resonance-window-final-structural-conclusions}）。

  \item \textbf{Dirichlet 实部坐标的漂移律：共振复极点条带向 \texorpdfstring{$\Re(s)=1$}{Re(s)=1} 的确定性逼近。}
  对任一根 \(\xi\) 定义 Dirichlet 实部坐标 \(\sigma_q(\xi):=\log|\xi|/\log r_q\)。
  在共振窗内，\(\sigma_q(\mu_q)>1/2\) 且单调上升（表 \ref{tab:fold_collision_shadow_spectral_packet_q9_17}），形成一条向 \(\Re(s)=1\) 逼近的非平凡极点条带；其漂移可分解为两步重整化比与幅值修正的和（注 \ref{rem:pom-shadow-spectral-packet-q9-17}）
  \[
  \sigma_q(\mu_q)=\frac{\log r_{q-2}+\log(|\mu_q|/r_{q-2})}{\log r_q}.
  \]
  因而在端点假设 \(|\mu_q|/r_{q-2}\to 1\) 与 \(R_q:=r_q/r_{q-2}\to\varphi\) 下，复模态缺口指数 \(g_q^{(c)}:=1-\sigma_q(\mu_q)\) 满足条件渐近 \(g_q^{(c)}\sim 2/q\)，给出一个由端点 \(\varphi\) 唯一决定的可证伪漂移速度常数。

  \item \textbf{判别式 \texorpdfstring{$-15$}{-15} 的模 \texorpdfstring{$p$}{p} 分裂选择律：相位锁定二次式的跨素数可证伪指纹。}
  二次式 \(2z^2+z+2\) 在 \(\mathbb{F}_p\) 上的分裂性由 \(\bigl(\frac{-15}{p}\bigr)\) 控制，并且在分歧素数 \(p\in\{3,5\}\) 处出现重根退化；
  因此若相位锁定二维块确由该二次式支配，则跨素数的模 \(p\) 因子分解指纹应呈系统性分裂/惰化二分（命题 \ref{prop:pom-phase-lock-discriminant-minus15-splitting} 与注 \ref{rem:pom-phase-lock-discriminant-minus15-prediction}）。

  \item \textbf{有限域嵌入旗标的表示论含义：低阶核作为高阶核的局部子动力学出现。}
  若存在 \(\chi_{A_2}(\lambda)\mid P_q(\lambda)\) 于 \(\mathbb{F}_p[\lambda]\)，则模 \(p\) 的递推状态空间中出现一个与 \(A_2\) 同型的嵌入子动力学（一个三维不变子表示；推论 \ref{cor:pom-charpoly-modp-a2-subdynamics}）。
  已核验切片给出在 \((q,p)=(14,7)\) 处的 \(A_2\)-影子嵌入旗标（命题 \ref{prop:pom-charpoly-modp-a2-embedding}），从而提供一个可用于局部—全局比较的有限域表示论接口。

  \item \textbf{判别式素支撑的硬指纹：\(\{2,7,11\}\) 的强制整除谱。}
  对 \(P_q(\lambda)=\lambda^{a_0}P_q^\ast(\lambda)\)（\(P_q^\ast(0)\ne 0\)）定义 \(\mathrm{Disc}^\ast(P_q):=\mathrm{Disc}(P_q^\ast)\)。
  在共振窗的已核验切片内，由 repeat? 非平方自由旗标可直接推出 \(p\mid \mathrm{Disc}^\ast(P_q)\) 的强制整除性；其中模 \(2\) 影子给出全窗 \(2\mid \mathrm{Disc}^\ast(P_q)\)，而模 \(7\)、模 \(11\) 的切片分别在若干 \(q\) 上强制出现（注 \ref{rem:pom-discriminant-prime-support}）。

  \item \textbf{复合模数下的周期坍缩：\texorpdfstring{$q=13$}{q=13} 在模 \texorpdfstring{$182$}{182} 下的最终周期整除 \texorpdfstring{$168$}{168}。}
  由模 \(2/7/13\) 的周期上界与中国剩余结构，在对齐共同瞬态后有 \(\mathrm{Per}_{182}(q=13)\mid 168\)（注 \ref{rem:pom-mod182-period-collapse-q13}）。

  \item \textbf{有限域因子抽取定律：\texorpdfstring{$\chi\mid P_q$}{chi|Pq} \texorpdfstring{$\Rightarrow$}{=>} 存在可显式构造的 \texorpdfstring{$\chi$}{chi}-子递推。}
  若在 \(\mathbb{F}_p[\lambda]\) 中 \(P_q\equiv \chi Q\)，则对任一解序列 \(S_q\) 取 \(S_q^{(\chi)}:=Q(E)S_q\) 即满足 \(\chi(E)S_q^{(\chi)}\equiv 0\pmod p\)（命题 \ref{prop:pom-modp-factor-extraction}）。
  对嵌入旗标 \(\chi_{A_2}\mid P_q\) 可据此抽取 \(A_2\)-影子子递推（注 \ref{rem:pom-a2-subrecurrence-extraction}）。

  \item \textbf{模 \texorpdfstring{$11$}{11} 的 Jordan-阶分裂：不可约调制阶在 \texorpdfstring{$q=12$}{q=12} 与 \texorpdfstring{$q=13,15$}{q=13,15} 间严格不同。}
  由表 \ref{tab:fold_collision_charpoly_mod11_fingerprint_q12_15}，\(q=12\) 含二次不可约因子的平方而 \(q=13,15\) 仅含线性平方因子；因此二者的 Jordan 调制在不可约层面具有不同的最小次数下界，并可在不求复谱的条件下作为共振窗内部细分判据（注 \ref{rem:pom-charpoly-mod11-modulation}）。

  \item \textbf{厚尾亏损函数 \texorpdfstring{$\delta(q)$}{delta(q)} 的二阶常数含义：由 \texorpdfstring{$r_q$}{rq} 的次主项决定的全局标量不变量。}
  Paley--Zygmund 下包络自然诱导非负缺口
  \[
  \delta(q):=\log\varphi+\log r_{2q}-2\log r_q\ge 0
  \]
  （推论 \ref{cor:pom-spectrum-intermittency-deficit}）。
  若端点展开 \(\log r_q=\frac{q}{2}\log\varphi+a+o(1)\) 成立，则由代入直接得到
  \[
  \lim_{q\to\infty}\delta(q)=\log\varphi-a,
  \]
  因此 \(\delta(q)\) 不仅度量“厚尾保证指数”的亏损，还反演出碰撞谱端点的二阶自由能常数 \(a\)，从而构成一个可跨版本对比的全局标量指纹。

  \item \textbf{prime-orbit Mertens 常数的谱敏感性：由 \texorpdfstring{$\chi_A'(r)$}{chi'(r)} 的乘积结构给出的放大量纲。}
  对有限状态碰撞核 \(A\)（SFT 邻接矩阵），其 Abel-归一化主极点留数常数满足闭式
  \[
  C(A)=\frac{r^{d-1}}{\chi_A'(r)}=\prod_{\lambda_j\in\sigma(A)\setminus\{r\}}\Bigl(1-\frac{\lambda_j}{r}\Bigr)^{-1},
  \]
  从而 \(C(A)\) 仅由 \((\chi_A,r)\) 决定并在任何保持同一 \(\chi_A\) 的非负整数 realization 下保持不变；并且对 \(A_2,A_3,A_4\) 可将 \(C(A_q)\) 升格为整系数最小多项式的唯一实根证书，从而形成“留数常数 \(\Rightarrow\) 代数方程 \(\Rightarrow\) 可审计证书链”的闭环（见 \S\ref{subsec:pom-s4}）。

  \item \textbf{一个完全可积的 Abel-归一化 \texorpdfstring{$\zeta$}{zeta} 核：极点集合落在 \texorpdfstring{$\QQ(\sqrt5)$}{Q(sqrt5)} 单位群轨道上。}
  在真实输入 \(40\) 状态核的 Abel-归一化坐标 \(r=\lambda z\) 下，归一化 \(\widetilde{\zeta}_M(r)\) 具有完全分解的行列式表达，其全部奇点可被显式枚举为
  \(
  r\in\{1,\pm\lambda,\lambda^2,\varphi^3,-\varphi\}
  \)
  且除 \(r=1\) 外均满足 \(|r|>1\)，因此在 \(|r|<1\) 内无任何次级极点障碍；并且主极点留数常数给出闭式代数值
  \(
  C=(47+21\sqrt5)/100
  \)
  （\S\ref{app:real-input-40-abel}）。
\end{enumerate}

\paragraph{端点簇—中心截面—谱凝聚：从四分支单值化到 \texorpdfstring{$p^{-6}$}{p^-6} 衰减律}
为将“端点退化 jet—中心截面根盘—高阶碰撞核谱凝聚”三类证书统一压缩为可直接复用的审计接口，下列断言给出一组互补的可核验命题（证书分布于同步核加权消元附录与碰撞谱表）：
\begin{enumerate}
  \item \textbf{端点簇的四分支单值化与局部二面体型单群。}
  在同步核速率曲线的退化端 \(u\to 0\) 上，特征曲线满足 \(F(\lambda,0)=\lambda^5(\lambda-1)\)，其五重零根包刚性分裂为四条 Puiseux 支与一条解析小根支（命题 \ref{prop:rate-resultant-uadic-split-7}）：
  \[
  \lambda(u)=\zeta\,u^{1/4}+O(u^{1/2})\quad(\zeta^4=1),
  \qquad
  \lambda(u)=-u+O(u^2).
  \]
  因此围绕 \(u=0\) 的解析延拓对四条 Puiseux 支诱导 \(u^{1/4}\mapsto i\,u^{1/4}\) 的四循环单值化；并且侧翼端点簇的正确统一化变量必须取 \(u=s^4\)（\(s=u^{1/4}\)），从而对同一 \(u>0\) 形成四页覆盖 \(s\in\{\pm u^{1/4},\pm i\,u^{1/4}\}\)（注 \ref{rem:rate-algebraic-elimination-structure}）。
  结合复共轭对分支的作用，端点簇在局部上携带包含 \(C_4\) 的二面体型单群（局部 monodromy 至少为 \(D_4\)）；因此任何把该端点簇误判为“平方根分歧”的实现，必在环绕 \(u=0\) 时产生错误的支切换。

  \item \textbf{中心截面无实奇点迫使最近奇点纯复；最近根给出可审计收敛盘与系数振荡指纹。}
  中心截面分解式 \(R(\tfrac12,u)=-(u-1)^6Q(u)/64\) 中的 \(Q(u)\in\ZZ[u]\) 具有严格回文自倒易性 \(Q(u)=u^{16}Q(1/u)\)，并且 Sturm 计数给出 \(Q\) 在实轴上无零点，从而 \(Q(u)>0\) 对所有 \(u\in\RR\) 成立（命题 \ref{prop:rate-center-Q-reciprocal-no-real}）。
  因此围绕 \(u=1\) 的局部解析展开，其收敛障碍必由复奇点对决定；最近奇点距离
  \(
  R_\star=\min\{|u-1|:\ Q(u)=0\}
  \)
  与最近根 \(u_\star\) 已给出可复核数值证书（推论 \ref{cor:rate-center-Rstar}）：
  \[
  R_\star\approx 1.0507607822,
  \qquad
  u_\star\approx -0.0492815273\pm 0.0557359652\,\mathrm{i}.
  \]
  按标准奇点分析/系数转移原理，任何以该中心截面（或由其构造的完成化行列式）生成的泰勒系数序列在高阶上必出现由 \(\arg(u_\star-1)\) 锁定的准周期振荡；其频率与衰减率共同构成“最近奇点”口径下可审计的谱指纹。

  \item \textbf{最小模数障碍：\texorpdfstring{$p=5$}{p=5} 为中心截面局部展开的唯一素数排除点。}
  令 \(\omega_p=e^{2\pi i/p}\)，则 \(|\omega_p-1|=2\sin(\pi/p)\)。
  由最近奇点距离证书 \(R_\star\) 可得
  \(
  |\omega_5-1|>R_\star
  \)
  而
  \(
  |\omega_7-1|<R_\star
  \)，并由单调性推出对所有 \(p\ge 7\) 有 \(|\omega_p-1|<R_\star\)（推论 \ref{cor:rate-center-omega-threshold}）。
  因此任一仅依赖 \(u=1\) 的局部解析展开在所有素数 \(p\ge 7\) 的单位根扭曲处必处于严格收敛域内，而在 \(p=5\) 上原则上不可由该局部口径控制，从而给出一个最小且唯一的素数排除点。

  \item \textbf{高阶碰撞核的 \texorpdfstring{$q$}{q}-根尺度谱凝聚：\texorpdfstring{$\Lambda_q^{1/q}$}{Lambda^(1/q)} 与 \texorpdfstring{$r_q^{1/q}$}{r^(1/q)} 对 \texorpdfstring{$\sqrt{\varphi}$}{sqrt(phi)} 的夹逼。}
  对碰撞 moment-kernel \(A_q\) 的 Perron 根 \(r_q\) 与次谱模 \(\Lambda_q\)（表 \ref{tab:fold_collision_moment_spectrum_k18_23} 与表 \ref{tab:fold_collision_kernel_rh_scan_q18_23}），定义
  \[
  b_q:=r_q^{1/q},\qquad a_q:=\Lambda_q^{1/q}.
  \]
  由极端阶端点 \(\lim_{q\to\infty}r_q^{1/q}=\sqrt{\varphi}\)（命题 \ref{prop:pom-rq-qinfty-endpoint}）可将 \(\sqrt{\varphi}\) 视为谱端点常数。
  在已核验区间 \(q=18,\dots,23\) 内，数值证书给出严格夹逼 \(a_q<\sqrt{\varphi}<b_q\)；例如
  \[
  (a_{18},b_{18})\approx(1.26368,1.28566),\qquad
  (a_{23},b_{23})\approx(1.26740,1.28011),
  \qquad
  \sqrt{\varphi}\approx 1.27202,
  \]
  显示二者在 \(q\)-根尺度上对 \(\sqrt{\varphi}\) 作单调收缩夹逼，从而把“越线谱相变”与“端点常数”统一到同一条谱凝聚轨道上。

  \item \textbf{谱凝聚指数的可审计单调律：\texorpdfstring{$\Lambda_q/r_q$}{Lambda/r} 以 \texorpdfstring{$\exp(-qg_q^{(\Lambda)})$}{exp(-qg)} 贴近 \(1\)。}
  在同一已核验区间内，定义谱凝聚指数
  \[
  g_q^{(\Lambda)}:=\frac{1}{q}\log\frac{r_q}{\Lambda_q}\,,
  \qquad
  \frac{\Lambda_q}{r_q}=\exp\!\bigl(-q\,g_q^{(\Lambda)}\bigr).
  \]
  由表 \ref{tab:fold_collision_kernel_rh_scan_q18_23} 的精确数据直接读出 \(g_q^{(\Lambda)}\) 在 \(q=18,\dots,23\) 上严格下降（从 \(g_{18}^{(\Lambda)}\approx 0.017241\) 降至 \(g_{23}^{(\Lambda)}\approx 0.009979\)），等价地 \(\Lambda_q/r_q\) 在该区间内由 \(0.733\) 单调上升至 \(0.795\)。
  因此“次模态贴 Perron”不仅可由 \(\log\Lambda_q/\log r_q\uparrow 1\) 的 Dirichlet 实部口径描述，还可由一条显式可核验的 \(g_q^{(\Lambda)}\) 序列量化；其单调衰减提供了判断谱凝聚是否进入近简并区的最低成本证书。

  \item \textbf{矩阶选择的离散变分阈值：最优 \texorpdfstring{$q^\star(c)$}{q*(c)} 的区间不等式刻画。}
  令 \(P(q)=\log r_q\) 并记割线斜率 \(\beta_q=(P(q)-P(1))/(q-1)\)。
  在指数风险预算 \(\varepsilon_m=e^{-cm}\) 下，证书指数可写为
  \[
  \gamma^{\mathrm{cert}}(c):=\inf_{q\ge 2}\left(\beta_q+\frac{c}{q-1}\right),
  \]
  并由离散凸性导出 \(q\mapsto\beta_q\) 单调不减（命题 \ref{prop:pom-pressure-convexity}；并见注 \ref{rem:pom-quantile-budget-computable}）。
  任一最优矩阶 \(q^\star(c)\) 满足离散 Euler--Lagrange 型阈值不等式（命题 \ref{prop:pom-qstar-discrete-variational-threshold}）：
  \[
  (q^\star-1)(q^\star-2)\bigl(\beta_{q^\star}-\beta_{q^\star-1}\bigr)\le c\le q^\star(q^\star-1)\bigl(\beta_{q^\star+1}-\beta_{q^\star}\bigr),
  \]
  从而 \(c\mapsto q^\star(c)\) 为分段常值且整体单调不减的“相图”。该刻画把“选 \(q\)”从启发式参数调度提升为完全由序列 \(\{r_q\}\) 决定的确定性阈值律。

  \item \textbf{回文多项式的 Chebyshev 降维：单位根扭曲被压缩为 \texorpdfstring{$[-2,2]$}{[-2,2]} 上的一变量实多项式。}
  由 \(Q(u)=u^{16}Q(1/u)\)（命题 \ref{prop:rate-center-Q-reciprocal-no-real}），可知 \(u^{-8}Q(u)\) 为反演 \(u\mapsto 1/u\) 的不变量；由 Laurent 环不变量子环恒等式 \(\ZZ[u,u^{-1}]^{u\mapsto 1/u}=\ZZ[u+u^{-1}]\)（命题 \ref{prop:completion-subring-unique-angle}）得存在唯一 \(\widetilde Q\in\ZZ[v]\)（\(\deg\widetilde Q=8\)）使
  \[
  Q(u)=u^8\,\widetilde Q\!\left(u+\frac1u\right).
  \]
  又因 \(Q(u)>0\) 对所有实 \(u\) 成立，得 \(\widetilde Q(v)>0\) 对所有 \(v\in(-\infty,-2]\cup[2,\infty)\) 成立。
  特别地，对任意单位根 \(\omega_p^k\) 有
  \[
  Q(\omega_p^k)=\omega_p^{8k}\,\widetilde Q\!\left(2\cos\frac{2\pi k}{p}\right),
  \]
  从而“单位根扭曲下的中心截面非退化性/奇点接近性”被严格压缩为 \(\widetilde Q\) 在区间 \([-2,2]\) 上的取值问题，给出一条最短可计算的 local--global 算术分层接口。

  \item \textbf{中心六重平坦性导出 \texorpdfstring{$p^{-6}$}{p^-6} 衰减律：单位根扭曲的不可见阶数与最小可辨尺度。}
  中心截面分解 \(R(\tfrac12,u)=-(u-1)^6Q(u)/64\) 与 \(Q(1)=8515584\neq 0\)（推论 \ref{cor:rate-center-unique-real} 与注 \ref{rem:rate-center-min-audit}）表明 \(u=1\) 处前五阶导数严格湮灭。
  因此对 \(u=\omega_p\) 且 \(p\to\infty\) 有主项
  \[
  R\!\left(\tfrac12,\omega_p\right)
  =-\frac{Q(1)}{64}\,(\omega_p-1)^6+O(p^{-7}),
  \qquad
  \bigl|R\!\left(\tfrac12,\omega_p\right)\bigr|
  =\frac{|Q(1)|}{64}\,(2\sin(\pi/p))^6+O(p^{-7})\asymp p^{-6}.
  \]
  等价地，单位根扭曲在五阶及以下的任何局部展开中均不可被分辨；“六阶”因此成为中心截面识别扭曲的最小可辨阶数与最小尺度律，并直接主导相关的数值稳定性与剥离策略（注 \ref{rem:rate-algebraic-elimination-structure}）。
\end{enumerate}

\paragraph{有限部—层析—端点喷射：从 \texorpdfstring{$O(\log\varepsilon^{-1})$}{O(log eps^-1)} 截断律到非平方自由幂零旗标}
为将“有限部常数的可计算性—最小数据依赖的 Lucas 层析—自对偶对称的偶喷射—模 \(p\) 影子的幂零通道”统一压缩为可直接复用的审计接口，下列断言给出一组互补的可核验命题（证书分布于有限部附录与 \(\xi\)/同步核/共振窗章节）：
\begin{enumerate}
  \item \textbf{Möbius--Abel 有限部常数的对数精度截断律：复杂度 \texorpdfstring{$O(\log\varepsilon^{-1})$}{O(log eps^-1)} 的可证书算法。}
  若 \(\zeta_M(z)=\det(I-zM)^{-1}\) 在闭圆盘 \(|z|\le \lambda^{-2}\) 内解析且 \(\zeta_M(0)=1\)，并令
  \[
  A:=\sup_{0<|z|\le \lambda^{-2}}\frac{|\log\zeta_M(z)|}{|z|}<\infty,
  \qquad
  \log\mathfrak{M}:=\log C+\sum_{m\ge 2}\frac{\mu(m)}{m}\log\zeta_M(\lambda^{-m}),
  \]
  则对任意 \(K\ge 2\) 的截断 \(\log\mathfrak{M}^{(K)}\) 都有显式尾项界（命题 \ref{prop:real-input-40-logM-trunc-bound}）：
  \[
  |\log\mathfrak{M}-\log\mathfrak{M}^{(K)}|
  \le \frac{A}{(K+1)(\lambda-1)}\,\lambda^{-K}.
  \]
  因此给定 \(\varepsilon\in(0,1)\)，取
  \[
  K\ \ge\ \left\lceil \frac{\log\!\bigl(\tfrac{A}{\varepsilon}(\lambda-1)\bigr)}{\log \lambda}\right\rceil
  \]
  即可保证 \(|\log\mathfrak{M}-\log\mathfrak{M}^{(K)}|\le \varepsilon\)；
  从而有限部常数在“对数精度”意义下可由 \(K\asymp \log(1/\varepsilon)\) 项完成证书化计算，而不依赖任何 primitive 轨道枚举或高维谱细节。

  \item \textbf{三点评估的 Lucas 层析：仅用迹样本重建 Abel 常数差并给出端到端误差封口。}
  在真实输入核的 prime-splitting 口径下，Lucas 迹样本 \(L_n=\Tr(F^n)\)（满足 \(L_n=L_{n-1}+L_{n-2}\)）经 Möbius/Witt 反演给出
  \(
  \pi_n=p_n(F)
  \)
  与
  \(
  r_n=p_n(F\otimes F)
  \)；
  再由模 \(4\) 选择律得到 \(p_n(M)\)（注 \ref{rem:real-input-40-lucas-pipeline}），并在奇素数长度上进一步坍缩为乘积律 \(p_p(M)=L_p\pi_p\)（推论 \ref{cor:real-input-40-prime-length-product}）。
  因此 Abel 常数差 \(\log\mathfrak{M}_M-\log\mathfrak{M}_{\mathrm{in}}\) 可仅由 \(P_F(z)=\sum_{n\ge 1}\pi_n z^n\) 在三点 \(z\in\{-\lambda^{-1},\pm\lambda^{-2}\}\) 的取值重建（推论 \ref{cor:real-input-40-logM-fast-prime}），并且当把 \(P_F\) 截断到 \(N\) 项时，总误差由显式尾项泛函 \(\mathrm{Tail}(N;|z|)\) 统一封口（注 \ref{rem:real-input-40-lucas-pipeline}），从而得到一条“最小数据依赖”的端到端可审计编译管线。

  \item \textbf{对称重尾的定量指纹：primitive 层标准化峰度与尾部增强阈值。}
  在自对偶完成化对称下，primitive 分层计数的中心化尺度具有显式累积量密度（定理 \ref{thm:sync-kernel-even-edgeworth-c4}）：
  \[
  \sigma^2=P''(0)=\frac{11}{102},\qquad
  \kappa_4=P^{(4)}(0)=\frac{1559}{58956},
  \qquad
  c_4=\frac{\kappa_4}{24\sigma^4}=\frac{1559}{16456}.
  \]
  因而标准化四阶比为 \(\kappa_4/\sigma^4=4677/2057\)，标准化峰度满足
  \[
  \frac{\mathbb{E}[X^4]}{\sigma^4}=3+\frac{\kappa_4}{\sigma^4}=\frac{10848}{2057}\approx 5.274>3,
  \]
  从而在奇阶中心矩严格为零的同时呈现严格的尖峰—厚尾偏离。
  并且由 Edgeworth 一阶修正多项式 \(H_4(x)=x^4-6x^2+3\)（定理 \ref{thm:sync-kernel-even-edgeworth-c4}）与其对尾概率的导数恒等式，可得首阶尾部增强在 \(|x|>\sqrt3\) 时不可避免出现（符号阈值由 \(H_3(x)=x^3-3x\) 锁定）。

  \item \textbf{中偏差“无奇项”相图：率函数喷射仅含偶次幂且二次曲率常数为 \texorpdfstring{$51/11$}{51/11}。}
  令 \(\alpha(\theta)=P'(\theta)\) 并以 \(I(\alpha)=\sup_\theta(\theta\alpha-P(\theta)+P(0))\) 定义率函数。
  自对偶偶性强制 \(I(\alpha)=I(1-\alpha)\)，从而在中心点 \(\alpha_\star=\tfrac12\) 的 Taylor 喷射仅含偶次幂（命题 \ref{prop:sync-kernel-moderate-deviation-cauchy-certified}）。
  特别地，
  \[
  I(\alpha)=\frac{(\alpha-\tfrac12)^2}{2P''(0)}+O\!\left((\alpha-\tfrac12)^4\right)
  =\frac{51}{11}\Bigl(\alpha-\tfrac12\Bigr)^2+O\!\left((\alpha-\tfrac12)^4\right),
  \]
  其中二次系数由 \(P''(0)=11/102\)（定理 \ref{thm:sync-kernel-even-edgeworth-c4}）唯一决定。
  并且在中偏差标度 \(|k-n/2|\le n^{2/3}\) 下，喷射截断的指数尺度误差可由真实解析半径 \(R_\theta\) 的 Cauchy 证书统一封口（命题 \ref{prop:sync-kernel-moderate-deviation-cauchy-certified} 与推论 \ref{cor:pressure-analytic-radius}），无需额外自由参数。

  \item \textbf{零温极限的双相谱隙常数：\texorpdfstring{$\lambda_{\mathrm{sub}}(q)/\lambda(q)\to\sqrt{2/3}$}{lambda_sub/lambda -> sqrt(2/3)} 给出一致慢模基准。}
  对真实输入核的 arity-charge 加权矩阵族 \(M(q)\)（\(q>0\)），其非平凡谱由次数 \(7\) 的代数曲线 \(P_7(\Lambda,q)=0\) 控制（定理 \ref{thm:real-input-40-arity-charge-det-closed}），从而 \(q\mapsto\lambda(q)=\rho(M(q))\) 为一条解析代数分支。
  当 \(q\to+\infty\)（零温端）时，Perron 分支满足
  \(
  \lambda(q)=\sqrt{3q}+O(1)
  \)
  （命题 \ref{prop:real-input-40-arity-charge-zero-temp-expansion}），并且同一代数曲线还存在一条正实增长分支
  \(
  \lambda_{\mathrm{sub}}(q)=\sqrt{2q}+O(1)
  \)
  （注 \ref{rem:real-input-40-arity-charge-large-bias-mechanism}）。
  因此
  \[
  \lim_{q\to\infty}\frac{\lambda_{\mathrm{sub}}(q)}{\lambda(q)}=\sqrt{\frac{2}{3}}<1,
  \]
  给出零温区间内仍与 \(1\) 保持正距离的统一慢模尺度常数，可作为谱隙/混合速率的一致性基准。

  \item \textbf{模 \texorpdfstring{$p$}{p} 非平方自由旗标的动力学含义：重复因子 \texorpdfstring{$\Leftrightarrow$}{<->} nilpotent 误差通道。}
  在有限域切片中，将递推特征多项式分解为
  \(P_q(\lambda)\equiv\prod_i f_i(\lambda)^{e_i}\)。
  若出现非零平方因子（存在 \(e_i\ge 2\) 且 \(f_i\neq\lambda\)），则对应的 companion 动力学含非半单 Jordan 块，从而解空间必出现 \((n^j\lambda^n)\) 型的“周期项叠加低阶多项式漂移”（注 \ref{rem:pom-charpoly-modp-audit-method}；注 \ref{rem:pom-charpoly-mod11-modulation}）。
  在已核验表中，模 \(7\) 的 \(q\in\{13,14,15\}\) 与模 \(11\) 的 \(q=12\) 均出现 \(\mathrm{repeat?}=1\)（表 \ref{tab:fold_collision_charpoly_mod7_selection_q13_15} 与表 \ref{tab:fold_collision_charpoly_mod11_fingerprint_q12_15}），并在 \((q,p)=(14,7)\)、\((12,11)\) 处同时出现 \(\chi_{A_2}\mid P_q\) 的嵌入旗标（推论 \ref{cor:pom-charpoly-modp-a2-subdynamics}），从而把“幂零误差通道”与“低阶影子子动力学”两类局部指纹在同一模 \(p\) 审计口径下并置。

  \item \textbf{Blaschke 端点喷射的可识别性：有限阶导数给出离临界线偏移与高度的加权矩。}
  在有限 Blaschke 设定下，端点径向吸收剖面 \(\mathcal{A}_B\) 在 \(r=1\) 的低阶喷射满足（定理 \ref{thm:xi-endpoint-absorption-jet-moments}）：
  \[
  \mathcal{A}_B(1)=\sum_k \delta_k,\qquad
  \mathcal{A}_B'(1)=\sum_k \delta_k(\delta_k-1),\qquad
  \mathcal{A}_B''(1)=\frac12\sum_k \delta_k\bigl(3(\delta_k-1)^2-\gamma_k^2\bigr).
  \]
  当 \(\delta_k\in(0,\tfrac12)\) 时有 \(\mathcal{A}_B'(1)<0\)，而二阶项中的 \(-\tfrac12\sum_k\delta_k\gamma_k^2\) 给出对大高度的二阶抑制响应。
  因此 \((\mathcal{A}_B(1),\mathcal{A}_B'(1),\mathcal{A}_B''(1))\) 构成一个有限维“统计截面”：它同时区分偏移总量、偏移离散度与高度平方的加权惩罚，从而把零点集合数据压缩为可复核的有限维测度不变量。

  \item \textbf{单位根扭曲的最小可审计阶：\texorpdfstring{$p_\star=6$}{p*=6} 的非素数阈值与素数 \texorpdfstring{$p=5$}{p=5} 的孤立性。}
  对单位根扭曲 \(u=\omega_p=e^{2\pi i/p}\)，以 \(\theta=\log u\) 为局部参数的 Taylor/累积展开在该点可作为收敛算法使用的充要必要条件为 \(2\pi/p<R_\theta\)（推论 \ref{cor:sync-kernel-unit-root-twist-threshold}）。
  在同步核输出位势加权中，最近分歧点证书给出 \(R_\theta\approx 1.0709533953\) 并推出最小阈值 \(p_\star=6\)（推论 \ref{cor:pressure-unit-root-modulus-threshold}）。
  该阈值不是素数门槛；因此在素数扭曲中，\(p=5\) 成为首个原则上无法由该局部口径保证收敛的孤立点，而 \(p\ge 7\) 自动落入可审计收敛域（并与中心截面根盘阈值链 \(\ref{cor:rate-center-omega-threshold}\) 一致）。
\end{enumerate}

\paragraph{有限部常数与扭曲指数的谱闭式：导数重建、Möbius--Abel 截断与 cyclotomic 消元的统一证书链}
在有限状态核的行列式口径下，常数层有限部与模 \(m\) 扭曲的最坏混合指数均可被压缩为整系数代数对象：主极点留数可由特征多项式导数完全重建，Möbius--Abel 级数给出可认证的截断误差界，而 completion 与 cyclotomic resultant 则把单位根扭曲的谱半径问题还原为一元整系数多项式的最小模根判定。其结论性接口可概括为：
\begin{enumerate}
  \item \textbf{主极点留数常数的谱闭式与代数审计：由 \texorpdfstring{$\chi_A'(r)$}{chi'(r)} 完全重建。}
  对任意原始非负邻接矩阵 \(A\)，设 \(\chi_A(\lambda)=\det(\lambda I-A)\)（首一，\(\deg\chi_A=d\)），且 Perron 根 \(r=\rho(A)\) 为单根；令 \(\zeta_A(z)=\det(I-zA)^{-1}\)、\(z_\star=1/r\)，并定义
  \(
  C(A):=\lim_{z\to z_\star}(1-rz)\zeta_A(z)
  \)。
  则
  \[
  \boxed{\;
  C(A)=\frac{r^{d-1}}{\chi_A'(r)}
  =\prod_{\lambda\in\sigma(A)\setminus\{r\}}\left(1-\frac{\lambda}{r}\right)^{-1}\;}
  \]
  （命题 \ref{prop:finite-part-residue-reduced-determinant}；亦见 \S\ref{subsec:pom-s4}）。
  因而 \(C(A)\) 在任何保持同一 \(\chi_A\) 的非负整数 realization 下不变；并且对 \(A_2,A_3,A_4\) 的可复现算例，\(C(A_q)\) 可进一步被升格为整系数最小多项式的唯一实根证书（\S\ref{subsec:pom-s4}），从而把常数输出封装为纯代数审计对象。

  \item \textbf{Möbius--Abel 快速公式与指数精度截断证书：有限部常数可计算且可认证。}
  在 \(\zeta_M\) 于 \(|z|\le z_\star^{\,2}\) 内解析且 \(\zeta_M(0)=1\) 的一般设定下，有限部常数的 Möbius--Abel 表示与显式尾项界见命题 \ref{prop:real-input-40-logM-trunc-bound}；其要点在于：仅需对收敛圆盘内部的有限采样即可同时给出高精度数值与可审计的误差预算，从而把 \(\mathsf{M}=\gamma+\log\mathfrak{M}\) 提升为稳定的标量证书输出。

  \item \textbf{常数层 Artin 分解：平凡因子可完全显式剥离，有限部机制仅依赖非平凡谱块。}
  若 \(\Delta(z)=\det(I-zM)\) 分解为 \(\Delta(z)=\Delta_{\mathrm{triv}}(z)\Delta_{\mathrm{nt}}(z)\)，其中 \(\Delta_{\mathrm{triv}}(z)=(1-z)^a(1+z)^b\) 且主极点 \(z_\star\in(0,1)\) 来自 \(\Delta_{\mathrm{nt}}\) 的简单零点，则 Abel 有限部常数满足严格加性分解
  \[
  \boxed{\ \log\mathfrak{M}=\log\mathfrak{M}_{\mathrm{nt}}+(a-b)z_\star+b z_\star^2\ }
  \]
  （命题 \ref{prop:real-input-40-reduced-mertens}）。
  在真实输入 \(40\) 状态核上，该剥离进一步细化为 even/twist/triv 的三因子常数分解（注 \ref{rem:real-input-40-logM-artin-factorization}），从而把常数层差异精确归因到表示通道与非平凡谱块的有限维机制。

  \item \textbf{加权 Ihara--Witt Euler 乘积与 Dwork 型同余：prime-shadow 生成整系数 Witt 向量。}
  对带权 Hashimoto 算子 \(B(u)\) 的二变量 Ihara \(\zeta\)-函数 \(Z_K(u,t)=\det(I-tB(u))^{-1}\)，存在整系数 primitive Witt 坐标 \(q_n(u)\in\ZZ[u]\) 使
  \(
  Z_K(u,t)=\prod_{n\ge 1}(1-q_n(u)t^n)^{-1}
  \)
  严格成立；并且 ghost 分量 \(b_n(u)=\Tr(B(u)^n)\) 同时满足 necklace 整除律与 Dwork 型同余
  \(b_{p^k}(u)\equiv b_{p^{k-1}}(u^p)\ (\mathrm{mod}\ p^k)\)
  （注 \ref{cor:big-witt-audit-guardrails}；命题 \ref{prop:witt-necklace-dwork-R}）。
  因此 prime-like（Euler 乘积）侧与迹/谱侧在整系数层面被统一为可机核验的交换图，并为后续 \(p\)-进分析对象提供结构性入口。
  另一方面，循环块迹类提升给出 Euler 产品到行列式 \(\zeta\) 的构造性满射（推论 \ref{cor:cyclic-euler-product}），因此仅凭“可写作 determinant/Euler product”的形式并不推出任何 RH-类谱约束；非平凡结构只能来自 Perron--Frobenius 正性、局部转移约束与可审计的通道选择律等附加输入。

  \item \textbf{自对偶消元的结构化降维：六重退化点与 \texorpdfstring{$1/6$}{1/6} 临界反演指数。}
  同步核速率曲线的规范化消元多项式满足自对偶对称 \(R(\alpha,u)=u^{22}R(1-\alpha,1/u)\)，并在 \(\alpha=\tfrac12\) 截面出现六重聚合
  \(R(\tfrac12,u)=-(u-1)^6Q(u)/64\)（推论 \ref{cor:rate-center-unique-real}）。
  因而反解 \(u=u(\alpha)\) 在中心点附近具有标度
  \(
  |u-1|\asymp|\alpha-\tfrac12|^{1/6}
  \)
  （注 \ref{rem:unit-root-final-spec}；推论 \ref{cor:rate-duality-center-6-3x2}），并可在对称坐标 \(X=(2\alpha-1)^2\)、\(s=u+1/u\) 下实现“关于 \(X\) 的三次 + 二次根式延拓”的降维（推论 \ref{cor:rate-duality-center-6-3x2}）。

  \item \textbf{cyclotomic elimination：模 \texorpdfstring{$m$}{m} residue-mixing 指数的 resultant 证书与精确次数律。}
  对模 \(m\) 的 residue mixing，primitive 计数可由单位根离散 Fourier 采样表示为
  \(
  p_{n,a}^{(m)}=\frac1m\sum_{r=0}^{m-1}\omega_m^{-ar}\,p_n(\omega_m^r),\qquad \omega_m:=e^{2\pi i/m}
  \)；
  其误差指数由最坏扭曲 Perron 根 \(\rho_m=\max_{1\le r\le m-1}|\lambda(\omega_m^r)|\) 控制，并可由 cyclotomic resultant 多项式 \(R_m(w)\in\ZZ[w]\) 精确刻画为
  \(
  \rho_m=1/\min\{|w|:R_m(w)=0\}
  \)
  （命题 \ref{prop:sync-cyclotomic-elimination}；表 \ref{tab:sync_kernel_cyclotomic_elimination_summary}）。
  当 \(m\ge 4\) 时，\(R_m\) 的次数满足精确公式 \(\deg_w R_m=3\varphi(2m)\)（定理 \ref{thm:sync-cyclotomic-elimination-degree-law}）。

  \item \textbf{单位圆扭曲的真实解析延拓半径：最近复分岔点给出非微扰收敛域护栏。}
  对单位圆扭曲 \(\theta\mapsto\log|\lambda(e^{i\theta})|\) 的局部展开，真实解析半径由最近复分岔点钉死（注 \ref{rem:sync-hatdelta-analytic-radius-complex}），其角平面坐标满足
  \(
  \theta_\star\approx 1.0556869306\pm 0.1801840156\,i
  \)
  并诱导半角坐标 \(x=\tan(\theta/2)\) 的最近奇点模长 \(|x_\star|\approx 0.4916\)（表 \ref{tab:sync_kernel_hatdelta_nearest_complex_branch_point_tan_half_angle}）。
  该半径为基于 Taylor/Pad\'e 的相图外推与误差指数预测提供可审计的“解析护栏”。

  \item \textbf{window-\texorpdfstring{$6$}{6} 的中心结构与 Levi 刚性：\texorpdfstring{$(\ZZ_2)^3$}{(Z2)^3} 的规范嵌入、内生 \texorpdfstring{$\ZZ_6$}{Z6} 商与 SM Levi 的最小闭包强制。}
  window-\(6\) 的 boundary dyadic 读出给出规范中心嵌入 \((\ZZ_2)^3\hookrightarrow Z(\mathcal{G}^{\mathrm{bin}}_6)\)，并对任何连续包络候选施加可证伪的中心秩约束（注 \ref{rem:bdry-center-z2-global-compatibility}）。
  与 cyclic 扇区的 \(\ZZ_3\) 源并置可构造内生 \(\ZZ_6\) 商，从而把标准模型常用的全局形式 \((SU(3)\times SU(2)\times U(1))/\ZZ_6\) 解释为离散输出强制的最小一致性后果；并且当 orbit--Levi 骨架与阈值轻/重分裂要求在同一连续包络内兼容闭合时，将强制出现 SM 型 Levi 子结构（命题 \ref{prop:window6-common-refinement-sm-levi}；亦见猜想 \ref{conj:window6-sm-gauge-as-minimal-compact-observable}）。
\end{enumerate}

\paragraph{表示通道的剥离原则与完成化刚性：toy-RH、常数层与二变量谱的可审计归约}
在有限核口径下，若不先剥离由表示通道强制嵌入的显式谱支与短周期平凡修饰，则若干表观“达界/越界”现象中的平方根级误差项与常数级漂移将不可避免地干扰 RH/Ramanujan 型判别；而一旦以行列式因子化与完成化同余对这些通道作结构性剥离，剩余部分在严格意义上落入可审计的谱界与有限生成证书闭包。
\begin{enumerate}
  \item \textbf{toy-RH 的唯一达界来自符号表示通道，而非内部共振。}
  对真实输入 \(40\) 状态核，行列式在张量骨架与符号通道之间存在精确三因子分解（命题 \ref{prop:real-input-40-A0-tensor-hidden}）：
  \[
  \det(I-zM)=(1-z)^2(1+z)\,\det\!\bigl(I-z(A_0\otimes A_0)\bigr)\,\det\!\bigl(I-z(-A_0)\bigr).
  \]
  其中 \(A_0\) 为黄金均值骨架，且唯一触及平方根界 \(|\mu|=\sqrt{\lambda}=\varphi\) 的非 Perron 模态由 \(\det(I-z(-A_0))=1+z-z^2\) 强制产生（注 \ref{rem:finite-rh-40-hidden-minusA0}），而 \(A_0\otimes A_0\) 仅贡献模长 \(1\) 与 \(\varphi^{-2}\) 的次谱，从而严格满足 \(|\nu|<\sqrt{\lambda}\)。
  因此“临界线达界”并非高维核结构的偶然共振，而是符号表示通道的结构性嵌入；在通道投影意义下剥离 \(\zeta_{-A_0}\) 后，剩余因子满足严格 Ramanujan 比率 \(<1\)，并在可观测量上消去平方根级振荡项。

  \item \textbf{短周期平凡因子只在 \texorpdfstring{$n=1,2$}{n=1,2} 留下 primitive 痕迹。}
  平凡周期因子 \(\zeta_{\mathrm{triv}}=(1-z)^{-2}(1+z)^{-3}\) 的 Möbius--\(\log\) 变换给出严格闭式
  \(P_{\mathrm{triv}}(z)=-z+3z^2\)，从而 \(p^{\mathrm{triv}}_n=0\ (n\ge 3)\)（命题 \ref{prop:triv-factor-primitive-polynomial}）。
  因此对所有长度 \(n\ge 3\)，primitive 轨道谱在严格意义上与平凡呈示无关；相应地，剥离 \(\zeta_{\mathrm{triv}}\) 后的非平凡部分在 primitive 层面自动落入 \(\zeta_{\mathrm{even}}\) 与符号因子 \(\zeta_{\chi}\) 的精确二因子结构（注 \ref{rem:real-input-40-logM-artin-factorization}）。

  \item \textbf{符号通道的 Abel/Mertens 常数是与内部动力学无关的普适偏移量。}
  对 twist 因子 \(\zeta_{-F}(z)=(1+z-z^2)^{-1}\) 在主极点 \(z_\star=\varphi^{-2}\) 处定义的有限部常数
  \(K_{-F}=\sum_{m\ge 1}\frac{\mu(m)}{m}\log\zeta_{-F}(z_\star^{\,m})\)
  仅依赖该符号通道本身，并在常数层满足
  \(\log\mathfrak{M}=\log\mathfrak{M}_{\mathrm{even}}+K_{-F}+P_{\mathrm{triv}}(z_\star)\)（注 \ref{rem:real-input-40-logM-artin-factorization}；定义 \ref{def:real-input-40-artin-sign-mertens}）。
  其数值见表 \ref{tab:real_input_40_artin_sign_mertens}（\(K_{-F}\approx -0.1348652875\ldots\)），从而为“符号表示通道是否存在”提供一条无需枚举长轨的常数级诊断口径。

  \item \textbf{完成化 ghost 的 Frobenius 刚性与 \texorpdfstring{$\ZZ[x]$}{Z[x]} 全有理闭包。}
  在完成化坐标下，素数长度同余给出 Teichm\"uller 型刚性：\(\widehat a_p(s)\equiv s^p\ (\mathrm{mod}\ p)\)（推论 \ref{cor:completed-teichmueller-congruence}）。
  同时，unitary-slice 的倍角/幂半群作用在迹坐标上由 Chebyshev--Adams 多项式 \(C_d\) 统一本地化，并可在 Cayley 有理参数 \(x\) 上以 \(\ZZ[x]\) 的纯乘加递推实现（命题 \ref{prop:chebyshev-adams-fully-rational}）；并与 primitive 完成化的倍角律一起，在 \(\ZZ[s]\) 上给出幂替换的 Adams（\(\lambda\)-环）实现（推论 \ref{cor:primitive-completion-hatp}）。

  \item \textbf{arity-charge 二变量 \texorpdfstring{$\zeta$}{zeta} 的显式谱支、零荷熵锚点与有限生成证书。}
  对 arity-charge 加权核 \(M(q)\)，行列式具有闭式分解 \(\det(I-zM(q))=(1+z)(1-qz^2)\,Q_7(z,q)\)（定理 \ref{thm:real-input-40-arity-charge-det-closed}），从而谱中对任意 \(q\) 都包含显式特征值 \(\pm\sqrt q\) 与固定特征值 \(-1\)，并诱导完全显式的偶奇振荡项（推论 \ref{cor:real-input-40-arity-charge-sqrtq}）。
  因此关于平方根界/类 RH 的谱判别在该族中应先剥离因子 \((1+z)(1-qz^2)\)（等价 \((\Lambda+1)(\Lambda^2-q)\)）；剥离后由次数 \(7\) 的部分承载非平凡谱机制。
  零电荷子系统的 \(\zeta_0\) 亦为闭式有理函数，其 Perron 根 \(\kappa\) 给出约束内增长率的代数锚点（命题 \ref{prop:real-input-40-arity-charge-zero-zeta}）。
  进一步，\(\det(I-zM(q))\) 在 \(z\) 方向有效次数为 \(10\)、在 \(q\) 方向次数为 \(3\)，并可由有限段 primitive 数据经 Newton 递推闭合恢复（注 \ref{rem:real-input-40-arity-charge-det-q-cubic} 与注 \ref{rem:real-input-40-arity-charge-newton}），从而把二变量谱对象压缩为有限审计证书。
\end{enumerate}

\paragraph{纤维拓扑二分、序格完备表示与模轴束/视界正性的统一证书链：从 \texorpdfstring{$\tau(x)$}{tau(x)} 到 Chebotarev--Plancherel 与 Toeplitz--PSD}
在本文固定的折叠—投影—扭曲接口下，纤维侧的拓扑/序/统计结构与同余轴束侧的扭曲谱隙、profinite 极限及视界正性可被压缩为一条有限维、整系数、可审计的证书链：纤维的同伦型与格结构由分量长度多重集 \(\{\ell_i\}\) 完全决定；最坏纤维与碰撞缺口由非平凡 Fourier 能量定量驱动；而模 \(q\) 等分布、加性能量主项与 profinite Haar 二分律则由扭曲 Perron 根的统一谱隙统摄；在视界接口上，Carath\'eodory/Herglotz 正性进一步等价于一族 Toeplitz--PSD 层级，从而把无限对象的正性判别落地为可逐级核验的有限维可行性约束。
\begin{enumerate}
  \item \textbf{纤维独立集复形的“球面—可缩”二分刚性由单一整数 \texorpdfstring{$\tau(x)$}{tau(x)} 完全参数化。}
  对任意稳定类型 \(x\in X_m\)，候选位点诱导图 \(\Gamma_x\) 的路径分量长度 \((\ell_i)\) 给出闭式参数
  \(\tau(x)=\sum_i\lfloor(\ell_i+1)/3\rfloor\)，并强制独立集复形 \(|K_x|\) 的同伦型满足严格二分：要么可缩、要么同伦等价于单一球面 \(S^{\tau(x)-1}\)；且非可缩情形下约化同调仅在唯一维度上为 \(\ZZ\)（定理 \ref{thm:pom-fiber-independence-complex-classification}；推论 \ref{cor:pom-fiber-independence-complex-homology}）。

  \item \textbf{高多重纤维指数稀疏：达到固定比例最坏值的类型集合在指数尺度上必稀薄。}
  令 \(D_m=\max_{x\in X_m}d_m(x)\)，并对 \(\alpha\in(0,1]\) 定义
  \(S_\alpha(m)=\{x:\ d_m(x)\ge \alpha D_m\}\)。
  则
  \[
  |S_\alpha(m)|\le \frac{2^m}{\alpha D_m},
  \qquad
  \frac{|S_\alpha(m)|}{|X_m|}=O\!\left(\frac{1}{\alpha}\left(\frac{2}{\varphi^{3/2}}\right)^m\right)
  \]
  （推论 \ref{cor:pom-high-multiplicity-sparse}）。
  同时，最坏/平均放大比的指数律与奇偶乘性常数可被完全显式化（注 \ref{rem:pom-defect-gap}），从而“极端空间损失”与“其承载集合的稀疏性”可被并列审计。

  \item \textbf{纤维内部的格结构具有 Birkhoff 完备表示：锯齿理想格模型给出 join/meet 与区间结构的全分类。}
  在自然的“单次逆代换开关可达性”偏序下，每条纤维 \(F_x\) 具有规范分布式格结构 \(\mathcal L_x\)，并存在同构
  \[
  \mathcal L_x\ \cong\ J\!\Big(\bigsqcup_{i=1}^r Z_{\ell_i}\Big),
  \]
  其中 \(Z_{\ell}\) 为锯齿（fence）偏序、\(J(P)\) 为理想格；并且 join-不可约元偏序同构于 \(\bigsqcup_i Z_{\ell_i}\)，从而 \(\{\ell_i\}\) 完整决定纤维的合取/析取与区间结构（定理 \ref{thm:pom-fiber-birkhoff-fence-ideal-lattice}）。

  \item \textbf{最大纤维与碰撞缺口由非平凡 Fourier 能量定量驱动，并给出常数级偏置下界。}
  令 \(d_m(r)\) 为模 \(F_{m+2}\) 的多重度剖面，\(\widehat d_m(k)\) 为其 Fourier 幅度，\(\overline d_m=2^m/F_{m+2}\) 为均值，且 \(M_m=\max_r d_m(r)\)。
  则
  \[
  M_m\ge \overline d_m+\frac{1}{F_{m+2}}\Bigl(\sum_{k\ne 0}|\widehat d_m(k)|^2\Bigr)^{1/2},
  \qquad
  M_m\ge \overline d_m+\frac{1}{F_{m+2}}\max_{k\ne 0}|\widehat d_m(k)|
  \]
  （定理 \ref{thm:fold-max-fiber-spectrum} 与定理 \ref{thm:fold-max-fiber-fourier}）。
  并且折叠碰撞概率满足
  \(
  \mathsf{Col}_m\ge M_m^2/2^{2m}
  \)
  （定理 \ref{thm:fold-collision-max-fiber}）。
  进一步，由临界共振常数 \(C_\varphi=\prod_{n\ge 2}|\cos(\pi\varphi^{-n})|\) 的极限化（定理 \ref{thm:fold-critical-resonance-constant}），得到
  \(\liminf_{m\to\infty}M_m/\overline d_m\ge 1+C_\varphi\) 的常数级偏置下界（推论 \ref{cor:fold-max-fiber-constant-bias}）。

  \item \textbf{几何侧浓缩指数与统计侧最优可分性指数严格同一化为一枚代数常数。}
  最大/平均纤维浓缩指数满足
  \[
  \boxed{\ \lim_{m\to\infty}\frac{1}{m}\log_2\frac{D_m}{\overline d_m}=C_\varphi^{(\mathrm{Ch})}:=\frac32\log_2\varphi-1\ }
  \]
  （推论 \ref{cor:pom-D-over-avg-exponent-equals-chernoff}），从而把“最坏/平均浓缩”的几何指数与二相位识别的 Chernoff 常数在指数尺度上严格对齐。

  \item \textbf{扭曲 primitive 抽取具有 Möbius--Adams 同步性：时间除子与 \texorpdfstring{$(\chi,u)$}{(chi,u)} 伸缩强耦合。}
  在角色轴 \(\chi\) 与能量参数 \(u\) 的扭曲 \(\zeta\) 框架中，primitive 抽取满足带伸缩的 Möbius 公式
  \[
  B_n(\chi;u)=\frac{1}{n}\sum_{d\mid n}\mu(d)\,a_{n/d}(\chi^d;u^d),
  \]
  其中 \((\chi,u)\mapsto(\chi^d,u^d)\) 同步出现（命题 \ref{prop:kernel-chebotarev-mobius-adams}；并见注 \ref{rem:kernel-chebotarev-mobius-adams-meaning} 与命题 \ref{prop:witt-renorm-monoid} 的半群化表述）。

  \item \textbf{Chebotarev--Plancherel 能量恒等式与商单调性：空间侧 \texorpdfstring{$L^2$}{L2} 偏差精确等于角色侧能量和。}
  在有限阿贝尔群标签分层中，均方偏差满足精确恒等式
  \(
  \sum_{c\in G}|\varepsilon_n(c;u)|^2=\frac{1}{|G|}\sum_{\chi\ne 1}|B_n(\chi;u)|^2
  \)
  （推论 \ref{cor:kernel-chebotarev-plancherel-energy}）。
  在一般有限群口径下，上式提升为类函数空间的 Parseval/Plancherel 恒等式，并给出对商映射的偏差单调性（命题 \ref{prop:kernel-chebotarev-plancherel-nonabelian}）。

  \item \textbf{扭曲谱隙 \texorpdfstring{$\Rightarrow$}{=>} 模 \texorpdfstring{$q$}{q} 指数等分布与加性能量主项：误差由最坏扭曲 Perron 根统一支配。}
  在模 \(q\) 的扭曲 Perron 谱隙条件下，剩余类计数具有指数误差项（定理 \ref{thm:gm-modq-equidistribution-spectral-gap}），并且模加性能量满足
  \[
  E_m(q)=\frac{|A_m|^4}{q}+O\!\bigl(q\,(\rho_q^\ast)^{4m}\bigr)
  \]
  （命题 \ref{prop:gm-modq-additive-energy-main-error}），从而把均匀分布与能量估计统一为一条可计算的谱链条。

  \item \textbf{Profinite Haar 二分律：要么 Haar 极限，要么支撑退化到真闭子群陪集。}
  由有限商 \(\ZZ/q\ZZ\) 上计数的相容性可在 \(\widehat{\ZZ}\) 上构造测度列 \(\mu_m\)；若对每个固定 \(q\ge 2\) 均满足非退化谱隙，则 \(\mu_m\Rightarrow \mu_{\mathrm{Haar}}\)；
  反之若某个 \(q\) 上存在等半径共振，则 \(A_m\bmod q\) 终将落入某个真子群陪集，并诱导 \(\widehat{\ZZ}\) 上的闭子群陪集支撑退化（定理 \ref{thm:gm-profinite-haar-dichotomy}）。

  \item \textbf{主弧集中 \texorpdfstring{$\Longleftrightarrow$}{<->} 支撑退化：有理频谱共振与 profinite 非 Haar 的双向判别。}
  主弧端点上出现常数级指数和（主弧集中）当且仅当存在 cyclotomic 障碍/共振角色（定理 \ref{thm:gm-majorarc-energy-obstruction-loop}），而后者等价于 \(\widehat{\ZZ}\) 上非 Haar 的陪集支撑退化（定理 \ref{thm:gm-profinite-haar-dichotomy}）。
  因此“有理频谱大值”与“算术支撑退化”在同一有限维谱证书口径下互为可执行判别。

  \item \textbf{profinite 角色对角化与 Hilbert--Artin 分解：最坏扭曲谱半径可由算子限制精确读出。}
  在 \(\mathcal{H}=\CC^V\otimes L^2(\widehat{\ZZ})\) 上的群扩张邻接算子 \(\mathcal{M}_u\) 可按角色直和分解为有限维扭曲块 \(M_\chi(u)\)，并满足
  \[
  \rho(\mathcal{M}_u|_{\mathcal{H}_\perp})=\sup_{\chi\ne 1}\rho(M_\chi(u)),
  \qquad
  \rho(\mathcal{M}_u|_{\mathcal{H}_0})=\rho(M_1(u))=:\lambda_1(u),
  \]
  从而把“最坏扭曲谱半径”提升为算子化读出（命题 \ref{prop:pom-profinite-artin-hilbert}）。

  \item \textbf{视界 Toeplitz--PSD 的有限证书化：Carath\'eodory 正性 \texorpdfstring{$\Leftrightarrow$}{<->} PSD 层级，且共终端子序列核验与全层级等价。}
  视界 Carath\'eodory 对数导数的正性等价于 Toeplitz 主子矩阵族 \(\{\mathbf{T}_N\succeq 0\}_{N\ge 0}\)，并在 \(\zeta\) 情形下与 RH 等价（定理 \ref{thm:app-horizon-toeplitz-lmi}）。
  进一步，Toeplitz--PSD 的全层级核验与任意共终端子序列（如 \(N_k=2^k\)）核验严格等价，从而给出一条固定阶可嵌入的稀疏化/逼近机制（定理 \ref{thm:xi-toeplitz-psd-cofinal-sparsification}）。
\end{enumerate}


\paragraph{模障碍的导子化闭包、共振角色群的同调分类与 Ihara--Chebotarev 误差字典：从 cyclotomic 因子到 \texorpdfstring{$p$}{p}-典型整性与零温线性响应}
在有限状态加权图/同步核的扭曲谱框架中，“模障碍/共振”并非散点经验现象，而可被压缩为一组有限维、整系数、可计算的同调—谱证书：同调导子 \(d\) 决定全部例外模数；共振角色构成有限阿贝尔群并可由 Smith 标准形一次性输出；非共振通道上存在显式相位失配下界给出统一的扭曲谱降；一旦共振发生，行列式必含 cyclotomic 因子并诱导有限阶旋转模态，从而在时间域出现不可消去的严格周期振荡校正。
与之并行，Ihara--Chebotarev 的平方根级误差与谱隙条件互为充要，而 Ihara--Mertens 常数与局部 \(p\)-因子均可经导数/同余闭包被压缩为可审计的有限维对象；在零温端，Perron 根的一阶微扰系数进一步给出一条与谱隙正交的“易激发性”标量序。
\begin{enumerate}
  \item \textbf{有限导子原理：全部模障碍由单一整数 \texorpdfstring{$d$}{d} 完全控制。}
  对原始有向图 \(G\) 的整数边权 \(w\)，回路权和集合 \(\mathcal{C}\) 的最大公因子
  \(d=\gcd(\mathcal{C})\)
  给出模共边界的充要判别：存在 \(g:V\to\ZZ/q\ZZ\) 使
  \(w(e)\equiv g(\mathrm{head}(e))-g(\mathrm{tail}(e))\ (\mathrm{mod}\ q)\)
  当且仅当 \(q\mid d\)（命题 \ref{prop:gm-cycle-gcd-mod-obstruction}）。
  因此所谓“无限多模数共振”必等价于一个固定的同调障碍 \(d>1\)，而非解析偶然。

  \item \textbf{统一扭曲谱降的定量下界：非障碍模上谱隙由最小相位失配显式给出。}
  若 \(q\nmid d\)，则扭曲邻接算子满足
  \[
  \boxed{\ \rho\!\bigl(A(\chi_q)\bigr)\ \le\ \bigl(1-\kappa_G\eta(q)\bigr)\rho\ },
  \qquad
  \eta(q)=\min_{1\le t<q}\sin^2\!\Bigl(\frac{\pi t}{q}\Bigr)\asymp q^{-2}
  \]
  其中 \(\kappa_G>0\) 仅依赖底图的原始性常数（定理 \ref{thm:gm-uniform-twist-gap-from-gcd}）。
  该不等式把“主弧衰减/谱压制”的来源压缩为有限维谱半径比较与显式三角函数下界的组合。

  \item \textbf{共振角色群的完整同调分类与可计算性：共振不是散点而是有限阿贝尔群。}
  共振角色集合
  \(\mathcal{X}_{\mathrm{res}}=\{\chi:\rho(A(\chi))=\rho(A(1))\}\)
  构成有限阿贝尔群，并同构于同调商挠性部分的 Pontryagin 对偶
  \[
  \boxed{\ \mathcal{X}_{\mathrm{res}}\ \cong\ \mathrm{Hom}\!\Bigl(\mathrm{tors}\bigl(H_1(G,\ZZ)/\langle[w]\rangle\bigr),\ \QQ/\ZZ\Bigr)\ }
  \]
  （定理 \ref{thm:gm-resonant-character-group-homology}）。
  其不变因子分解与生成角色表可由整数矩阵的 Smith 标准形在多项式时间内输出（命题 \ref{prop:gm-smith-compute-obstructions}），从而把“共振结构”提升为一次性可审计的离散输出。

  \item \textbf{cyclotomic 因子 \texorpdfstring{$\Longleftrightarrow$}{<->} 等半径共振：共振必然对应有限阶旋转模态并诱导严格周期振荡项。}
  在图 \(\zeta\)/Ihara--Artin \(\zeta\) 的行列式框架中，存在非平凡共振通道当且仅当扭曲 \(\zeta\)（或等价的 \(\det(I-uA(\chi))\)、\(\det(I-tB_\rho(u))\)）含 cyclotomic 因子；并且该情形等价于某个非平凡扭曲满足谱半径等半径
  \(\rho(A(\chi))=\rho(A(1))\)
  或
  \(\rho(B_\rho(u))=\lambda_1(u)\)
  （定理 \ref{thm:gm-graph-zeta-cyclotomic-classification}；定理 \ref{thm:conclusion75-ihara-cyclotomic-factor-classification}）。
  cyclotomic 因子对应单位圆上的单位根，从而把谱边界上的特征值强制为有限阶旋转；其时间域后果是：主项之外出现与 \(\lambda_1(u)^n/n\) 同阶的周期振荡校正项，其周期为相应单位根的阶。

  \item \textbf{统一扭曲谱隙的平凡共振判别与可计算闭包。}
  定义“统一扭曲谱隙”性质为：存在常数 \(\delta>0\) 使对一切 \(\chi\neq \mathbf 1\) 有
  \[
  \rho(A(\chi))\le (1-\delta)\rho(A(1)).
  \]
  则该性质当且仅当 \(\mathcal X_{\mathrm{res}}=\{\mathbf 1\}\)。
  事实上，若统一谱隙成立则不存在等半径共振角色；
  反之若 \(\mathcal X_{\mathrm{res}}=\{\mathbf 1\}\)，则由上一条的 cyclotomic 分类与障碍闭子群分解，非平凡扭曲被分为“障碍上非平凡”的无穷族与“障碍上平凡”的有限族：
  前者由定理 \ref{thm:gm-uniform-twist-gap-from-gcd} 给出统一谱降常数；
  后者由于有限且无等半径点，取最大谱半径比即可得到 \(\delta_1>0\)，合并得统一 \(\delta=\min(\delta_0,\delta_1)\)。
  由结论 \ref{con:xi-resonant-character-snf-computable} 的 SNF 可计算性，\(\mathcal X_{\mathrm{res}}\) 的平凡性可在多项式时间内审计，从而“统一扭曲谱隙是否成立”在该离散口径下属于确定性多项式时间可判定问题。

  \item \textbf{共振见证的素阶压缩：最小反证据必为素数阶。}
  若 \(\mathcal X_{\mathrm{res}}\neq\{\mathbf 1\}\)，则存在某个 \(\chi\in\mathcal X_{\mathrm{res}}\setminus\{\mathbf 1\}\) 其阶为素数。
  证明只需在有限阿贝尔群 \(\mathcal X_{\mathrm{res}}\) 中取阶最小的非平凡元：若其阶为合成数 \(m=ab\) 且 \(1<a<m\)，则 \(\chi^a\neq \mathbf 1\) 且 \(\mathrm{ord}(\chi^a)=b<m\)，与极小性矛盾，故该最小阶必为素数。
  因此一旦统一谱隙失败，总可在某个素阶通道上给出等半径共振的有限核验证据，从而把“失败见证”的搜索空间压缩到素数模数族。

  \item \textbf{Ihara--Chebotarev 的谱隙充要性：平方根级误差与 Ramanujan 型谱界互为等价判别。}
  对按 Frobenius 共轭类分层的 Ihara-primitive 计数 \(q_{n,C}(u)\)，平方根级一致误差
  \[
  q_{n,C}(u)=\frac{|C|}{|G|}\cdot\frac{\lambda_1(u)^n}{n}
  +O\!\left(\frac{\lambda_1(u)^{n/2}}{n}\right)
  \qquad(\text{对 }C\text{ 一致})
  \]
  与谱隙条件 \(\Lambda(u)\le \lambda_1(u)^{1/2}\) 互为充要（推论 \ref{cor:ihara-chebotarev-grh-criterion}）。
  因此误差尺度本身可作为谱隙质量的判别器，而无需先验知道全部扭曲谱。

  \item \textbf{Ihara--Mertens 常数的 Möbius--Abel 公式：主极点留数可由行列式导数给出有限维证书。}
  在 Perron--Frobenius 谱隙条件下，Ihara 口径的 Abel 有限部常数 \(\log\mathfrak{M}^{\mathrm{Ih}}_K(u)\) 具有与未回溯常数层同型的 Möbius--Abel 快速配方，并且主极点留数常数由
  \(
  \partial_t\det(I-tB(u))|_{t=t_\star(u)}
  \)
  的有限维导数证书完全决定（命题 \ref{prop:ihara-mertens-constant}）。

  \item \textbf{\texorpdfstring{$p$}{p}-典型 \texorpdfstring{$p$}{p}-进整性闭包：从 Dwork 同余到局部 \texorpdfstring{$p$}{p}-因子并自动稳定于多势参数。}
  Dwork 型同余与 Dieudonn\'e--Dwork 整性判据给出 \(p\)-典型局部因子
  \(
  Z_p(T;u)=\exp(\sum_{k\ge 0}a_{p^k}(u)T^{p^k}/p^k)
  \)
  在 \(1+T\ZZ_p[u][[T]]\) 内闭合（定理 \ref{thm:witt-dwork-zp-integrality}；引理 \ref{lem:witt-dieudonne-dwork}）。
  同时，多维势下的 necklace 整除与多维 Dwork 同余逐字成立（\S\ref{app:multi-dim-witt-guardrails}），因此“局部 \(p\)-因子”结构在系数环升级时不引入额外审计代价。

  \item \textbf{零温线性响应的代数刚性：一阶 carry 密度完全由 carry-free 骨架决定并导出可比标量序。}
  若 carry-free 骨架 \(A_0\) primitive 且 Perron 根 \(\lambda_0\) 简单，则谱半径在 \(u=0\) 处解析并满足
  \[
  \lambda_K(u)=\lambda_0+c_1u+O(u^2),
  \qquad
  c_1=\frac{\ell^{\mathsf T}A_1 r}{\ell^{\mathsf T}r},
  \qquad
  \mathbb{E}_u[\kappa]=\frac{c_1}{\lambda_0}u+O(u^2)
  \]
  （命题 \ref{prop:lambda-linear-response} 与推论 \ref{cor:kappa-linear-response-order}）。
  在 \(K_9/K_{13}/K_{21}\) 三核算例中，该一阶斜率给出零温附近的“抗 carry”排序，并且 \(K_{21}\) 的一阶响应常数在代数上塌缩为 \(c_1=\sqrt5\)（推论 \ref{cor:k21-c1-sqrt5}）。
\end{enumerate}

\begin{theorem}[carry 势零温极限的支撑刚性]\label{thm:conclusion-carry-zero-temp-support-collapse}
设 \(M_K(u)\) 为有限图上的 carry 加权转移矩阵，并取局部常势
\[
\varphi_\beta:=-\beta\kappa,
\qquad
\beta\to+\infty.
\]
令 \(\mu_\beta\) 为相应 Gibbs 平衡态。则任意零温极限点 \(\mu_0\) 都满足
\[
\boxed{\;\mu_0(\kappa=0)=1.\;}
\]
亦即，零温极限的全部质量必定落在 carry-free 子图生成的有限并 SFT 上。若 carry-free 子图仅有唯一的本质强连通分量，则 \(\mu_0\) 唯一，且正是该分量的 Parry 测度。
\end{theorem}
\begin{proof}
正文已建立谱半径导数与平衡态 carry 密度之间的恒等式
\[
\partial_\theta\log\lambda_K(e^\theta)=\mathbb E_{\mu_{e^\theta}}[\kappa].
\]
同时，关于零温极限的结论表明：当 \(\beta\to+\infty\) 时，任意极限平衡态都必须支撑在使势函数 \(-\kappa\) 取最大值的轨道上；而这恰等价于
\[
\kappa=0.
\]
故任意零温极限 \(\mu_0\) 都满足 \(\mu_0(\kappa=0)=1\)，从而其支撑落在 carry-free 子图所生成的有限并 SFT 上。

若 carry-free 子图只有一个本质强连通分量，则正文的零温选择定理进一步表明极限平衡态唯一，并等于该 primitive 分量的 Parry 最大熵测度。证毕。
\end{proof}


\paragraph{整数审计的失明、稀疏化流与常数层规约：从 cyclotomic 不可见性到相位闭包的离散化制度}
在本文的“有限证书 + 离线核验”协议口径下，\(\bmod p^k\) 的整性护栏、常数层 strictification 与 unitary-slice 的相位闭包核验共同形成一条与阿基米德谱半径/极点数据正交的离散约束链；其关键后果可概括为：
\begin{enumerate}
  \item \textbf{Cyclotomic 不可见性：\texorpdfstring{$p$}{p}-幂层的整数审计对单位根扭曲失明。}
  若迹 ghost 序列满足 \(p\)-幂 Dwork 同余，则在任意奇素数 \(p\) 下，对所有 \(u^p=1\) 与一切 \(k\ge 1\) 都有
  \[
  a_{p^k}(u)\equiv a_{p^k}(1)\pmod{p^k}
  \]
  （推论 \ref{cor:witt-cyclotomic-invisibility}）。
  因此所有 \(p\)-幂长度层的 \(\bmod p^k\) 检验无法区分“扭曲”与“无扭曲”；扭曲差异只能在阿基米德谱半径与常数漂移层显影。

  \item \textbf{Frobenius 稀疏化的重整化流：高层信息在固定精度下被压缩到低层的 Frobenius 推前。}
  Dwork 同余逐系数比较强制 \(p^k\)-稀疏化：模 \(p^k\) 的多项式只落在 \((\ZZ/p^k\ZZ)[u^p]\)（命题 \ref{prop:witt-pk-sparsification}）；
  更进一步，对任意 \(1\le \ell\le k\) 存在迭代下降
  \(a_{p^k}(u)\equiv a_{p^{\ell-1}}(u^{p^{k-\ell+1}})\ (\mathrm{mod}\ p^\ell)\)
  （命题 \ref{prop:witt-frobenius-iterated-descent}），从而在固定精度 \(p^\ell\) 下“高层 \(p^k\)”的全部信息仅通过幂替换回推到低层。

  \item \textbf{常数级异常的可规约性：漂移向量形成 \texorpdfstring{$\mathcal{A}_G$}{AG}-torsor。}
  在有限阿贝尔 lift 群 \(G\) 上，纯常数 counterterm 门 \(\mathrm{CT}_\eta\) 保持 pole 数据不变，仅对有限部漂移做加性平移（定义 \ref{def:pom-finite-part-counterterm-gate}；命题 \ref{prop:pom-counterterm-anom-cancel}），
  因而在固定 pole 数据前提下，允许的漂移向量族构成以 \(\mathcal{A}_G\) 为结构群的仿射纤维（注 \ref{rem:xi-finite-part-drift-torsor}）。

  \item \textbf{相位闭包核验的强离散化：Tate 障碍退化为单一 Maslov 指标并制度性排除连续残差轴。}
  对自对偶 bounded type 函数的 Hardy/Smirnov 模型，等变 Tate 障碍分裂为点状整数部分与边界连续部分；若制度要求 unitary-slice 的相位闭包核验，则必要求奇异内因子 \(S_F\equiv 1\)，从而障碍退化为单一整数 \(\deg(B_F)\)（命题 \ref{prop:selfdual-maslov-singular-dichotomy}，并见定理 \ref{thm:xi-tate-maslov-pontryagin-unification}）。

  \item \textbf{Bartholdi fugacity 的无穷远根簇机制：\texorpdfstring{$t=1$}{t=1} 处出现 \texorpdfstring{$13$}{13} 条必然逸散支。}
  在真实输入核的 Bartholdi 二变量行列式闭式中，\(F(z,t)\) 在 \(t\to 1\) 时发生 \(z\)-次数塌缩 \(20\to 7\)，从而在射影意义下必有 \(13\) 条根支向无穷远逸散（注 \ref{rem:real-input-40-bartholdi-infinite-root-cluster}；定理 \ref{thm:real-input-40-bartholdi-det}）。
  更精细地，\((t-1)\)-adic Newton--Puiseux 多边形在 \(t=1\) 处强制出现三段负斜率边，给出逸散指数谱 \(\{1,\tfrac12,\tfrac13\}\) 及首项初形方程，从而把逸散条数与首项代数域提升为可审计的纯代数输出（命题 \ref{prop:real-input-40-bartholdi-newton-puiseux-spectrum}、\ref{prop:real-input-40-bartholdi-escape-initial-forms}；注 \ref{rem:real-input-40-bartholdi-graded-spectrum}）。

  \item \textbf{最坏角色的二次型判别律：选择律判别量 \texorpdfstring{$g(p)$}{g(p)} 的极限常数由 \texorpdfstring{$\Sigma_{33}$}{Sigma33} 单独给出。}
  在离散角色扫描中，令 \(\Sigma=\nabla^2P(0)\) 为平衡态 Fisher 二次型，则
  \[
  g(p):=\left(1-\frac{\rho_{3,3,p}}{\lambda}\right)p^2
  \xrightarrow[p\to\infty]{}\ 2\pi^2\,\Sigma_{33}
  \]
  （注 \ref{rem:arity-335-selection-law-gp}），从而把“最坏角色”的大模极限缩放常数压缩为单一协方差矩元的可证伪定律。

  \item \textbf{Green 核的整格刚性：谱中 \texorpdfstring{$4\sin^2$}{4 sin^2} 的系数是不可标度的算术归一化常数。}
  锯齿 Green 核 \(K_k(i,j)=\min(i,j)\) 具有整数 Gram 分解与 \(\det(K_k)=1\)（引理 \ref{lem:pom-Kk-gram-det}），从而强制
  \(\prod_{p=1}^k 4\sin^2(\varphi_p)=1\)
  与相应求和恒等式成立（推论 \ref{cor:pom-Kk-sine-product-sum}）；
  其中系数 \(4\) 由整格归一化刚性钉死，不能被任意标度吸收。
\end{enumerate}

\begin{proposition}[三阶结合缺陷的 \(2\)-群封装与 strictification 的上同调判别]\label{prop:conclusion-2group-postnikov-strictification}
\leanverified{paper\_conclusion\_2group\_postnikov\_strictification}
令 \(G\) 为群、\(A\) 为 \(G\)-模（交换群），并取一个 \(3\)-上同调类 \([\omega]\in H^3(G;A)\)。
则存在一个单对象单胚群oid \(\mathcal{G}_\omega\)（亦可视为 skeletal \(2\)-群），其不变量满足
\[
\pi_0(\mathcal{G}_\omega)\cong G,\qquad \pi_1(\mathcal{G}_\omega)\cong A,
\]
且其结合子由 \([\omega]\) 完全决定；并且 \([\omega]=0\) 当且仅当 \(\mathcal{G}_\omega\) 与某个严格 \(2\)-群单胚等价（即可 strictification）。
\par
进一步，若存在系数扩张 \(0\to A\to \widehat A\to B\to 0\) 使得 \([\omega]\) 在自然映射 \(H^3(G;A)\to H^3(G;\widehat A)\) 下像为零，则由 \(\widehat A\) 作为新的 \(1\)-自同态系数可实现 strictification：对应的 \(\widehat{\mathcal{G}}\) 具有平凡 Postnikov 类。
\end{proposition}
\begin{proof}[证明要点]
对任意 \(3\)-余环 \(\omega\in Z^3(G;A)\)，可取对象集为 \(G\)，并令 \(g\to g\) 的自同态集为 \(A\)（以加法记），张量积在对象上为群乘法、在自同态上为 \(A\) 的加法；把结合子指定为 \(\omega(g,h,k)\) 所给出的自然同构即可得到单对象单胚群oid。
不同余环相差 \(3\)-余边界当且仅当所得单胚结构单胚等价，故等价类仅由 \([\omega]\in H^3(G;A)\) 决定。
\par
若 \([\omega]=0\)，则存在 \(2\)-余链 \(\beta\) 使 \(\omega=\delta\beta\)，用 \(\beta\) 扭转结合子即可消去结合缺陷，从而 strictification。
反之，若 strictification 存在，则结合子可被单胚等价拉回为恒等，因此其 Postnikov 类必为零。
系数扩张的断言由自然性与“像为零 \(\Rightarrow\) 可选取代表为余边界”同理推出。证毕。
\end{proof}

\begin{conjecture}[POM 复制谱根与范畴维数的对齐字典]\label{conj:conclusion-pom-fpdim-dictionary}
在本文的 POM 编译口径下，设 \(\{r_q\}_{q\ge 2}\) 为 moment-kernel 的 Perron 根族（定理 \ref{thm:pom-collision-kernel-ak}），并设 \(\mathcal{A}\) 为由 strictification 商与异常数据所诱导的某个有限范畴化对称对象（例如由命题 \ref{prop:conclusion-2group-postnikov-strictification} 的 \(2\)-群封装所生成的表示/缺陷范畴）。
猜想存在一个使两类“复制”操作相容的对齐：\(q\) 次 MOM 放大对应于范畴中的 \(q\) 次张量幂，且 \(r_q\) 与相应张量幂的 Frobenius--Perron 量纲在同一归一化下给出一致的主量纲读数。
该对齐若成立，将把“投影词—异常—谱指纹”的审计链条提升为一个可比较的范畴维数接口。
\end{conjecture}

\paragraph{异常轴的两点有限化与 RH-sharp 梳齿的有限探针极限：互素判别、规避机制与分辨率标度}
在本文的 \(\mathbf{PW}\) 语义中，异常类 \([\mathrm{Anom}]\) 既是重写闭环的路径不变量，也是“无限族约束是否可有限证书化”的逻辑触发轴；在 RH-sharp 的循环 lift 闭包中，临界线极点的虚部结构进一步提供了一个可量化的“有限采样”极限。其结论性接口可概括为：
\begin{enumerate}
  \item \textbf{互素两点判别（无需 torsion-free）：全矩阶零异常的两点完备性。}
  在异常上同调群 \(H\) 中，矩放大律
  \(
  [\mathrm{Anom}(\mathrm{MOM}_q\circ w)]=q\,[\mathrm{Anom}(w)]
  \)
  使得任意两点互素零检验已足以推出 \([\mathrm{Anom}(w)]=0\)（定理 \ref{thm:pom-anom-coprime-two-point-test}），并等价实现“对所有 \(q\ge 2\)”的全矩阶零异常命题的有限化（定理 \ref{thm:pom-anom-all-moments-collapse-to-two}；推论 \ref{cor:pom-anom-oracle-collapse-two-point-trigger}）。

  \item \textbf{有限探针不可证伪：循环 lift 可规避任意有限采样集。}
  若存在达界模态 \(\mu=\sqrt{\lambda}e^{i\theta}\)，则循环 lift \(M^{\langle q\rangle}\) 在临界线生成 \(q\)-等分梳齿；
  除去平凡对齐特例外，任意有限采样集 \(S\) 都可被无穷多个 lift 阶 \(q\) 同时规避（定理 \ref{thm:zeta-cyclic-lift-finite-probe-evasion}）。
  因而“有限探针审计”在 lift 闭包意义下无法排除 RH-sharp 的临界线极点存在性。

  \item \textbf{可寻址网格的分辨率：一般相位为 \texorpdfstring{$Q^{-1}$}{Q^-1}，对齐相位可改进为 \texorpdfstring{$Q^{-2}$}{Q^-2}。}
  令 \(Q\) 为允许的循环 lift 阶上界，则 RH-sharp 梳齿并集在基本周期上诱导的最坏分辨率 \(\varepsilon_Q\) 满足
  \(
  \varepsilon_Q\asymp (Q\log\lambda)^{-1}
  \)
  的一般尺度（推论 \ref{cor:zeta-cyclic-lift-resolution-law}），并且点数计数给出普适的 \(Q^{-2}\) 级下界（命题 \ref{prop:zeta-cyclic-lift-covering-radius-counting-lb}）。
  在对齐情形 \(\theta\equiv 0\) 下，网格进一步退化为标准 Farey 型有理网格，从而最坏覆盖半径可达 \(\Theta(Q^{-2})\) 的改进（注 \ref{rem:zeta-cyclic-lift-aligned-farey-q2}）。
\end{enumerate}


\paragraph{纤维谱的有理完全性与交换对称 MOM 坐标：有限重建、统一逼近、零温集中与 Euler 缺陷能量}
在固定分辨率 \(m\) 的层面，本文的碰撞矩与 escort 机制不仅提供指数尺度的压力/谱读出，也给出一组严格的“有限证书化”接口：它们把纤维多重度的离散谱、交换对称可观测量的可学习空间以及 Euler 缺陷的交叉能量统一压缩为可计算且可审计的有限对象。其结论性接口可概括为：
\begin{enumerate}
  \item \textbf{纤维谱 resolvent 的有理结构与完全重建：有限个幂和矩决定全部极点与重数。}
  固定 \(m\) 后，幂和矩生成函数 \(H_m(u)=\sum_{q\ge 0}S_q(m)u^q\) 具有严格部分分式展开，其极点集合与留数分别编码纤维多重度的不同取值 \(\{a_j\}\) 与其重数 \(\{c_j\}\)（定理 \ref{thm:pom-fiber-spectrum-resolvent-rational}）。
  进一步，若 \(R\) 为该层的不同多重度取值个数，则至多 \(2R\) 个幂和数据 \(\{S_q(m)\}_{q=0}^{2R-1}\) 即可唯一重建 \(H_m\)，并且在一般情形下少一项将失去唯一性（定理 \ref{thm:pom-fiber-spectrum-finite-reconstruction-sharp}）。

  \item \textbf{交换对称可观测量的统一逼近：\texorpdfstring{$\{p_q\}$}{pq} 生成全部可学习对称统计量。}
  在概率单纯形 \(\Delta(X_m)\) 上，幂和特征 \(p_q(w)=\sum_x w(x)^q\)（\(2\le q\le |X_m|\)）生成的代数在一致范数下稠密于所有交换对称连续可观测量空间（定理 \ref{thm:pom-mom-symmetric-observables-stone-weierstrass}）。
  因而在“对类型标签不敏感”的学习/审计口径下，有限个碰撞矩特征已足以统一逼近任意对称统计量。

  \item \textbf{碰撞矩坐标图：一般位置下有限幂和给出对称分布空间的局部坐标。}
  在单纯形内点且权重两两不等的开稠密集合上，映射 \(w\mapsto(p_2(w),\dots,p_n(w))\) 的 Jacobian 由 Vandermonde 行列式控制，从而在交换对称意义下成为局部坐标图（定理 \ref{thm:pom-power-sum-vandermonde-jacobian}；推论 \ref{cor:pom-power-sum-local-chart}）。

  \item \textbf{零温极限定量律：escort 在 \texorpdfstring{$q\sim m$}{q~m} 标度下呈现 \texorpdfstring{$m^2$}{m^2} 级超指数集中。}
  escort 分布 \(w_{m,q}(x)\propto d_m(x)^q\) 对“低于最坏指数 \(\alpha_m-\varepsilon\)”的集合质量满足显式上界 \(w_{m,q}(A_m(\varepsilon))\le |X_m|e^{-q\varepsilon m}\)，从而当 \(q=\Theta(m)\) 时出现 \(\exp(-\Theta(m^2))\) 的超指数收缩（推论 \ref{cor:pom-escort-zero-temperature-superexp-concentration}）。

  \item \textbf{预言机反演的预算曲线：成功指数、容量指数与热相切割的闭式读出。}
  在 \(w_m\) 下，受限比特预算的最优反演成功率的指数主项由 \(I_w\) 通过折叠变分式严格给出，并可改写为双支极值（定理 \ref{thm:pom-oracle-success-variational-laplace}）。
  相应的容量指数满足 \(\sup_{\alpha}(f(\alpha)+\min(\alpha,a))\) 的变分表述，并在热相区间内由单层 \(\alpha=a\) 主导，从而多重分形计数谱可由容量曲线直接读出（定理 \ref{thm:pom-oracle-capacity-thermal-cut}；推论 \ref{cor:pom-oracle-capacity-slope-readout}）。
  同时微正则条件化与 escort 在层集上具有相对熵率为零的信息等价，临界预算窗满足高斯极限并给出二阶反演公式（定理 \ref{thm:pom-microcanonical-escort-kl-rate-zero}；定理 \ref{thm:pom-oracle-critical-window-gaussian}；推论 \ref{cor:pom-oracle-critical-window-bit-inversion}）。

  \item \textbf{Euler 缺陷的正交能量：交叉项 \texorpdfstring{$L^2$}{L2} 能量等于到可加子空间的最小距离。}
  在 \(\widehat G=\prod_i\widehat{G_i}\) 的均匀测度下，坐标平均算子诱导 Hoeffding 型正交分解并给出 Pythagoras 恒等式（定理 \ref{thm:pom-euler-defect-orthogonal-pythagoras}）。
  相应地，二阶及以上交叉分量的总能量严格等于到“常数 \(\,+\)”轴可加子空间 \(\mathcal S_{\mathrm{add}}\) 的最小 \(L^2\) 距离（定理 \ref{thm:pom-euler-defect-min-action-gauge}），从而提供一条规范化的 Euler 缺陷能量 gauge。
\end{enumerate}


\paragraph{乘法链式律与熵亏损场：压力谱的乘法分解、几何望远镜与冻结判别}
在 escort 口径下，整数压力谱 \(\{P_q\}_{q\ge 1}\)（其中 \(P_q=\lim_{m\to\infty}\frac{1}{m}\log S_q(m)\)）与 R\'enyi 熵率之间存在一组严格的“乘法一致性”恒等式：它们把乘法尺度 \(q\mapsto ab\) 上的压力跃迁精确分解为熵亏损项，并由此导出几何望远镜收敛、冻结判别与可实现谱的约束锥。其结论性接口可概括为：
\begin{enumerate}
  \item \textbf{乘法链式律：压力谱在乘法尺度上的精确分解等式。}
  有限层面上，escort 分布 \(\pi_{m,a}=w_{m,a}\) 的 R\'enyi 熵给出严格恒等式
  \(\log S_{ab}(m)=b\log S_a(m)-(b-1)H_b(\pi_{m,a})\)（定理 \ref{thm:pom-multiplicative-chain-rule-finite-m}）；
  在极限层面则得到乘法分解
  \(P_{ab}=bP_a-(b-1)h^{\mathrm{esc}}_b(a)\)（定理 \ref{thm:pom-pressure-entropy-factorization-multiplicative}）。

  \item \textbf{乘法细化单调性：\texorpdfstring{$\alpha_\ast$}{alpha*} 的几何望远镜收敛与显式误差界。}
  乘法细化序列 \(\alpha_{a,b}(n)=P_{ab^n}/(ab^n)\) 单调不增，并具有望远镜展开与可计算误差上界
  \(0\le \alpha_{a,b}(N)-\alpha_\ast\le (\log\varphi)/(ab^N)\)（定理 \ref{thm:pom-multiplicative-refinement-telescope-alpha}）。

  \item \textbf{可观测性：整数压力谱完全决定全部 escort R\'enyi 熵率。}
  熵率曲面可由压力谱逐点反演：
  \(h^{\mathrm{esc}}_b(a)=(bP_a-P_{ab})/(b-1)\)，且基准分布满足 \(h^{\mathrm{esc}}_b(1)=(b\log 2-P_b)/(b-1)\)（推论 \ref{cor:pom-escort-renyi-rate-inversion}）。

  \item \textbf{熵亏损 \(1\)-cocycle：乘法半群上的平直结构与分解公式。}
  熵亏损场 \(\mathcal{H}_{\mathrm{loss}}(a,b):=bP_a-P_{ab}\) 满足严格 \(1\)-cocycle 方程
  \(\mathcal{H}_{\mathrm{loss}}(a,bc)=c\mathcal{H}_{\mathrm{loss}}(a,b)+\mathcal{H}_{\mathrm{loss}}(ab,c)\)，并诱导因子链展开与交换一致性的平直性恒等式（定义 \ref{def:pom-entropy-loss-cocycle}；定理 \ref{thm:pom-entropy-loss-1cocycle-flat}；推论 \ref{cor:pom-entropy-loss-factor-chain-expansion}）。

  \item \textbf{零温度恒等式：escort 的 min-entropy 率等于 \texorpdfstring{$P_a-a\alpha_\ast$}{Pa-a alpha*}。}
  escort 的 min-entropy 熵率满足 \(h_\infty^{\mathrm{esc}}(a)=P_a-a\alpha_\ast\)，并给出冻结判别：\(h_\infty^{\mathrm{esc}}(a)=0\Longleftrightarrow P_{ab}=bP_a\)（定理 \ref{thm:pom-escort-minentropy-zero-temperature-identity}；推论 \ref{cor:pom-escort-freezing-criterion}）。

  \item \textbf{熵谱可实现性的约束锥：由界、平直性与冻结极限导出非平凡可行域。}
  \(\mathcal{H}_{\mathrm{loss}}\) 的非负/上界、\(1\)-cocycle 平直性以及冻结极限兼容性共同定义压力谱的可行域 \(\mathfrak{C}\)，从而对任何候选压力谱提供可证伪的必要约束（定理 \ref{thm:pom-entropy-loss-feasibility-cone}）。
\end{enumerate}


\paragraph{对称幂升级的平方根谱界不稳定：第一变分放大、线性激发过滤与预言机化选择轴}
在本文的 moment-kernel 编译口径下，“占据数对称幂商”是一条自然的多项式态压缩路线（命题 \ref{prop:pom-moment-symmetric-power}；命题 \ref{prop:pom-factorized-sympower-upper-bound}）。
然而，对称幂在谱界上并不保守：它会把任何非 Perron 模态沿乘法尺度系统性放大到最终穿越平方根界，从而把“核版 RH”从一个可验证约束提升为一个\emph{需要额外选择轴才能维持}的结构性工程问题。其结论性接口可概括为：
\begin{enumerate}
  \item \textbf{对称幂的平方根谱界不稳定性：除谱秩一外，任意原始核终将违反核版 RH。}
  对原始非负矩阵 \(A\)，若非 Perron 谱半径 \(\eta>0\)，则存在 \(b_0\) 使得对所有 \(b\ge b_0\)，\(\mathrm{Sym}^b(A)\) 的非 Perron 谱半径严格大于 \(\sqrt{\rho(A)^b}\)，从而对称幂最终必违反核版 RH 的平方根谱界（定理 \ref{thm:pom-sympower-sqrt-bound-instability}）。
  并且“对所有 \(b\ge 2\) 始终满足平方根谱界”当且仅当 \(\eta=0\)，亦即非零谱仅含 Perron 根的谱秩一退化（推论 \ref{cor:pom-sympower-sqrt-bound-stable-iff-spectral-rankone}）。

  \item \textbf{第一变分保真 \texorpdfstring{$\Rightarrow$}{=>} 高阶核版 RH 不可能：MOM 乘法升级的致命放大。}
  若乘法升级 \(a\mapsto ab\) 的编译存在满射 intertwiner \(\Pi_{a,b}:\mathrm{Sym}^b(A_a)\twoheadrightarrow A_{ab}\) 且在第一变分扇区上不被消灭（第一变分保真），则当 \(b\) 足够大时 \(A_{ab}\) 必出现 \(\eta_{ab}>\sqrt{\rho_{ab}}\) 的越线（定义 \ref{def:pom-first-variation-sector}；定理 \ref{thm:pom-first-variation-fidelity-rh-breaking}）。

  \item \textbf{线性“激发过滤”下界：维持核版 RH 必须消灭线性比例的低激发扇区。}
  对任意 \(\mathrm{Sym}^b(A)\) 的非零子商 \(B_b\)，若其满足核版 RH，则最小存活激发数 \(k_b\) 必满足显式线性下界
  \(k_b\ge \frac{b}{2}\frac{\log\rho}{\log(\rho/\eta)}\)，在 RH-sharp 情形进一步强制 \(k_b\ge b\)（定理 \ref{thm:pom-excitation-filter-linear-bound}）。

  \item \textbf{对称幂商编译的不可避免障碍：无限核 RH 必然要求额外选择子，等价于预言机轴。}
  若 \(A_{ab}\) 的 realization 仅由占据数对称幂商给出（不引入随 \(b\) 变化的额外投影/选择子去消灭低激发扇区），则高阶核版 RH 必破裂；
  反之，若要求对固定 \(a\) 的全部 \(b\ge 2\) 同时维持核版 RH，则除谱秩一退化外，必须引入满足线性阈值的激发过滤选择子族 \(P_{a,b}\)，其在协议语义中表现为一条外置的“选择/预言机”编译轴（定理 \ref{thm:pom-occupancy-symquotient-rh-obstruction}）。
\end{enumerate}


\paragraph{异常同调、合流模异常与零温序门的可审计编译接口}
以下条目作为对 \S\ref{subsubsec:xi-horizon-ramanujan-feasible-reduction} 的“六点定位”与 \S\ref{subsubsec:xi-ramanujan-horizon-ledger-framework} 的“六模块闭环”之外的\emph{补充层}：它们把“不可交换/不可并行”的残差提升为可计算的 \(2\)-上同调类（曲率），把优先级与控制流解释为零温序投影的硬化极限，并给出一条将计算完备性与谱/证书语言对接的最小编译通路。与视界 Toeplitz/SDP 证书链已在上述两处完成总体定位的部分，此处仅给出可直接引用的闭合陈述与编号指引，不再重复其“六点/六模块”的总览表述。
\begin{enumerate}
  \item \textbf{异常的 \(2\)-上同调定域化与线性放大：从“不可交换”到可审计曲率。}
  在三生成元片段上，异常代表元随规约证书规范的改变仅差一个 coboundary，从而诱导出路径不变量的异常上同调类 \([\mathrm{Anom}(w)]\)，并对水平复合满足加法律（命题 \ref{prop:pom-anom-additive-up-to-boundary}；推论 \ref{cor:pom-anom-cohomology-class}）。
  在含 \(\mathrm{MOM}_q\) 的扩展片段中，moment 化对异常类满足严格的 \(q\)-倍线性放大，从而可将任意非零残差提升到任意预设审计阈值而不离开有限核可编译的语义域（命题 \ref{prop:pom-anom-moment-amplification}；命题 \ref{prop:pom-anom-gap-amplification}）。

  \item \textbf{合流模异常与算法等价：投影词判等由 \texorpdfstring{$\mathrm{NF}$}{NF} 与 \texorpdfstring{$[\mathrm{Anom}]$}{[Anom]} 完全决定。}
  在有界审计口径下，“绝对合流”被替换为“合流模异常”：任意临界对的两支规约可并到同一接口正规形，而残差差异必落在边界项中（命题 \ref{prop:pom-confluence-mod-anom-audited}；亦见注 \ref{rem:pom-residual-2confluence}）。
  因而算法等价可取为“正规形一致 + 异常类一致”，并在该片段内给出可终止的判等流程（定义 \ref{def:pom-algorithmic-equivalence-nf-anom}；推论 \ref{cor:pom-three-gen-normal-form-decidable}）。

  \item \textbf{零温序门的计算完备性接口：FRACTRAN 作为受序投影门控的素数纤维平移。}
  FRACTRAN 的一步语义可在素数指数向量纤维 \(\mathcal{P}\) 上严格实现为“偏序可行性 \(\Rightarrow\) first-fit 选择 \(\Rightarrow\) 确定性平移”，从而把控制流外置为序投影的选择机制（定义 \ref{def:pom-fractran}；命题 \ref{prop:pom-fractran-prime-translation}）。
  并且 first-fit 的确定性优先级可解释为温度化序投影 \(\mathrm{PROJ}_u\) 在 \(u\to 0\) 的零温硬化极限（定义 \ref{def:pom-proj-temp}；注 \ref{rem:pom-order-as-zero-temp}）。

  \item \textbf{Ihara--Artin 通道化下 cyclotomic 因子刻画模障碍：共振来源的完全压缩。}
  在扭曲 Ihara--Artin \(\zeta\) 的行列式因子化框架中，非平凡表示通道达到 Perron 半径（模障碍/共振通道出现）当且仅当相应因子含 cyclotomic 多项式因子；从而“谱隙失败/慢模来源”被压缩为一族可枚举的单位根障碍（定理 \ref{thm:conclusion75-ihara-cyclotomic-factor-classification}；亦见定理 \ref{thm:gm-graph-zeta-cyclotomic-classification}）。

  \item \textbf{异常 strictification 与散射绕数的指数同一：代数组合残差 \(=\) 拓扑障碍。}
  在 unitary slice 的可见通道口径下，异常严格为零等价于散射回路可取高能归一化且绕数为零；反之，若全部可见通道绕数可被归一化为零，则闭环接口必因子化通过普适 strictification 商 \(G^{\mathrm{BC}}_m\) 并在其上实现严格交换（定理 \ref{thm:xi-anom-winding-obstruction}；命题 \ref{prop:pom-bc-quotient-strictify}、\ref{prop:pom-bc-quotient-universal}）。

  \item \textbf{扫描剖面的逆问题：轴外缺陷参数的唯一重建。}
  在有限缺陷口径下，扫描剖面（Poisson/Herglotz 型叠加）作为实解析函数，其亚纯延拓的极点集合按重数唯一编码全部轴外参数；因此给定任意非平凡实区间上的剖面数据即可唯一重建缺陷参数多重集（定理 \ref{thm:xi-offcritical-poisson-herglotz-inversion}）。

  \item \textbf{端点热核探针：有限 Toeplitz 数据对端点原子质量的单调抽取。}
  对 Carath\'eodory/Herglotz 表示下的 Toeplitz 主子矩阵 \(\mathbf{T}_N\)，二项式探针 \(a_N\) 具有严格热核恒等式并单调收敛到端点原子质量 \(\mu(\{-1\})\)，从而端点阈值承载可由有限阶 Toeplitz 数据在不引入额外连续轴的前提下单调抽取（定理 \ref{thm:xi-endpoint-heat-kernel-probe}）。

  \item \textbf{核版 (G)RH 的区间半正定证书与 SDP 归约：从通道多项式到可审计 LMI。}
  在圆盘完成化与自倒数多项式口径下，核版 \(\Xi\)-RH 等价于一元多项式根落入区间 \([-2,2]\)（命题 \ref{prop:xi-rh-interval-criterion}），并可通过 Schur--Cohn Hermitian 形式导出显式 LMI/SDP 证书；当系数落入 \(\QQ[u+u^{-1}]\) 时进一步压缩为区间 \(x\in[-2,2]\) 上的区间半正定约束，并可由 SOS/半定规划给出可审计证书（注 \ref{rem:xi-schur-cohn-sdp-interface}）。

  \item \textbf{素数矩半代数锥：每一级 Toeplitz--PSD 约束的绝对收敛化。}
  在 \(|w|<1/3\) 的绝对收敛补丁内，Toeplitz 矩数据可由绝对收敛的 von Mangoldt 素数矩生成，从而每一级 Toeplitz--PSD 约束可被等价转写为由绝对收敛素数矩生成的有限半代数锥不等式系统（定理 \ref{thm:xi-prime-moment-semialgebraic-cone}）。

  \item \textbf{strictification 群骨架的逆极限：可见对称性的 profinite 完成。}
  分辨率塔诱导 strictification 商群的逆系统 \(\{G^{\mathrm{BC}}_m\}\)，并定义其逆极限 \(G^{\mathrm{BC}}_\infty=\varprojlim_m G^{\mathrm{BC}}_m\)，从而为“无限分辨率下仍与折叠严格相容的可见寄存器骨架”提供内禀 profinite 容器（推论 \ref{cor:pom-bc-inverse-system}）。

  \item \textbf{负平方指数的稳定读出：Toeplitz 负惯性 \(=\) Blaschke 缺陷 \(=\) 真共振计数。}
  在有限型点状缺陷口径下，广义 Carath\'eodory/Schur 的负平方指数可由高阶 Toeplitz 截断的负惯性稳定值完全审计读出，从而盘内点状共振（真共振）计数可由有限维线性代数在足够高阶截断上稳定提取（定理 \ref{thm:xi-toeplitz-inertia-stabilizes-to-kappa}；亦见定理 \ref{thm:xi-toeplitz-negative-squares-equals-blaschke-defect}）。
\end{enumerate}


\paragraph{二幂 Chebyshev--Dwork 闭包、确定性强 Szeg\H{o} 分解与主弧 Gram 谱判定：从 \texorpdfstring{$p$}{p}-进指纹到瓶颈的单一谱证书}
下述结论把三条表面上异质的“可审计计算”机制统一为可引用接口：二幂塔上的同余闭包与对数深度、离散 Green/Laplacian 的对象级行列式恒等式与黄金点算术化、以及主弧集中与有限维 Gram 谱半径判定的逆向刻画。
\begin{enumerate}
  \item \textbf{\texorpdfstring{$p=2$}{p=2} 的 Frobenius 在不变量环 \texorpdfstring{$\ZZ[t]$}{Z[t]} 上闭合。}
  完成化 Dwork 同余在不变量坐标 \(t=u+u^{-1}\) 上严格闭合为 Chebyshev--Dwork 同余链
  \(
  A_{2^k}(t)\equiv A_{2^{k-1}}(t^2-2)\ (\mathrm{mod}\ 2^k)
  \)，从而把 \(u\mapsto u^2\) 的 \(2\)-进算术兼容性精确转译为 \(\ZZ[t]\) 上的倍角动力学（定理 \ref{thm:chebyshev-dwork-congruence-chain}）。

  \item \textbf{对数深度的矩阵自由审计：固定精度 \texorpdfstring{$2^k$}{2^k} 的计算可在单变量动力学上完成。}
  在二幂塔上，关于 \(B(u)^{2^k}\) 的同余核验可被替换为在 \((\ZZ/2^k\ZZ)[t]\) 中迭代代入 \(t\mapsto t^2-2\) 并取模；递推步数为 \(k\)，相对于长度 \(n=2^k\) 呈对数深度（推论 \ref{cor:chebyshev-dwork-matrixfree}）。

  \item \textbf{稳定导数残差谱：由 \texorpdfstring{$p$}{p}-进缩放律生成一族无限 \texorpdfstring{$\mathbb{F}_p$}{F_p} 指纹。}
  以 \(u=e^\theta\) 定义 \(\theta\)-导数矩 \(A^{(r)}_n\)，Dwork 同余推出逐阶缩放同余
  \(A^{(r)}_{p^k}\equiv p^rA^{(r)}_{p^{k-1}}\ (\mathrm{mod}\ p^k)\)
  （命题 \ref{prop:witt-theta-derivative-filter}），并因此得到稳定首位残差
  \(\alpha_{p,r}\in\mathbb{F}_p\) 以及 \(p=2\) 的稳定二进制指纹 \(\varepsilon_r\in\{0,1\}\)（推论 \ref{cor:witt-theta-stable-residue} 与推论 \ref{cor:witt-theta-binary-fingerprint}）。

  \item \textbf{强 Szeg\H{o} 分解的严格恒等式版本：面积项、常数项与指数小修正全显式封闭。}
  对混合边界 Laplacian \(L_k\) 与 \(t>0\)，存在对象级恒等式
  \[
  \det(L_k+tI)=\frac{\cosh((2k+1)\eta(t))}{\cosh(\eta(t))}
  =q(t)^{-k}\frac{1+q(t)^{2k+1}}{1+q(t)},
  \]
  并给出 \(\log\det\) 的“面积项 + 常数项 + 指数小余项”严格分解与显式上界证书（推论 \ref{cor:pom-Lk-area-law-qform}；并见命题 \ref{prop:pom-Kk-strong-szego-bulk-surface}）。

  \item \textbf{黄金耦合的唯一性与质量隙锁定：\texorpdfstring{$t=1$}{t=1} 唯一把自由能密度锁到 \texorpdfstring{$2\log\varphi$}{2 log phi}。}
  方程 \(\eta(t)=\log\varphi\) 在 \(t>0\) 上至多一解，且与 \(q(t)=\varphi^{-2}\) 及
  \(\lim_{k\to\infty}\frac1k\log\det(L_k+tI)=2\log\varphi\)
  严格等价，从而给出同一唯一正解 \(t_\star=1\)（推论 \ref{cor:pom-Lk-golden-coupling-unique}）。

  \item \textbf{黄金点的全算术化：分母统一为 \texorpdfstring{$F_{2k+1}$}{F_{2k+1}} 的两点 Green 关联闭式。}
  在 \(t=1\) 下，\(A_k=L_k+I\) 的行列式与逆矩阵元完全算术化：
  \(\det(A_k)=F_{2k+1}\) 且
  \((A_k^{-1})_{ij}=\frac{F_{2i}F_{2(k-j)+1}}{F_{2k+1}}\)
  （推论 \ref{cor:pom-Lk-t1-fibonacci-det-green}；并见注 \ref{rem:pom-Lk-fibonacci-area-law-exact} 对面积项/常数项/指数修正的 Fibonacci 恒等式实现）。

  \item \textbf{边界 Riccati 流的精确指数收敛与黄金点的 Fibonacci 迭代。}
  反射系数 \(q_k(t)=(L_k+tI)^{-1}_{11}\) 满足 Schur--Riccati 递推并具有完全显式误差形式
  \(0<q_k(t)-q(t)=\frac{(1-q(t)^2)q(t)^{2k}}{1+q(t)^{2k+1}}\)；
  在 \(t=1\) 下退化为 Fibonacci 比
  \(q_k(1)=F_{2k-1}/F_{2k+1}\to\varphi^{-2}\)
  （命题 \ref{prop:pom-Lk-boundary-riccati-recursion}；并见推论 \ref{cor:pom-Lk-boundary-fixed-point-riccati}）。

  \item \textbf{边界热力学响应与 Parry 质量的同源闭合：谱侧导数等于测度侧期望。}
  表面自由能 \(g(t)=-\log(1+q(t))\) 满足
  \(g'(t)=\frac{q(t)}{1+q(t)}\cdot\frac{1}{\sqrt{t(4+t)}}\)
  （推论 \ref{cor:pom-Lk-surface-free-energy}）；
  在 \(t=1\) 时有
  \(\frac{q(1)}{1+q(1)}=\pi(1)\) 且 \(g'(1)=\pi(1)/\sqrt5\)，其中 \(\pi(1)\) 为黄金均值 shift 的 Parry 最大熵测度在符号 \(1\) 上的质量（推论 \ref{cor:pom-Lk-t1-parry-surface-derivative}）。

  \item \textbf{负谱行列式的完备不变量：Poisson 权重与伪双曲 Vandermonde 的乘积律。}
  点状缺陷的 Pick--Poisson Gram 行列式具有闭式分解
  \(
  \det\mathsf P=(\prod_j P_{a_j}(-1))(\prod_{j<k}\rho(a_j,a_k)^2)
  \)，并与 Toeplitz 证书的负特征值乘积极限严格对齐：
  \(\prod_{j=1}^{\kappa}(-\lambda_j^-(\mathbf{T}_N))\to\det\mathsf P\)
  （命题 \ref{prop:xi-pick-poisson-det-factorization}）。

  \item \textbf{\texorpdfstring{$M^4$}{M^4} 逆原理与统一瓶颈的单一谱判定化：改进性归约为有理频率 Gram 谱半径的数值衰减。}
  仿射均匀化残差的 \(M^4\) 近饱和当且仅当 Fourier 质量在某一有理主弧上集中（定理 \ref{thm:gm-affine-inverse-major-arc}），并强制产生高密度近仿射自相关结构（命题 \ref{prop:gm-major-arc-rigidity-affine-autocorr}）。
  同时残差谱范数在常数意义下等价于显式有限维大筛 Gram 矩阵 \(K_M\) 的最大特征值（定理 \ref{thm:gm-residual-opnorm-gram-equivalence}），从而所有幂次改进空间统一归约为单一可计算判定
  \(
  \|K_M\|_{\mathrm{op}}\ll M^{-\vartheta}
  \)
  是否成立（推论 \ref{cor:gm-KM-computable-vtheta}；并见二分律定理 \ref{thm:gm-affine-improvement-dichotomy}）。
\end{enumerate}


\paragraph{有限域影子与 Jordan 毛发：共振窗内递推动力学的最终多项式律与可计算周期证书}
在共振窗的最小整系数递推已被审计固定后（定理 \ref{thm:pom-moment-resonance}），特征多项式 \(P_q(\lambda)\) 在有限域 \(\mathbb{F}_p\) 上的约化与因子指数把序列 \(S_q(m)\) 的长期动力学压缩为一组完全内部、可复核的指纹接口：它们既不依赖额外的连续极限假设，也不需要引入外部近似模型，而直接由可计算的有限维代数对象给出。与正文的模态共振链条（注 \ref{rem:pom-charpoly-mod2-shadow} 与注 \ref{rem:pom-charpoly-mod3-selection}）相衔接，可将该类接口概括为以下五条互补断言。
\begin{enumerate}
  \item \textbf{模 \texorpdfstring{$2$}{2} 的最终多项式律：差分湮灭把奇偶动力学压缩为有限阶离散多项式。}
  在共振窗口 \(q=9,\dots,17\) 上，模 \(2\) 的影子坍缩给出
  \(P_q(\lambda)\equiv \lambda^{a_q}(\lambda+1)^{b_q}\ (\mathrm{mod}\ 2)\)，从而移位算子满足
  \(E^{a_q}(E+1)^{b_q}S_q\equiv 0\ (\mathrm{mod}\ 2)\)。
  记 \(\mathsf{D}:=E+1\)（特征 \(2\) 下与前向差分等同），则在递推生效区间内有
  \(\mathsf{D}^{b_q}S_q(m+a_q)\equiv 0\ (\mathrm{mod}\ 2)\)，因此 \(S_q(m+a_q)\bmod 2\) 在长期段退化为次数 \(<b_q\) 的离散多项式（注 \ref{rem:pom-charpoly-mod2-shadow}；表 \ref{tab:fold_collision_charpoly_mod2_shadow_q9_17}）。
  由此 \((a_q,b_q)\) 不仅是形式分解的指数对，也可作为“最终多项式阶/瞬态阶”的完整不变量，用以刻画奇偶动力学在足够后段的确定性退化结构。
  在更低阶的已核验窗口 \(4\le q\le 11\) 内，指数对进一步凝结为闭式
  \(b_q=2\lfloor (q+1)/4\rfloor\)、\(a_q=3+2\,\ind\{q\ \mathrm{even},\ q\ge 6\}\)，并呈现 \(q\mapsto q+4\) 的台阶相（命题 \ref{prop:pom-charpoly-mod2-closed-exponent-law-q4-11}）。

  \begin{proposition}[\(\FF_2\) 上差分幂零所强制的二项式基多项式性]\label{prop:conclusion-f2-binomial-basis-from-delta-nilpotent}
  令 \((a_m)_{m\ge 0}\) 为取值于 \(\FF_2\) 的序列，并记前移算子 \(E(a)_m=a_{m+1}\) 与差分算子 \(\Delta:=E+1\)（特征 \(2\) 下亦可写为 \(E-1\)）。
  若存在整数 \(b\ge 1\) 使得对充分大的 \(m\) 有
  \[
  \Delta^{b}a_m=0,
  \]
  则存在常数 \(c_0,\dots,c_{b-1}\in\FF_2\) 使得对充分大的 \(m\) 有
  \[
  a_m=\sum_{j=0}^{b-1}c_j\binom{m}{j}\quad(\mathrm{mod}\ 2).
  \]
  \leanverified{paper\_conclusion\_f2\_binomial\_basis\_from\_delta\_nilpotent}
  \end{proposition}

  \begin{proof}
  组合恒等式给出 \(\Delta\binom{m}{j}=\binom{m}{j-1}\)，从而 \(\Delta^b\binom{m}{j}=0\) 当且仅当 \(j<b\)。
  若 \(\Delta^b a_m=0\) 在后段成立，则 \(\Delta^{b-1}a_m\) 在后段为常值，记为 \(c_{b-1}\)。
  令 \(a^{(1)}_m:=a_m-c_{b-1}\binom{m}{b-1}\)，则 \(\Delta^{b-1}a^{(1)}_m=0\) 在后段成立。
  递归下降直至阶数 \(0\)，得到所示表示。证毕。
  \end{proof}

  \item \textbf{模 \texorpdfstring{$3$}{3} 的两扇区分解：常相多项式与交替相多项式的直和压缩。}
  当 \(P_q(\lambda)\equiv \lambda^2(\lambda-1)^4(\lambda+1)^5\ (\mathrm{mod}\ 3)\)（例如 \(q=13,15\)；注 \ref{rem:pom-charpoly-mod3-selection}）时，对 \(\widetilde S_q:=E^2S_q\) 有
  \((E-1)^4(E+1)^5\widetilde S_q\equiv 0\ (\mathrm{mod}\ 3)\)，从而存在次数分别受 \((4,5)\) 指数对控制的两类广义相位模态：一类为关于 \(m\) 的次数 \(\le 3\) 多项式，另一类为 \((-1)^m\) 乘以次数 \(\le 4\) 多项式（命题 \ref{prop:pom-moment-mod3-two-phase}；表 \ref{tab:fold_collision_charpoly_mod3_selection_q13_15}）。
  该两扇区分解把模 \(3\) 世界中的全部复杂性压缩为“常相 + 交替相”两块，并为后续的 Jordan 缺陷检测与嵌入事件定位提供基础层约束。

  \item \textbf{\texorpdfstring{$A_2$}{A2} 嵌入与判别式零化的并置：局部谱分歧点作为有限域诊断。}
  当 \(\chi_{A_2}(\lambda)\mid P_q(\lambda)\) 于 \(\mathbb{F}_p[\lambda]\) 成立时，模 \(p\) 的递推状态空间中出现一个三维不变子动力学，其 companion 作用的特征多项式恰为 \(\chi_{A_2}\)（命题 \ref{prop:pom-charpoly-modp-a2-embedding}；推论 \ref{cor:pom-charpoly-modp-a2-subdynamics}）。
  在已核验切片 \((q,p)=(14,7)\)、\((12,11)\) 中，该嵌入旗标与 \(\mathrm{repeat?}=1\) 同时出现，从而 \(\widetilde P_q:=P_q/\lambda^{a_0}\) 非平方自由并满足 \(\gcd(\widetilde P_q,\partial_\lambda\widetilde P_q)\neq 1\)；等价地，其判别式在模 \(p\) 意义下为零（注 \ref{rem:pom-a2-embedding-local-ramification}）。
  因此嵌入事件可被识别为发生在伴随 Jordan 调制通道的局部谱分歧点处，形成“子表示嵌入 \(\times\) 分歧/非半单性”的联合指纹。

  \item \textbf{判别式黑名单 \texorpdfstring{$\mathcal{B}_q$}{Bq} 与全局置换对称：异常素数集的完全操作化与可检验频率预言。}
  在 \(q=9,10,11,13\) 上，可观测最小多项式 \(P_q\) 首次显露固定亏格的阶塌缩并保持不可约（命题 \ref{prop:pom-observable-minpoly-order-collapse-q9-13}），其整数判别式具有可审计的结构化分解（式 \eqref{eq:fold_collision_observable_minpoly_discriminants_q9_13}）。
  因而分歧素数黑名单
  \(\mathcal{B}_q=\{p:\ p\mid \mathrm{Disc}(P_q)\}\)
  可直接由 \(\mathrm{Disc}(P_q)\) 生成，并且对任何 \(p\notin\mathcal{B}_q\) 都严格排除局部重根/谱分歧（命题 \ref{prop:pom-observable-minpoly-discriminant-ramification-q9-13}），从而把“异常素数集”完全提升为有限可审计集合。
  更进一步地，在同一切片上 \(\mathrm{Gal}(P_q/\QQ)\) 达到最大对称群 \(S_d\)（命题 \ref{prop:pom-observable-minpoly-galois-sd-q9-13}），于是 Chebotarev 密度给出分裂型的严格频率律，尤其不可约密度为 \(1/d\)（推论 \ref{cor:pom-observable-minpoly-chebotarev-density-q9-13}）；并且 Frobenius 的奇偶性可由判别式平方类的 Kronecker 符号 \(\left(\frac{\Delta_q^{\mathrm{disc}}}{p}\right)\) 直接读出（推论 \ref{cor:pom-observable-minpoly-frob-parity-q9-13}）。

  \item \textbf{period bound 的统一公式与额外 \texorpdfstring{$p$}{p} 因子判据：Jordan 非半单性可无需完全因子分解检测。}
  设
  \(P_q(\lambda)\equiv \lambda^{a_0}\prod_j f_j(\lambda)^{e_j}\ (\mathrm{mod}\ p)\)，其中 \(f_j\neq \lambda\) 不可约且 \(d_j=\deg f_j\)，并令 \(\varepsilon_{q,p}=\ind\{\exists j:\ e_j\ge 2\}\)。
  令 \(\widetilde P_q(\lambda):=P_q(\lambda)/\lambda^{a_0}\) 表示非零谱部分。
  表 \ref{tab:fold_collision_charpoly_mod7_selection_q13_15}--\ref{tab:fold_collision_charpoly_mod13_fingerprint_q13} 的 \emph{period bound} 列采用统一粗界
  \[
  \mathrm{PerBound}_p(q):=p^{\varepsilon_{q,p}}\cdot \mathrm{lcm}_j(p^{d_j}-1),
  \]
  它只区分“是否出现平方因子”（即 \(\varepsilon_{q,p}\in\{0,1\}\)），并将该非半单性以额外因子 \(p\) 的形式纳入最终周期上界。
  其中 \(\varepsilon_{q,p}=1\) 与 \(\mathrm{repeat?}=1\) 等价，并且可通过 \(\gcd(\widetilde P_q,\partial_\lambda\widetilde P_q)\neq 1\) 直接读出，而无需对 \(P_q(\lambda)\ (\mathrm{mod}\ p)\) 完整因子分解。
  更精细地，Jordan 指数 \(e_j\) 的数量级会产生 \(p^{\lceil\log_p e_j\rceil}\) 的 \(p\)-进贡献；相应的精确周期上界见命题 \ref{prop:conclusion-finitefield-jordan-exponent-period-bound}。
\end{enumerate}

\begin{proposition}[Jordan 指数的 \(p\)-进贡献与最终周期上界]\label{prop:conclusion-finitefield-jordan-exponent-period-bound}
\leanverified{paper\_conclusion\_finitefield\_jordan\_exponent\_period\_bound}
设 \(p\) 为素数，\((a_m)_{m\ge 0}\) 为取值于 \(\FF_p\) 的序列。
设存在多项式 \(P(E)\in\FF_p[E]\) 与整数 \(m_0\) 使得对所有 \(m\ge m_0\) 有
\[
(P(E)a)_m=0,
\qquad
(Ea)_m:=a_{m+1}.
\]
将 \(P(E)\) 写成
\[
P(E)=E^{a_0}\prod_{j=1}^{J}Q_j(E)^{e_j},
\]
其中 \(Q_j\) 两两互素、首一不可约，且 \(e_j\ge 1\)。
令 \(d_j:=\deg Q_j\) 并记 \(e_{\max}:=\max_j e_j\)。
则序列 \(a\) 在后段为最终周期序列，其最终周期 \(\mathrm{per}(a)\) 满足整除关系
\[
\boxed{\;
\mathrm{per}(a)\ \mid\ \mathrm{lcm}_{1\le j\le J}\Bigl(p^{\lceil\log_p e_j\rceil}\,(p^{d_j}-1)\Bigr)\; }.
\]
从而亦有粗界
\[
\boxed{\;
\mathrm{per}(a)\ \mid\ p^{\lceil\log_p e_{\max}\rceil}\cdot \mathrm{lcm}_{1\le j\le J}(p^{d_j}-1).\;}
\]
\end{proposition}
\begin{proof}
对充分大的 \(m\)，递推由 companion 动力系统实现：取状态向量 \(v_m\in\FF_p^{\deg P-a_0}\)（由有限个连续项组成）使 \(v_{m+1}=Av_m\)，其中 \(A\) 为相应 companion 矩阵。
因 \(E^{a_0}\) 仅产生有限瞬态，可将其吸收进起始段并忽略于最终周期的估计中。
\par
把 \(\FF_p^{\deg P-a_0}\) 视为 \(\FF_p[E]\)-模（令 \(E\) 作用为 \(A\)），由主分解定理可将其分解为各 \(Q_j\)-初等因子主分量的直和；因此 \(A\) 的最终周期整除各主分量上 \(A\) 的最终周期的最小公倍数。
\par
固定某个 \(j\)。在分裂域 \(\FF_{p^{d_j}}\) 上，对应主分量分解为若干 Jordan 块的直和，每个 Jordan 块形如
\(
J=\lambda(I+N)
\)
其中 \(\lambda\in\FF_{p^{d_j}}^\times\) 且 \(N\) 为幂零矩阵，满足 \(N^{e_j}=0\)。
取 \(t_j:=\lceil\log_p e_j\rceil\) 则 \(p^{t_j}\ge e_j\)，从而 \(N^{p^{t_j}}=0\)。
在特征 \(p\) 下由二项式恒等式（Frobenius 取幂）得
\[
(I+N)^{p^{t_j}}=I+N^{p^{t_j}}=I,
\qquad\text{从而}\qquad
J^{p^{t_j}}=\lambda^{p^{t_j}}I.
\]
又 \(\lambda^{p^{d_j}-1}=1\)，故
\[
J^{p^{t_j}(p^{d_j}-1)}=(J^{p^{t_j}})^{p^{d_j}-1}=(\lambda^{p^{t_j}}I)^{p^{d_j}-1}=I.
\]
因此该主分量上 \(A\) 的周期整除 \(p^{t_j}(p^{d_j}-1)\)。
对 \(j\) 取最小公倍数得到第一条整除关系；第二条由 \(p^{\lceil\log_p e_j\rceil}\le p^{\lceil\log_p e_{\max}\rceil}\) 直接推出。证毕。
\end{proof}

\begin{corollary}[模 \(2\) 的 \texorpdfstring{$(E+1)^e$}{(E+1)^e} 周期界]\label{cor:conclusion-mod2-eplus1-power-period}
设 \((a_m)\) 为取值于 \(\FF_2\) 的序列，并存在整数 \(e\ge 1\) 使得对充分大的 \(m\) 有
\[
(E+1)^e a_m=0.
\]
则 \(a\) 在后段的最终周期整除 \(2^{\lceil\log_2 e\rceil}\)。
特别地，当 \(e=8\) 时，最终周期整除 \(8\)。
\end{corollary}
\begin{proof}
对 \(t=\lceil\log_2 e\rceil\) 有 \(e\le 2^{t}\)，因而 \((E+1)^{2^{t}}a=0\) 在后段成立。
在特征 \(2\) 下有 \((E+1)^{2^{t}}=E^{2^{t}}+1\)，从而 \(E^{2^{t}}a=a\) 在后段成立，得到周期整除 \(2^{t}\)。
最后一条取 \(e=8\) 得 \(t=3\)。证毕。
\end{proof}
\leanverified{paper\_conclusion\_mod2\_eplus1\_power\_period}

\begin{theorem}[共振窗奇偶动力学的 \(2\)-进幺正完备分解]\label{thm:conclusion-resonance-z2-unipotent-complete-parity}
在共振窗 \(q=9,\dots,17\) 上，取定理 \ref{thm:pom-resonance-hensel-splitting-z2} 的唯一 Hensel 分解
\[
P_q(\lambda)=U_q(\lambda)V_q(\lambda),
\qquad
U_q(\lambda)\equiv \lambda^{a_q}\pmod 2,
\qquad
V_q(\lambda)\equiv (\lambda+1)^{b_q}\pmod 2,
\]
以及相应的环分裂
\[
\ZZ_2[\lambda]/(P_q)
\cong
\ZZ_2[\lambda]/(U_q)\times \ZZ_2[\lambda]/(V_q).
\]
对任意满足 \(P_q(E)S=0\) 的 \(2\)-进序列 \(S\)，令
\[
S=S_{\mathrm{nil}}+S_{\mathrm{unip}}
\]
为推论 \ref{cor:pom-resonance-hensel-splitting-sequence-module} 的典范分裂，则模 \(2\) 的长期轨道在至多 \(a_q\) 步之后完全由 \(\overline{S}_{\mathrm{unip}}\) 决定，而与 \(\overline{S}_{\mathrm{nil}}\) 无关。更精确地，
\[
\overline{S}_{\mathrm{nil}}(m+a_q)=0
\qquad\text{对充分大的 }m,
\]
并且
\[
(E+1)^{b_q}\overline{S}_{\mathrm{unip}}=0.
\]
因此 \(\overline{S}_{\mathrm{unip}}\) 在后段必为次数 \(<b_q\) 的离散二项式多项式流，其最终周期满足
\[
\mathrm{per}\bigl(\overline{S}_{\mathrm{unip}}\bigr)\mid 2^{\lceil\log_2 b_q\rceil}.
\]
换言之，长期奇偶信息由 \(2\)-进幺正因子完备决定，而幂零因子仅贡献有限瞬态。
\end{theorem}
\begin{proof}
唯一 Hensel 分解与环分裂正是定理 \ref{thm:pom-resonance-hensel-splitting-z2} 的内容。对任意满足 \(P_q(E)S=0\) 的 \(2\)-进序列 \(S\)，推论 \ref{cor:pom-resonance-hensel-splitting-sequence-module} 给出典范直和分裂
\[
S=S_{\mathrm{nil}}+S_{\mathrm{unip}},
\qquad
U_q(E)S_{\mathrm{nil}}=0,
\qquad
V_q(E)S_{\mathrm{unip}}=0.
\]
将这两式模 \(2\) 约化。由于
\[
U_q(\lambda)\equiv \lambda^{a_q}\pmod 2,
\]
便得到
\[
E^{a_q}\,\overline{S}_{\mathrm{nil}}=0,
\]
故 \(\overline{S}_{\mathrm{nil}}(m+a_q)=0\) 在后段成立。于是总序列 \(\overline S\) 的长期轨道与 \(\overline{S}_{\mathrm{unip}}\) 完全一致。

另一方面，由
\[
V_q(\lambda)\equiv (\lambda+1)^{b_q}\pmod 2
\]
及 \(V_q(E)S_{\mathrm{unip}}=0\) 的模 \(2\) 约化，得到
\[
(E+1)^{b_q}\overline{S}_{\mathrm{unip}}=0.
\]
于是由命题 \ref{prop:conclusion-f2-binomial-basis-from-delta-nilpotent}，\(\overline{S}_{\mathrm{unip}}\) 在后段必为次数 \(<b_q\) 的离散二项式多项式流；再由推论 \ref{cor:conclusion-mod2-eplus1-power-period}，
\[
\mathrm{per}\bigl(\overline{S}_{\mathrm{unip}}\bigr)\mid 2^{\lceil\log_2 b_q\rceil}.
\]
故长期奇偶信息由幺正分量完备决定，而幂零分量只产生至多 \(a_q\) 步的有限瞬态。证毕。
\end{proof}
\leanverified{paper\_conclusion\_resonance\_z2\_unipotent\_complete\_parity}

\begin{theorem}[共振窗碰撞矩的 \texorpdfstring{$2$}{2}-进 Hensel 尾周期定理]\label{thm:conclusion-resonance-window-z2-hensel-tail-periodicity}
\leanverified{paper\_conclusion\_resonance\_window\_z2\_hensel\_tail\_periodicity}
沿用表 \ref{tab:fold_collision_charpoly_mod2_shadow_q9_17} 的指数对 \((a_q,b_q)\)。在定理 \ref{thm:pom-moment-resonance} 的递推生效区间内，对每个 \(q\in\{9,\dots,17\}\) 与每个整数 \(k\ge 1\)，序列
\[
S_q(m)\pmod{2^k}
\]
都在去掉至多 \(a_qk\) 项前导之后进入纯周期壳。更精确地，存在整数
\[
N_{q,k}\le a_qk,
\qquad
T_{q,k}\mid 2^{\,k+\lceil\log_2(b_qk)\rceil-1}
\]
使得对充分大的 \(m\) 都有
\[
S_q(m+N_{q,k}+T_{q,k})
\equiv
S_q(m+N_{q,k})
\pmod{2^k}.
\]
因此，共振窗的 dyadic 尾动力学在任意有限精度 \(2^k\) 上都由“有限步幂零消失 + 单一 \((-1)\)-扇区上的 unipotent 厚化”完全控制。
\end{theorem}
\begin{proof}
把整数值序列 \(S_q\) 视为 \(2\)-进序列，则 \(S_q\in \mathscr M_q^{(2)}\)。由推论 \ref{cor:pom-resonance-hensel-splitting-sequence-module}，
\[
S_q=S_{q,\mathrm{nil}}+S_{q,\mathrm{unip}},
\qquad
U_q(E)S_{q,\mathrm{nil}}=0,
\qquad
V_q(E)S_{q,\mathrm{unip}}=0.
\]

先看幂零部分。由
\[
U_q(\lambda)\equiv \lambda^{a_q}\pmod 2
\]
可写
\[
U_q(\lambda)=\lambda^{a_q}+2R_q(\lambda)
\qquad
\bigl(R_q(\lambda)\in \ZZ_2[\lambda]\bigr).
\]
于是
\[
E^{a_q}S_{q,\mathrm{nil}}
=
-2R_q(E)S_{q,\mathrm{nil}}
\in
2\,\mathscr M_q^{(2)}.
\]
归纳得到
\[
E^{a_qk}S_{q,\mathrm{nil}}\in 2^k\,\mathscr M_q^{(2)}.
\]
故幂零分量在模 \(2^k\) 下至多经过 \(a_qk\) 步便完全消失。

再看幺正部分。令
\[
D:=E+1.
\]
由
\[
V_q(\lambda)\equiv (\lambda+1)^{b_q}\pmod 2
\]
可写
\[
V_q(E)=D^{b_q}+2W_q(E)
\qquad
\bigl(W_q(E)\in \ZZ_2[E]\bigr).
\]
因此
\[
D^{b_q}S_{q,\mathrm{unip}}
=
-2W_q(E)S_{q,\mathrm{unip}}
\in
2\,\mathscr M_q^{(2)}.
\]
同样归纳得到
\[
D^{b_qk}S_{q,\mathrm{unip}}\in 2^k\,\mathscr M_q^{(2)}.
\]
故在模 \(2^k\) 下，若记
\[
N:=b_qk,
\qquad
s:=k+\lceil\log_2 N\rceil-1,
\]
则 \(D^N=0\) 且
\[
E^{2^s}=(-1+D)^{2^s}
=
1+\sum_{t=1}^{2^s}(-1)^t\binom{2^s}{t}D^t.
\]
对 \(1\le t<N\)，经典赋值公式给出
\[
v_2\!\binom{2^s}{t}=s-v_2(t)\ge k,
\]
因为 \(v_2(t)\le \lceil\log_2 N\rceil-1\)。而对 \(t\ge N\)，已有 \(D^t=0\) 模 \(2^k\)。故
\[
E^{2^s}S_{q,\mathrm{unip}}\equiv S_{q,\mathrm{unip}}\pmod{2^k},
\]
即幺正分量的周期整除 \(2^s\)。

将两部分合并：在去掉至多 \(a_qk\) 项前导之后，幂零分量消失，而剩余幺正分量以周期整除 \(2^s=2^{k+\lceil\log_2(b_qk)\rceil-1}\) 运行。证毕。
\end{proof}

\begin{corollary}[共振窗的统一 dyadic 周期壳与终端幺正深度极值点]\label{cor:conclusion-resonance-window-unified-dyadic-shell-extremal-endpoint}
\leanverified{paper\_conclusion\_resonance\_window\_unified\_dyadic\_shell\_extremal\_endpoint}
由表 \ref{tab:fold_collision_charpoly_mod2_shadow_q9_17} 读得
\[
(a_q,b_q)=
(3,4),(5,4),(3,6),(7,6),(3,8),(7,6),(3,8),(7,6),(3,10)
\]
分别对应 \(q=9,\dots,17\)。因此对整个共振窗统一有
\[
N_{q,k}\le 7k,
\qquad
T_{q,k}\mid 2^{\,k+\lceil\log_2(10k)\rceil-1}.
\]
并且 \(q=17\) 是当前审计窗口内唯一满足 \(b_q=10\) 的点；等价地，它是已核验高阶窗内唯一达到最大 dyadic 幺正厚度的终端点。
\end{corollary}
\begin{proof}
前两式由定理 \ref{thm:conclusion-resonance-window-z2-hensel-tail-periodicity} 与表 \ref{tab:fold_collision_charpoly_mod2_shadow_q9_17} 的逐项代入直接得到。最后一句只是指出该表中 \(b_q\) 的最大值 \(10\) 仅在 \(q=17\) 处出现。证毕。
\end{proof}

\begin{corollary}[模 \(2\) 影子差分阶数的台阶跃迁与碰撞矩奇偶动力学的可检验分级]\label{cor:conclusion-fold-collision-mod2-step-jump-q7}
\leanverified{paper\_conclusion\_fold\_collision\_mod2\_step\_jump\_q7}
在模 \(2\) 影子分解
\(
P_q(\lambda)\equiv \lambda^{a_q}(\lambda+1)^{b_q}\ (\mathrm{mod}\ 2)
\)
与湮灭式 \(E^{a_q}(E+1)^{b_q}S_q\equiv 0\ (\mathrm{mod}\ 2)\) 的设定下，
对每个固定 \(q\) 记 \(\Delta:=E+1\)。
则 \(S_q(m+a_q)\bmod 2\) 在后段满足 \(\Delta^{b_q}S_q(m+a_q)=0\)，从而由命题 \ref{prop:conclusion-f2-binomial-basis-from-delta-nilpotent} 得到二项式基多项式表示，其次数严格小于 \(b_q\)。
特别地，在已核验切片中 \(q\in\{4,5,6\}\) 对应 \(b_q=2\)，而从 \(q=7\) 起（在所给审计范围内）出现 \(b_q=4\)；
因此模 \(2\) 的差分幂零阶数在 \(q=7\) 处发生可复核的台阶跃迁：由最终仿射相跃迁到最终至多三次相。
\end{corollary}

\begin{conjecture}[模 \(2\) 塌缩指数的幂次塔：差分幂零深度的无界性]\label{conj:conclusion-fold-collision-mod2-power-tower}
令 \(b_q\) 为模 \(2\) 影子分解 \(P_q(\lambda)\equiv \lambda^{a_q}(\lambda+1)^{b_q}\ (\mathrm{mod}\ 2)\) 的指数。
猜想：
\begin{enumerate}
  \item \(b_q\) 仅取 \(2\) 的幂，并随 \(q\) 非减；
  \item 存在无穷增序列 \(q_1<q_2<\cdots\) 使 \(b_{q_{n+1}}=2b_{q_n}\)；
  \item 对每个 \(q\)，\(b_q\) 同时给出 \(S_q(m)\bmod 2\) 的最小差分幂零阶：在后段有 \(\Delta^{b_q}S_q=0\) 且 \(\Delta^{b_q/2}S_q\not\equiv 0\)。
\end{enumerate}
\end{conjecture}

\begin{theorem}[共振窗模 \(2\) 影子的低阶二项式基塌缩]\label{thm:conclusion-resonance-window-mod2-binomial-collapse}
\leanverified{paper\_conclusion\_resonance\_window\_mod2\_binomial\_collapse}
在共振窗 \(q=9,\dots,17\) 上，模 \(2\) 约化后的特征多项式满足
\[
P_q(\lambda)\equiv \lambda^{a_q}(\lambda+1)^{b_q}\pmod 2,
\]
因此对应碰撞矩奇偶序列满足
\[
E^{a_q}(E+1)^{b_q}S_q\equiv 0\pmod 2.
\]
记 \(\Delta:=E+1\)。则对充分大的 \(m\) 有
\[
\Delta^{b_q}S_q(m+a_q)\equiv 0\pmod 2.
\]
从而在 \(\FF_2\) 上存在唯一系数 \(c_{q,0},\dots,c_{q,b_q-1}\in\FF_2\)，使得
\[
S_q(m+a_q)\equiv \sum_{j=0}^{b_q-1}c_{q,j}\binom{m}{j}\pmod 2
\]
在后段成立；并且其最终周期整除
\[
T_q:=2^{\lceil\log_2 b_q\rceil}\le 16.
\]
更具体地，在该审计窗中仅出现
\[
T_q\mid 4,\qquad T_q\mid 8,\qquad T_q\mid 16
\]
三种周期级别。故整个 resonance window 的模 \(2\) 影子已塌缩为一族低阶二项式基多项式序列。
\end{theorem}

\begin{proof}
模 \(2\) 影子分解与湮灭关系由注 \ref{rem:pom-charpoly-mod2-shadow} 给出，而定理 \ref{thm:conclusion-resonance-z2-unipotent-complete-parity} 已将其进一步提升为
\[
\Delta^{b_q}S_q(m+a_q)\equiv 0\pmod 2
\]
的长期幂零律。
于是由命题 \ref{prop:conclusion-f2-binomial-basis-from-delta-nilpotent}，\(S_q(m+a_q)\bmod 2\) 在后段唯一表示为次数 \(<b_q\) 的二项式基多项式。
再由推论 \ref{cor:conclusion-mod2-eplus1-power-period}，最终周期整除 \(2^{\lceil\log_2 b_q\rceil}\)。
最后把表 \ref{tab:fold_collision_charpoly_mod2_shadow_q9_17} 中的 \(b_q\in\{4,6,8,10\}\) 逐项代入，即得 \(T_q\) 只可能落在 \(4,8,16\) 三种级别。证毕。
\end{proof}

\leanverified{paper\_conclusion\_mod2\_shadow\_insufficient\_for\_archimedean\_instability}
\begin{corollary}[固定模 \(2\) 影子不足以见证 archimedean 越线强度]\label{cor:conclusion-mod2-shadow-insufficient-for-archimedean-instability}
在同一共振窗内，存在不同的 \(q\) 具有完全相同的模 \(2\) 影子，但具有不同的 archimedean 越线指数 \(\delta_q=\log(\Lambda_q^2/r_q)\)。例如，
\[
P_{12}(\lambda)\equiv P_{14}(\lambda)\equiv P_{16}(\lambda)\equiv \lambda^7(\lambda+1)^6\pmod 2,
\]
\[
P_{13}(\lambda)\equiv P_{15}(\lambda)\equiv \lambda^3(\lambda+1)^8\pmod 2,
\]
而相应地
\[
\delta_{12}\approx 2.2776615,\qquad
\delta_{14}\approx 2.8248141,\qquad
\delta_{16}\approx 3.3661669,
\]
\[
\delta_{13}\approx 2.5520014,\qquad
\delta_{15}\approx 3.0961823.
\]
故任何仅经由 \(P_q(\lambda)\bmod 2\) 或其等价湮灭子类而因子化的审计器，都不可能恢复、排序或完备监测真正的 archimedean 越线强度。
\end{corollary}

\begin{proof}
表 \ref{tab:fold_collision_charpoly_mod2_shadow_q9_17} 直接给出上述模 \(2\) 影子重合关系；而表 \ref{tab:fold_collision_kernel_rh_scan_q2_17} 与表 \ref{tab:fold_collision_kernel_rh_scan_q18_23} 给出对应的 \(\delta_q\) 数值，且它们在同一 shadow 类中彼此不同。
因此映射
\[
q\longmapsto P_q(\lambda)\bmod 2
\]
不是关于 \(\delta_q\) 的充分统计量。若某个审计器完全经由该影子因子化，则对同一 shadow 类中的不同 \(q\) 必给出相同输出，因而不可能完备见证 archimedean 越线强度。证毕。
\end{proof}

\begin{theorem}[共振窗奇偶尾喷流的有限维完全性]\label{thm:conclusion-resonance-window-mod2-tail-jet-completeness}
\leanverified{paper\_conclusion\_resonance\_window\_mod2\_tail\_jet\_completeness}
对每个 \(q\in\{9,\dots,17\}\)，定义模 \(2\) 尾序列
\[
s_q(n):=S_q(n+a_q)\bmod 2,
\qquad
n\ge 0.
\]
再记
\[
\mathscr J_q:=
\{s:\NN\to\FF_2:\ \Delta^{b_q}s=0\},
\qquad
\Delta:=E+1.
\]
则尾喷流映射
\[
J_q:\mathscr J_q\longrightarrow \FF_2^{\,b_q},
\qquad
J_q(s):=(s(0),\dots,s(b_q-1))
\]
是线性同构。因而：
\begin{enumerate}
\normalfont
  \item 一切可实现的奇偶尾序列恰构成一个 \(b_q\)-维 \(\FF_2\)-线性空间；
  \item 任意 \(s_q\) 都由前 \(b_q\) 个值唯一决定；
  \item 存在一个至多含 \(2^{b_q}\) 个状态的确定性有限状态监测器，足以生成全部共振窗奇偶尾部动力学。
\end{enumerate}
\end{theorem}
\begin{proof}
由定理 \ref{thm:conclusion-resonance-window-mod2-binomial-collapse} 与命题 \ref{prop:conclusion-f2-binomial-basis-from-delta-nilpotent}，
\[
\mathscr J_q
=
\operatorname{span}_{\FF_2}
\left\{
\binom{m}{0},\dots,\binom{m}{b_q-1}
\right\},
\]
故 \(\dim_{\FF_2}\mathscr J_q=b_q\)。

再看评价矩阵
\[
\left(\binom{i}{j}\right)_{0\le i,j\le b_q-1}.
\]
它是下三角矩阵，且对角线上全为 \(1\)，故在 \(\FF_2\) 上可逆。因此前 \(b_q\) 个值与二项式基系数一一对应，\(J_q\) 遂为线性同构。三条结论随即成立。证毕。
\end{proof}

\begin{theorem}[共振窗奇偶动力学的三层 dyadic 室分解]\label{thm:conclusion-resonance-window-dyadic-chambers}
\leanverified{paper\_conclusion\_resonance\_window\_dyadic\_chambers}
定义共振窗的奇偶周期室
\[
\mathcal C_2:=\{9,10\},
\qquad
\mathcal C_3:=\{11,\dots,16\},
\qquad
\mathcal C_4:=\{17\}.
\]
则对每个 \(q\in\mathcal C_r\)，其尾序列 \(s_q\) 的最终周期都整除 \(2^r\)。因此，共振窗 \(q=9,\dots,17\) 的模 \(2\) 尾动力学精确分裂为三类 dyadic 室，而不存在第四种已审计周期级别。
\end{theorem}
\begin{proof}
定理 \ref{thm:conclusion-resonance-window-mod2-binomial-collapse} 已给出最终周期上界
\[
T_q\mid 4\quad(q=9,10),
\qquad
T_q\mid 8\quad(q=11,\dots,16),
\qquad
T_q\mid 16\quad(q=17).
\]
这恰对应三类 dyadic 周期室 \(\mathcal C_2,\mathcal C_3,\mathcal C_4\)。由于该窗口内所有 \(q\) 已被这三类穷尽，故不存在第四种已审计周期级别。证毕。
\end{proof}

\begin{theorem}[模 \(2\) 平方自由化的单根坍缩]\label{thm:conclusion-resonance-window-mod2-single-root-sqfree-collapse}
对每个 \(q\in\{9,\dots,17\}\)，设
\[
P_q^{(2)}(\lambda):=P_q(\lambda)\bmod 2
=
\lambda^{a_q}(\lambda+1)^{b_q},
\]
其中在当前共振窗审计中 \(a_q\) 为奇数而 \(b_q\) 为偶数。则有
\[
\gcd\!\bigl(P_q^{(2)},(P_q^{(2)})'\bigr)
=
\lambda^{a_q-1}(\lambda+1)^{b_q},
\]
从而平方自由化满足
\[
\operatorname{sqfree}\!\bigl(P_q^{(2)}\bigr)=\lambda.
\]
因此，在该窗口中：
\begin{enumerate}
\normalfont
  \item 唯一的简单谱点是 \(\lambda=0\)；
  \item 点 \(\lambda=1\) 只以纯非半单方式出现；
  \item 模 \(2\) 光谱的全部平方自由信息退化为单个一次因子。
\end{enumerate}
\end{theorem}
\begin{proof}
在特征 \(2\) 下，
\[
(P_q^{(2)})'(\lambda)
=
a_q\lambda^{a_q-1}(\lambda+1)^{b_q}
+
b_q\lambda^{a_q}(\lambda+1)^{b_q-1}.
\]
由于 \(a_q\) 为奇数而 \(b_q\) 为偶数，遂在 \(\FF_2\) 中化为
\[
(P_q^{(2)})'(\lambda)=\lambda^{a_q-1}(\lambda+1)^{b_q}.
\]
故
\[
\gcd\!\bigl(P_q^{(2)},(P_q^{(2)})'\bigr)
=
\lambda^{a_q-1}(\lambda+1)^{b_q}.
\]
再将 \(P_q^{(2)}\) 除以此最大公因子，便得
\[
\operatorname{sqfree}\!\bigl(P_q^{(2)}\bigr)=\lambda.
\]
三条结论立即随之成立。证毕。
\end{proof}

\begin{theorem}[最大模 \(2\) 共根度与判别式强制整除]\label{thm:conclusion-resonance-window-mod2-maximal-coroot-discriminant}
对每个 \(q\in\{9,\dots,17\}\)，都有
\[
\deg \gcd\!\bigl(P_q^{(2)},(P_q^{(2)})'\bigr)=d_q-1,
\qquad
2^{d_q-1}\mid \Disc(P_q).
\]
其中第一条在所有次数为 \(d_q\) 且导数不恒零的首一模 \(2\) 多项式中，已经达到最大可能值。
\end{theorem}
\begin{proof}
由定理 \ref{thm:conclusion-resonance-window-mod2-single-root-sqfree-collapse}，
\[
\deg \gcd\!\bigl(P_q^{(2)},(P_q^{(2)})'\bigr)
=
(a_q-1)+b_q
=
d_q-1.
\]
另一方面，对任一次数为 \(d_q\) 的首一多项式 \(f\) 且 \(f'\neq 0\)，恒有
\[
\deg \gcd(f,f')\le \deg f'=d_q-1.
\]
故这里确已达到最大允许共根度。

至于判别式整除，主稿对共振窗的 Sylvester 结果式核维已给出
\[
2^{d_q-1}\mid \Disc(P_q).
\]
证毕。
\end{proof}

\begin{theorem}[共振窗 Hankel-null 模的全尺度主理想化]\label{thm:conclusion-resonance-window-hankel-null-principalization}
对每个 \(q\in\{9,\dots,17\}\)，记
\[
\operatorname{NullModes}^{(q)}(m):=\ker_{\ZZ}\!\bigl((H^{(m)})^\top\bigr),
\]
\[
M^{(m)}(P_q):=
\left\{
\operatorname{coeff}_{<m}\!\bigl(Q(\lambda)P_q(\lambda)\bigr):
Q\in\ZZ[\lambda],\ \deg Q<m-d_q
\right\}\subset \ZZ^m.
\]
则对一切 \(m\ge d_q+1\) 都有严格恒等式
\[
\operatorname{NullModes}^{(q)}(m)=M^{(m)}(P_q),
\qquad
\operatorname{rank}\operatorname{NullModes}^{(q)}(m)=m-d_q.
\]
特别地，在已审计 primitive 口径下，
\[
\Delta_{\mathrm{prim}}(q)=0
\qquad
(q=9,\dots,17),
\]
故不存在任何非倍乘型的 primitive null mode。
\end{theorem}
\begin{proof}
主稿对共振窗已审计给出：对任意 \(m\ge d_q+1\)，全尺度 Hankel-null 模满足
\[
\operatorname{NullModes}^{(q)}(m)=M^{(m)}(P_q),
\qquad
\operatorname{rank}\operatorname{NullModes}^{(q)}(m)=m-d_q.
\]
同时，同一审计链又给出
\[
\Delta_{\mathrm{prim}}(q)=0
\qquad
(q=9,\dots,17).
\]
这意味着 primitive quotient
\[
\operatorname{NullModes}^{(q)}(m)\big/M^{(m)}(P_q)
\]
在当前窗口内恒为零，故每一个 null certificate 都已经是 \(P_q\) 的倍乘产物。证毕。
\end{proof}

\begin{theorem}[奇素数不可约性与模 \(2\) 贫乏性的正交分离]\label{thm:conclusion-resonance-window-oddprime-fulldegree-vs-mod2-poverty}
对每个 \(q\in\{9,\dots,17\}\)，存在素数 \(p_{\mathrm{irr}}(q)\) 使得
\[
P_q(\lambda)\bmod p_{\mathrm{irr}}(q)
\]
在 \(\FF_{p_{\mathrm{irr}}(q)}[\lambda]\) 中不可约，且次数恰为 \(d_q\)。在当前审计窗中，可以取
\[
p_{\mathrm{irr}}(q)=11,17,37,29,29,37,17,239,31
\]
分别对应
\[
q=9,10,11,12,13,14,15,16,17.
\]
然而同一族多项式在模 \(2\) 下却同时满足：
\begin{enumerate}
\normalfont
  \item 平方自由化恒为 \(\lambda\)；
  \item 尾奇偶动力学只分裂成三种 dyadic 周期室；
  \item 所有 Hankel-null 模均由 \(P_q\) 的倍乘主理想完全生成。
\end{enumerate}
因此，共振窗中存在严格的双重分离：奇素数方向保持 full-degree 不可约性，而模 \(2\) 方向同时呈现 phase-poor、sqfree-poor 与 principal-ideal 贫乏性。
\end{theorem}
\begin{proof}
主稿的奇素数审计表明：对
\[
q=9,10,11,12,13,14,15,16,17
\]
分别存在
\[
11,17,37,29,29,37,17,239,31
\]
使得 \(P_q(\lambda)\bmod p_{\mathrm{irr}}(q)\) 不可约，因此保留完整次数 \(d_q\)。

另一方面，定理 \ref{thm:conclusion-resonance-window-mod2-single-root-sqfree-collapse}、定理 \ref{thm:conclusion-resonance-window-dyadic-chambers} 与定理 \ref{thm:conclusion-resonance-window-hankel-null-principalization} 已分别给出：模 \(2\) 世界中平方自由化仅剩 \(\lambda\)，尾动力学只剩三类 dyadic 周期室，而全部 Hankel-null 证书都由单一主理想生成。故奇素数方向的 degree-full 不可约性与模 \(2\) 方向的结构贫乏之间不存在单调投影关系，它们是彼此正交的两种代数复杂度。证毕。
\end{proof}

\begin{conjecture}[两条实谱核与无界 dyadic 相位塔]\label{conj:conclusion-resonance-window-two-real-cores-unbounded-dyadic-tower}
结合主稿关于实谱核稳定分裂的猜想与当前已审计的模 \(2\) 台阶律，设应存在阈值 \(q_0\) 与两常数
\[
d_{\mathrm{even}},d_{\mathrm{odd}}\in\{15,17\},
\]
使得对充分大的 \(n\) 有
\[
d_{2n}=d_{\mathrm{even}},
\qquad
d_{2n+1}=d_{\mathrm{odd}},
\]
并且相应的实谱 Perron 根分别稳定落在两个固定全实数域中；与此同时，模 \(2\) 影子的 \((\lambda+1)\)-重数 \(b_q\) 形成一个无界的 dyadic 台阶塔。等价地，长程共振塔在特征 \(0\) 上仅剩两条有限次数实谱核，而在特征 \(2\) 上却经历无限多个 dyadic 相位跳迁。
\end{conjecture}

\noindent\emph{证据.}
当前共振窗审计已经表明：一方面，模 \(2\) 影子在 \(q=9,\dots,17\) 上只剩有限喷流、单根平方自由化与三层 dyadic 周期室；另一方面，奇素数方向仍保持 full-degree 不可约性，而更高阶的实谱审计则已经出现奇偶两支 \(15/17\) 的准稳定分裂。由此看，特征 \(0\) 上的长程对象更像两条稳定实谱核，而特征 \(2\) 上的长程对象则更像一座持续增高的 dyadic 相位塔。


\begin{theorem}[秩缺口与 characteristic-\texorpdfstring{$2$}{2} 终端记忆深度的独立性]\label{thm:conclusion-resonance-window-gap-vs-mod2-terminal-memory-independent}
在共振窗 \(q=9,\dots,17\) 中，设
\[
\Delta(q):=\bigl(2\lfloor q/2\rfloor+1\bigr)-d_q
\]
为秩缺口，并记
\[
\operatorname{Per}_2^{\max}(q)
\]
为定理 \ref{thm:conclusion-resonance-window-dyadic-chambers} 给出的模 \(2\) 最终周期上界。则不存在函数
\[
F:\ZZ_{\ge 0}\to \ZZ_{\ge 1}
\]
使得
\[
\operatorname{Per}_2^{\max}(q)=F(\Delta(q))
\qquad (q=9,\dots,17).
\]
换言之，\(\Delta(q)\) 与 characteristic-\(2\) 的终端记忆深度是相互独立的两条终端轴。
\leanverified{resonance\_window\_gap\_not\_terminal\_type\_function}
\leanverified{resonance\_window\_delta2\_witnesses}
\leanverified{resonance\_window\_five\_terminal\_types\_distinct}
\leanverified{paper\_resonance\_window\_terminal\_extended}
\end{theorem}

\begin{proof}
表 \ref{tab:fold_collision_resonance_gap_delta_q_9_17} 给出
\[
\Delta(9)=\Delta(10)=\Delta(11)=\Delta(13)=\Delta(14)=2,
\qquad
\Delta(12)=0,
\]
\[
\Delta(15)=\Delta(16)=\Delta(17)=4.
\]
另一方面，定理 \ref{thm:conclusion-resonance-window-dyadic-chambers} 给出
\[
\operatorname{Per}_2^{\max}(q)\mid 4 \quad (q=9,10),
\]
\[
\operatorname{Per}_2^{\max}(q)\mid 8 \quad (q=11,\dots,16),
\]
\[
\operatorname{Per}_2^{\max}(q)\mid 16 \quad (q=17).
\]
因此同一缺口层 \(\Delta=2\) 内同时出现周期级别 \(4\) 与 \(8\)，而同一缺口层 \(\Delta=4\) 内同时出现 \(8\) 与 \(16\)。故 \(\operatorname{Per}_2^{\max}(q)\) 不可能仅由 \(\Delta(q)\) 决定。证毕。
\end{proof}

\begin{corollary}[共振窗至少具有五种有限特征终端型]\label{cor:conclusion-resonance-window-five-terminal-types}
若以有序对
\[
\bigl(\Delta(q),\operatorname{Per}_2^{\max}(q)\bigr)
\]
作为共振窗 \(q=9,\dots,17\) 的有限特征终端型，则至少出现下列五种互异类型：
\[
(2,4),\qquad
(2,8),\qquad
(0,8),\qquad
(4,8),\qquad
(4,16).
\]
更具体地，
\[
q=9,10 \leadsto (2,4),\qquad
q=11,13,14 \leadsto (2,8),\qquad
q=12 \leadsto (0,8),
\]
\[
q=15,16 \leadsto (4,8),\qquad
q=17 \leadsto (4,16).
\]
\leanverified{resonance\_window\_five\_types\_pairwise\_distinct}
\end{corollary}

\begin{proof}
这只是定理 \ref{thm:conclusion-resonance-window-gap-vs-mod2-terminal-memory-independent} 证明中两张已审计表格的逐项合并。五种不同的有序对均在共振窗内实际出现，故至少有五种互异的有限特征终端型。证毕。
\end{proof}

\begin{theorem}[mixed-residue 终端尾在模 \texorpdfstring{$6$}{6} 下的 \texorpdfstring{$72$}{72} 周期壳]\label{thm:conclusion-resonance-window-q13-q15-mod6-period72}
对 \(q\in\{13,15\}\)，记
\[
\overline S_q(m):=S_q(m)\bmod 6.
\]
则 \(\overline S_q(m)\) 在充分大的 \(m\) 上最终周期满足
\[
\operatorname{Per}_6(\overline S_q)\mid 72.
\]
更精确地，
\[
\operatorname{Per}_2(\overline S_q)\mid 8,
\qquad
\operatorname{Per}_3(\overline S_q)\mid 18,
\qquad
\operatorname{Per}_6(\overline S_q)\mid \operatorname{lcm}(8,18)=72.
\]
\leanverified{resonance\_window\_lcm\_8\_18\_eq\_72}
\leanverified{resonance\_window\_eight\_dvd\_72}
\leanverified{resonance\_window\_eighteen\_dvd\_72}
\leanverified{resonance\_window\_lcm\_universal}
\leanverified{paper\_resonance\_window\_mod6\_period\_witness}
\end{theorem}

\begin{proof}
对 \(q\in\{13,15\}\)，定理 \ref{thm:conclusion-resonance-window-dyadic-chambers} 已给出模 \(2\) 终端周期整除 \(8\)。另一方面，推论 \ref{cor:pom-mod6-period-collapse-72-q13-15} 证明了模 \(3\) 的两相尾在最终阶段周期整除 \(18\)。由中国剩余分解
\[
\ZZ/6\ZZ\cong \ZZ/2\ZZ\times \ZZ/3\ZZ,
\]
模 \(6\) 终端尾的最终周期遂整除
\[
\operatorname{lcm}(8,18)=72.
\]
证毕。
\end{proof}

\begin{conjecture}[有限特征终端指纹的完整性]\label{conj:conclusion-resonance-window-finite-characteristic-terminal-fingerprint}
在共振窗 \(q=9,\dots,17\) 内，定义有限特征终端指纹
\[
\mathfrak F_q:=
\bigl(
P_q,\,
\mathfrak S_q^{\mathrm{dist}},\,
\Pi_q^{(2)},\,
\Pi_q^{(3)}
\bigr),
\]
其中
\[
P_q=\text{首一最小递推特征多项式},
\qquad
\mathfrak S_q^{\mathrm{dist}}=\text{推论 \ref{cor:conclusion-hankel-good-prime-rigidity-finite-distortion-support} 中的有限失真素数集},
\]
\[
\Pi_q^{(2)}=\text{定理 \ref{thm:conclusion-resonance-window-dyadic-chambers} 所给的 characteristic-}2\text{ 终端类},
\qquad
\Pi_q^{(3)}=
\begin{cases}
\text{定理 \ref{thm:conclusion-resonance-window-q13-q15-mod6-period72} 所对应的 characteristic-}3\text{ 两相终端类}, & q\in\{13,15\},\\[1mm]
\varnothing, & \text{其余 }q.
\end{cases}
\]
则该指纹应足以决定 \(q\)-moment 的有理 Hankel 等价类。

换言之，共振窗中的有限特征信息应恰好分裂为三层彼此正交的对象：有理主骨架 \(P_q\)、有限素数失真账本 \(\mathfrak S_q^{\mathrm{dist}}\)、以及终端局部相 \((\Pi_q^{(2)},\Pi_q^{(3)})\)。
\end{conjecture}


\begin{conjecture}[固定特征影子与 archimedean 越线的双层分裂]\label{conj:conclusion-fixed-characteristic-shadow-archimedean-splitting}
对每个固定素数 \(p\)，记 \(P_q(\lambda)\bmod p\) 为第 \(q\) 阶碰撞核的有限特征影子，而
\[
\delta_q=\log\!\left(\frac{\Lambda_q^2}{r_q}\right),
\qquad
\beta_q^{\mathrm{RH}}=\frac{\log\Lambda_q}{\log r_q}
\]
为 archimedean 越线参数。则应存在如下普遍二分：
\begin{enumerate}
  \item 对固定 \(p\)，\(q\mapsto P_q(\lambda)\bmod p\) 的长期复杂度只在有限个有限状态扇区间跳变；等价地，其诱导的 \(p\)-进差分动力学始终可由有界阶的离散多项式或幂零--幺正规约描述。
  \item 独立于此，\(\delta_q\) 在高阶方向继续上升，而 \(\beta_q^{\mathrm{RH}}\) 逼近 \(1\)；真正的灾变只发生在实谱的二次竞争指数中，而不发生在固定特征影子的语法复杂度中。
\end{enumerate}
本段的模 \(2\) 共振窗现象只给出了该猜想的第一个可审计实例：有限特征世界保持有限状态，而 archimedean 世界已经进入单调放大的越线相。
\end{conjecture}

\paragraph{视界纤维容量与 Rényi-\texorpdfstring{$4$}{4} 碰撞谱的临界贴线：面积律饱和与平方根稳定阈值的有限维桥接}
在本文的外部化/视界投影语义中，黑洞熵被严格放置为外部接口约束下典型纤维的对数容量：在固定外部代数与边界几何（面积 \(A\)）下，\(\pi_{\mathrm{ext}}\) 的典型纤维对数规模存在最大允许上界，且在 Planck 单位下由
\(
S_{\mathrm{BH}}=A/4
\)
给出规范化的饱和界（命题 \ref{prop:dephys-horizon-fiber-capacity-area-law}；式 \eqref{eq:dephys-bh-entropy-area-quarter}）。
在算子代数口径下，该容量由有限指数条件期望的对数 \(\log\mathrm{Ind}(E)\) 唯一锁定：它等于外部化可造成的最坏相对熵损失，并且仅依赖外部接口与嵌入结构而与具体态无关（命题 \ref{prop:op-algebra-dmax-capacity-equals-log-index}；亦见命题 \ref{prop:fold-condexp-index-maxfiber} 的折叠离散实现）。
\par
与之并行，碰撞谱的 \(q\)-阶 moment-kernel 为同一纤维结构在 Rényi 方向上的可计算探针：以 \(q=4\) 为例，
\(
S_4(m)=\sum_{x\in X_m} d_m(x)^4
\)
同时等价于“四条独立均匀窗口折叠后完全碰撞”的概率 \(S_4(m)/16^m\)，并由五态非负整数核 \(A_4\) 精确实现，其指数增长率由 Perron 根 \(r_4=\rho(A_4)\) 控制（命题 \ref{prop:pom-s4-recurrence}）。
在核版 RH 的平方根稳定界 \(\Lambda\le \sqrt{\rho}\)（定义 \ref{def:kernel-rh}）口径下，\(A_4\) 首次出现贴线级越线：\(\Lambda_4>\sqrt{r_4}\) 且越线量级为 \(10^{-3}\)，并由稳定的负实慢模态 \(\mu_2<0\) 主导（注 \ref{rem:pom-collision-rh-break-a4-spectrum}）。
该负实次谱模态不仅控制越线量级，也在离散轨道计数侧导出一个可证书化的二步交替指纹：命题 \ref{prop:pom-a4-trace-primitive-two-term} 给出
\(\Tr(A_4^n)=r_4^n+(-\Lambda_4)^n+O(\rho_4^n)\)，从而 primitive 轨道数的偶/奇长差异在指数尺度上由 \(\Lambda_4^n/n\) 主导并符号严格交替；该二周期慢模可视为 \(q=4\) 贴线破缺的最小离散序参量。
进一步把 \(A_4\) 的多重边权视为连续参数后，该贴线破缺可被组织为一条 Newman 型阈值并由 resultant 消元给出完全代数证书（注 \ref{rem:pom-collision-rh-break-a4-threshold}），其阈值语义与“视界作为扰乱器/编码器所诱导的信息回收阈值”在制度解释上同型（注 \ref{rem:dephys-horizon-scrambler-recovery-threshold}）。
在同一核族上，偶谱五次曲线 \(E_t(u)=\det(I-uA_4(t)^2)=0\) 给出一个最短的偶扇区闭式入口；其 \(u\)-判别式分解排除 \(t\ge 0\)（除孤立点 \(t=1\) 外）上的偶谱退化，并将阈值证书 \(P_4(t)\) 重新刻画为 Perron 曲线与偶谱嵌入曲线的交点消元结果式（命题 \ref{prop:pom-a4t-even-zeta-quintic} 与命题 \ref{prop:pom-a4t-newman-resultant-curve-intersection}）。
并且在判别式分辨率上，四阶核自身的特征多项式诱导伴随虚二次域 \(K_4^{\mathrm{disc}}=\QQ(\sqrt{-985219})\)，其类数为 \(195\)，从而把 near-RH 的 \(q=4\) 贴线相与一条可审计的常数级算术不变量链条并置（命题 \ref{prop:pom-a4-disc-quadratic-class-number}）。
\par
因此，\(q=4\) 的 near-RH 破缺可被理解为：当以四阶 Rényi 碰撞探针审计外部化纤维结构时，系统第一次在谱侧显影出临界慢模通道，从而为“面积律饱和的纤维容量”与“临界稳定阈值贴线”之间提供一条可审计的有限维桥接。
\par
在同一复制统计的自由能记号下，令 \(F_q:=\log r_q\)、\(G_q:=\log\Lambda_q\) 并记
\(\delta_q:=2G_q-F_q=\log(\Lambda_q^2/r_q)\)（见注 \ref{rem:pom-collision-kernel-rh-phase}；亦见命题 \ref{prop:pom-deltaq-spectral-dimension-identity} 与命题 \ref{prop:pom-deltaq-quadratic-observable-max-exponent}）。
若令折叠推前概率为 \(w_m(x)=d_m(x)/2^m\)，则
\(
S_q(m)/2^{qm}=\sum_{x\in X_m}w_m(x)^q
\)
给出一个完全离散、可审计的整数复制配分函数实现。
在该口径中，\(F_q\) 给出“断开”主贡献的作用密度，而由次谱模态生成的连通二点修正在指数尺度上以 \(\exp(m\cdot 2G_q)\) 计量；因此相对断开主项 \(\exp(mF_q)\) 的竞争强度由 \(\exp(m\delta_q)\) 控制。
由表 \ref{tab:fold_collision_kernel_rh_scan_q2_8} 的符号变化 \(\delta_2<0,\delta_3<0,\delta_4>0\) 可将 \(q_c=4\) 识别为 Fold 碰撞谱在复制阶方向上的最小越线点。
同一表亦给出 \(\beta_q:=\log\Lambda_q/\log r_q\) 在 \(q=4\) 首次满足 \(\beta_q>\tfrac12\)，从而“平方根压制”在该阶首次被允许在指数层面被慢模通道突破；该首次穿越可视为一类半熵临界。
在黑洞 Rényi 熵的类比中，上述越线可被解释为连通鞍点（复制虫洞；replica wormhole）与断开鞍点在作用密度上的首次等量竞争；并且该竞争的判据完全落入有限维可计算谱不变量 \(\delta_q\)，而不依赖任何外部连续极限或热流极限假设。
推论 \ref{cor:pom-deltaq-sample-complexity-threshold} 进一步在可采样口径下将该越线指数读作最小样本指数阈值。
\par
最后，定义临界可放大性常数 \(\kappa_4:=|\delta_4|^{-1}\)。由于 \(\delta_4\) 为 \(10^{-3}\) 量级（表 \ref{tab:fold_collision_kernel_rh_scan_q2_8}），\(\kappa_4\) 落在 \(10^2\) 量级并给出“复制演化深度”方向上相对权重发生量级跃迁的最小尺度：当 \(m\asymp \kappa_4\) 时，连通/断开通道的相对权重在指数意义下发生量级跃迁；并且对任意谱微扰 \(\delta_4\mapsto \delta_4+\Delta\)，该相对权重的乘性漂移精确为 \(\exp(m\Delta)\)，从而 near-RH 贴线区在深度方向的结构稳定半径被强制压缩到 \(O(1/m)\) 级别。该现象把“谱贴线”从静态不等式提升为可审计的动态放大机制：临界区并非仅是“接近”，而呈现“高增益”响应。
另一方面，上述由 \(\kappa_4\) 所度量的谱侧临界深度与 Chernoff 分位证书的最优矩阶断点并不等价：在 \(\varepsilon=10^{-6}\) 下，表 \ref{tab:fold_collision_qstar_breakpoints_eps1e6_q2_20} 给出 \(4\leftrightarrow 5\) 的断点 \(m\approx 1241.89\)，而 \(3\leftrightarrow 4\) 的断点为 \(m\approx 2128.30\)（注 \ref{rem:pom-qstar-phase-diagram-breakpoints}）。因此在可审计参数下出现严格错位：\(q=4\) 的谱贴线可在 \(m\sim\kappa_4\) 的深度尺度上被放大，而在“最优尾界/扫描策略”口径下 \(q=4\) 成为最优的深度却显著后移；该错位表明谱相变与策略相变由两套不可约的临界结构分别控制，并可在同一有限窗口上被直接核验。


\paragraph{零温端分支常数、完成化 Frobenius 塔与端点层析的有限证书化接口}
除去本文各处已强调的“主尺度（Perron 根/压力）决定指数增长”之外，若进一步以代数分支的端点展开、完成化同余塔与端点层析反演为锚，则若干看似解析性的常数/稳定性结论可被统一压缩为有限维、整系数、可审计的证书链：
\begin{enumerate}
  \item \textbf{arity-charge 零温端的 Puiseux 精细展开把主项提升到常数项证书。}
  对 \(P_7(\Lambda,q)=0\) 的最大正实根分支，零温端展开不仅给出 \(\lambda(q)\sim \sqrt{3q}\)，其常数项与更高阶校正均由代数曲线在无穷远处的分支结构强制锁定（命题 \ref{prop:real-input-40-arity-charge-zero-temp-expansion}）。
  同一代数方程的次正实增长分支亦具有可审计的三项闭式（命题 \ref{prop:real-input-40-arity-charge-subbranch-puiseux}），从而把“统一慢模基准”从比例极限细化到可比较的端点喷射指纹。

  \item \textbf{零荷端的解析起跳斜率提供一个无拟合的微扰强度常数。}
  由 \(P_7(\Lambda,0)\) 的显式因子化与隐函数定理，\(q\mapsto\lambda(q)\) 在 \(0\) 邻域解析且 \(\lambda'(0)\) 具有闭式代数表达（命题 \ref{prop:real-input-40-arity-charge-derivative-q0}），从而“从零荷到正荷”的谱响应速度被压缩为可离线核验的一枚标量证书。

  \item \textbf{完成化 Dwork 同余在 Chebyshev 形式下闭合，并在 \texorpdfstring{$\bmod p$}{mod p} 上塌缩为纯 Frobenius 迭代。}
  completed ghost 的 \(p^k\)-级同余塔以 \(S\mapsto C_p(S)\) 的 Chebyshev 形式严格闭合（定理 \ref{thm:discussion-completed-dwork-chebyshev}），并进一步在 \(\bmod p\) 层面塌缩为 \(\widehat a_{p^k}(S)\equiv S^{p^k}\) 的 Frobenius 刚性（推论 \ref{cor:completed-dwork-frobenius-tower}）。

  \item \textbf{端点剖面是 C-finite 的：Hankel 秩与 Prony 反演把离线零点参数化为有限样本证书。}
  在有限 Blaschke 口径下，端点径向吸收剖面在 \(q\)-坐标中为有理函数并给出极点对的显式位置（定理 \ref{thm:xi-endpoint-absorption-inversion-uniqueness}），从而其 Taylor 系数满足常系数递推并诱导 Hankel 秩证书（推论 \ref{cor:xi-endpoint-profile-cfinite-hankel-rank}）；进一步，有限个系数即可在原则上通过 Prony/核空间法恢复全部 \((\gamma_k,\delta_k)\)（推论 \ref{cor:xi-endpoint-profile-prony-invertibility}）。
\end{enumerate}


\paragraph{有限分辨率可见层的布尔化、真实规范不变量与逆极限异常障碍}
把附录 \ref{subsec:op-algebra-fold-quantum-channel} 的折叠量子信道、\S\ref{subsec:pom-pw-tannaka-krein-reconstruction} 的 profinite 可见对称、以及 \S\ref{subsec:cdim-hardy-finite-mode-projection} 的 Hardy 有限模态投影拼接后，可把本项目中的“有限分辨率量子性”压缩为一条更细的分层：真实 fold-gauge 不变量先裂解为一个可见相干矩阵扇区与一组暗标量账本；在此之后，折叠可见信道才把该相干扇区压到布尔对角。量子资源仍由纤维多重度场精确控制，而任何真正超出有限层 gauge-commutant 的全局非经典性则只能滞留在 inverse-limit 异常残差中。

\begin{theorem}[有限分辨率可见层的布尔化与量子性的残差外置]\label{thm:conclusion-fold-quantum-visible-booleanization}
\leanverified{paper\_conclusion\_fold\_quantum\_visible\_booleanization}
固定分辨率 \(m\ge 1\)，令 \(\mathcal H_m=\ell^2(\Omega_m)\)，并记
\[
\mathcal H_m=\bigoplus_{x\in X_m}\mathcal H_x,
\qquad
\mathcal H_x:=\Pi_x\mathcal H_m,
\qquad
d_m(x):=\dim \mathcal H_x.
\]
对折叠量子信道
\[
\mathcal E_m(A)=\sum_{x\in X_m}\frac{\Tr(\Pi_xA)}{d_m(x)}\,\Pi_x
\qquad (A\in\mathcal B(\mathcal H_m))
\]
定义可见代数
\[
\mathcal Z_m:=\bigoplus_{x\in X_m}\CC\,\Pi_x\ \cong\ C(X_m).
\]
则：
\begin{enumerate}
  \item \(\mathrm{Im}(\mathcal E_m)=\mathcal Z_m\)
  \item \(\mathcal E_m\) 的全部可见尖事件恰为
  \[
  \Pi_S:=\sum_{x\in S}\Pi_x
  \qquad (S\subseteq X_m)
  \]
  \item 因而有限分辨率可见层的尖事件格规范同构于幂集 \(\mathcal P(X_m)\)，特别地为布尔格。
\end{enumerate}
因此，有限分辨率读出中不存在内生的非布尔量子逻辑；所有非交换性只能居留于被 \(\mathcal E_m\) 抹去的残差/扩张层。
\end{theorem}
\begin{proof}
对任意 \(A\in\mathcal B(\mathcal H_m)\)，由定义立即有
\[
\mathcal E_m(A)=\sum_{x\in X_m}\frac{\Tr(\Pi_xA)}{d_m(x)}\,\Pi_x\in \mathcal Z_m,
\]
故 \(\mathrm{Im}(\mathcal E_m)\subseteq \mathcal Z_m\)。
反之若 \(A=\sum_x a_x\Pi_x\in\mathcal Z_m\)，则
\[
\Tr(\Pi_xA)=a_x\,\Tr(\Pi_x)=a_x\,d_m(x),
\]
从而 \(\mathcal E_m(A)=A\)，故 \(\mathcal Z_m\subseteq \mathrm{Im}(\mathcal E_m)\)。
于是 \(\mathrm{Im}(\mathcal E_m)=\mathcal Z_m\)。
第二条正是命题 \ref{prop:op-algebra-fold-quantum-channel-projection-01-rigidity} 的内容；第三条随之立即成立。证毕。
\end{proof}

\begin{proposition}[射影纯态的规范商恰为经典概率单纯形]\label{prop:conclusion-fold-quantum-projective-simplex-quotient}
\leanverified{paper\_conclusion\_fold\_quantum\_projective\_simplex\_quotient}
令
\[
U_m:=\prod_{x\in X_m}U(\mathcal H_x)
\]
按块对角方式作用于 \(\mathbb P(\mathcal H_m)\)。
则存在规范同构
\[
\mathbb P(\mathcal H_m)/U_m\ \cong\ \Delta(X_m),
\qquad
[\psi]\longmapsto p_\psi,
\qquad
p_\psi(x):=\frac{\|\Pi_x\psi\|^2}{\|\psi\|^2}.
\]
换言之，有限分辨率可见纯态模空间并不产生新的量子几何，而严格退化为经典概率单纯形。
\end{proposition}
\begin{proof}
映射 \( [\psi]\mapsto p_\psi \) 对射影缩放与块内酉作用都不变，故先降到商空间。
任取 \(p=(p_x)_{x\in X_m}\in \Delta(X_m)\)，对每个 \(x\) 任选单位向量 \(\xi_x\in \mathcal H_x\)，令
\[
\psi:=\bigoplus_{x\in X_m}\sqrt{p_x}\,\xi_x.
\]
则 \(\|\psi\|=1\) 且 \(p_\psi=p\)，故满射成立。

若 \(p_\psi=p_\phi\)，归一化后可设 \(\|\psi\|=\|\phi\|=1\)。
记 \(\psi_x:=\Pi_x\psi\)、\(\phi_x:=\Pi_x\phi\)，则对每个 \(x\) 有 \(\|\psi_x\|=\|\phi_x\|\)。
若 \(\psi_x=\phi_x=0\) 则任取 \(U_x\in U(\mathcal H_x)\)；若二者非零，则存在 \(U_x\in U(\mathcal H_x)\) 使 \(U_x\psi_x=\phi_x\)。
令 \(U:=\bigoplus_x U_x\in U_m\)，则 \(U\psi=\phi\)，故二者在 \(U_m\)-商中等价。
因此该映射双射。证毕。
\end{proof}

\begin{theorem}[纤维多重度多重集给出折叠量子信道的有限分辨率完全分类]\label{thm:conclusion-fold-quantum-channel-multiplicity-classification}
\leanverified{paper\_conclusion\_fold\_quantum\_channel\_multiplicity\_classification}
设
\[
\mathcal H=\bigoplus_{x\in X}\mathcal H_x,
\qquad
\mathcal H'=\bigoplus_{y\in Y}\mathcal H'_y
\]
是两组有限维正交分解，并分别定义测量--制备型折叠信道
\[
\mathcal E(A)=\sum_{x\in X}\frac{\Tr(P_xA)}{d_x}\,P_x,
\qquad
\mathcal E'(A')=\sum_{y\in Y}\frac{\Tr(P'_yA')}{d'_y}\,P'_y,
\]
其中 \(P_x,P'_y\) 为块投影，\(d_x=\dim\mathcal H_x\)，\(d'_y=\dim\mathcal H'_y\)。
若 \(\dim\mathcal H=\dim\mathcal H'=n\)，则 \(\mathcal E\) 与 \(\mathcal E'\) 酉共轭，当且仅当其纤维多重度多重集
\[
\mathcal D(\mathcal E):=\{d_x:\ x\in X\},
\qquad
\mathcal D(\mathcal E'):=\{d'_y:\ y\in Y\}
\]
相同。
\end{theorem}
\begin{proof}
若 \(\mathcal D(\mathcal E)=\mathcal D(\mathcal E')\)，则可取双射 \(\sigma:X\to Y\) 使 \(d_x=d'_{\sigma(x)}\)。
对每个 \(x\) 选酉同构 \(V_x:\mathcal H_x\to \mathcal H'_{\sigma(x)}\)，令 \(V:=\bigoplus_x V_x\)。
则 \(VP_xV^\ast=P'_{\sigma(x)}\)，故对任意 \(A\in\mathcal B(\mathcal H)\) 有
\[
V\mathcal E(A)V^\ast=\mathcal E'(VAV^\ast).
\]
从而两信道酉共轭。

反之，若二者酉共轭，则其归一化 Choi 态谱相同。
由定理 \ref{thm:op-algebra-fold-quantum-channel-choi-spectrum-mi}，非零特征值恰为
\[
\lambda_d=\frac{1}{nd},
\]
而若某个整数 \(d\) 在多重度多重集中出现 \(N_d\) 次，则 \(\lambda_d\) 的总代数重数为 \(N_d d^2\)。
因此由 Choi 谱可唯一恢复每个 \(d\) 与其出现次数 \(N_d\)，从而唯一恢复 \(\mathcal D(\mathcal E)\)。
故酉共轭必迫使两多重度多重集相同。证毕。
\end{proof}

\begin{corollary}[时间纤维 Fibonacci 立方谱即量子资源的精确配分函数]\label{cor:conclusion-fold-quantum-fiber-spectrum-resource-partition}
\leanverified{paper\_conclusion\_fold\_quantum\_fiber\_spectrum\_resource\_partition}
若对每个 \(x\in X_m\)，纤维重构图满足
\[
\mathcal F_x\cong \prod_{i=1}^{r(x)}\Gamma_{\ell_i(x)}
\]
（结论 \ref{con:xi-terminal-fiber-reconstruction-cartesian}），则
\[
d_m(x)=|\mathcal F_x|=\prod_{i=1}^{r(x)}F_{\ell_i(x)+2}.
\]
于是折叠量子信道 \(\mathcal E_m\) 的资源量完全由该时间纤维立方谱给出：
\begin{enumerate}
  \item 最小 Stinespring 环境维数为
  \[
  \min\dim(\mathcal K_{\mathrm{env}})
  =
  \sum_{x\in X_m}\prod_{i=1}^{r(x)}F_{\ell_i(x)+2}^{\,2};
  \]
  \item 归一化 Choi 态 \(\rho_{RB}=J(\mathcal E_m)/2^m\) 的正特征值对每个 \(x\) 为
  \[
  \lambda_x
  =
  2^{-m}\prod_{i=1}^{r(x)}F_{\ell_i(x)+2}^{-1},
  \]
  其重数恰为
  \[
  \prod_{i=1}^{r(x)}F_{\ell_i(x)+2}^{\,2};
  \]
  \item Choi 互信息满足
  \[
  I(R\!:\!B)_{\rho_{RB}}=H(\pi_m),
  \qquad
  \pi_m(x)=2^{-m}\prod_{i=1}^{r(x)}F_{\ell_i(x)+2}.
  \]
\end{enumerate}
因此有限分辨率量子资源并非额外原语，而是时间纤维立方谱的精确配分函数。
\end{corollary}
\begin{proof}
由结论 \ref{con:xi-terminal-fiber-reconstruction-cartesian}，\(\mathcal F_x\) 为 Fibonacci 立方的 Cartesian 乘积。
而 \(|\Gamma_\ell|=F_{\ell+2}\)，故
\[
d_m(x)=|\mathcal F_x|=\prod_i |\Gamma_{\ell_i(x)}|
=\prod_i F_{\ell_i(x)+2}.
\]
将该表达代入定理 \ref{thm:op-algebra-fold-quantum-channel-minimal-environment-equals-s2} 与定理 \ref{thm:op-algebra-fold-quantum-channel-choi-spectrum-mi} 即得三式。证毕。
\end{proof}

\begin{corollary}[profinite 可见极限上的有限维量子逻辑不可能性]\label{cor:conclusion-profinite-visible-logic-finite-dim-boolean}
\leanverified{paper\_conclusion\_profinite\_visible\_logic\_finite\_dim\_boolean}
令 \(G_\infty^{\mathrm{BC}}\cong \mathrm{Homeo}_{\mathrm{cyl}}(X_\infty)\) 为定理 \ref{thm:pom-bc-profinite-reconstruction-homeo-cyl} 给出的柱语义保持 profinite 群。
则任意有限维连续表示
\[
\pi:G_\infty^{\mathrm{BC}}\to \GL(V)
\]
都经某一有限层 \(G_m^{\mathrm{BC}}\) 因子化；并且一旦在该有限层施加与折叠读出相容的可见条件期望，其对应的尖事件代数必落在某个交换代数 \(\mathcal Z_m\) 中。
因而其尖事件格必为布尔格。
特别地，inverse-limit 可见理论中的任何真正非布尔量子逻辑，都不能被有限维连续表象捕获。
\end{corollary}
\begin{proof}
定理 \ref{thm:pom-repcont-profinite-colimit} 表明任意有限维连续表示都经某个有限商 \(G_m^{\mathrm{BC}}\) 因子化。
在该有限层上，可见读出由折叠条件期望与投影族 \(\{\Pi_x\}_{x\in X_m}\) 给出，而定理 \ref{thm:conclusion-fold-quantum-visible-booleanization} 已证明相应可见代数恰为交换代数
\(
\mathcal Z_m=\bigoplus_{x\in X_m}\CC\,\Pi_x
\)
且其投影格布尔。
故任何有限维连续扇区都只能看到这一布尔骨架，而不能承载真正的非布尔尖事件逻辑。证毕。
\end{proof}

\begin{proposition}[Born 权重不是附加公理，而是 Hardy 投影与折叠条件期望的派生量]\label{prop:conclusion-born-weight-derived-from-hardy-and-fold-expectation}
\leanverified{paper\_conclusion\_born\_weight\_derived\_from\_hardy\_and\_fold\_expectation}
令 \(\Pi_N\) 为命题 \ref{prop:unit-circle-phase-hardy-kernel-rational} 与 \S\ref{subsec:cdim-hardy-finite-mode-projection} 中的有限模态 Hardy 投影。
对任意 \(f\in L^2(\RR,\dd\mu)\) 令
\[
\psi_N:=\Pi_N f\neq 0,
\qquad
\rho_{N,f}:=\frac{\lvert\psi_N\rangle\langle\psi_N\rvert}{\|\psi_N\|^2}.
\]
则有限分辨率可见概率分布由严格恒等式给出
\[
\mu_{N,m}^f(x)
:=
\Tr\!\bigl(\Pi_x\,\mathcal E_m(\rho_{N,f})\bigr)
=
\frac{\|\Pi_x\psi_N\|^2}{\|\psi_N\|^2}.
\]
又由
\[
(\Pi_N f)(t)=\int_{\RR}K_N(t,y)\,f(y)\,\dd\mu(y)
\]
且 \(K_N\) 为显式有理核（命题 \ref{prop:unit-circle-phase-hardy-kernel-rational}），\(\mu_{N,m}^f(x)\) 遂成为输入振幅 \(f\) 的显式有理积分函数。
因此，在该框架内有限分辨率 Born 权重不是附加公理，而是 \(\Pi_N\) 与 \(\mathcal E_m\) 的派生定理。
\end{proposition}
\begin{proof}
由 \(\Pi_x\Pi_y=\delta_{xy}\Pi_x\) 对任意态 \(\rho\) 有
\[
\Tr\!\bigl(\Pi_x\mathcal E_m(\rho)\bigr)
=
\sum_{y\in X_m}\frac{\Tr(\Pi_y\rho)}{d_m(y)}\,\Tr(\Pi_x\Pi_y)
=
\Tr(\Pi_x\rho).
\]
代入 \(\rho=\rho_{N,f}\) 得
\[
\mu_{N,m}^f(x)
=
\Tr(\Pi_x\rho_{N,f})
=
\frac{\langle \psi_N,\Pi_x\psi_N\rangle}{\|\psi_N\|^2}
=
\frac{\|\Pi_x\psi_N\|^2}{\|\psi_N\|^2}.
\]
最后一条核表示正是命题 \ref{prop:unit-circle-phase-hardy-kernel-rational} 的结论。证毕。
\end{proof}

\begin{theorem}[真实 fold-gauge 不变量代数的矩阵扇区裂解]\label{thm:conclusion-fold-gauge-commutant-matrix-sector}
\leanverified{paper\_conclusion\_fold\_gauge\_commutant\_matrix\_sector}
对固定分辨率 \(m\ge 1\)，令
\[
F_x:=\Fold_m^{-1}(x),
\qquad
d_x:=|F_x|,
\qquad
\mathcal H_m=\ell^2(\Omega_m)=\bigoplus_{x\in X_m}\mathcal H_x,
\qquad
\mathcal H_x:=\ell^2(F_x).
\]
定义纤维置换规范群
\[
G_m^{\mathrm{fib}}:=\prod_{x\in X_m}\Sym(F_x),
\]
并记其在 \(\mathcal H_m\) 上的自然置换表示为 \(\rho_m\)。
对每个 \(x\in X_m\)，令
\[
u_x:=d_x^{-1/2}\sum_{\omega\in F_x}\delta_\omega,
\qquad
P_x^+:=\lvert u_x\rangle\langle u_x\rvert,
\qquad
P_x^-:=\Pi_x-P_x^+.
\]
再定义
\[
Q_m:=\sum_{x\in X_m}P_x^+,
\qquad
U_m:=Q_m\mathcal H_m=\Span\{u_x:\ x\in X_m\},
\qquad
W_m:=\bigoplus_{x\in X_m}P_x^-\mathcal H_m.
\]
则存在自然 \(^*\)-代数同构
\[
\rho_m(G_m^{\mathrm{fib}})'
\cong
\mathcal B(U_m)\ \oplus\ \bigoplus_{\substack{x\in X_m\\ d_x\ge 2}}\CC\,P_x^-.
\]
特别地，
\[
Q_m\,\rho_m(G_m^{\mathrm{fib}})'\,Q_m=\mathcal B(U_m)\cong M_{|X_m|}(\CC).
\]
因此真实 fold-gauge 的固定点代数并不交换；它内生地包含一个维数 \(|X_m|\) 的完整量子矩阵扇区。
\end{theorem}
\begin{proof}
对每个 \(x\in X_m\)，有正交分解
\[
\mathcal H_x=\CC u_x\oplus W_x,
\qquad
W_x:=u_x^\perp=P_x^-\mathcal H_m.
\]
其中 \(\Sym(F_x)\) 在 \(\CC u_x\) 上平凡作用，在 \(W_x\) 上给出标准表示；当 \(d_x\ge 2\) 时，该标准表示在 \(\CC\) 上不可约。
于是
\[
\mathcal H_m
=
U_m\oplus \bigoplus_{\substack{x\in X_m\\ d_x\ge 2}}W_x
\]
是 \(G_m^{\mathrm{fib}}\) 的等型分解，并且 \(U_m\) 是平凡表示的 \(|X_m|\) 重等型块。

由于 \(G_m^{\mathrm{fib}}\) 在 \(U_m\) 上平凡作用，故 \(\mathcal B(U_m)\subset \rho_m(G_m^{\mathrm{fib}})'\)。
另一方面，若 \(x\neq y\) 且 \(T:W_x\to W_y\) 与 \(G_m^{\mathrm{fib}}\) 对易，则把 \(T\) 限制到第 \(x\) 个因子 \(\Sym(F_x)\) 可见：\(W_x\) 上该因子作用非平凡，而 \(W_y\) 上该因子作用平凡，故
\[
T\in \Hom_{\Sym(F_x)}(W_x,\mathbf 1)=0.
\]
同理，
\[
\Hom_{G_m^{\mathrm{fib}}}(U_m,W_x)
=
\Hom_{G_m^{\mathrm{fib}}}(W_x,U_m)
=0.
\]
又由 Schur 引理，
\[
\End_{G_m^{\mathrm{fib}}}(W_x)=\CC\,I_{W_x}=\CC\,P_x^-.
\]
综上，交换子代数恰为
\[
\mathcal B(U_m)\oplus \bigoplus_{\substack{x\in X_m\\ d_x\ge 2}}\CC\,P_x^-.
\]
再乘以 \(Q_m\) 即得
\(
Q_m\rho_m(G_m^{\mathrm{fib}})'Q_m=\mathcal B(U_m)
\)。
证毕。
\end{proof}

\begin{corollary}[真实规范纯态空间复现为射影空间]\label{cor:conclusion-fold-gauge-purestate-projective-space}
\leanverified{paper\_conclusion\_fold\_gauge\_purestate\_projective\_space}
存在规范酉同构
\[
U_m\cong \ell^2(X_m),
\qquad
u_x\longmapsto e_x.
\]
因此真实 fold-gauge 不变量代数中的相干纯态空间同构于
\[
\mathbb P(U_m)\cong \mathbb P^{|X_m|-1}(\CC).
\]
并且
\[
\dim U_m=|X_m|,
\qquad
\dim W_m=\sum_{x\in X_m}(d_x-1)=2^m-|X_m|.
\]
当 \(m=6\) 时，由推论 \ref{cor:fold-bin-m6-fiber-hist} 的直方图
\[
2{:}8,\qquad 3{:}4,\qquad 4{:}9
\]
可得
\[
|X_6|=21,\qquad \dim U_6=21,\qquad \dim W_6=43,
\]
并且
\[
\rho_6(G_6^{\mathrm{fib}})'
\cong
M_{21}(\CC)\oplus \CC^{21}.
\]
\end{corollary}
\begin{proof}
\(\{u_x\}_{x\in X_m}\) 两两正交且均为单位向量，故构成 \(U_m\) 的一组标准正交基，从而给出所示酉同构。
维数公式由
\[
\sum_{x\in X_m}d_x=|\Omega_m|=2^m
\]
立得。
纯态空间结论是有限维 Hilbert 空间的射影化。
最后，当 \(m=6\) 时，由推论 \ref{cor:fold-bin-m6-fiber-hist} 得
\[
|X_6|=8+4+9=21,
\qquad
\dim W_6=8(2-1)+4(3-1)+9(4-1)=43.
\]
再代入定理 \ref{thm:conclusion-fold-gauge-commutant-matrix-sector} 即得
\[
\rho_6(G_6^{\mathrm{fib}})'\cong M_{21}(\CC)\oplus \CC^{21}.
\]
这也解释了命题 \ref{prop:conclusion-fold-quantum-projective-simplex-quotient} 的单纯形并非原初对象，而是对更大的块内酉群 \(\prod_{x\in X_m}U(\mathcal H_x)\) 进一步商化后的结果。证毕。
\end{proof}

\begin{theorem}[可见布尔读出是矩阵扇区退相干与暗多重度坍缩的复合]\label{thm:conclusion-fold-gauge-average-visible-decoherence}
在定理 \ref{thm:conclusion-fold-gauge-commutant-matrix-sector} 的设定下，定义群平均
\[
\mathfrak T_m(A):=\frac{1}{|G_m^{\mathrm{fib}}|}\sum_{g\in G_m^{\mathrm{fib}}}\rho_m(g)A\rho_m(g)^*
\qquad
(A\in \mathcal B(\mathcal H_m)).
\]
则对任意 \(A\in\mathcal B(\mathcal H_m)\)，有显式公式
\[
\mathfrak T_m(A)
=
Q_mAQ_m
+\sum_{\substack{x\in X_m\\ d_x\ge 2}}
\frac{\Tr(P_x^-A)}{d_x-1}\,P_x^-.
\]
另一方面，折叠可见信道满足
\[
\mathcal E_m(A)
=
\sum_{x\in X_m}
\frac{\langle u_x,Au_x\rangle+\Tr(P_x^-A)}{d_x}\,\Pi_x.
\]
因此 \(\mathcal E_m\) 的布尔化并非原初结构，而是先保留整个 \(\mathcal B(U_m)\) 量子矩阵块，再沿基 \(\{u_x\}\) 完全退相干，并把每条纤维的暗标准分量仅以总迹的形式坍缩到相应对角项之后的复合结果。
\end{theorem}
\begin{proof}
群平均 \(\mathfrak T_m\) 是到交换子代数 \(\rho_m(G_m^{\mathrm{fib}})'\) 的条件期望。
由定理 \ref{thm:conclusion-fold-gauge-commutant-matrix-sector}，
\(\rho_m(G_m^{\mathrm{fib}})'\) 在分解
\[
\mathcal H_m=U_m\oplus \bigoplus_{x:d_x\ge 2}W_x
\]
上恰为
\(
\mathcal B(U_m)\oplus \bigoplus_x \CC P_x^-
\)。
故 \(\mathfrak T_m\) 在平凡表示块 \(U_m\) 上不做压缩，保留为 \(Q_mAQ_m\)；
而在每个不可约块 \(W_x\) 上只能留下标量部分，其系数由迹归一化给出：
\[
\frac{\Tr(P_x^-A)}{\dim W_x}
=
\frac{\Tr(P_x^-A)}{d_x-1}.
\]
于是得到第一式。

第二式由
\[
\Pi_x=P_x^++P_x^-,
\qquad
\Tr(P_x^+A)=\langle u_x,Au_x\rangle
\]
直接推出：
\[
\Tr(\Pi_xA)=\langle u_x,Au_x\rangle+\Tr(P_x^-A).
\]
代回 \(\mathcal E_m\) 的定义即得。
因此 \(\mathcal E_m\) 先读取 \(U_m\) 上完整的相干块，再把它退相干到 \(\{u_x\}\) 的对角，并把暗块只保留其总迹权重。证毕。
\leanverified{paper\_conclusion\_fold\_gauge\_average\_visible\_decoherence}
\end{proof}

\begin{theorem}[可见布尔化核的显式分解与精确维数]\label{thm:conclusion-fold-quantum-visible-kernel-explicit-decomposition}
\leanverified{paper\_conclusion\_fold\_quantum\_visible\_kernel\_explicit\_decomposition}
在定理 \ref{thm:conclusion-fold-gauge-average-visible-decoherence} 的设定下，记
\[
A_m^{\mathrm{vis}}:=\rho_m(G_m^{\mathrm{fib}})',
\qquad
E_m^{\mathrm{vis}}:=\mathcal E_m|_{A_m^{\mathrm{vis}}}:A_m^{\mathrm{vis}}\to Z_m,
\]
\[
N_{\ge 2}(m):=\#\{x\in X_m:\ d_x\ge 2\},
\qquad
\mathcal O_m:=\{B\in \mathcal B(U_m):\ \langle u_x,Bu_x\rangle=0\ \forall x\in X_m\}.
\]
则 \(E_m^{\mathrm{vis}}\) 为满射，且其核具有显式直和分解
\[
\ker E_m^{\mathrm{vis}}
=
\mathcal O_m
\oplus
\bigoplus_{\substack{x\in X_m\\ d_x\ge 2}}
\CC\Bigl(P_x^--(d_x-1)P_x^+\Bigr).
\]
特别地，
\[
\dim_{\CC}\ker E_m^{\mathrm{vis}}
=
|X_m|(|X_m|-1)+N_{\ge 2}(m).
\]
因此，可见布尔化在固定分辨率上擦除的是一个精确可计数的非交换核，而不是一团不可分辨的“量子剩余”。
\end{theorem}
\begin{proof}
由定理 \ref{thm:conclusion-fold-gauge-commutant-matrix-sector}，
任意 \(A\in A_m^{\mathrm{vis}}\) 都可唯一写成
\[
A=B+\sum_{\substack{x\in X_m\\ d_x\ge 2}}\lambda_xP_x^-,
\qquad
B\in\mathcal B(U_m).
\]
再由定理 \ref{thm:conclusion-fold-gauge-average-visible-decoherence}，
\[
E_m^{\mathrm{vis}}(A)
=
\sum_{x\in X_m}
\frac{\langle u_x,Bu_x\rangle+(d_x-1)\lambda_x}{d_x}\,\Pi_x,
\]
其中当 \(d_x=1\) 时约定 \(\lambda_x=0\)。

先证满射：对任意
\[
Z=\sum_{x\in X_m}c_x\Pi_x\in Z_m
\]
取
\[
B:=\sum_{x\in X_m}d_xc_x\,P_x^+\in \mathcal B(U_m),
\qquad
\lambda_x:=0,
\]
则 \(E_m^{\mathrm{vis}}(B)=Z\)。

为求核，把 \(B\) 在基 \(\{u_x\}_{x\in X_m}\) 下分解为
\[
B=B_0+\sum_{x\in X_m}b_xP_x^+,
\qquad
B_0\in\mathcal O_m.
\]
于是 \(A\in\ker E_m^{\mathrm{vis}}\) 当且仅当对每个 \(x\in X_m\) 有
\[
b_x+(d_x-1)\lambda_x=0.
\]
故
\[
A
=
B_0+\sum_{\substack{x\in X_m\\ d_x\ge 2}}\lambda_x\Bigl(P_x^--(d_x-1)P_x^+\Bigr),
\]
从而得到所示直和分解。最后，
\[
\dim_{\CC}\mathcal O_m
=
\dim_{\CC}\mathcal B(U_m)-|X_m|
=
|X_m|^2-|X_m|,
\]
每个 \(d_x\ge 2\) 再贡献一个独立一维核方向，故
\[
\dim_{\CC}\ker E_m^{\mathrm{vis}}
=
|X_m|^2-|X_m|+N_{\ge 2}(m).
\]
证毕。
\end{proof}


\leanverified{paper\_conclusion\_fold\_capacity\_wedderburn\_exact\_bit\_threshold}
\begin{corollary}[容量曲线、Wedderburn 型与精确反演门槛的同一性]\label{cor:conclusion-fold-capacity-wedderburn-exact-bit-threshold}
固定分辨率 \(m\ge 1\)。
连续容量曲线
\[
\mathcal C_m^{\mathrm{cont}}(T):=\sum_{x\in X_m}\min(d_x,T)
\qquad (T\ge 0)
\]
唯一决定纤维直方图
\[
n_k(m):=\#\{x\in X_m:\ d_x=k\},
\]
从而唯一决定群胚代数的 Wedderburn 型
\[
\CC[\mathcal R_m^{\mathrm{bin}}]
\cong
\bigoplus_{k\ge 1}M_k(\CC)^{\oplus n_k(m)}.
\]
此外，达到精确反演
\[
\sum_{x\in X_m}\min(d_x,2^B)=2^m
\]
所需的最小侧信息比特数满足
\[
B_{\mathrm{exact}}(m)
=
\Bigl\lceil \log_2 d_{\max}(m)\Bigr\rceil
=
\Bigl\lceil \log_2 \operatorname{Ind}_{\mathrm{PP}}(E_m^Z)\Bigr\rceil,
\]
其中 \(E_m^Z\) 为定理 \ref{thm:xi-time-part9ze-center-condexp-pimsner-popa-index} 的中心条件期望。
特别地，在 window-\(6\) 上
\[
d_{\max}(6)=4,
\qquad
B_{\mathrm{exact}}(6)=2.
\]
\end{corollary}
\begin{proof}
由定理 \ref{thm:xi-time-part60aaa-wedderburn-information-completeness}，
\(\mathcal C_m^{\mathrm{cont}}\)、纤维直方图 \(\{n_k(m)\}_{k\ge 1}\) 与半单代数同构型彼此唯一决定，故第一部分成立。

另一方面，
\[
\sum_{x\in X_m}\min(d_x,2^B)=2^m=\sum_{x\in X_m}d_x
\]
成立，当且仅当 \(2^B\ge d_x\) 对全部 \(x\) 同时成立；
故最小 \(B\) 正是 \(\lceil\log_2 d_{\max}(m)\rceil\)。
再由定理 \ref{thm:xi-time-part9ze-center-condexp-pimsner-popa-index}，
\[
\operatorname{Ind}_{\mathrm{PP}}(E_m^Z)=d_{\max}(m),
\]
遂得第二式。
对 \(m=6\)，由推论 \ref{cor:fold-bin-m6-fiber-hist} 得 \(d_{\max}(6)=4\)，从而 \(B_{\mathrm{exact}}(6)=2\)。证毕。
\end{proof}

\begin{corollary}[bin-fold 的全部正阶 Rényi 熵率坍缩到 \texorpdfstring{$\log\varphi$}{log phi}]\label{cor:conclusion-binfold-renyi-rate-collapse}
\leanverified{paper\_conclusion\_binfold\_renyi\_rate\_collapse}
在二进制区间口径下，令
\[
\mu_m(x):=\frac{d_m^{\mathrm{bin}}(x)}{2^m}
\qquad(x\in X_m),
\]
并对任意 \(q>0\) 定义 Rényi 熵
\[
H_q(\mu_m):=
\begin{cases}
\dfrac{1}{1-q}\log\sum_x \mu_m(x)^q,& q\neq 1,\\[1ex]
-\sum_x \mu_m(x)\log\mu_m(x),& q=1.
\end{cases}
\]
则对一切 \(q>0\) 都有
\[
\lim_{m\to\infty}\frac{1}{m}H_q(\mu_m)=\log\varphi.
\]
因此 bin-fold 的全部正阶 Rényi 熵率共享同一指数斜率，非均匀性只能停留在 \(O(1)\) 的亚指数层。
\end{corollary}
\begin{proof}
当 \(q\neq 1\) 时，定理 \ref{thm:xi-fold-renyi-constant-spectrum-closed-form} 给出
\[
H_q(\mu_m)=m\log\varphi+c(q)+O\!\Bigl(\Bigl(\frac{\varphi}{2}\Bigr)^m\Bigr),
\]
故
\[
\frac{1}{m}H_q(\mu_m)\to \log\varphi.
\]
当 \(q=1\) 时，同一定理给出
\[
H(\mu_m)=m\log\varphi+\frac{\log\varphi}{\varphi\sqrt 5}+o(1),
\]
故极限仍为 \(\log\varphi\)。
于是对全部 \(q>0\) 结论统一成立。证毕。
\end{proof}

\begin{conjecture}[inverse-limit 异常 \(2\)-类是全局量子性的唯一障碍]\label{conj:conclusion-inverse-limit-anomaly-global-quantumity}
记 \(\mathfrak a_\infty^{\mathbf{PW}}\) 为投影词 \(2\)-范畴 \(\mathbf{PW}\) 在 inverse-limit 可见塔上的兼容异常类极限，即由各有限切片中的异常上同调类 \([\mathrm{Anom}]\) 所拼接的残差 \(2\)-类。
则 inverse-limit 可见理论存在真正非布尔的完备尖事件逻辑，当且仅当 \(\mathfrak a_\infty^{\mathbf{PW}}\) 在所有有限截断上都不能被消去。
换言之，若全局层仍保留不可约量子性，则其唯一持续来源应是贯穿全部有限截断而不消失的异常 \(2\)-类，而非任何有限维可见读出本身。
\end{conjecture}

\begin{remark}[支持性理由]\label{rem:conclusion-inverse-limit-anomaly-global-quantumity-support}
定理 \ref{thm:conclusion-fold-quantum-visible-booleanization} 已表明每个有限分辨率可见层都严格布尔化；
推论 \ref{cor:conclusion-profinite-visible-logic-finite-dim-boolean} 进一步说明每个有限维连续对称扇区都只能看到该布尔骨架。
另一方面，\(\mathbf{PW}\) 中的异常是否为零由上同调类控制（推论 \ref{cor:pom-anom-cohomology-class}），而在有界切片上 strictification 的存在性可被归约为有限异常消失约束的可满足性（推论 \ref{cor:pom-strictification-existence-decidable-bounded-slice}）。
特别地，在 window-\(6\) 的最小非平凡锚点上，异常标签可压缩为 \((\ZZ_2)^{21}\) 上的有限线性可满足性问题，并在 boundary 中进一步投影到规范的 \((\ZZ_2)^3\) 子格（推论 \ref{cor:fold-bin-gauge-m6} 与推论 \ref{cor:fold-bin-boundary-center-z2}）。
因此，若 inverse limit 处仍有真正不可压平的非经典性，它就不可能来自任何单个有限层或任何有限维连续扇区，而只能来自这一异常残差类在极限中的持续存活。
\end{remark}

% --- Distillation writeback ---
% --- Distillation writeback: Wang-Zahl fold-aware coherent grain audit ---
\begin{theorem}[fold-aware coherent grains inside the visible matrix sector]\label{distill:wang_zahl-grain_prism_structure_dichotomy-conclusion-visible-matrix-grains}
\textbf{日期时间：2026-04-23T06:15:38+08:00. 依赖状态：rigidity.}
沿用定理 \ref{thm:conclusion-fold-gauge-commutant-matrix-sector} 的记号。对 \(x,y\in X_m\)，令
\[
E_{xy}:=|u_x\rangle\langle u_y|\in\mathcal B(U_m),
\qquad E_{xx}=P_x^+.
\]
设 \(\mathcal A\subset\mathcal B(U_m)\) 是一个 unital \(^\ast\)-子代数，并假设
\[
P_x^+\in\mathcal A\qquad(x\in X_m).
\]
定义图 \(\Gamma_{\mathcal A}\) 的顶点集为 \(X_m\)，且当 \(x\ne y\) 且
\[
P_x^+\mathcal A P_y^+\ne 0
\]
时连接 \(x\) 与 \(y\)。令 \(\mathcal C(\mathcal A)\) 为 \(\Gamma_{\mathcal A}\) 的连通分支，并记
\[
U_C:=\Span\{u_x:x\in C\}.
\]
则
\[
\mathcal A=\bigoplus_{C\in\mathcal C(\mathcal A)}\mathcal B(U_C)
\]
作为 \(U_m=\bigoplus_C U_C\) 上的块对角 \(^\ast\)-代数。

特别地，\(\mathcal C(\mathcal A)\) 正是该 visible audit algebra 的 maximal coherent grains。若 \(Y\in\mathcal C(\mathcal A)\)，而非空集合 \(P\subset X_m\setminus Y\) 满足：对每个 \(y\in P\) 存在 \(x_y\in Y\) 使
\[
E_{x_y y}\in\mathcal A,
\]
则
\[
\mathcal B(U_{Y\cup P})\subset\mathcal A,
\]
从而 \(Y\) 不可能仍为 maximal coherent grain。也就是说，在真实 fold-gauge 可见矩阵扇区内，一个 fully occupied thin prism 一旦携带通向 maximal grain 的相干 matrix-unit attachment，就必然被吸收到更大的 grain；它不能作为独立 residual prism 存活。
\end{theorem}
\begin{proof}
首先注意，对任意 \(T\in\mathcal B(U_m)\)，由于 \(P_x^+\) 与 \(P_y^+\) 均为一维正交投影，存在唯一标量 \(\lambda_{xy}(T)\in\CC\) 使
\[
P_x^+TP_y^+=\lambda_{xy}(T)E_{xy}.
\]
因此若 \(P_x^+\mathcal A P_y^+\ne0\)，则某个 \(T\in\mathcal A\) 满足 \(\lambda_{xy}(T)\ne0\)，从而
\[
E_{xy}=\lambda_{xy}(T)^{-1}P_x^+TP_y^+\in\mathcal A,
\]
因为 \(P_x^+,P_y^+\in\mathcal A\)。反向显然：若 \(E_{xy}\in\mathcal A\)，则 \(P_x^+\mathcal A P_y^+\ne0\)。

若 \(x=x_0,x_1,\ldots,x_r=y\) 是 \(\Gamma_{\mathcal A}\) 中的一条路径，则由上一段可知每个 \(E_{x_{j-1}x_j}\in\mathcal A\)，并且
\[
E_{x_0x_1}E_{x_1x_2}\cdots E_{x_{r-1}x_r}=E_{xy}.
\]
故同一连通分支内的全部 matrix units \(E_{xy}\) 都属于 \(\mathcal A\)。于是
\[
\mathcal B(U_C)=\Span\{E_{xy}:x,y\in C\}\subset\mathcal A
\]
对每个分支 \(C\) 成立。

另一方面，若 \(x,y\) 属于不同连通分支，则按图的定义必有 \(P_x^+\mathcal A P_y^+=0\)。因此任意 \(T\in\mathcal A\) 在不同分支之间的矩阵块都为零，故
\[
\mathcal A\subset\bigoplus_{C\in\mathcal C(\mathcal A)}\mathcal B(U_C).
\]
与上一段的反向包含合并，得到所示块分解。

最后设 \(Y\in\mathcal C(\mathcal A)\)，且 \(P\subset X_m\setminus Y\) 满足题设 attachment 条件。对每个 \(y\in P\)，由 \(E_{x_y y}\in\mathcal A\) 可知 \(x_y\) 与 \(y\) 在 \(\Gamma_{\mathcal A}\) 中相邻；而 \(x_y\in Y\)，故 \(y\) 与整个分支 \(Y\) 位于同一个连通分支。于是 \(Y\cup P\) 包含在某个连通分支 \(C'\) 中。由已证的分解，\(\mathcal B(U_{C'})\subset\mathcal A\)，从而特别有 \(\mathcal B(U_{Y\cup P})\subset\mathcal A\)。若 \(P\ne\varnothing\)，则 \(C'\) 严格大于 \(Y\)，这与 \(Y\) 作为 maximal coherent grain 相矛盾。

该论证只使用 fold-aware 的纤维平均向量 \(u_x\) 与真实可见矩阵扇区 \(U_m\)。当某些 \(d_x=1\) 时，\(P_x^-=0\) 仅意味着相应暗标准分量为空；上述 rank-one matrix-unit 论证不受影响。证毕。
\end{proof}
% --- End distillation writeback ---


\paragraph{调和异常刚性、可见 Walsh 条码、非可缩零温筛选与有限核心坍缩}
在前述 strictification 上同调判别、fold-gauge 可见量子扇区、非可缩最大纤维相图与 Toeplitz--PSD 量词消去机制已经闭合之后，还可把若干原本分散在不同章节中的后继效应压缩为一条更高阶的终端链：异常的最优残留并非任意 counterterm 之后的剩余，而严格是调和分量；低阶 Walsh 可见谱一旦缺陷全灭，就会收缩为一条由 Gram 秩向量完全控制的精确谱条码；非可缩拓扑筛选不会改变零温主指数，只会留下模 \(6\) 的算术前因子；而在一致 realization 维数上界下，双参数证书逆系统又会坍缩为固定有限 PSD 核心之上的任意共终深度链。最后，window-\(6\) 还把“精确反演”与“异常账本完备读出”分裂为两个严格不同的信息尺度。

\begin{theorem}[异常类的调和刚性与最小残留 holonomy]\label{thm:conclusion-anomaly-harmonic-rigidity}
设 \(\mathcal K\) 为有限二维 CW 复形，\(V\) 为有限维实内积空间，且 \(a\in C^2(\mathcal K;V)\) 为异常 \(2\)-上链。
在正交分解
\[
C^2(\mathcal K;V)=\operatorname{im}\delta\oplus\ker\delta^\ast
\]
下写
\[
a=\delta\eta_0+h,
\qquad
h\in\ker\delta^\ast.
\]
则有：
\begin{enumerate}
  \item \(a\) 可 strictify 当且仅当 \(h=0\)；
  \item 对任意 \(\eta\in C^1(\mathcal K;V)\)，恒有精确恒等式
  \[
  \|a-\delta\eta\|_2^2
  =
  \|h\|_2^2+\|\delta(\eta-\eta_0)\|_2^2;
  \]
  \item 因而所有 counterterm 中的最小残留异常范数恰为 \(\|h\|_2\)；
  \item 对任意 \(2\)-循环 \(\Sigma\in Z_2(\mathcal K)\)，有
  \[
  \langle a-\delta\eta_0,\Sigma\rangle
  =
  \langle h,\Sigma\rangle.
  \]
\end{enumerate}
换言之，strictification 后仍不可消去的全部 holonomy，严格等于异常类调和代表在循环上的配对。
\leanverified{norm\_add\_sq\_of\_inner\_zero}
\leanverified{norm\_sub\_sq\_of\_inner\_zero}
\leanverified{anomaly\_pythagoras}
\leanverified{min\_residual\_eq\_h\_norm\_sq}
\leanverified{strictify\_iff\_h\_zero}
\leanverified{min\_residual\_anomaly}
\leanverified{paper\_conclusion\_anomaly\_harmonic\_rigidity}
\end{theorem}
\begin{proof}
由于 \(\mathcal K\) 有限，\(C^2(\mathcal K;V)\) 为有限维内积空间，故有正交分解
\[
C^2(\mathcal K;V)=\operatorname{im}\delta\oplus(\operatorname{im}\delta)^\perp
=\operatorname{im}\delta\oplus\ker\delta^\ast.
\]
于是可唯一写成 \(a=\delta\eta_0+h\)，其中 \(h\perp\operatorname{im}\delta\)。
对任意 \(\eta\in C^1(\mathcal K;V)\)，
\[
a-\delta\eta
=
h+\delta(\eta_0-\eta),
\]
且两项正交，故
\[
\|a-\delta\eta\|_2^2
=
\|h\|_2^2+\|\delta(\eta-\eta_0)\|_2^2.
\]
这立即给出第三条。又 \(a\) 可 strictify 当且仅当 \(a\in\operatorname{im}\delta\)，亦即当且仅当 \(h=0\)，得到第一条。

最后，由定理 \ref{thm:pom-pw-noncommutative-stokes}，对任意 \(2\)-循环 \(\Sigma\) 与任意 \(\xi\in C^1(\mathcal K;V)\) 都有
\[
\langle \delta\xi,\Sigma\rangle=0,
\]
于是
\[
\langle a-\delta\eta_0,\Sigma\rangle
=
\langle h,\Sigma\rangle.
\]
证毕。
\end{proof}

\begin{corollary}[二维异常的双互素矩层析完备性]\label{cor:conclusion-anomaly-two-coprime-layer-tomography}
沿用定理 \ref{thm:conclusion-anomaly-harmonic-rigidity} 的设定，并设
\[
[a]\in H^2(\mathcal K;V)
\]
为原始异常类。
取任意两整数 \(q_1,q_2\ge 2\) 且 \(\gcd(q_1,q_2)=1\)，并记
\[
a^{(q_1)},\qquad a^{(q_2)}
\]
为文稿中的两条 \(q\)-矩提升异常。
若
\[
[\Sigma_1],\dots,[\Sigma_{\beta_2}]
\]
为 \(H_2(\mathcal K;\ZZ)/\mathrm{tors}\) 的一组基，则以下命题等价：
\[
[a]=0\in H^2(\mathcal K;V);
\]
\[
\langle [a^{(q_1)}],[\Sigma_j]\rangle
=
\langle [a^{(q_2)}],[\Sigma_j]\rangle
=0
\qquad (j=1,\dots,\beta_2).
\]
因此，二维异常的 strictifiability 可由仅仅 \(2\beta_2\) 个向量配对完全判定。
\end{corollary}
\begin{proof}
因 \(V\) 为实向量空间，存在自然同构
\[
H^2(\mathcal K;V)
\cong
\operatorname{Hom}\bigl(H_2(\mathcal K;\ZZ)/\mathrm{tors},V\bigr).
\]
故一个 \(2\)-上同调类为零，当且仅当它在基 \([\Sigma_1],\dots,[\Sigma_{\beta_2}]\) 上的全部配对同时为零。
于是上式等价于
\[
[a^{(q_1)}]=[a^{(q_2)}]=0.
\]
再由定理 \ref{thm:pom-anom-coprime-two-point-test} 与定理 \ref{thm:pom-anom-all-moments-collapse-to-two}，两个互素矩提升同时消失当且仅当原始异常类消失。
反向蕴含显然。证毕。
\end{proof}

\begin{theorem}[低阶 Walsh 可见谱条码]\label{thm:conclusion-visible-walsh-low-order-barcode}
\leanverified{paper\_conclusion\_visible\_walsh\_low\_order\_barcode}
沿用 \S\ref{subsubsec:pom-visible-walsh-factor-chain-commutator} 的记号。对每个 \(k\in\{0,\dots,m\}\)，定义
\[
\mathcal G_k:=\operatorname{span}\{g_I:\ |I|=k\},
\qquad
M_k:=\operatorname{rank}\bigl(\langle g_I,g_J\rangle_{\pi_m}\bigr)_{|I|=|J|=k}.
\]
若存在 \(K\le m\) 使得对所有 \(|I|\le K\) 都有
\[
U_m^\ast[P_m^\square,E_m]\chi_I=0,
\]
则可见因子链在低能可见子空间
\[
\mathcal G_{\le K}:=\sum_{k=0}^{K}\mathcal G_k
\]
上具有精确谱分解
\[
\overline P_m^\square\big|_{\mathcal G_{\le K}}
\cong
\bigoplus_{k=0}^{K}\left(1-\frac{2k}{m}\right)I_{M_k}.
\]
特别地，
\[
\det\!\left(I-z\,\overline P_m^\square\big|_{\mathcal G_{\le K}}\right)
=
\prod_{k=0}^{K}\left(1-z\left(1-\frac{2k}{m}\right)\right)^{M_k},
\]
并且对任意 \(t\in\NN\)，
\[
\Tr\!\left(\left(\overline P_m^\square\big|_{\mathcal G_{\le K}}\right)^t\right)
=
\sum_{k=0}^{K}M_k\left(1-\frac{2k}{m}\right)^t.
\]
\end{theorem}
\begin{proof}
由定理 \ref{thm:pom-visible-walsh-commutator-defect}，当 \(|I|\le K\) 时
\[
(\overline P_m^\square-\lambda_I\mathrm{Id})g_I=0,
\qquad
\lambda_I=1-\frac{2|I|}{m}.
\]
因此对每个固定 \(k\le K\)，\(\mathcal G_k\) 中的全部生成元都对应同一特征值
\[
\lambda_k:=1-\frac{2k}{m}.
\]
由推论 \ref{cor:pom-visible-walsh-sectorwise-inheritance}，
\[
\overline P_m^\square\big|_{\mathcal G_k}
=
\lambda_k I.
\]
而 \(\dim\mathcal G_k\) 恰为相应 Gram 块的秩 \(M_k\)。
不同 \(k\) 对应不同特征值，故这些特征子空间之和为直和，从而得到所述谱分解。
行列式与迹公式即由有限维谱分解直接推出。证毕。
\end{proof}

\begin{theorem}[非可缩筛选不改变零温主指数，只留下模 \texorpdfstring{$6$}{6} 前因子记忆]\label{thm:conclusion-noncontractible-zero-temp-exponent-mod6-memory}
沿用 \S\ref{subsec:conclusion-max-noncontractible-fiber-zero-block-golden-coupling} 的记号。定义
\[
Z_m^{\mathrm{all}}(\beta):=\sum_{x\in X_m}d_m(x)^\beta,
\qquad
Z_m^{\mathrm{nc}}(\beta):=\sum_{\substack{x\in X_m\\ |K_x|\not\simeq *}}d_m(x)^\beta
\qquad (\beta>0).
\]
再记
\[
D_m:=\max_{x\in X_m}d_m(x),
\qquad
\widetilde D_m:=\max\{d_m(x):x\in X_m,\ |K_x|\not\simeq *\}.
\]
则
\[
\lim_{\beta\to\infty}\bigl(Z_m^{\mathrm{all}}(\beta)\bigr)^{1/\beta}=D_m,
\qquad
\lim_{\beta\to\infty}\bigl(Z_m^{\mathrm{nc}}(\beta)\bigr)^{1/\beta}=\widetilde D_m.
\]
并且
\[
\lim_{m\to\infty}\frac{1}{m}\log D_m
=
\lim_{m\to\infty}\frac{1}{m}\log \widetilde D_m
=
\frac12\log\varphi.
\]
更精确地，
\[
\lim_{\substack{m\to\infty\\ m\equiv r\!\!\!\!\pmod 6}}
\frac{\widetilde D_m}{D_m}
=
c_r,
\]
其中
\[
c_0=c_4=1,\qquad
c_1=c_5=1-\frac12\varphi^{-5},\qquad
c_2=1-\varphi^{-5},\qquad
c_3=1-\frac12\bigl(\varphi^{-5}+\varphi^{-9}\bigr).
\]
因此，非可缩拓扑筛选在零温相中并不触动指数级自由能，而只留下严格算术化的模 \(6\) 记忆。
\end{theorem}
\begin{proof}
有限集合上的 \(\ell^\beta\)-范数极限给出
\[
\lim_{\beta\to\infty}\left(\sum_i a_i^\beta\right)^{1/\beta}=\max_i a_i
\qquad (a_i\ge 0),
\]
故第一组极限立即成立。
再由推论 \ref{cor:conclusion-max-noncontractible-fiber-exponential-rate}，
\[
\lim_{m\to\infty}\frac1m\log D_m
=
\lim_{m\to\infty}\frac1m\log \widetilde D_m
=
\frac12\log\varphi.
\]
最后，模 \(6\) 比值极限正是定理 \ref{thm:conclusion-max-noncontractible-fiber-mod6-golden-ratios} 的四条公式。证毕。
\end{proof}

\begin{theorem}[一致维数上界下的有限核心坍缩]\label{thm:conclusion-finite-core-collapse-uniform-dimension}
记 \(\mathscr C_{b,N}(F)\) 为深度 \(b\) 且 Toeplitz--PSD 仅核验至阶 \(N\) 的证书对象类。
若所考察对象族的最小 realization 维数由某个常数 \(D\) 一致上界，则存在
\[
N_0\le 2D
\]
使得双参数逆系统
\[
\{\mathscr C_{b,N}(F)\}_{b\in\NN,\ N\in\NN}
\]
与单参数深度链
\[
\{\mathscr C_{b,N_0}(F)\}_{b\in\NN}
\]
在 pro-意义下同构。
更强地，对任意共终深度子序列 \(b_1<b_2<\cdots\to\infty\)，都有
\[
\varprojlim_{b,N}\mathscr C_{b,N}(F)
\cong
\varprojlim_{j}\mathscr C_{b_j,N_0}(F).
\]
因此，在一致维数上界制度下，无穷多 PSD 量词与无穷多深度量词同时坍缩为一个有限核心 \(N_0\) 与任意共终深度链。
\end{theorem}
\begin{proof}
由定理 \ref{thm:xi-oracle-collapse-toeplitz-psd-finite-truncation}，存在统一阈值 \(N_0\le 2D\)，使得当 \(N\ge N_0\) 时，“对所有 Toeplitz 截断 PSD”的无穷约束与“仅至 \(N_0\) 阶 PSD”的有限核验严格等价；更高阶截断均由 \(N_0\) 阶数据经合同压缩得到。
故 \(N\)-方向在 \(N_0\) 后不再产生新信息。

另一方面，由命题 \ref{prop:cdim-dyadic-cofinal-sparsification}，任意共终深度子序列都给出同一逆极限点对象，因此 \(b\)-方向对共终抽稀不敏感。
两者合并，即得整个双参数逆系统与固定核心 \(N_0\) 的单参数深度链在 pro-意义下同构，并给出所述逆极限同构。证毕。
\end{proof}

\begin{theorem}[window-\texorpdfstring{$6$}{6} 中精确反演与异常完备读出的信息分离]\label{thm:conclusion-window6-inversion-anomaly-information-separation}
在二进制区间折叠
\[
\Fold_6^{\mathrm{bin}}:\Omega_6\to X_6
\]
上，考虑两类辅助寄存器任务：
\begin{enumerate}
  \item \textbf{精确反演任务：}给定 \(w=\Fold_6^{\mathrm{bin}}(\omega)\) 及辅助寄存器内容 \(r\)，可唯一恢复微态 \(\omega\)；
  \item \textbf{异常完备任务：}辅助寄存器能够注入地编码
  \[
  \bigl(G_6^{\mathrm{bin}}\bigr)^{\mathrm{ab}}
  \cong
  (\ZZ/2\ZZ)^{21}.
  \]
\end{enumerate}
则最小所需二进制位数分别为
\[
b_{\mathrm{inv}}=2,
\qquad
b_{\mathrm{anom}}=21.
\]
从而任何试图同时完成两项任务的统一辅助寄存器都必满足
\[
b_{\mathrm{joint}}\ge 21,
\qquad
b_{\mathrm{joint}}-b_{\mathrm{inv}}\ge 19.
\]
因此，在 window-\(6\) 中，精确恢复微态与完整保留异常账本并不处于同一资源尺度：前者是局部纤维索引问题，后者则是全局规范交换化的独立方向计数。
\end{theorem}
\begin{proof}
由推论 \ref{cor:fold-bin-m6-fiber-hist}，window-\(6\) 的纤维直方图为
\[
2:8,\qquad 3:4,\qquad 4:9,
\]
故最大纤维大小为 \(4\)。
若辅助寄存器仅用于精确反演，则给定 \(w\) 后要区分同一纤维内的全部微态，必要条件为 \(2^b\ge 4\)，即 \(b\ge 2\)。
反之，对每条纤维任选一组局部编号
\[
\iota_w:\Fold_6^{\mathrm{bin}}{}^{-1}(w)\hookrightarrow \{0,1\}^2,
\]
即可用 \(2\) 位完成精确反演，故 \(b_{\mathrm{inv}}=2\)。

另一方面，由推论 \ref{cor:fold-bin-gauge-m6}，
\[
\bigl(G_6^{\mathrm{bin}}\bigr)^{\mathrm{ab}}\cong (\ZZ/2\ZZ)^{21},
\]
其元素总数为 \(2^{21}\)。
若用 \(b\) 位二进制寄存器注入地区分全部异常类，则必须 \(2^b\ge 2^{21}\)，即 \(b\ge 21\)；故 \(b_{\mathrm{anom}}=21\)。
联合任务至少与异常完备任务一样困难，从而
\[
b_{\mathrm{joint}}\ge 21.
\]
证毕。
\end{proof}

综上，有限分辨率量子扇区的出现本身还不是终点。真正控制该框架能否闭合为“无需额外物理公理”的，是四类更深的纯数学条件：调和异常是否可灭、低阶 Walsh 缺陷是否消失、非可缩基态是否保持主导、以及无限证书链是否已经坍缩为有限核心。它们分别控制 strictification、可见谱、零温拓扑与离线可判定性；而 window-\(6\) 的 \(2\) 比特/ \(21\) 比特裂分则说明，即便在最小非平凡锚点上，局部微态反演与全局异常账本保真也已经严格分属于两个不可混同的资源尺度。

\paragraph{屏幕余秩、Gödel 病态分裂、张量 Joukowsky 盲性与碰撞扭量}
在观察者时空 forcing 骨架、部分屏幕观测的相对顶维同调刻画、边界 Gödelization 的可计算可逆性以及 Joukowsky 椭圆族的解析矩盲性都已经建立之后，还可把这几条原本分散的接口再压缩成一条更高阶的结构链：屏幕补全成本实际上是一个 graphic matroid 的余秩函数；边界 Gödel 全息在算法复杂度上是充分的，却在度量意义上指数病态；Joukowsky 乘积族的全部全纯矩对尺度参数完全失明，而最小解除盲性的外置维数恰等于参数维数；最后，碰撞型 \(\mathsf{NULL}\) 不是布尔空值，而是由屏幕核所张成的自由阿贝尔扭量。

\begin{theorem}[屏幕审计成本的图拟阵余秩表示与超模性]\label{thm:conclusion-screen-graphic-matroid-corank-supermodularity}
固定分辨率 \(m\) 与维数 \(n\)。设 \(\Gamma=(V,E)\) 为定理 \ref{thm:spg-screen-kernel-connected-components} 中由全部顶维胞腔与外部顶点 \(\infty\) 组成的完整屏幕图。对任意可见屏幕 \(S\subseteq E\)，记 \(\Gamma_S:=(V,S)\)，并定义其最小审计补全成本
\[
\lambda(S):=\min_{(R,r^S)}\cdim(R).
\]
再记 \(r_{\mathrm{gr}}\) 为 \(\Gamma\) 的 graphic matroid 在边集 \(E\) 上的秩函数。则有精确公式
\[
\lambda(S)=c(\Gamma_S)-1=|V|-1-r_{\mathrm{gr}}(S).
\]
因而：
\begin{enumerate}
  \item \(\lambda(\varnothing)=|V|-1\)，\(\lambda(E)=0\)；
  \item 若 \(S\subseteq T\subseteq E\)，则 \(\lambda(T)\le \lambda(S)\)；
  \item \(\lambda\) 为超模函数，即对任意 \(S,T\subseteq E\) 都有
        \[
        \lambda(S\cap T)+\lambda(S\cup T)\ge \lambda(S)+\lambda(T).
        \]
\end{enumerate}
\end{theorem}

\begin{proof}
由推论 \ref{cor:spg-screen-injectivity-and-audit-cost-components}，
\[
\lambda(S)=c(\Gamma_S)-1.
\]
另一方面，graphic matroid 的秩公式给出
\[
r_{\mathrm{gr}}(S)=|V|-c(\Gamma_S),
\]
故
\[
\lambda(S)=|V|-1-r_{\mathrm{gr}}(S).
\]
第 \(1\) 条对应于 \(\Gamma_\varnothing\) 有 \(|V|\) 个分量而 \(\Gamma_E=\Gamma\) 连通。第 \(2\) 条由 \(r_{\mathrm{gr}}\) 的单调性立得。第 \(3\) 条则由拟阵秩函数的次模性
\[
r_{\mathrm{gr}}(S)+r_{\mathrm{gr}}(T)\ge r_{\mathrm{gr}}(S\cap T)+r_{\mathrm{gr}}(S\cup T)
\]
代入上述余秩表示后立即得到。证毕。
\end{proof}
\leanverified{paper\_conclusion\_screen\_graphic\_matroid\_corank\_supermodularity}

\begin{theorem}[延迟证书包的相对秩削减律与闭包判据]\label{thm:conclusion-screen-rank-gain-closure-diminishing-returns}
\leanverified{paper\_conclusion\_screen\_rank\_gain\_closure\_diminishing\_returns}
沿用定理 \ref{thm:conclusion-screen-graphic-matroid-corank-supermodularity} 的记号。对任意 \(S\subseteq E\) 与补充屏幕包 \(A\subseteq E\setminus S\)，定义补全收益
\[
\Delta_A\lambda(S):=\lambda(S)-\lambda(S\cup A).
\]
则有精确公式
\[
\Delta_A\lambda(S)=r_{\mathrm{gr}}(S\cup A)-r_{\mathrm{gr}}(S).
\]
特别地：
\begin{enumerate}
  \item 对单边 \(e\in E\setminus S\)，有
        \[
        \Delta_{\{e\}}\lambda(S)\in\{0,1\};
        \]
  \item \(\Delta_{\{e\}}\lambda(S)=1\) 当且仅当 \(e\) 的两个端点位于 \(\Gamma_S\) 的不同连通分量；
  \item 定义
        \[
        \operatorname{cl}_{\lambda}(S):=\{e\in E:\lambda(S\cup\{e\})=\lambda(S)\},
        \]
        则 \(\operatorname{cl}_{\lambda}(S)\) 恰为 graphic matroid 对 \(S\) 的闭包；
  \item 若 \(S\subseteq T\subseteq E\) 且 \(e\notin T\)，则有边际收益递减律
        \[
        \Delta_{\{e\}}\lambda(S)\ge \Delta_{\{e\}}\lambda(T).
        \]
\end{enumerate}
\end{theorem}

\begin{proof}
由定理 \ref{thm:conclusion-screen-graphic-matroid-corank-supermodularity}，
\[
\lambda(S)=|V|-1-r_{\mathrm{gr}}(S),
\]
故
\[
\Delta_A\lambda(S)=r_{\mathrm{gr}}(S\cup A)-r_{\mathrm{gr}}(S).
\]
这就给出了主公式。

对单边情形，graphic matroid 的秩增量至多为 \(1\)，故只能取 \(0\) 或 \(1\)。并且增量为 \(1\) 当且仅当 \(e\) 不在 \(S\) 的闭包中；对图拟阵来说，这等价于 \(e\) 连接了 \(\Gamma_S\) 的两个不同连通分量。第 \(3\) 条只是闭包定义的直接翻译。最后，第 \(4\) 条由拟阵秩的次模性等价推出，即边际秩增量沿包含关系单调递减。证毕。
\end{proof}

\begin{theorem}[算法充分性与度量稳定性的指数分离]\label{thm:conclusion-godel-algorithmic-sufficiency-metric-instability-separation}
\leanverified{paper\_conclusion\_godel\_algorithmic\_sufficiency\_metric\_instability\_separation}
固定 \(n\ge 1\)。对每个分辨率 \(m\)，存在一族边界 Gödel 编码
\[
\Phi_m:C_n(\mathcal Q_m;\FF_2)\to\NN
\]
满足：
\begin{enumerate}
  \item \(\Phi_m\) 在其像上可计算可逆；
  \item 存在常数 \(C>0\)，与 \(m\) 及输入无关，使得对任意 \(u\in C_n(\mathcal Q_m;\FF_2)\)，都有
        \[
        \bigl|K(\Phi_m(u)\mid n,m)-K(u\mid n,m)\bigr|\le C,
        \qquad
        K(u\mid \Phi_m(u),n,m)\le C.
        \]
\end{enumerate}
然而，对任意线性左逆
\[
L_m:C_{n-1}(\mathcal Q_m;\RR)\to C_n(\mathcal Q_m;\RR),
\qquad
L_m\partial_n=\mathrm{Id},
\]
都必有算子范数下界
\[
\|L_m\|_{2\to 2}\ge \frac{1}{\sqrt{2n}}\,2^{m/2}.
\]
因此，不存在一族同时满足“精确无损”“可计算反演”与“多项式条件数”的线性全息反演器。
\end{theorem}

\begin{proof}
前半部正是定理 \ref{thm:spg-stokes-godel-algorithmic-holographic-completeness} 的结论：边界 Gödelization 在固定 \((m,n)\) 下给出一个可计算双向字典，并把 Kolmogorov 复杂度守恒到 \(O(1)\) 常数。

后半部则由定理 \ref{thm:spg-dyadic-top-dimensional-holographic-inversion-exponential-ill-conditioning} 直接给出。该定理表明，任意线性左逆都满足
\[
\|L_m\|_{2\to 2}\ge \frac{1}{\sqrt{2n}}\,2^{m/2}.
\]
若反设存在 \(\|L_m\|_{2\to 2}\le m^C\) 的多项式条件数族，则与此指数下界矛盾。故算法上的可逆充分性与度量上的稳定可逆性严格分离。证毕。
\end{proof}

\begin{corollary}[噪声阈值下的统一线性重构不可能性]\label{cor:conclusion-godel-noise-threshold-no-uniform-linear-reconstruction}
\leanverified{paper\_conclusion\_godel\_noise\_threshold\_no\_uniform\_linear\_reconstruction}
在定理 \ref{thm:conclusion-godel-algorithmic-sufficiency-metric-instability-separation} 的设定下，对任意固定目标精度 \(\delta>0\)，若观测噪声满足
\[
\|\eta_m\|_2\gg 2^{-m/2},
\]
则不存在统一的线性重构规则能够在全部分辨率上保持
\[
\|L_m\eta_m\|_2\le \delta.
\]
\end{corollary}

\begin{proof}
由定理 \ref{thm:conclusion-godel-algorithmic-sufficiency-metric-instability-separation}，
\[
\|L_m\eta_m\|_2\ge \frac{1}{\sqrt{2n}}\,2^{m/2}\|\eta_m\|_2.
\]
若 \(\|\eta_m\|_2\gg 2^{-m/2}\)，则右端不能在 \(m\to\infty\) 下被统一压到任意固定的 \(\delta>0\) 之内。故此类统一线性重构不可能存在。证毕。
\end{proof}

\begin{theorem}[边界 Gödel 精确可逆性与实线性稳定不可逆性的严格分裂]\label{thm:conclusion-godel-arithmetic-exact-vs-linear-stable-separation}
\leanverified{paper\_conclusion\_godel\_arithmetic\_exact\_vs\_linear\_stable\_separation}
固定 \(m,n\ge 1\)，记 \(\mathcal Q_m\) 为 \(I^n=[0,1]^n\) 的 \(m\)-dyadic 立方剖分，并令
\[
\mathcal B_{m,n}:=\partial_n\bigl(C_n(\mathcal Q_m;\{0,1\})\bigr)
\subset C_{n-1}(\mathcal Q_m;\ZZ).
\]
则同时成立：
\begin{enumerate}
  \item 存在一个非线性、算术型、精确译码器
        \[
        D_m:\bigl\{H_m(U):\ U\subset I^n\ \text{为 dyadic polycube}\bigr\}
        \longrightarrow
        C_n(\mathcal Q_m;\{0,1\}),
        \]
        它由边界 Gödel 整数的素因子分解与树上传播组成，并满足
        \[
        D_m\bigl(H_m(U)\bigr)=[U].
        \]
        该译码可在与 dyadic 邻接图规模线性同阶的时间内完成；
  \item 对任意实线性算子
        \[
        L_m:C_{n-1}(\mathcal Q_m;\RR)\longrightarrow C_n(\mathcal Q_m;\RR),
        \]
        若对所有
        \[
        x\in C_n(\mathcal Q_m;\{0,1\})
        \]
        都满足
        \[
        L_m(\partial_n x)=x,
        \]
        则必有
        \[
        \|L_m\|_{2\to 2}\ge \frac{2^{m/2}}{\sqrt{2n}}.
        \]
\end{enumerate}
因此，在 dyadic 全息体系中，精确算术可逆性与实线性数值稳定性并不兼容：前者确实存在，而后者不可能具有分辨率无关的条件数。
\end{theorem}
\begin{proof}
对第一条，由定理 \ref{thm:spg-boundary-godelization-holographic-dictionary}，边界 Gödel 整数 \(H_m(U)\) 的素因子分解唯一确定边界链 \(\partial_n[U]\)。随后应用定理 \ref{thm:spg-linear-time-bulk-decode-from-boundary} 的树上传播译码，即可由该边界链在线性时间内恢复唯一的顶维 \(0/1\)-链 \([U]\)。这便定义了所需的算术型精确译码器 \(D_m\)。

对第二条，注意 \(\{0,1\}\)-链基向量包含全部顶维胞腔基向量 \(e_Q\)。由假设，
\[
L_m(\partial_n e_Q)=e_Q
\qquad
(\forall\,Q\in \mathcal Q_m^{(n)}).
\]
由于 \(L_m\) 与 \(\partial_n\) 都是实线性的，这立即推出
\[
L_m\partial_n=\mathrm{Id}
\qquad\text{on } C_n(\mathcal Q_m;\RR).
\]
故 \(L_m\) 是 \(\partial_n\) 的一个实线性左逆。于是可直接应用定理 \ref{thm:spg-dyadic-top-dimensional-holographic-inversion-exponential-ill-conditioning}，得到
\[
\|L_m\|_{2\to 2}\ge \frac{2^{m/2}}{\sqrt{2n}}.
\]
证毕。
\end{proof}

\leanverified{paper\_conclusion\_joukowsky\_product\_analytic\_blindness\_minimal\_radial\_sufficiency}
\begin{theorem}[Joukowsky 乘积族的张量盲性与最小径向充分性]\label{thm:conclusion-joukowsky-product-analytic-blindness-minimal-radial-sufficiency}
对任意整数 \(k\ge 1\)，取参数
\[
\mathbf r=(r_1,\dots,r_k)\in(1,\infty)^k,
\qquad
\mu_{\mathbf r}:=\mu_{r_1}\otimes\cdots\otimes\mu_{r_k},
\]
其中 \(\mu_{r_j}\) 为一维 Joukowsky 椭圆推前测度。则：
\begin{enumerate}
  \item 对任意全纯多项式
        \[
        p\in\CC[w_1,\dots,w_k],
        \]
        其矩泛函
        \[
        \varphi_{\mathbf r}(p):=\int_{\CC^k}p\,\dd\mu_{\mathbf r}
        \]
        与 \(\mathbf r\) 无关；
  \item 映射
        \[
        \mathbf r\longmapsto \mathbf R_2(\mathbf r):=
        \bigl(r_1^2+r_1^{-2},\dots,r_k^2+r_k^{-2}\bigr)
        \]
        在 \((1,\infty)^k\) 上为单射，且每个坐标可由
        \[
        r_j=
        \left(
        \frac{R_{2,j}+\sqrt{R_{2,j}^2-4}}{2}
        \right)^{1/2}
        \]
        唯一恢复；
  \item 任意企图在全部全纯矩之外仅再加入 \(d<k\) 个连续实标量观测而分离整个参数族的方案都不可能；
  \item 因而，恰需 \(k\) 个独立径向标量观测，方可极小地解除该盲性。
\end{enumerate}
\end{theorem}

\begin{proof}
若
\[
p(w)=w_1^{a_1}\cdots w_k^{a_k},
\]
则
\[
\varphi_{\mathbf r}(p)
=
\prod_{j=1}^{k}\int_{\CC} w_j^{a_j}\,\dd\mu_{r_j}(w_j).
\]
由推论 \ref{cor:prime-register-joukowsky-analytic-moments}，每个一维解析矩都与 \(r_j\) 无关，故乘积亦与 \(\mathbf r\) 无关；线性延拓即得第 \(1\) 条。

第 \(2\) 条逐坐标由定理 \ref{thm:conclusion-analytic-holography-inversion-rigidity} 给出：一维参数 \(r_j\) 可由二阶径向矩 \(R_{2,j}=r_j^2+r_j^{-2}\) 唯一恢复。因而 \(\mathbf R_2\) 在 \((1,\infty)^k\) 上单射。

若存在 \(d<k\) 个连续实标量观测 \(F_1,\dots,F_d\) 能分离整个参数族，则得到连续单射
\[
F:(1,\infty)^k\to\RR^d,
\qquad
F(\mathbf r):=(F_1(\mathbf r),\dots,F_d(\mathbf r)).
\]
但 \((1,\infty)^k\) 是 \(k\)-维流形，而连续单射到 \(\RR^d\)（\(d<k\)）与拓扑维数单调性矛盾。故第 \(3\) 条成立。

最后，第 \(2\) 条已经构造出 \(k\) 个连续径向观测 \((R_{2,1},\dots,R_{2,k})\) 并实现了参数的唯一反演，故第 \(4\) 条成立。证毕。
\end{proof}

\begin{corollary}[乘积 Joukowsky 族的最小分离 \texorpdfstring{$*$}{*}-代数扩张]\label{cor:conclusion-joukowsky-product-minimal-star-algebra-resolution}
\leanverified{paper\_conclusion\_joukowsky\_product\_minimal\_star\_algebra\_resolution}
令
\[
\mathcal H_k:=\CC[w_1,\dots,w_k].
\]
考虑所有有限生成交换 \(^*\)-代数
\[
\mathcal H_k\subseteq \mathcal A\subseteq C(\CC^k)
\]
并要求状态族 \(\{\varphi_{\mathbf r}\}_{\mathbf r\in(1,\infty)^k}\) 在 \(\mathcal A\) 上可分离。则：
\begin{enumerate}
  \item 任意这样的 \(\mathcal A\) 至少需要 \(k\) 个额外自伴生成元；
  \item 代数
        \[
        \mathcal A_{\min}:=\CC\bigl[w_1,\dots,w_k,\ |w_1|^2,\dots,|w_k|^2\bigr]
        \]
        达到该下界；
  \item 因而 \(\mathcal A_{\min}\) 是乘积 Joukowsky 参数族的最小分离扩张。
\end{enumerate}
\end{corollary}

\begin{proof}
全纯部分 \(\mathcal H_k\) 按定理 \ref{thm:conclusion-joukowsky-product-analytic-blindness-minimal-radial-sufficiency} 第 \(1\) 条对参数完全失明，故全部分离信息只能由额外自伴生成元提供。若其数目少于 \(k\)，则它们给出到某个 \(\RR^d\)（\(d<k\)）的连续分离映射，与定理 \ref{thm:conclusion-joukowsky-product-analytic-blindness-minimal-radial-sufficiency} 第 \(3\) 条矛盾。

另一方面，\((|w_1|^2,\dots,|w_k|^2)\) 在状态 \(\mu_{\mathbf r}\) 下的期望正是 \(\mathbf R_2(\mathbf r)\)，并由定理 \ref{thm:conclusion-joukowsky-product-analytic-blindness-minimal-radial-sufficiency} 第 \(2\) 条唯一恢复 \(\mathbf r\)。故 \(\mathcal A_{\min}\) 确实分离该参数族，并达到最小生成元数。证毕。
\end{proof}

\begin{theorem}[碰撞型 \texorpdfstring{$\mathsf{NULL}$}{NULL} 的拓扑化：完成纤维的同调扭量]\label{thm:conclusion-collision-null-topological-torsor}
\leanverified{paper\_conclusion\_collision\_null\_topological\_torsor}
固定分辨率 \(m,n\) 与部分屏幕 \(S\subseteq E\)。设
\[
f_S:C_n(\mathcal Q_m;\ZZ)\to \ZZ^S
\]
为定理 \ref{thm:spg-partial-boundary-screen-kernel-relative-homology} 的部分边界观测算子。对任意可见边界数据 \(y\in\operatorname{im}f_S\)，定义其完成纤维
\[
\mathcal C_S(y):=f_S^{-1}(y).
\]
则：
\begin{enumerate}
  \item 若 \(\mathcal C_S(y)\neq\varnothing\)，则 \(\mathcal C_S(y)\) 是 \(\ker f_S\) 的主齐次空间；
  \item 存在自然同构
        \[
        \ker f_S\cong H_n(\mathcal Q_m,A_{\neg S};\ZZ)\cong \ZZ^{\,c(\Gamma_S)-1},
        \]
        因而 \(\mathcal C_S(y)\) 是一个自由阿贝尔扭量；
  \item 在推论 \ref{cor:null-trichotomy-local-section} 的三分解语言中，一旦语义空值与协议空值都已排除，剩余的碰撞型 \(\mathsf{NULL}\) 正对应于该扭量的非平凡性；
  \item 碰撞型 \(\mathsf{NULL}\) 消失当且仅当 \(\Gamma_S\) 连通。
\end{enumerate}
\end{theorem}

\begin{proof}
若 \(\mathcal C_S(y)\neq\varnothing\)，取任意 \(x_0\in\mathcal C_S(y)\)。则
\[
\mathcal C_S(y)=x_0+\ker f_S,
\]
故 \(\mathcal C_S(y)\) 是 \(\ker f_S\) 的主齐次空间。第一条得证。

第二条由定理 \ref{thm:spg-partial-boundary-screen-kernel-relative-homology} 与定理 \ref{thm:spg-screen-kernel-connected-components} 复合即得：
\[
\ker f_S\cong H_n(\mathcal Q_m,A_{\neg S};\ZZ)\cong \ZZ^{\,c(\Gamma_S)-1}.
\]
于是 \(\mathcal C_S(y)\) 不是单纯的“多解集”，而是一个自由阿贝尔扭量。

对第 \(3\) 条，推论 \ref{cor:null-trichotomy-local-section} 已把 \(\mathsf{NULL}\) 分为语义、协议与碰撞三类。当前在 \(y\in\operatorname{im}f_S\) 的前提下，语义未定义与协议不相容都已被排除，于是剩余的失败只能来自同一可见数据允许多个 bulk completion；这恰等价于 \(\mathcal C_S(y)\) 作为 torsor 非平凡，亦即 \(\ker f_S\neq 0\)。

最后，由 \(\ker f_S\cong \ZZ^{\,c(\Gamma_S)-1}\)，碰撞型 \(\mathsf{NULL}\) 消失当且仅当 \(\ker f_S=0\)，亦即当且仅当 \(c(\Gamma_S)=1\)，等价于 \(\Gamma_S\) 连通。证毕。
\end{proof}

\begin{theorem}[延迟补全的碰撞扭量商与秩削减律]\label{thm:conclusion-delayed-completion-collision-torsor-rank-law}
\leanverified{paper\_conclusion\_delayed\_completion\_collision\_torsor\_rank\_law}
沿用定理 \ref{thm:conclusion-screen-graphic-matroid-corank-supermodularity} 与定理 \ref{thm:conclusion-collision-null-topological-torsor} 的记号。设 \(S\subseteq T\subseteq E\)，且 \(y_T\in\operatorname{im}f_T\) 在限制映射下投影为 \(y_S\in\operatorname{im}f_S\)。则存在自然嵌入
\[
\mathcal C_T(y_T)\hookrightarrow \mathcal C_S(y_S),
\]
并且有短正合列
\[
0\longrightarrow \ker f_T\longrightarrow \ker f_S\longrightarrow \ker f_S/\ker f_T\longrightarrow 0,
\qquad
\ker f_S/\ker f_T\cong \ZZ^{\,\lambda(S)-\lambda(T)}.
\]
特别地，若 \(T=S\cup A\)，则
\[
\operatorname{rank}\bigl(\ker f_S/\ker f_T\bigr)
=
\lambda(S)-\lambda(T)
=
r_{\mathrm{gr}}(T)-r_{\mathrm{gr}}(S).
\]
换言之，延迟送达的证书包 \(A\) 所消灭的并非某种模糊“未知性”，而恰好是
\[
r_{\mathrm{gr}}(T)-r_{\mathrm{gr}}(S)
\]
个彼此独立的碰撞同调方向。
\end{theorem}

\begin{proof}
因为 \(S\subseteq T\)，算子 \(f_T\) 比 \(f_S\) 观测更多坐标，故
\[
\ker f_T\subseteq \ker f_S.
\]
若 \(x\in\mathcal C_T(y_T)\)，则 \(f_T(x)=y_T\) 立即推出 \(f_S(x)=y_S\)，故存在自然嵌入
\[
\mathcal C_T(y_T)\hookrightarrow \mathcal C_S(y_S).
\]

两核都是自由阿贝尔群，其商 \(\ker f_S/\ker f_T\) 亦自由阿贝尔；其秩等于秩差：
\[
\operatorname{rank}(\ker f_S/\ker f_T)
=
\operatorname{rank}\ker f_S-\operatorname{rank}\ker f_T.
\]
由推论 \ref{cor:spg-partial-boundary-min-audit-cost-kernel-rank} 与定理 \ref{thm:conclusion-screen-graphic-matroid-corank-supermodularity}，
\[
\operatorname{rank}\ker f_S=\lambda(S),
\qquad
\operatorname{rank}\ker f_T=\lambda(T).
\]
故
\[
\ker f_S/\ker f_T\cong \ZZ^{\,\lambda(S)-\lambda(T)}.
\]
最后把定理 \ref{thm:conclusion-screen-graphic-matroid-corank-supermodularity} 中
\[
\lambda(U)=|V|-1-r_{\mathrm{gr}}(U)
\]
分别代入 \(U=S,T\)，便得
\[
\lambda(S)-\lambda(T)=r_{\mathrm{gr}}(T)-r_{\mathrm{gr}}(S).
\]
证毕。
\end{proof}

\begin{theorem}[离散同调歧义与连续尺度盲性的加性最小外置定理]\label{thm:conclusion-screen-joukowsky-additive-minimal-externalization}
\leanverified{paper\_conclusion\_screen\_joukowsky\_additive\_minimal\_externalization}
固定一个可接受屏幕 \(S\) 与一条非空纤维 \(f_S^{-1}(y)\)。设
\[
\beta:=\beta_n(\mathcal Q_m,A_{\neg S};\FF_2),
\qquad
\mathbf r=(r_1,\dots,r_k)\in (1,\infty)^k.
\]
考虑复合隐藏状态空间
\[
\mathcal H_{S,k}(y):=f_S^{-1}(y)\times (1,\infty)^k.
\]
称一组可分离的分离式审计架构由一对映射
\[
a:f_S^{-1}(y)\to \{0,1\}^b,
\qquad
b:(1,\infty)^k\to \RR^d
\]
给出，若乘积映射
\[
\Psi(x,\mathbf r):=(a(x),b(\mathbf r))
\]
在 \(\mathcal H_{S,k}(y)\) 上单射。则：
\begin{enumerate}
  \item 任意此类架构都必须满足
        \[
        b\ge \beta,
        \qquad
        d\ge k;
        \]
  \item 该下界是锐的：取 \(b=\beta\)、\(d=k\) 可以实现单射；
  \item 若进一步要求 \(b(\mathbf r)\) 完全因子化经由 \(k\)-重 Joukowsky 乘积族的全纯多项式矩，则对任意 \(b,d\) 都不存在可分离架构。
\end{enumerate}
\end{theorem}

\begin{proof}
先固定任意 \(\mathbf r_0\in (1,\infty)^k\)。若 \(\Psi\) 在 \(\mathcal H_{S,k}(y)\) 上单射，则
\[
x\longmapsto \Psi(x,\mathbf r_0)=(a(x),b(\mathbf r_0))
\]
在有限集合 \(f_S^{-1}(y)\) 上也必须单射，故 \(a\) 本身必须单射。由定理 \ref{thm:conclusion-screen-exact-binary-audit-homological-bit-bound}，这立即推出
\[
b\ge \beta.
\]

再固定任意 \(x_0\in f_S^{-1}(y)\)。由同样的单射性，
\[
\mathbf r\longmapsto \Psi(x_0,\mathbf r)=(a(x_0),b(\mathbf r))
\]
在 \((1,\infty)^k\) 上必须单射，因此 \(b:(1,\infty)^k\to \RR^d\) 为连续单射。由于 \((1,\infty)^k\) 是 \(k\)-维流形，按拓扑维数单调性，不存在到 \(\RR^d\) 的连续单射当 \(d<k\)。故
\[
d\ge k.
\]

这两个下界都是可达的。离散部分令 \(a\) 取定理 \ref{thm:conclusion-screen-exact-binary-audit-homological-bit-bound} 中的最优 \(\beta\)-比特核坐标编码；连续部分取
\[
b(\mathbf r):=\bigl(r_1^2+r_1^{-2},\dots,r_k^2+r_k^{-2}\bigr).
\]
由推论 \ref{cor:conclusion-joukowsky-product-minimal-star-algebra-resolution} 背后的逐坐标径向反演公式，每个 \(r_j\) 都可由对应分量唯一恢复，故该 \(\Psi\) 的确是单射。

最后，若 \(b(\mathbf r)\) 完全因子化经由 \(k\)-重 Joukowsky 乘积族的全纯多项式矩，则由定理 \ref{thm:conclusion-joukowsky-product-analytic-blindness-minimal-radial-sufficiency}，所有这类矩对 \(\mathbf r\) 都与 \(\mathbf r\) 无关，于是 \(b\) 必为常值映射，不可能在 \((1,\infty)^k\) 上单射。故此类可分离架构不存在。证毕。
\end{proof}

\endinput

\paragraph{逆极限地址上的有限前缀决定性与解析尺度的离散退场}
在 dyadic 证书链的逆极限地址空间、前缀超度量与 Joukowsky 乘积族的解析矩盲性都已经建立之后，还可把两条原本分立的链条压缩成一个更紧的离散稳定性原则：任意连续的离散值稳定可观测量都必在某个有限前缀层级上终止，而一旦其连续尺度侧又坚持纯全纯矩口径，则全部连续参数都会整体退场，最终剩下的可见信息只能是一条有限深度的 dyadic 地址函数。

\begin{theorem}[逆极限地址上的有限前缀决定性]\label{thm:conclusion-inverse-limit-address-finite-prefix-determinacy}
\leanverified{paper\_conclusion\_inverse\_limit\_address\_finite\_prefix\_determinacy}
设
\[
A_\infty=\varprojlim (A_b,\pi_{b'\to b})
\]
为 dyadic 证书链的逆极限地址空间，并赋以前缀超度量 \(d_{\mathrm{dy}}\)。令 \(Y\) 为离散空间。则任意连续映射
\[
v:A_\infty\longrightarrow Y
\]
都存在某个有限层级 \(b\) 以及映射
\[
v_b:A_b\longrightarrow Y
\]
使得
\[
v=v_b\circ \pi_b.
\]
换言之，任意连续的离散值稳定可观测量皆由有限前缀完全决定。
\end{theorem}

\begin{proof}
由于 \(A_\infty\) 是紧空间而 \(Y\) 离散，其连续像 \(v(A_\infty)\) 必为有限集。对每个 \(y\in v(A_\infty)\)，纤维
\[
U_y:=v^{-1}(\{y\})
\]
是 \(A_\infty\) 中的紧开集。

由前缀超度量的构造，柱集
\[
[a]_b:=\{a'\in A_\infty:\pi_b(a')=\pi_b(a)\}
\]
构成 \(A_\infty\) 的 clopen 基。故对每个 \(U_y\)，可取有限柱覆盖
\[
U_y=\bigcup_{\ell=1}^{N_y}[a_{y,\ell}]_{b_{y,\ell}}.
\]
令
\[
b:=\max_{y,\ell} b_{y,\ell}.
\]
把每个较浅柱进一步细分为深度 \(b\) 的柱后，可得 \(A_\infty\) 被若干深度 \(b\) 的柱集分割，而 \(v\) 在每一柱集上恒为常值。于是可定义
\[
v_b(\alpha):=v(a)
\qquad
(\pi_b(a)=\alpha),
\]
该定义与代表元 \(a\) 无关，并满足
\[
v=v_b\circ \pi_b.
\]
证毕。
\end{proof}

\begin{corollary}[全纯矩盲性下的地址--解析分离塌缩]\label{cor:conclusion-address-analytic-separation-collapse}
设 \(k\ge 1\)，并令
\[
X:=A_\infty\times (1,\infty)^k,
\]
其中第一因子为 dyadic 逆极限地址空间，第二因子为 \(k\)-参数 Joukowsky 尺度空间。设 \(Y\) 为离散空间，且
\[
v:X\longrightarrow Y
\]
连续。若对每个固定地址 \(a\in A_\infty\)，切片映射
\[
\mathbf r\longmapsto v(a,\mathbf r)
\]
完全因子化经由 \(k\)-重 Joukowsky 乘积族的全纯多项式矩，则存在某一有限层级 \(b\) 与映射
\[
\widetilde v_b:A_b\longrightarrow Y
\]
使得
\[
v(a,\mathbf r)=\widetilde v_b(\pi_b(a))
\qquad
(\forall (a,\mathbf r)\in X).
\]
特别地，此时 \(v\) 对全部连续尺度参数 \(\mathbf r\) 完全盲化；其全部可见内容仅存于有限深度的 dyadic 地址前缀。
\end{corollary}

\begin{proof}
由定理 \ref{thm:conclusion-joukowsky-product-analytic-blindness-minimal-radial-sufficiency}，\(k\)-重 Joukowsky 乘积族的一切全纯多项式矩都与 \(\mathbf r\) 无关。故对每个固定 \(a\)，既然切片 \(\mathbf r\mapsto v(a,\mathbf r)\) 因子化经由这些矩，它便只能是常值。于是存在映射
\[
u:A_\infty\to Y
\]
使得
\[
v(a,\mathbf r)=u(a).
\]
由 \(v\) 的连续性可知 \(u\) 也连续。再应用定理 \ref{thm:conclusion-inverse-limit-address-finite-prefix-determinacy}，便存在有限层级 \(b\) 与映射 \(\widetilde v_b:A_b\to Y\)，满足
\[
u=\widetilde v_b\circ \pi_b.
\]
把它代回上式即得
\[
v(a,\mathbf r)=\widetilde v_b(\pi_b(a)).
\]
证毕。
\end{proof}
\leanverified{paper\_conclusion\_address\_analytic\_separation\_collapse}

\endinput

\paragraph{cut-space 全息、精确屏幕亏损与边界算术坐标}
在 dyadic 体--界全息、部分屏幕的相对顶维同调刻画、边界 Gödel 码的可计算可逆性以及 dyadic 边界像的 LDPC 结构都已经建立之后，还可把这几条链进一步压缩为一组纯组合与纯算术的终态结论：部分可见性的全部障碍完全退化为屏幕图上的 cut/cycle 理论；典型 bulk 在部分屏幕下真正损失的条件信息精确等于屏幕缺口 \(c(\Gamma_S)-1\)；固定分辨率的边界像不仅构成正码率、常最小距离的精确全息码，而且边界 Gödel 整数在单一平方自由整型坐标内已经压缩了全部最小完备 tensor moment box，而任何坚持各向同性总次数截断的矩系都必须支付 \(2^m\) 级的次数门槛。

\begin{theorem}[partial screen trace 的 cut-space 识别]\label{thm:conclusion-partial-screen-trace-cut-space-identification}
固定 \(m,n\ge 1\)。沿用定理 \ref{thm:spg-partial-boundary-screen-kernel-relative-homology} 的部分边界观测算子与定理 \ref{thm:spg-screen-kernel-connected-components} 的屏幕图 \(\Gamma_S\) 记号，取系数域 \(\FF_2\)，并把任意
\[
u\in C_n(\mathcal Q_m;\FF_2)
\]
识别为 \(\Gamma_S\) 的顶点标记函数
\[
\widetilde u:V(\Gamma_S)\to \FF_2,
\qquad
\widetilde u(\infty):=0.
\]
再记
\[
d_{\Gamma_S}:C^0(\Gamma_S;\FF_2)\to C^1(\Gamma_S;\FF_2)
\]
为图 coboundary。则部分边界观测算子
\[
f_S:C_n(\mathcal Q_m;\FF_2)\to \FF_2^{\,S}
\]
与图 coboundary 精确等同：
\[
f_S(u)=d_{\Gamma_S}\widetilde u.
\]
特别地，
\[
\operatorname{im}f_S=B^1(\Gamma_S;\FF_2),
\]
即恰为屏幕图 \(\Gamma_S\) 的 cut space。
\end{theorem}
\leanverified{paper\_conclusion\_partial\_screen\_trace\_cut\_space\_identification}

\begin{proof}
在 \(\FF_2\) 系数下，取向符号消失，故 \(f_S(u)\) 在每条可见 \((n-1)\)-面上的值只取决于相邻顶维胞腔的占据奇偶。

若 \(e\in S\) 是内部可见面，其两侧 \(n\)-胞腔为 \(Q,Q'\)，则
\[
f_S(u)(e)=u(Q)+u(Q')=d_{\Gamma_S}\widetilde u(e).
\]
若 \(e\in S\) 是外边界可见面，其唯一相邻胞腔为 \(Q\)，则
\[
f_S(u)(e)=u(Q)=u(Q)+\widetilde u(\infty)=d_{\Gamma_S}\widetilde u(e).
\]
因此 \(f_S=d_{\Gamma_S}\)。

由于 cut space 按定义就是
\[
B^1(\Gamma_S;\FF_2)=\operatorname{im}(d_{\Gamma_S}),
\]
遂得
\[
\operatorname{im}f_S=B^1(\Gamma_S;\FF_2).
\]
证毕。
\end{proof}

\begin{corollary}[screen realizability 的唯一障碍即图环空间]\label{cor:conclusion-screen-realizability-cycle-space-obstruction}
沿用定理 \ref{thm:conclusion-partial-screen-trace-cut-space-identification} 的记号。记 \(V(\Gamma_S)\)、\(E(\Gamma_S)\) 分别为屏幕图的顶点与边集，\(c(\Gamma_S)\) 为其连通分量数，\(\beta_1(\Gamma_S)\) 为其第一 Betti 数。则
\[
\dim_{\FF_2}\operatorname{im}f_S=|V(\Gamma_S)|-c(\Gamma_S),
\]
并且
\[
\operatorname{codim}_{\FF_2^{\,S}}\operatorname{im}f_S
=
|E(\Gamma_S)|-|V(\Gamma_S)|+c(\Gamma_S)
=
\beta_1(\Gamma_S).
\]
从而
\[
f_S\text{ 为满射}
\iff
\beta_1(\Gamma_S)=0
\iff
\Gamma_S\text{ 为森林}.
\]
\leanverified{paper\_conclusion\_screen\_realizability\_cycle\_space\_obstruction}
\end{corollary}

\begin{proof}
由定理 \ref{thm:conclusion-partial-screen-trace-cut-space-identification}，
\[
\operatorname{im}f_S=B^1(\Gamma_S;\FF_2).
\]
有限图的 cut space 维数满足
\[
\dim_{\FF_2}B^1(\Gamma_S;\FF_2)=|V(\Gamma_S)|-c(\Gamma_S).
\]
又因 \(\FF_2^{\,S}\cong C^1(\Gamma_S;\FF_2)\) 的维数等于 \(|E(\Gamma_S)|\)，故余维为
\[
|E(\Gamma_S)|-|V(\Gamma_S)|+c(\Gamma_S)=\beta_1(\Gamma_S).
\]
满射性等价于余维为零，而有限图满足 \(\beta_1(\Gamma_S)=0\) 当且仅当 \(\Gamma_S\) 为森林。证毕。
\end{proof}

\begin{theorem}[全内部屏幕的二值剩余律：bulk 仅余整体补元]\label{thm:conclusion-full-internal-screen-global-complement-law}
令
\[
S_{\mathrm{int}}\subset \overrightarrow{\mathcal F_m^{n-1}}
\]
为全部内部取向 \((n-1)\)-面所构成的屏幕，即只观测内部界面而忽略全部外边界面。则
\[
\ker f_{S_{\mathrm{int}}}
=
\langle \mathbf 1\rangle
\cong
\FF_2,
\]
其中 \(\mathbf 1\) 表示在全部 \(n\)-胞腔上恒等于 \(1\) 的常值链。等价地，对任意 \(u\in C_n(\mathcal Q_m;\FF_2)\)，其内部屏幕迹唯一确定
\[
u
\quad\text{与}\quad
u+\mathbf 1
\]
这一对互补 bulk，而不存在其他歧义。
\end{theorem}

\begin{proof}
屏幕 \(S_{\mathrm{int}}\) 所对应的图 \(\Gamma_{S_{\mathrm{int}}}\) 正是全部 \(n\)-胞腔的邻接图外加一个孤立顶点 \(\infty\)。
其中胞腔邻接图连通，而 \(\infty\) 因未观测任何外边界面而自成单独分量，因此
\[
c(\Gamma_{S_{\mathrm{int}}})=2.
\]
由定理 \ref{thm:spg-screen-kernel-connected-components} 在 \(\FF_2\) 系数下的同样论证可知，
\[
\ker f_{S_{\mathrm{int}}}
\]
由“在每个连通分量上常值、并在 \(\infty\) 分量上取 \(0\)”的顶维链组成。于是胞腔分量上只剩一个自由常值位，恰由 \(\mathbf 1\) 生成。证毕。
\end{proof}
\leanverified{paper\_conclusion\_full\_internal\_screen\_global\_complement\_law}

\begin{theorem}[uniform bulk 的 cut-space 均匀化]\label{thm:conclusion-uniform-bulk-cut-space-uniformization}
在 \(C_n(\mathcal Q_m;\FF_2)\) 上取均匀分布。则对任意屏幕 \(S\)，其屏幕迹 \(f_S(u)\) 在
\[
B^1(\Gamma_S;\FF_2)
\]
上严格均匀分布。更精确地，对任意 \(y\in \FF_2^{\,S}\) 都有
\[
\PP\bigl(f_S(u)=y\bigr)
=
2^{-(|V(\Gamma_S)|-c(\Gamma_S))}
\mathbf 1_{\,B^1(\Gamma_S;\FF_2)}(y).
\]
特别地，若 \(\Gamma_S\) 为森林，则 \(f_S(u)\) 在整个 \(\FF_2^{\,S}\) 上均匀分布；因而全部可见屏幕位恰为彼此独立的公平 Bernoulli 比特。
\leanverified{paper\_conclusion\_uniform\_bulk\_cut\_space\_uniformization}
\end{theorem}

\begin{proof}
映射
\[
f_S:C_n(\mathcal Q_m;\FF_2)\to \FF_2^{\,S}
\]
是有限维 \(\FF_2\)-线性映射。故对任意 \(y\in\operatorname{im}f_S\)，其纤维 \(f_S^{-1}(y)\) 都是 \(\ker f_S\) 的仿射平移，因而具有相同基数
\[
\#f_S^{-1}(y)=2^{\dim_{\FF_2}\ker f_S}.
\]
由均匀源通过等纤维线性映射的推前仍在像空间上均匀，得到
\[
\PP\bigl(f_S(u)=y\bigr)=|\operatorname{im}f_S|^{-1}\mathbf 1_{\operatorname{im}f_S}(y).
\]
再由推论 \ref{cor:conclusion-screen-realizability-cycle-space-obstruction}，
\[
|\operatorname{im}f_S|
=
2^{\dim_{\FF_2}\operatorname{im}f_S}
=
2^{|V(\Gamma_S)|-c(\Gamma_S)},
\]
便得所示公式。

若 \(\Gamma_S\) 为森林，则由推论 \ref{cor:conclusion-screen-realizability-cycle-space-obstruction}，\(\operatorname{im}f_S=\FF_2^{\,S}\)。有限积空间上的均匀分布正是独立公平 Bernoulli 比特的乘积测度。证毕。
\end{proof}

\begin{theorem}[screen deficit 的算法信息 sharp 性]\label{thm:conclusion-screen-deficit-algorithmic-sharpness}
\leanverified{paper\_conclusion\_screen\_deficit\_algorithmic\_sharpness}
设
\[
r_S:=c(\Gamma_S)-1.
\]
则存在与通用前缀机无关的常数，使得对任意 \(u\in C_n(\mathcal Q_m;\FF_2)\)，都有
\[
K\bigl(u\mid f_S(u),S,m,n\bigr)
\le
r_S+O(\log(r_S+1)).
\]
反之，在 \(C_n(\mathcal Q_m;\FF_2)\) 的均匀分布下，对任意 \(t\ge 0\)，有
\[
\PP\!\left(
K\bigl(u\mid f_S(u),S,m,n\bigr)\le r_S-t
\right)
\le
2^{-t+O(1)}.
\]
因此，对典型 bulk 而言，partial screen 真正丢失的条件复杂度正好是
\[
r_S=c(\Gamma_S)-1
\]
比特。
\end{theorem}

\begin{proof}
由定理 \ref{thm:spg-screen-kernel-connected-components} 在 \(\FF_2\) 系数下的版本，
\[
\dim_{\FF_2}\ker f_S=r_S.
\]
固定 \(S,m,n\) 后，可用有限维线性代数可计算地求出 \(\ker f_S\) 的一组基
\[
k_1,\dots,k_{r_S},
\]
并选取一个可计算线性截面
\[
s:\operatorname{im}f_S\to C_n(\mathcal Q_m;\FF_2).
\]
于是任意 \(u\) 都可唯一写成
\[
u=s(f_S(u))+\sum_{j=1}^{r_S}\varepsilon_j k_j,
\qquad
\varepsilon_j\in\FF_2.
\]
给定 \(f_S(u),S,m,n\) 后，只需再给出二进制向量 \((\varepsilon_1,\dots,\varepsilon_{r_S})\) 即可恢复 \(u\)。采用前缀型自定界编码，其长度为
\[
r_S+O(\log(r_S+1)),
\]
从而得到上界。

再证下界。对任意固定 \(y\in\operatorname{im}f_S\)，纤维 \(f_S^{-1}(y)\) 的大小恰为 \(2^{r_S}\)。而长度不超过 \(r_S-t-c\) 的前缀描述至多有 \(2^{r_S-t-c+1}\) 个，其中 \(c\) 为通用机常数。故在该纤维中，满足
\[
K(u\mid y,S,m,n)\le r_S-t
\]
的点至多占 \(2^{-t+O(1)}\) 的比例。由定理 \ref{thm:conclusion-uniform-bulk-cut-space-uniformization}，均匀 bulk 在各纤维上条件均匀，故对全部 \(y\) 平均后仍得
\[
\PP\!\left(
K\bigl(u\mid f_S(u),S,m,n\bigr)\le r_S-t
\right)\le 2^{-t+O(1)}.
\]
证毕。
\end{proof}

\begin{corollary}[轴向屏幕的精确直和分裂与 visible/hidden 完全独立]\label{cor:conclusion-axial-screen-direct-sum-visible-hidden-split}
\leanverified{paper\_conclusion\_axial\_screen\_direct\_sum\_visible\_hidden\_split}
取命题 \ref{prop:spg-axial-screen-area-law-audit-cost} 中的轴向屏幕 \(S^{(1)}\)。则 \(\Gamma_{S^{(1)}}\) 为森林，且
\[
c(\Gamma_{S^{(1)}})=2^{m(n-1)}+1.
\]
因此
\[
f_{S^{(1)}}:C_n(\mathcal Q_m;\FF_2)\to \FF_2^{\,S^{(1)}}
\]
为满射，核维数为 \(2^{m(n-1)}\)。特别地，存在线性同构
\[
C_n(\mathcal Q_m;\FF_2)
\cong
\FF_2^{(2^m-1)2^{m(n-1)}}
\oplus
\FF_2^{2^{m(n-1)}}.
\]
其中第一因子正是轴向可见屏幕位，第二因子是最小 hidden completion bits；并且在均匀 bulk 分布下，这两个因子严格独立且各自均匀。
\end{corollary}

\begin{proof}
命题 \ref{prop:spg-axial-screen-area-law-audit-cost} 已证明
\[
c(\Gamma_{S^{(1)}})=2^{m(n-1)}+1,
\qquad
\rank_{\FF_2}\ker f_{S^{(1)}}=2^{m(n-1)}.
\]
由于 \(\Gamma_{S^{(1)}}\) 由 \(2^{m(n-1)}\) 条彼此不交的路径外加孤立顶点 \(\infty\) 组成，故它确为森林。于是由推论 \ref{cor:conclusion-screen-realizability-cycle-space-obstruction}，\(f_{S^{(1)}}\) 为满射。

再数可见坐标。对每个固定的 \((i_2,\dots,i_n)\in\{0,\dots,2^m-1\}^{n-1}\)，沿 \(x_1\) 方向恰有 \(2^m-1\) 条内部可见面，故
\[
|S^{(1)}|=(2^m-1)2^{m(n-1)}.
\]
于是短正合列
\[
0\to \ker f_{S^{(1)}}\to C_n(\mathcal Q_m;\FF_2)\to \FF_2^{\,S^{(1)}}\to 0
\]
在域 \(\FF_2\) 上必分裂，从而得到
\[
C_n(\mathcal Q_m;\FF_2)
\cong
\FF_2^{|S^{(1)}|}
\oplus
\FF_2^{2^{m(n-1)}}.
\]
代入 \(|S^{(1)}|\) 即得所示维数分解。

最后，在有限向量空间的直和分解下，均匀分布分解为两个因子的独立均匀分布，故 visible/hidden 两部分严格独立。证毕。
\end{proof}

\begin{corollary}[fixed-resolution hologram 的精确纠错半径]\label{cor:conclusion-fixed-resolution-hologram-exact-error-radius}
记
\[
\mathcal C_{m,n}:=\operatorname{im}\bigl(\partial_n:C_n(\mathcal Q_m;\FF_2)\to C_{n-1}(\mathcal Q_m;\FF_2)\bigr)
\]
为 dyadic 边界像码。则：
\begin{enumerate}
  \item 任意不超过 \(n-1\) 个取向 \((n-1)\)-面上的任意错误，都可被唯一纠正；
  \item 任意不超过 \(2n-1\) 个取向 \((n-1)\)-面上的任意错误，都可被确定检测；
  \item 该码同时具有常数校验度与正码率
        \[
        R_{m,n}=\frac{2^m}{n(2^m+1)}\xrightarrow[m\to\infty]{}\frac1n,
        \]
        因而 fixed-resolution hologram 构成一族常局部校验、常纠错半径而正码率的精确全息码。
\end{enumerate}
\end{corollary}
\leanverified{paper\_conclusion\_fixed\_resolution\_hologram\_exact\_error\_radius}

\begin{proof}
由定理 \ref{thm:spg-dyadic-boundary-image-ldpc}，\(\mathcal C_{m,n}\) 是局部校验度为 \(4\)、变量度为 \(2(n-1)\) 的规则 LDPC 码，且码率如上。
又由推论 \ref{cor:spg-dyadic-boundary-image-min-distance}，
\[
d_{\min}(\mathcal C_{m,n})=2n.
\]
编码论的标准结论遂给出唯一纠错半径
\[
\left\lfloor\frac{d_{\min}-1}{2}\right\rfloor=n-1
\]
以及检测半径
\[
d_{\min}-1=2n-1.
\]
证毕。
\end{proof}

\begin{theorem}[boundary arithmetic distance 对完整 tensor moment box 的面积尺度收缩]\label{thm:conclusion-boundary-arithmetic-distance-tensor-box-contraction}
对 fixed-resolution dyadic polycube \(U\subset I^n\)，定义完整 tensor moment box
\[
\mathcal M_m(U):=\bigl(\mu_\alpha(U)\bigr)_{0\le \alpha_i\le 2^m-1},
\qquad
\mu_\alpha(U):=\int_U x^\alpha\,\dd x.
\]
再定义加权上确界范数
\[
\|\Delta\|_{\star,m}
:=
\sup_{0\le \alpha_i\le 2^m-1}
(\alpha_1+1)\,|\Delta_\alpha|.
\]
则对任意 dyadic polycube \(U,V\subset I^n\)，都有
\[
\|\mathcal M_m(U)-\mathcal M_m(V)\|_{\star,m}
\le
2^{-m(n-1)}\,\operatorname{Dist}_m(U,V),
\]
其中 \(\operatorname{Dist}_m\) 为定理 \ref{thm:spg-godel-boundary-gcd-lipschitz-stability} 的 \(\gcd\)-算术边界距离。
\leanverified{paper\_conclusion\_boundary\_arithmetic\_distance\_tensor\_box\_contraction}
\end{theorem}

\begin{proof}
对任意满足 \(0\le \alpha_i\le 2^m-1\) 的多重指标 \(\alpha\)，推论 \ref{cor:spg-godel-moment-robust-gcd-bound} 给出
\[
|\mu_\alpha(U)-\mu_\alpha(V)|
\le
\frac{2^{-m(n-1)}}{\alpha_1+1}\,\operatorname{Dist}_m(U,V).
\]
两边同乘 \(\alpha_1+1\)，再对完整 tensor moment box 取上确界，即得
\[
\|\mathcal M_m(U)-\mathcal M_m(V)\|_{\star,m}
\le
2^{-m(n-1)}\,\operatorname{Dist}_m(U,V).
\]
证毕。
\end{proof}

\begin{theorem}[exact holography 与 isotropic polynomial tomography 之间的指数各向异性鸿沟]\label{thm:conclusion-exact-holography-vs-isotropic-tomography-anisotropy-gap}
固定 \(m,n\ge 1\)，并令 \(N:=2^m\)。则：
\begin{enumerate}
  \item 边界 Gödel 整数 \(H_m(U)\) 唯一决定最小完备的 tensor moment box
        \[
        \{\mu_\alpha(U):0\le \alpha_i\le N-1\},
        \]
        其坐标数恰为
        \[
        N^n=2^{mn};
        \]
  \item 若仅允许 isotropic 总次数矩族
        \[
        \{\mu_\alpha(U):|\alpha|\le r\},
        \]
        则要在线性意义上达到完备的必要门槛，必须满足
        \[
        \binom{n+r}{n}\ge 2^{mn},
        \]
        从而在 \(n\) 固定、\(m\to\infty\) 时必有
        \[
        r\ge (n!)^{1/n}2^m(1+o(1)).
        \]
\end{enumerate}
因此，fixed-resolution exact holography 与 isotropic polynomial tomography 之间存在一条不可消去的指数各向异性鸿沟：前者以单个平方自由整型边界坐标压缩了全部最小完备 tensor 数据；后者若坚持总次数各向同性，则不可避免地支付 \(2^m\) 级的次数爆炸。
\leanverified{paper\_conclusion\_exact\_holography\_vs\_isotropic\_tomography\_anisotropy\_gap}
\end{theorem}

\begin{proof}
由推论 \ref{cor:spg-boundary-godel-finite-moment-completeness}，边界 Gödel 整数 \(H_m(U)\) 在固定分辨率类内唯一决定 \(U\)；而同一推论与定理 \ref{thm:spg-dyadic-finite-moment-completeness} 一并表明，只需 tensor moment box
\[
\{\mu_\alpha(U):0\le \alpha_i\le N-1\}
\]
就已足以唯一恢复 \(U\)。其坐标数显然为
\[
N^n=2^{mn}.
\]

另一方面，推论 \ref{cor:spg-total-degree-moment-threshold-exponential-scale} 已给出：若总次数截断矩族
\[
\{\mu_\alpha:|\alpha|\le r\}
\]
要满足线性完备的必要门槛，则必须有
\[
\binom{n+r}{n}\ge 2^{mn},
\]
并因此推出
\[
r\ge (n!)^{1/n}2^m(1+o(1)).
\]
两部分合并即得结论。证毕。
\end{proof}

\begin{corollary}[边界 Gödel 完备性与各向同性矩全息之间的指数不可压缩分裂]\label{cor:conclusion-boundary-godel-vs-isotropic-moment-exponential-incompressibility}
\leanverified{paper\_conclusion\_boundary\_godel\_vs\_isotropic\_moment\_exponential\_incompressibility}
固定 \(m,n\ge 1\)，记 \(\mathcal P_{m,n}\) 为 \(I^n\) 中全体 \(m\)-dyadic polycube 的集合。对每个 \(U\in \mathcal P_{m,n}\)，令
\[
\mathcal M_r(U):=\bigl(\mu_\alpha(U)\bigr)_{|\alpha|\le r}
\]
为总次数不超过 \(r\) 的各向同性单项式矩截断。若存在映射
\[
J_m=\Psi_m\circ \mathcal M_{r_m}:\mathcal P_{m,n}\to Y_m
\]
并且 \(J_m\) 在 \(\mathcal P_{m,n}\) 上单射，则必有
\[
r_m\ge (n!)^{1/n}2^m(1+o(1))
\qquad (m\to\infty,\ n\ \text{固定}).
\]
与此相反，边界 Gödel 整数 \(H_m(U)\) 已在固定分辨率类 \(\mathcal P_{m,n}\) 内对 \(U\) 完全单射。故边界算术全息的完备性不能被任何满足 \(r_m=o(2^m)\) 的各向同性线性矩草图所承载。
\end{corollary}

\begin{proof}
若 \(J_m\) 在 \(\mathcal P_{m,n}\) 上单射，则其前置映射 \(\mathcal M_{r_m}\) 也必须在 \(\mathcal P_{m,n}\) 上单射。于是定理 \ref{thm:conclusion-exact-holography-vs-isotropic-tomography-anisotropy-gap} 的第二部分立即给出必要门槛
\[
r_m\ge (n!)^{1/n}2^m(1+o(1)).
\]
另一方面，同一定理的第一部分已表明 \(H_m(U)\) 在固定分辨率类内唯一决定 \(U\)，故 boundary Gödel holography 的完备性确实成立。两者合并即得“精确边界算术全息 / 亚临界各向同性矩草图”之间的指数不可压缩分裂。证毕。
\end{proof}

\begin{corollary}[低阶 Stokes/\texorpdfstring{$\mu_\alpha$}{mu alpha} 证书族在跨分辨率上必然失完备]\label{cor:conclusion-subcritical-isotropic-moment-certificates-incomplete-across-resolutions}
对任意满足 \(r_m=o(2^m)\) 的整数列，都存在两列不同的 \(m\)-dyadic polycube
\[
U_m\neq V_m,
\qquad
U_m,V_m\in \mathcal P_{m,n},
\]
使得
\[
\mu_\alpha(U_m)=\mu_\alpha(V_m)
\qquad (|\alpha|\le r_m),
\]
但
\[
H_m(U_m)\neq H_m(V_m).
\]
因此，任何只依赖于亚临界总次数矩的有限证书族，都不能在全部分辨率上充当边界 Gödel 全息的完备替身。
\end{corollary}

\begin{proof}
若对某个 \(r_m=o(2^m)\) 不存在这样的 \(U_m,V_m\)，则 \(\mathcal M_{r_m}\) 在 \(\mathcal P_{m,n}\) 上便是单射。这与推论 \ref{cor:conclusion-boundary-godel-vs-isotropic-moment-exponential-incompressibility} 的必要下界矛盾。故所述不同的 \(U_m,V_m\) 必然存在。证毕。
\end{proof}

\begin{theorem}[fixed-resolution 全息的线性--非线性严格分叉]\label{thm:conclusion-fixedresolution-linear-vs-nonlinear-holography-strict-bifurcation}
\leanverified{paper\_conclusion\_fixedresolution\_linear\_vs\_nonlinear\_holography\_strict\_bifurcation}
记
\[
\mathcal V_m:=\mathrm{span}\{1_{Q}:Q\in \mathcal Q_m^{(n)}\},
\qquad
\dim \mathcal V_m=2^{mn}.
\]
对任意有限线性测量族
\[
\Lambda_m=\{\ell_1,\dots,\ell_L\}\subset \mathcal V_m^\ast,
\]
定义测量映射
\[
E_{\Lambda_m}:\mathcal V_m\to \RR^L,
\qquad
f\longmapsto (\ell_1(f),\dots,\ell_L(f)).
\]
则：
\begin{enumerate}
  \item 若 \(E_{\Lambda_m}\) 在整个 \(\mathcal V_m\) 上单射，则必有
        \[
        L\ge 2^{mn}.
        \]
  \item 此下界由完整 tensor moment box 恰好达到；亦即
        \[
        \bigl(\mu_\alpha\bigr)_{0\le \alpha_i\le 2^m-1}
        \]
        给出一个坐标数精确为 \(2^{mn}\) 的线性完备全息。
  \item 相反，总次数截断矩族
        \[
        \bigl(\mu_\alpha\bigr)_{|\alpha|\le m-1}
        \]
        在 dyadic polycube 的示性函数子类上都不是单射。
  \item 然而，非线性的算术边界码
        \[
        U\longmapsto H_m(U)\in \mathfrak Q_m^\partial\subset \NN
        \]
        却已经对固定分辨率类给出精确恢复。
\end{enumerate}
因此，fixed-resolution holography 在线性范畴与算术范畴之间发生严格分叉：前者的完备全息不可压到少于 \(2^{mn}\) 个标量坐标，而后者在单个平方自由整数坐标内已经闭合。
\end{theorem}

\begin{proof}
若 \(E_{\Lambda_m}\) 在 \(\mathcal V_m\) 上单射，则有限维线性代数立刻给出
\[
L\ge \dim \mathcal V_m=2^{mn}.
\]
这证明了第一条。

第二条由定理 \ref{thm:conclusion-boundary-stokes-strict-linear-holography} 直接给出：完整 tensor moment box 在 fixed-resolution 类上构成与边界像等价的线性同构，因此正好以 \(2^{mn}\) 个坐标达到线性完备门槛。

对第三条，直接援引主稿中“固定分辨率 dyadic polycube 的线性矩全息最小维数”定理的第三部分：当
\[
r=m-1
\]
时，总次数截断矩族
\[
\bigl(\mu_\alpha\bigr)_{|\alpha|\le m-1}
\]
甚至在 dyadic polycube 的示性函数子类上都不是单射。故它当然更不可能充当 fixed-resolution 类上的完备线性全息。

第四条则由定理 \ref{thm:conclusion-fixedresolution-boundary-godel-momentbox-equivalence} 给出：边界 Gödel 整数 \(H_m(U)\) 与最小完备 tensor moment box 彼此可计算互逆，并在固定分辨率类内唯一决定 \(U\)。证毕。
\end{proof}

\begin{conjecture}[boundary Gödel 码的一坐标全息唯一性]\label{conj:conclusion-boundary-godel-single-coordinate-holography-uniqueness}
设
\[
\Psi_m:\mathcal P_{m,n}\to \NN
\]
为一个可计算编码，并满足以下三条：
\begin{enumerate}
  \item \(\Psi_m(U)\) 通过有限局部综合征即可判定是否来自某个合法边界链，且在合法情形下 admit 唯一 bulk 解码；
  \item \(\Psi_m\) 在某个算术距离 \(d_{\Psi}\) 下与边界对称差的 Hamming 距离等距，并嵌入某个有限算术中值空间；
  \item \(\Psi_m(U)\) 唯一决定 fixed-resolution 的完整 tensor moment box
        \[
        \{\mu_\alpha(U):0\le \alpha_i\le 2^m-1\}.
        \]
\end{enumerate}
则应存在一个由取向面重标记与素数重编号诱导的可计算算术超立方自同构
\[
\Theta_m:\mathfrak Q_m^\partial\to \mathfrak Q_m^\partial
\]
使得在把 \(\Psi_m(\mathcal P_{m,n})\) 识别为一个 squarefree boundary cube 之后，
\[
\Psi_m=\Theta_m\circ H_m.
\]
换言之，一旦同时要求局部 syndrome 可判定、算术中值可见与有限矩全息完备，则 boundary Gödel 码应当是 fixed-resolution dyadic polycube 类上一坐标全息的唯一极小实现。
\end{conjecture}

\begin{theorem}[全屏--部分屏幕的复杂度相变]\label{thm:conclusion-fullscreen-partialscreen-complexity-transition}
\leanverified{paper\_conclusion\_fullscreen\_partialscreen\_complexity\_transition}
固定 \(m,n\ge 1\)。对任一屏幕
\[
S\subset \overrightarrow{\mathcal F}^{\,n-1}_m,
\]
记
\[
f_S:=\Pi_S\circ \partial_n:C_n(\mathcal Q_m;\ZZ)\to \ZZ^S,
\]
并在其像上取任意可计算可逆的平方自由 Gödel 编码
\[
G_S:\operatorname{im}f_S\to \ZZ_{>0}.
\]
对每个 dyadic polycube \(U\subset I^n\)，定义部分读出
\[
\operatorname{Obs}_S(U):=G_S(f_S(u_U)).
\]
则存在统一的 \(O(1)\)-复杂度反演
\[
K\bigl(U\mid \operatorname{Obs}_S(U),m,n\bigr)=O(1)
\qquad\text{对全部 }U
\]
当且仅当
\[
c(\Gamma_S)=1.
\]
更一般地，使该反演成立所需的最小补全自由秩恰为
\[
\operatorname{Cost}(S)=c(\Gamma_S)-1.
\]
\end{theorem}
\begin{proof}
全边界情形中，boundary Gödel 全息坐标与 \(U\) 在给定 \((m,n)\) 后可计算可逆，故固定分辨率全息已经给出 \(O(1)\)-复杂度恢复。对一般屏幕 \(S\)，由于 \(G_S\) 在 \(\operatorname{im}f_S\) 上可计算可逆，\(\operatorname{Obs}_S(U)\) 与 \(f_S(u_U)\) 等价。

另一方面，主稿对部分屏幕的相对同调核已给出
\[
\ker f_S\cong H_n(\mathcal Q_m,A_{\neg S};\ZZ)\cong \ZZ^{\,c(\Gamma_S)-1},
\]
并且 sharp 补全成本公式正是
\[
\operatorname{Cost}(S)=\operatorname{rank}_{\ZZ}\ker f_S=c(\Gamma_S)-1.
\]
因此当 \(c(\Gamma_S)=1\) 时，\(\ker f_S=0\)，部分读出唯一决定 \(u_U\)，从而唯一决定 \(U\)，故有
\[
K\bigl(U\mid \operatorname{Obs}_S(U),m,n\bigr)=O(1).
\]
反之，若 \(c(\Gamma_S)>1\)，则恰有 \(c(\Gamma_S)-1\) 条独立整数赤字方向不能由屏幕读出恢复，因而任何可逆补全都至少需要该自由秩。证毕。
\end{proof}

\begin{corollary}[轴向屏幕赤字的指数律]\label{cor:conclusion-axial-screen-deficit-exponential-law}
对轴向屏幕 \(S^{(1)}\)，有精确公式
\[
\operatorname{Cost}(S^{(1)})=2^{m(n-1)}.
\]
因此，固定分辨率全息在全屏处是 \(O(1)\)-保真的；而一旦压缩为单轴屏幕，其补全自由秩立即跃迁到 codimension-one 面积的指数级。
\end{corollary}
\begin{proof}
由定理 \ref{thm:conclusion-fullscreen-partialscreen-complexity-transition}，
\[
\operatorname{Cost}(S)=c(\Gamma_S)-1.
\]
对轴向屏幕，主稿已给出
\[
c(\Gamma_{S^{(1)}})=2^{m(n-1)}+1.
\]
代入即得
\[
\operatorname{Cost}(S^{(1)})=2^{m(n-1)}.
\]
证毕。
\end{proof}
\leanverified{paper\_conclusion\_axial\_screen\_deficit\_exponential\_law}

\begin{conjecture}[固定分辨率完备审计四元组]\label{conj:conclusion-fixed-resolution-complete-audit-quadriga}
对固定分辨率折叠系统，存在最小完备审计函子
\[
\mathfrak A_{\mathrm{comp}}
\simeq
\mathfrak H_{\mathrm{bdry}}
\oplus
\mathfrak H_{\mathrm{rel}}(S)
\oplus
\mathfrak S_{\mathrm{tr}}
\oplus
\mathfrak T_{\mathrm{trans}},
\]
其中 \(\mathfrak H_{\mathrm{bdry}}\) 表示全边界全息坐标，\(\mathfrak H_{\mathrm{rel}}(S)\) 表示相对屏幕同调赤字，\(\mathfrak S_{\mathrm{tr}}\) 表示 ambiguity shell 的瞬态模块，而 \(\mathfrak T_{\mathrm{trans}}\) 表示算术--动力横截谱模块（例如 CRT 可见扇区与高阶耦合核心）。

并且，上述四个分量彼此不可由其余三者函子性重建。换言之，全边界全息不能替代部分屏幕补全，周期/Fredholm 证书不能替代瞬态壳，Watatani 指标不能替代屏幕赤字，而算术可见扇区与动力耦合核心之间也不存在统一的有理主分解。
\end{conjecture}


\endinput

\paragraph{边界 Stokes 全息、Gödel 距离等周律与最小观测维数}
在 fixed-resolution dyadic polycube 的有限矩完备性、边界 Gödel 码的有限矩全息完备性、边界像的 LDPC 阴影以及部分屏幕的 cut-space 几何都已固定之后，还可把边界侧的 Stokes 读出链进一步压缩为一组纯线性与纯组合的终态结论：完整张量积矩盒与边界像并不是两种不同的全息，而是同一个 \(2^{mn}\)-维线性对象的两套严格等价坐标；边界 Gödel 距离满足精确的离散等周律，并把“最小不可见障碍”唯一压缩为单个 dyadic 立方体；与此同时，低总次数 Stokes 观测族在固定分辨率上严格不完备，而完整张量积盒恰以最小维数达到边界线性全息的唯一完备门槛。

\begin{theorem}[Stokes--边界全息的严格线性同构]\label{thm:conclusion-boundary-stokes-strict-linear-holography}
\leanverified{paper\_conclusion\_boundary\_stokes\_strict\_linear\_holography}
固定 \(m,n\ge 1\)，记 \(N:=2^m\)，
\[
\mathcal A_m:=\{0,\dots,N-1\}^n,
\qquad
B_{m,n}^{\QQ}:=\operatorname{im}\!\bigl(\partial_n:C_n(\mathcal Q_m;\QQ)\to C_{n-1}(\mathcal Q_m;\QQ)\bigr).
\]
对每个 \(\alpha=(\alpha_1,\dots,\alpha_n)\in\mathcal A_m\)，定义
\[
\omega^{(1)}_{\alpha}
:=
\frac{x^{\alpha+e_1}}{\alpha_1+1}\,\dd x_2\wedge\cdots\wedge \dd x_n,
\qquad
\dd\omega^{(1)}_{\alpha}=x^\alpha\,\dd x_1\cdots \dd x_n.
\]
则线性映射
\[
\mathcal S_{m,n}:B_{m,n}^{\QQ}\longrightarrow \QQ^{\mathcal A_m},
\qquad
\mathcal S_{m,n}(b):=
\Bigl(\int_b\omega^{(1)}_{\alpha}\Bigr)_{\alpha\in\mathcal A_m}
\]
是同构。特别地，固定分辨率下的边界像与完整张量积矩盒并不是两种不同的全息，而是同一有限维对象的两套严格等价坐标；其模 \(2\) 阴影正是定理 \ref{thm:spg-dyadic-boundary-image-ldpc} 的 dyadic 边界 LDPC 码族。
\end{theorem}

\begin{proof}
由于 \(\mathcal Q_m\) 没有 \((n+1)\)-胞腔，顶维边界算子
\[
\partial_n:C_n(\mathcal Q_m;\QQ)\longrightarrow B_{m,n}^{\QQ}
\]
为线性同构。

把 \(C_n(\mathcal Q_m;\QQ)\) 与定理 \ref{thm:spg-dyadic-finite-moment-completeness} 证明中出现的分段常数空间 \(\mathcal V_m\) 通过标准 dyadic \(n\)-胞腔基底识别，并定义完整矩盒映射
\[
\mathcal M_m:C_n(\mathcal Q_m;\QQ)\longrightarrow \QQ^{\mathcal A_m},
\qquad
\mathcal M_m(c):=
\Bigl(\int_c x^\alpha\,\dd x_1\cdots \dd x_n\Bigr)_{\alpha\in\mathcal A_m}.
\]
由定理 \ref{thm:spg-dyadic-finite-moment-completeness} 的线性版本，\(\mathcal M_m\) 是同构。

对任意 \(b=\partial_n c\in B_{m,n}^{\QQ}\)，Stokes 定理给出
\[
\int_b\omega^{(1)}_{\alpha}
=
\int_c \dd\omega^{(1)}_{\alpha}
=
\int_c x^\alpha\,\dd x_1\cdots\dd x_n.
\]
因此
\[
\mathcal S_{m,n}
=
\mathcal M_m\circ (\partial_n)^{-1}.
\]
右侧是两个线性同构的复合，故 \(\mathcal S_{m,n}\) 也是线性同构。证毕。
\end{proof}

\begin{corollary}[边界线性泛函的截断多项式穷尽性]\label{cor:conclusion-boundary-linear-functionals-truncated-polynomial-exhaustion}
沿用定理 \ref{thm:conclusion-boundary-stokes-strict-linear-holography} 的记号。则对任意线性泛函
\[
\Lambda:B_{m,n}^{\QQ}\to \QQ
\]
都存在唯一截断多项式
\[
p_\Lambda(x)=\sum_{\alpha\in\mathcal A_m} c_\alpha x^\alpha
\]
使得对每个 dyadic polycube \(U\subset I^n\) 都有
\[
\Lambda(\partial_n[U])=\int_U p_\Lambda(x)\,\dd x.
\]
\leanverified{paper\_conclusion\_boundary\_linear\_functionals\_truncated\_polynomial\_exhaustion}
换言之，边界线性统计量恰由坐标逐项次数 \(<2^m\) 的多项式体积分全部穷尽。
\end{corollary}

\begin{proof}
由定理 \ref{thm:conclusion-boundary-stokes-strict-linear-holography}，\(\Lambda\) 唯一对应于一组系数 \((c_\alpha)_{\alpha\in\mathcal A_m}\)，使
\[
\Lambda(b)=\sum_{\alpha\in\mathcal A_m} c_\alpha \int_b\omega^{(1)}_{\alpha}
\qquad
(b\in B_{m,n}^{\QQ}).
\]
令
\[
p_\Lambda(x):=\sum_{\alpha\in\mathcal A_m} c_\alpha x^\alpha.
\]
对任意 dyadic polycube \(U\subset I^n\)，由 Stokes 定理得
\[
\Lambda(\partial_n[U])
=
\sum_{\alpha\in\mathcal A_m} c_\alpha \int_{\partial_n[U]}\omega^{(1)}_{\alpha}
=
\sum_{\alpha\in\mathcal A_m} c_\alpha \int_U x^\alpha\,\dd x
=
\int_U p_\Lambda(x)\,\dd x.
\]
唯一性由定理 \ref{thm:conclusion-boundary-stokes-strict-linear-holography} 的单射性立即推出。证毕。
\end{proof}

\begin{theorem}[边界 Gödel 距离的离散等周律]\label{thm:conclusion-boundary-godel-distance-discrete-isoperimetric-law}
设 \(n\ge 2\)。对任意同分辨率 dyadic polycube \(U,V\subset I^n\)，定义其边界 Gödel 对称差计数
\[
\operatorname{Dist}_m(U,V)
:=
\omega\!\left(\frac{H_m(U)H_m(V)}{\gcd(H_m(U),H_m(V))^2}\right),
\]
并记
\[
N_\Delta(U,V):=\#\bigl\{Q\in \mathcal Q_m:\ Q\subset U\triangle V\bigr\}.
\]
则有精确双边估计
\[
2n\,N_\Delta(U,V)^{\frac{n-1}{n}}
\le
\operatorname{Dist}_m(U,V)
\le
2n\,N_\Delta(U,V).
\]
\end{theorem}

\begin{proof}
令 \(W:=U\triangle V\)。由定理 \ref{thm:spg-godel-boundary-gcd-lipschitz-stability} 与推论 \ref{cor:spg-boundary-godel-ratio-squareclass-isometry}，
\[
\operatorname{Dist}_m(U,V)
=
\bigl\|\partial_n[U]-\partial_n[V]\bigr\|_{\ell^1}.
\]
而对 dyadic polycube 的 \(0/1\)-占据链，\(\partial_n[U]-\partial_n[V]\) 的支撑恰对应于 \(W=U\triangle V\) 的外露边界，故
\[
\operatorname{Dist}_m(U,V)=F(W),
\]
其中 \(F(W)\) 为 \(W\) 的边界面数。再由定理 \ref{thm:spg-dyadic-polyclube-discrete-isoperimetry}，
\[
2n\,N(W)^{\frac{n-1}{n}}
\le
F(W)
\le
2n\,N(W).
\]
由于 \(N(W)=N_\Delta(U,V)\)，代入即得结论。证毕。
\end{proof}
\leanverified{paper\_conclusion\_boundary\_godel\_distance\_discrete\_isoperimetric\_law}

\begin{corollary}[单立方体编辑刚性]\label{cor:conclusion-single-cube-edit-rigidity}
在定理 \ref{thm:conclusion-boundary-godel-distance-discrete-isoperimetric-law} 的设定下：
\[
\operatorname{Dist}_m(U,V)<2n \Longrightarrow U=V;
\]
并且若 \(U\neq V\) 且
\[
\operatorname{Dist}_m(U,V)=2n,
\]
则
\[
U\triangle V
\]
恰为一个唯一的 \(m\)-dyadic \(n\)-立方体。
\end{corollary}

\begin{proof}
若 \(\operatorname{Dist}_m(U,V)<2n\)，由定理 \ref{thm:conclusion-boundary-godel-distance-discrete-isoperimetric-law} 得
\[
2n\,N_\Delta(U,V)^{\frac{n-1}{n}}<2n,
\]
故 \(N_\Delta(U,V)=0\)，即 \(U=V\)。

若 \(U\neq V\) 且 \(\operatorname{Dist}_m(U,V)=2n\)，则同样由该定理得
\[
N_\Delta(U,V)^{\frac{n-1}{n}}\le 1.
\]
由于 \(U\neq V\)，必有 \(N_\Delta(U,V)\ge 1\)，故只能是 \(N_\Delta(U,V)=1\)。这正表示 \(U\triangle V\) 恰为一个唯一的 \(m\)-dyadic \(n\)-立方体。证毕。
\end{proof}
\leanpartial{paper\_conclusion\_single\_cube\_edit\_rigidity}{The stated Lean signature is false for $n=0$; a sound proof needs the ambient hypothesis $n\ge 1$.}

\begin{theorem}[隐藏面预算的尖锐阈值与极小障碍分类]\label{thm:conclusion-hidden-face-budget-sharp-threshold-classification}
设 \(n\ge 2\)，\(T\) 为一组被隐藏的取向 \((n-1)\)-面。若两个 dyadic polycube \(U,V\subset I^n\) 的边界读出在可见部分上一致，即
\[
\Pi_{\overrightarrow{\mathcal F}^{\,n-1}_m\setminus T}\,\partial_n[U]
=
\Pi_{\overrightarrow{\mathcal F}^{\,n-1}_m\setminus T}\,\partial_n[V],
\]
则成立：
\begin{enumerate}
  \item 若 \(|T|<2n\)，必有 \(U=V\)；
  \item 若 \(|T|=2n\) 且 \(U\neq V\)，则 \(U\triangle V\) 恰为一个唯一的 dyadic \(n\)-立方体 \(Q\)，并且
  \[
  T=\operatorname{supp}(\partial_n[Q]).
  \]
\end{enumerate}
反之，任取一个 dyadic \(n\)-立方体 \(Q\)，若取
\[
T=\operatorname{supp}(\partial_n[Q]),
\]
则 \(U\) 与 \(U\triangle Q\) 在可见部分上不可区分。
\end{theorem}

\begin{proof}
由可见部分一致，\(\partial_n[U]-\partial_n[V]\) 的支撑只能落在 \(T\) 内，故
\[
\operatorname{Dist}_m(U,V)\le |T|.
\]
若 \(|T|<2n\)，则由推论 \ref{cor:conclusion-single-cube-edit-rigidity} 立即得到 \(U=V\)。

若 \(|T|=2n\) 且 \(U\neq V\)，则仍由推论 \ref{cor:conclusion-single-cube-edit-rigidity}，\(U\triangle V\) 必为一个唯一 dyadic \(n\)-立方体 \(Q\)。于是
\[
\operatorname{supp}(\partial_n[Q])
=
\operatorname{supp}\bigl(\partial_n[U]-\partial_n[V]\bigr)
\subseteq T.
\]
两边基数同为 \(2n\)，故必有
\[
T=\operatorname{supp}(\partial_n[Q]).
\]
反向方向显然：若隐藏恰好覆盖 \(Q\) 的全部边界面，则翻转 \(Q\) 的占据状态不会改变任何可见坐标。证毕。
\end{proof}

\begin{proposition}[高阶碰撞矩奇偶影子的有限状态崩缩]\label{prop:conclusion-high-collision-moment-mod2-finite-state-collapse}
对更新稿审计窗口中的碰撞矩序列 \(S_q(m)\)（\(q=9,\dots,17\)），记前移算子 \(E\) 作用为 \((Ef)(m)=f(m+1)\)，并令
\[
D:=E+1
\]
为特征 \(2\) 下的前向差分。则每个 \(q\in\{9,\dots,17\}\) 都存在整数 \(a_q,b_q\) 使
\[
D^{\,b_q}S_q(m+a_q)\equiv 0 \pmod 2.
\]
因而 \(m\mapsto S_q(m)\bmod 2\) 的尾序列必为有限周期列；更精确地，其周期满足
\[
T_q\mid 4 \quad(q=9,10),\qquad
T_q\mid 8 \quad(q=11,\dots,16),\qquad
T_q\mid 16 \quad(q=17).
\]
\end{proposition}

\begin{proof}
更新稿在审计窗口 \(q=9,\dots,17\) 上已给出模 \(2\) 差分湮灭律
\[
D^{\,b_q}S_q(m+a_q)\equiv 0\pmod 2.
\]
在 \(\FF_2\) 上，这等价于尾序列 \(m\mapsto S_q(m+a_q)\) 是次数 \(<b_q\) 的离散多项式，故可写成
\[
S_q(m+a_q)\equiv \sum_{k=0}^{b_q-1} c_k\binom{m}{k}\pmod 2.
\]
由 Lucas 定理，\(\binom{m}{k}\bmod 2\) 仅依赖于 \(m\) 的有限低位二进制展开，故其周期整除
\[
2^{\lceil \log_2(k+1)\rceil}.
\]
于是整个和式的周期整除
\[
2^{\lceil \log_2 b_q\rceil}.
\]
再把更新稿表格中对应的 \(b_q\) 范围代入，即得
\[
T_q\mid 4 \quad(q=9,10),\qquad
T_q\mid 8 \quad(q=11,\dots,16),\qquad
T_q\mid 16 \quad(q=17).
\]
证毕。
\end{proof}

\begin{theorem}[低总次数 Stokes 家族在边界侧的不可分离性]\label{thm:conclusion-low-total-degree-stokes-family-nonseparation}
设
\[
\Omega_{\le r}:=\operatorname{span}\{\omega^{(1)}_{\alpha}:|\alpha|\le r\}\subset \Omega^{n-1}(I^n).
\]
当 \(r=m-1\) 时，映射
\[
B_{m,n}^{\QQ}\longrightarrow \Omega_{\le m-1}^{\ast},
\qquad
b\longmapsto \bigl(\omega\mapsto \int_b\omega\bigr)
\]
不是单射。更强地，存在互异 dyadic polycube \(U\neq V\) 使
\[
\partial_n[U]\neq \partial_n[V],
\qquad
\int_{\partial_n[U]}\omega=\int_{\partial_n[V]}\omega
\quad
(\forall\,\omega\in\Omega_{\le m-1}).
\]
\end{theorem}

\begin{proof}
由定理 \ref{thm:spg-prouhet-thue-morse-obstruction-dyadic-polyclube-flux-moments}，存在互异 \(m\)-级 dyadic polycube \(U\neq V\) 使得对一切 \(|\alpha|\le m-1\) 都有
\[
\mu_\alpha(U)=\mu_\alpha(V).
\]
由于 \(\partial_n\) 在顶维上单射，故
\[
\partial_n[U]\neq \partial_n[V].
\]
任取 \(\omega\in\Omega_{\le m-1}\)，可写成
\[
\omega=\sum_{|\alpha|\le m-1} c_\alpha\,\omega^{(1)}_{\alpha}.
\]
于是由 Stokes 定理，
\[
\int_{\partial_n[U]}\omega
=
\sum_{|\alpha|\le m-1} c_\alpha \int_{\partial_n[U]}\omega^{(1)}_{\alpha}
=
\sum_{|\alpha|\le m-1} c_\alpha \mu_\alpha(U)
=
\sum_{|\alpha|\le m-1} c_\alpha \mu_\alpha(V)
=
\int_{\partial_n[V]}\omega.
\]
故该观测族不能分离全部边界像。证毕。
\end{proof}
\leanverified{paper\_conclusion\_low\_total\_degree\_stokes\_family\_nonseparation}

\begin{theorem}[边界 Stokes 观测的最小维数]\label{thm:conclusion-boundary-stokes-observation-minimal-dimension}
在固定分辨率 \((m,n)\) 下，任何能够分离全部边界像
\[
B_{m,n}^{\QQ}
\]
的线性 Stokes 观测族，其基数至少为
\[
2^{mn}.
\]
此下界可达；恰由完整张量积盒
\[
\{\omega^{(1)}_{\alpha}:\alpha\in\{0,\dots,2^m-1\}^n\}
\]
实现。
\leanverified{paper\_conclusion\_boundary\_stokes\_observation\_minimal\_dimension\_seeds}
\end{theorem}

\begin{proof}
设 \(\ell_1,\dots,\ell_L\) 为一族线性 Stokes 观测，并令联合映射
\[
\ell:B_{m,n}^{\QQ}\to \QQ^L,
\qquad
\ell(b):=(\ell_1(b),\dots,\ell_L(b)).
\]
若 \(\ell\) 单射，则必有
\[
L\ge \dim_{\QQ} B_{m,n}^{\QQ}.
\]
而 \(\partial_n:C_n(\mathcal Q_m;\QQ)\to B_{m,n}^{\QQ}\) 为同构，所以
\[
\dim_{\QQ} B_{m,n}^{\QQ}
=
\dim_{\QQ} C_n(\mathcal Q_m;\QQ)
=
2^{mn}.
\]
故必有 \(L\ge 2^{mn}\)。

另一方面，定理 \ref{thm:conclusion-boundary-stokes-strict-linear-holography} 已证明，完整张量积盒所诱导的 Stokes 观测给出一个线性同构
\[
B_{m,n}^{\QQ}\xrightarrow{\sim}\QQ^{2^{mn}}.
\]
因此该下界可达。证毕。
\end{proof}

\begin{theorem}[次临界边界扰动下的二阶 Gödel--Stokes 刚性]\label{thm:conclusion-boundary-subcritical-perturbation-second-order-rigidity}
设 \(A\subset I^n\) 满足
\[
N_m(A)=M\,2^{md}+o(2^{md}),
\qquad
M\in(0,\infty),\ d\in(0,n],
\]
其中 \(N_m(A)\) 为 \(A\) 的 dyadic 外逼近计数，\(U_m(A)\) 为对应外逼近。
再设 \((V_m)_{m\ge 1}\) 为同分辨率 dyadic polycube 列，并定义
\[
N_m(V_m):=2^{mn}\operatorname{Vol}(V_m).
\]
若
\[
\operatorname{Dist}_m\bigl(U_m(A),V_m\bigr)=o\!\bigl(2^{m(d-1)}\bigr),
\]
则有
\[
N_m(V_m)=M\,2^{md}+o(2^{md}).
\]
特别地，\(V_m\) 与 \(U_m(A)\) 具有相同的 dyadic Minkowski 维数 \(d\) 与相同的 dyadic Minkowski 含量 \(M\)。等价地，由定理 \ref{thm:spg-godel-doublelog-second-order-minkowski-content} 恢复出的 primorial 双对数二阶常数在该阈值以下保持刚性。
\end{theorem}
\begin{proof}
由定理 \ref{thm:spg-godel-boundary-gcd-lipschitz-stability}(2)，
\[
\bigl|\operatorname{Vol}(V_m)-\operatorname{Vol}(U_m(A))\bigr|
\le
2^{-m(n-1)}\operatorname{Dist}_m\bigl(U_m(A),V_m\bigr)
=
o\!\bigl(2^{-m(n-d)}\bigr).
\]
两边同乘 \(2^{mn}\)，并注意
\[
N_m(A)=2^{mn}\operatorname{Vol}(U_m(A)),
\qquad
N_m(V_m)=2^{mn}\operatorname{Vol}(V_m),
\]
便得到
\[
|N_m(V_m)-N_m(A)|=o(2^{md}).
\]
再代入
\[
N_m(A)=M\,2^{md}+o(2^{md}),
\]
即得
\[
N_m(V_m)=M\,2^{md}+o(2^{md}).
\]
因此 \(V_m\) 具有与 \(A\) 相同的 dyadic Minkowski 维数与含量。最后应用定理 \ref{thm:spg-godel-doublelog-second-order-minkowski-content}，可知相应 primorial 双对数的一阶斜率与二阶正规化常数也完全一致。证毕。
\end{proof}

\begin{theorem}[维数抬升的必要边界税率]\label{thm:conclusion-boundary-dimension-lifting-necessary-tax}
设 \(A\subset I^n\) 的上 dyadic Minkowski 维严格小于 \(q\le n\)。等价地，
\[
N_m(A)=o(2^{mq}).
\]
令 \(U_m(A)\) 为其 dyadic 外逼近。若存在一列同分辨率 dyadic polycube \(V_m\) 与常数 \(c>0\)，使得沿某个无穷子列有
\[
\operatorname{Vol}(V_m)\ge c\,2^{-m(n-q)},
\]
则沿同一子列必有
\[
\operatorname{Dist}_m\bigl(U_m(A),V_m\bigr)=\Omega\!\bigl(2^{m(q-1)}\bigr).
\]
\end{theorem}
\begin{proof}
由
\[
N_m(A)=o(2^{mq})
\]
可得
\[
\operatorname{Vol}(U_m(A))
=
2^{-mn}N_m(A)
=
o\!\bigl(2^{-m(n-q)}\bigr).
\]
因此沿所给子列，对充分大的 \(m\) 有
\[
\bigl|\operatorname{Vol}(V_m)-\operatorname{Vol}(U_m(A))\bigr|
\ge
\frac{c}{2}\,2^{-m(n-q)}.
\]
再由定理 \ref{thm:spg-godel-boundary-gcd-lipschitz-stability}(2)，
\[
\bigl|\operatorname{Vol}(V_m)-\operatorname{Vol}(U_m(A))\bigr|
\le
2^{-m(n-1)}\operatorname{Dist}_m\bigl(U_m(A),V_m\bigr),
\]
遂得
\[
\operatorname{Dist}_m\bigl(U_m(A),V_m\bigr)
\ge
\frac{c}{2}\,2^{m(q-1)}.
\]
证毕。
\end{proof}

\begin{corollary}[固定维数下内容重整的必要边界税率]\label{cor:conclusion-boundary-content-renormalization-necessary-tax}
设 \(A\subset I^n\) 满足
\[
N_m(A)=M\,2^{md}+o(2^{md}),
\qquad
M>0,\ d\in(0,n].
\]
设另一列同分辨率 dyadic polycube \(V_m\) 满足
\[
N_m(V_m)=M'\,2^{md}+o(2^{md}),
\qquad
M'\neq M.
\]
则必有
\[
\operatorname{Dist}_m\bigl(U_m(A),V_m\bigr)=\Omega\!\bigl(2^{m(d-1)}\bigr).
\]
\end{corollary}
\begin{proof}
由两条渐近式相减，
\[
|N_m(V_m)-N_m(A)|
=
|M'-M|\,2^{md}+o(2^{md}).
\]
两边乘以 \(2^{-mn}\)，即得
\[
\bigl|\operatorname{Vol}(V_m)-\operatorname{Vol}(U_m(A))\bigr|
=
|M'-M|\,2^{-m(n-d)}+o\!\bigl(2^{-m(n-d)}\bigr).
\]
因此对充分大的 \(m\)，存在常数 \(c>0\) 使
\[
\bigl|\operatorname{Vol}(V_m)-\operatorname{Vol}(U_m(A))\bigr|
\ge
c\,2^{-m(n-d)}.
\]
再由定理 \ref{thm:spg-godel-boundary-gcd-lipschitz-stability}(2)，
\[
\operatorname{Dist}_m\bigl(U_m(A),V_m\bigr)
\ge
c\,2^{m(d-1)}.
\]
证毕。
\end{proof}

\begin{conjecture}[实维 surgery 的边界税率等号原理]\label{conj:conclusion-boundary-real-dimension-surgery-tax-equivalence}
对任意实数 \(s\in(0,n]\)，设 \((U_m)_{m\ge 1}\) 为一列同分辨率 dyadic polycube，并考虑所有满足
\[
\bigl|\operatorname{Vol}(V_m)-\operatorname{Vol}(U_m)\bigr|
\asymp
2^{-m(n-s)}
\]
的 dyadic surgery 列 \((V_m)_{m\ge 1}\)。
则其最小可能边界算术代价应满足
\[
\inf_{V_m}\operatorname{Dist}_m(U_m,V_m)\asymp 2^{m(s-1)}.
\]
亦即，指数 \(s-1\) 应是“以边界算术距离制造 \(s\)-维 bulk 质量”的普适税率。定理 \ref{thm:conclusion-boundary-dimension-lifting-necessary-tax} 与推论 \ref{cor:conclusion-boundary-content-renormalization-necessary-tax} 已经给出这一税率作为必要下界的普适性；剩余问题在于构造达到该指数级别的近极小 surgery 族。
\end{conjecture}

\paragraph{屏幕核的 solenoid 完备化、\texorpdfstring{$p$}{p}-进隐熵与 Kolmogorov 缺损的精确核坐标}
在部分屏幕的相对顶维同调核、cut-space 余秩闭式与完整边界 Gödel 全息都已经建立之后，屏幕缺损还可再压缩成一条更细的算术--算法链：屏幕核不仅给出一个自由阿贝尔群，而且天然携带一个由实相位部分与素数账本部分拼接而成的 solenoid 完备化；同一个秩 \(r_S=c(\Gamma_S)-1\) 在每个素数上都重新出现为严格的仿射纤维维数；并且在 \(\FF_2\) 口径下，部分屏幕真正丢失的 Kolmogorov 信息并不是某种模糊缺损，而精确等于屏幕核坐标的条件复杂度。于是，几何屏幕缺口、\(p\)-进隐熵与算法缺损三者被压缩为同一整数 \(r_S\) 的三种表现。

\leanverified{paper\_conclusion\_screen\_kernel\_solenoid\_bicompletion}
\begin{theorem}[屏幕核的规范双完备化与 solenoid 出现]\label{thm:conclusion-screen-kernel-solenoid-bicompletion}
\leanverified{paper\_conclusion\_screen\_kernel\_solenoid\_bicompletion}
固定分辨率 \(m\) 与维数 \(n\)，并沿用定理 \ref{thm:spg-partial-boundary-screen-kernel-relative-homology} 的部分边界观测算子
\[
f_S:C_n(\mathcal Q_m;\ZZ)\to \ZZ^S
\]
与定理 \ref{thm:spg-screen-kernel-connected-components} 的记号。记
\[
K_S:=\ker f_S,
\qquad
r_S:=c(\Gamma_S)-1.
\]
定义其规范双完备化
\[
\Sigma_S
:=
\bigl((K_S\otimes_{\ZZ}\RR)\times (K_S\otimes_{\ZZ}\widehat{\ZZ})\bigr)\big/K_S,
\qquad
\widehat{\ZZ}:=\prod_p \ZZ_p,
\]
其中 \(K_S\) 对角嵌入到两因子之中。则 \(\Sigma_S\) 是一个与基选择无关的紧交换群，并具有自然短正合列
\[
0\longrightarrow K_S\otimes_{\ZZ}\widehat{\ZZ}
\longrightarrow \Sigma_S
\longrightarrow (K_S\otimes_{\ZZ}\RR)/K_S
\longrightarrow 0.
\]
由于 \(K_S\cong \ZZ^{r_S}\)，该短正合列具体化为
\[
0\longrightarrow \prod_p \ZZ_p^{\,r_S}
\longrightarrow \Sigma_S
\longrightarrow \mathbb T^{r_S}
\longrightarrow 0,
\qquad
\mathbb T:=\RR/\ZZ.
\]
因此，每一个部分屏幕都内生地决定了一个 \(r_S\)-维 solenoid：其连通相位部分恰为 \(\mathbb T^{r_S}\)，其全不连通素数账本部分恰为 \(\prod_p\ZZ_p^{\,r_S}\)。
\end{theorem}

\begin{proof}
由定理 \ref{thm:spg-screen-kernel-connected-components}，
\[
K_S\cong \ZZ^{\,r_S}
\]
是自由阿贝尔群。故
\[
K_S\otimes_{\ZZ}\RR\cong \RR^{r_S},
\qquad
K_S\otimes_{\ZZ}\widehat{\ZZ}\cong \widehat{\ZZ}^{\,r_S}\cong \prod_p \ZZ_p^{\,r_S}.
\]
将 \(K_S\) 以对角方式嵌入
\[
K_S\hookrightarrow (K_S\otimes_{\ZZ}\RR)\times (K_S\otimes_{\ZZ}\widehat{\ZZ}),
\qquad
k\mapsto (k,k),
\]
并定义商群 \(\Sigma_S\) 如上。

考虑投影到第一因子所诱导的满射
\[
\pi_S:\Sigma_S\to (K_S\otimes_{\ZZ}\RR)/K_S,
\qquad
[(x,z)]\mapsto x+K_S.
\]
若 \(\pi_S([(x,z)])=0\)，则 \(x\in K_S\)。于是
\[
[(x,z)]=[(0,z-x)],
\]
故 \(\ker\pi_S\) 由形如 \([(0,z)]\) 的类组成。映射
\[
\iota_S:K_S\otimes_{\ZZ}\widehat{\ZZ}\to \Sigma_S,
\qquad
z\mapsto [(0,z)]
\]
显然为群同态。若 \([(0,z)]=[(0,z')]\)，则存在 \(k\in K_S\) 使
\[
(0,z')-(0,z)=(k,k).
\]
比较第一坐标得 \(k=0\)，因而 \(z=z'\)，故 \(\iota_S\) 为单射。由上面对 \(\ker\pi_S\) 的描述可知
\[
\ker\pi_S=\operatorname{im}\iota_S\cong K_S\otimes_{\ZZ}\widehat{\ZZ}.
\]
这就得到所述短正合列。

最后，由 \(K_S\cong \ZZ^{r_S}\) 可得
\[
(K_S\otimes_{\ZZ}\RR)/K_S\cong \RR^{r_S}/\ZZ^{r_S}\cong \mathbb T^{r_S}.
\]
若更换 \(K_S\) 的 \(\ZZ\)-基，则变化矩阵属于 \(\GL_{r_S}(\ZZ)\)，它在 \(\RR^{r_S}\) 与 \(\widehat{\ZZ}^{\,r_S}\) 上同时诱导自同构，故商群 \(\Sigma_S\) 的同构类型与基选择无关。证毕。
\end{proof}

\begin{corollary}[部分屏幕的 \texorpdfstring{$p$}{p}-进隐熵律]\label{cor:conclusion-screen-padic-hidden-entropy-law}
\leanverified{paper\_conclusion\_screen\_padic\_hidden\_entropy\_law}
沿用定理 \ref{thm:conclusion-screen-kernel-solenoid-bicompletion} 的记号。对任意素数 \(p\)，令
\[
\bar f_{S,p}:C_n(\mathcal Q_m;\FF_p)\to \FF_p^{\,S}
\]
为 \(f_S\) 的模 \(p\) 约化。则对每个可实现的屏幕综合征 \(y\in\operatorname{im}(\bar f_{S,p})\)，其纤维都是一个维数恰为 \(r_S\) 的仿射 \(\FF_p\)-空间；特别地，
\[
\#\bar f_{S,p}^{-1}(y)=p^{r_S}.
\]
因此，部分屏幕在每个素数 \(p\) 上留下的隐藏信息量恰为
\[
r_S\log p,
\]
并且该量不依赖于综合征 \(y\)，只依赖于屏幕图的连通分量数 \(c(\Gamma_S)\)。
\end{corollary}

\begin{proof}
由定理 \ref{thm:spg-screen-kernel-connected-components}，
\[
K_S\cong \ZZ^{\,r_S}.
\]
张量化到 \(\FF_p\) 后得到
\[
K_S\otimes_{\ZZ}\FF_p\cong \FF_p^{\,r_S}.
\]
另一方面，\(\bar f_{S,p}\) 的核恰为 \(f_S\) 在 \(\FF_p\) 上的基变换核，即
\[
\ker \bar f_{S,p}\cong K_S\otimes_{\ZZ}\FF_p\cong \FF_p^{\,r_S}.
\]
于是任一非空纤维 \(\bar f_{S,p}^{-1}(y)\) 都是 \(\ker\bar f_{S,p}\) 的平移，故为维数 \(r_S\) 的仿射 \(\FF_p\)-空间，其点数正是 \(p^{r_S}\)。证毕。
\end{proof}

\begin{theorem}[屏幕缺损的 Kolmogorov 精确分解]\label{thm:conclusion-screen-kolmogorov-deficit-exact-splitting}
\leanverified{paper\_conclusion\_screen\_kolmogorov\_deficit\_exact\_splitting}
沿用定理 \ref{thm:conclusion-screen-kernel-solenoid-bicompletion} 的记号，并将系数域取为 \(\FF_2\)。记
\[
\bar f_{S,2}:C_n(\mathcal Q_m;\FF_2)\to \FF_2^{\,S},
\qquad
r_S=c(\Gamma_S)-1.
\]
则存在可计算线性同构
\[
\Xi_S:
C_n(\mathcal Q_m;\FF_2)
\xrightarrow{\ \sim\ }
\operatorname{im}(\bar f_{S,2})\oplus \FF_2^{\,r_S},
\qquad
\Xi_S(x)=\bigl(\bar f_{S,2}(x),h_S(x)\bigr),
\]
其中 \(h_S:C_n(\mathcal Q_m;\FF_2)\to \FF_2^{\,r_S}\) 为某个可计算核坐标映射。再令 \(\Phi(x)\) 为定理 \ref{thm:spg-stokes-godel-algorithmic-holographic-completeness} 的完整边界 Gödel 码。则有
\[
\bigl|K(\Phi(x)\mid n,m,S)-K(\Xi_S(x)\mid n,m,S)\bigr|\le O(1),
\]
以及
\[
\bigl|K(x\mid \bar f_{S,2}(x),n,m,S)-K(h_S(x)\mid \bar f_{S,2}(x),n,m,S)\bigr|\le O(1).
\]
换言之，部分屏幕丢失的 Kolmogorov 信息，精确等于其核坐标的条件复杂度。
\end{theorem}

\begin{proof}
由于 \(C_n(\mathcal Q_m;\FF_2)\) 与 \(\FF_2^{\,S}\) 都是有限维 \(\FF_2\)-向量空间，可通过高斯消元可计算地求出
\[
K_{S,2}:=\ker \bar f_{S,2}\cong \FF_2^{\,r_S}
\]
的一组基，并构造一个补空间 \(V_S\subset C_n(\mathcal Q_m;\FF_2)\)，使得
\[
C_n(\mathcal Q_m;\FF_2)=V_S\oplus K_{S,2}
\]
且 \(\bar f_{S,2}|_{V_S}:V_S\to \operatorname{im}(\bar f_{S,2})\) 为线性同构。记其可计算逆为
\[
s_S:\operatorname{im}(\bar f_{S,2})\to V_S.
\]
再取可计算同构
\[
\vartheta_S:K_{S,2}\xrightarrow{\sim}\FF_2^{\,r_S},
\]
并定义
\[
h_S:=\vartheta_S\circ \pi_{K},
\]
其中 \(\pi_K\) 为上述直和分解对应到核的投影。于是对任意 \(x\in C_n(\mathcal Q_m;\FF_2)\)，都有唯一分解
\[
x=s_S(\bar f_{S,2}(x))+\vartheta_S^{-1}(h_S(x)).
\]
从而 \(\Xi_S\) 为可计算双射，其逆由
\[
(y,h)\mapsto s_S(y)+\vartheta_S^{-1}(h)
\]
给出。

另一方面，定理 \ref{thm:spg-stokes-godel-algorithmic-holographic-completeness} 已给出完整边界 Gödel 码 \(\Phi(x)\) 与 \(x\) 在固定 \((m,n)\) 下的可计算双向等价；加入条件 \(S\) 不改变该双向可计算性。故
\[
\Phi(x)\leftrightarrow x\leftrightarrow \Xi_S(x)
\]
均为可计算变换，从而
\[
\bigl|K(\Phi(x)\mid n,m,S)-K(\Xi_S(x)\mid n,m,S)\bigr|\le O(1).
\]

再在条件 \(y:=\bar f_{S,2}(x)\) 下，由上式中的唯一分解可知
\[
x\leftrightarrow h_S(x)
\]
仍是可计算双射，因此
\[
\bigl|K(x\mid y,n,m,S)-K(h_S(x)\mid y,n,m,S)\bigr|\le O(1).
\]
把 \(y\) 代回为 \(\bar f_{S,2}(x)\) 即得结论。证毕。
\end{proof}

\begin{corollary}[部分屏幕的最大条件复杂度恰等于其屏幕秩]\label{cor:conclusion-screen-max-conditional-complexity-equals-rank}
\leanverified{paper\_conclusion\_screen\_max\_conditional\_complexity\_equals\_rank}
在定理 \ref{thm:conclusion-screen-kolmogorov-deficit-exact-splitting} 的设定下，对任意可实现的屏幕输出 \(y\in \operatorname{im}(\bar f_{S,2})\)，都有
\[
\max_{x:\,\bar f_{S,2}(x)=y} K(x\mid y,n,m,S)=r_S+O(1).
\]
因此，屏幕 \(S\) 的最大条件 Kolmogorov 容量精确到 \(O(1)\) 地等于
\[
r_S=c(\Gamma_S)-1.
\]
换言之，相对顶维同调秩、屏幕核维数与最大条件算法信息量在这里严格重合。
\end{corollary}

\begin{proof}
由定理 \ref{thm:conclusion-screen-kolmogorov-deficit-exact-splitting}，
\[
K(x\mid y,n,m,S)=K(h_S(x)\mid y,n,m,S)+O(1).
\]
固定 \(y\) 后，\(h_S(x)\) 在纤维 \(\bar f_{S,2}^{-1}(y)\) 上恰取遍 \(\FF_2^{\,r_S}\)。因此必有统一上界
\[
K(x\mid y,n,m,S)\le r_S+O(1).
\]

另一方面，在 \(\FF_2^{\,r_S}\) 中必存在关于条件 \((y,n,m,S)\) 不可压缩的元素 \(h\)，满足
\[
K(h\mid y,n,m,S)\ge r_S-O(1).
\]
取与之对应的 \(x\in \bar f_{S,2}^{-1}(y)\)，再用定理 \ref{thm:conclusion-screen-kolmogorov-deficit-exact-splitting} 的条件复杂度等价，便得
\[
K(x\mid y,n,m,S)\ge r_S-O(1).
\]
合并上下界即得结论。证毕。
\end{proof}

\begin{theorem}[屏幕相位对象的普适 solenoid 分解与有理自同态环]\label{thm:conclusion-screen-phase-object-universal-solenoid-rational-endomorphisms}
\leanverified{paper\_conclusion\_screen\_phase\_object\_universal\_solenoid\_rational\_endomorphisms}
沿用定理 \ref{thm:conclusion-screen-kernel-solenoid-bicompletion} 的记号。设
\[
K_S\cong \ZZ^{\,r_S},
\qquad
\Sigma_S^{\mathrm{ph}}
:=
\bigl((K_S\otimes_{\ZZ}\RR)\times (K_S\otimes_{\ZZ}\widehat{\ZZ})\bigr)\big/K_S,
\qquad
\widehat{\ZZ}:=\prod_p\ZZ_p.
\]
再记一维普适 solenoid 为
\[
\mathbb U:= (\RR\times \widehat{\ZZ})/\ZZ.
\]
则存在自然拓扑群同构
\[
\Sigma_S^{\mathrm{ph}}\cong \mathbb U^{\,r_S}.
\]
特别地，
\[
\bigl(\Sigma_S^{\mathrm{ph}}\bigr)^\vee \cong \QQ^{\,r_S},
\qquad
\End_{\mathrm{cts}}\bigl(\Sigma_S^{\mathrm{ph}}\bigr)
\cong
M_{r_S}(\QQ).
\]
因此，只要 \(r_S\ge 2\)，\(\Sigma_S^{\mathrm{ph}}\) 就含有非平凡连续幂等自同态；换言之，屏幕相位对象在连续范畴中并非刚性不可分，而是精确地由一个有理矩阵代数控制。
\end{theorem}

\begin{proof}
由 \(K_S\cong \ZZ^{\,r_S}\) 可选取一组 \(\ZZ\)-基，从而得到
\[
K_S\otimes_{\ZZ}\RR\cong \RR^{r_S},
\qquad
K_S\otimes_{\ZZ}\widehat{\ZZ}\cong \widehat{\ZZ}^{\,r_S}.
\]
于是
\[
\Sigma_S^{\mathrm{ph}}
\cong
(\RR^{r_S}\times \widehat{\ZZ}^{\,r_S})/\ZZ^{r_S}.
\]
而 \(\ZZ^{r_S}\) 的对角作用逐坐标分离，因此商群自然同构于
\[
\bigl((\RR\times \widehat{\ZZ})/\ZZ\bigr)^{r_S}
=
\mathbb U^{\,r_S}.
\]

Pontryagin 对偶把有限直积送到有限直和，而一维普适 solenoid 的对偶恰为加法群 \(\QQ\)。故
\[
\bigl(\Sigma_S^{\mathrm{ph}}\bigr)^\vee
\cong
\bigl(\mathbb U^\vee\bigr)^{r_S}
\cong
\QQ^{\,r_S}.
\]
再由 Pontryagin 对偶的反变等价，
\[
\End_{\mathrm{cts}}\bigl(\Sigma_S^{\mathrm{ph}}\bigr)
\cong
\End_{\ZZ}\bigl(\QQ^{\,r_S}\bigr)
\cong
M_{r_S}(\QQ).
\]
当 \(r_S\ge 2\) 时，任取一个非平凡有理投影矩阵便给出非平凡幂等自同态。证毕。
\end{proof}

\begin{corollary}[非平凡屏幕相位对象不可有限素数局部化]\label{cor:conclusion-screen-phase-object-not-finitely-localizable}
设 \(T\subset \mathbb P\) 为有限素数集，记
\[
G_T:=\ZZ[T^{-1}] ,
\qquad
\widehat{G_T}
\]
为其 Pontryagin 对偶。则任意满足 \(r_S\ge 1\) 的屏幕相位对象 \(\Sigma_S^{\mathrm{ph}}\) 都不可能与 \(\widehat{G_T}\) 连续同构。更强地，\(\Sigma_S^{\mathrm{ph}}\) 也不可能与任意有限个有限素数局部化相位对象的有限直积连续同构。
\end{corollary}

\begin{proof}
由定理 \ref{thm:conclusion-screen-phase-object-universal-solenoid-rational-endomorphisms}，
\[
\bigl(\Sigma_S^{\mathrm{ph}}\bigr)^\vee\cong \QQ^{\,r_S}.
\]
若 \(r_S\neq 1\)，则其 \(\QQ\)-秩与 \(G_T\) 的 \(\QQ\)-秩 \(1\) 已经不同，故不可能同构。现只需考虑 \(r_S=1\) 的情形。

若此时有
\[
\Sigma_S^{\mathrm{ph}}\cong \widehat{G_T},
\]
则对偶离散群必须满足
\[
\QQ\cong G_T=\ZZ[T^{-1}].
\]
进而其挠商也应同构。但
\[
\QQ/\ZZ \cong \bigoplus_{p\in\mathbb P} C_{p^\infty},
\]
具有全素数支撑；另一方面，
\[
G_T/\ZZ \cong \bigoplus_{p\in T} C_{p^\infty},
\]
只具有有限素数支撑。二者不可能同构，矛盾。

有限直积情形同理：若 \(\Sigma_S^{\mathrm{ph}}\) 与若干 \(\widehat{G_{T_i}}\) 的有限直积同构，则其对偶将同构于若干 \(\ZZ[T_i^{-1}]\) 的有限直和，其挠商仍只具有有限素数支撑；而 \((\QQ/\ZZ)^{r_S}\) 始终具有全素数支撑。故亦不可能同构。证毕。
\end{proof}
\leanverified{paper\_conclusion\_screen\_phase\_object\_not\_finitely\_localizable}

\endinput

\paragraph{屏幕亏秩的熵恒等式、坐标束板层基与协议型 \texorpdfstring{$\mathsf{NULL}$}{NULL} 的相对同调实现}
在 screen deficit 的 cut-space 识别、屏幕核的 solenoid 完备化以及 Kolmogorov 缺损的精确核坐标都已经建立之后，还可把部分屏幕的几何亏秩进一步压成一条更紧的链：其一，均匀 bulk 在部分屏幕下所损失的 Shannon 条件熵，恰与最小二进审计位数完全重合；其二，内部坐标束屏幕的核不仅有闭式秩，而且具有由板层示性链给出的显式整数基，并在坐标幂集格上退化为精确模律；其三，在只允许访问可见屏幕而禁止隐式侧信息泄漏的类型化读出中，协议型 \(\mathsf{NULL}\) 恰对应相对顶维同调亏秩为正的情形。于是，entropy、audit、coordinate bundle 与 typed \(\mathsf{NULL}\) 被压缩为同一相对亏秩不变量的不同投影。

\leanverified{paper\_conclusion\_screen\_entropy\_audit\_identity}
\begin{theorem}[部分屏幕的精确熵--审计恒等式]\label{thm:conclusion-screen-entropy-audit-identity}
固定 \(m,n\ge 1\)，令 \(\mathcal Q_m\) 为 \(I^n\) 的 \(m\)-级 dyadic 立方复形，并取任意部分屏幕
\[
S\subseteq \overrightarrow{\mathcal F_m^{n-1}}.
\]
设
\[
f_S:C_n(\mathcal Q_m;\FF_2)\longrightarrow \FF_2^{\,S}
\]
为部分边界观测算子，\(X\) 为 \(C_n(\mathcal Q_m;\FF_2)\) 上的均匀分布随机变量，且 \(Y:=f_S(X)\)。再记
\[
r_S:=\beta_n(\mathcal Q_m,A_{\neg S};\FF_2)=c(\Gamma_S)-1.
\]
则有精确恒等式
\[
H(X\mid Y)=r_S.
\]
并且，这一共同整数恰等于使 \(f_S\) 成为精确读出的最小二进审计位数，即
\[
\min\Bigl\{
b:\exists\ \iota,\ \iota:C_n(\mathcal Q_m;\FF_2)\hookrightarrow \FF_2^{\,S}\times \FF_2^{\,b},\
\operatorname{pr}_1\circ \iota=f_S
\Bigr\}=r_S.
\]
\end{theorem}

\begin{proof}
由推论 \ref{cor:spg-partial-boundary-fiber-cardinality-f2}，
对任意 \(y\in \operatorname{im}(f_S)\)，其非空纤维大小恒为
\[
\#f_S^{-1}(y)=2^{\beta_n(\mathcal Q_m,A_{\neg S};\FF_2)}=2^{r_S}.
\]
因此，在均匀先验下，条件分布 \(X\mid(Y=y)\) 在该纤维上均匀，从而
\[
H(X\mid Y=y)=\log_2 \#f_S^{-1}(y)=r_S
\]
对一切 \(y\in\operatorname{im}(f_S)\) 恒成立。再对 \(Y\) 取平均，便得
\[
H(X\mid Y)=r_S.
\]

另一方面，若附加 \(b\) 个二进审计位并得到一个单射提升
\[
\iota:C_n(\mathcal Q_m;\FF_2)\hookrightarrow \FF_2^{\,S}\times \FF_2^{\,b},
\qquad
\operatorname{pr}_1\circ\iota=f_S,
\]
则在任一固定纤维 \(f_S^{-1}(y)\) 内，第二坐标至多可区分 \(2^b\) 个候选；故必须有
\[
2^b\ge \#f_S^{-1}(y)=2^{r_S},
\]
从而 \(b\ge r_S\)。反之，取 \(\ker f_S\) 的一组 \(\FF_2\)-基，把每个元素相对于某个直和分裂的核坐标输出为 \(r_S\) 个比特，便得到一个恰用 \(r_S\) 位的单射补全。故最小审计位数正是 \(r_S\)。证毕。
\end{proof}

\begin{theorem}[最优部分审计的 Kolmogorov 等距性]\label{thm:conclusion-screen-optimal-audit-kolmogorov-isometry}
\leanverified{paper\_conclusion\_screen\_optimal\_audit\_kolmogorov\_isometry}
沿用定理 \ref{thm:conclusion-screen-entropy-audit-identity} 的记号。固定 \(S\) 的一组可计算核基，并取一个可计算直和分解
\[
C_n(\mathcal Q_m;\FF_2)=L_S\oplus \ker f_S.
\]
定义最优二进审计编码
\[
\mathrm{Enc}_S(x):=\bigl(f_S(x),\kappa_S(x)\bigr),
\]
其中 \(\kappa_S(x)\in \FF_2^{\,r_S}\) 表示 \(x\) 在上述分解下的核坐标。则存在仅依赖于 \(S,m,n\) 与所选通用机/分解的常数 \(C_S\)，使得对任意 \(x\in C_n(\mathcal Q_m;\FF_2)\) 都有
\[
\bigl|K(x\mid m,n,S)-K(\mathrm{Enc}_S(x)\mid m,n,S)\bigr|\le C_S,
\]
并且
\[
K\bigl(x\mid \mathrm{Enc}_S(x),m,n,S\bigr)\le C_S.
\]
反之，任何少于 \(r_S\) 位的可计算补全，都不可能在全域上同时满足单射性与上述 \(O(1)\)-级可逆复杂度。
\end{theorem}

\begin{proof}
由构造，\(\mathrm{Enc}_S\) 是可计算映射，故
\[
K(\mathrm{Enc}_S(x)\mid m,n,S)\le K(x\mid m,n,S)+O(1).
\]
另一方面，由直和分解，任意 \(x\) 都可唯一写成
\[
x=s_S\bigl(f_S(x)\bigr)+\sum_{j=1}^{r_S}\varepsilon_j k_j,
\]
其中 \(s_S:\operatorname{im}(f_S)\to L_S\) 为可计算截面，\(\{k_j\}_{j=1}^{r_S}\) 为所固定的核基，而
\[
\kappa_S(x)=(\varepsilon_1,\dots,\varepsilon_{r_S}).
\]
故从 \(\mathrm{Enc}_S(x)\) 恢复 \(x\) 只需执行固定的有限维线性代数运算，因而其逆映射在像上可计算。于是
\[
K(x\mid m,n,S)\le K(\mathrm{Enc}_S(x)\mid m,n,S)+O(1),
\qquad
K\bigl(x\mid \mathrm{Enc}_S(x),m,n,S\bigr)\le O(1).
\]
合并即得前两式。

若存在某个使用 \(b<r_S\) 位的可计算补全在全域上仍为单射，则这会给出定理 \ref{thm:conclusion-screen-entropy-audit-identity} 中更短的单射提升，与最小审计位数等于 \(r_S\) 相矛盾。因而此类更短补全不可能同时拥有全域单射性与 \(O(1)\)-级逆复杂度。证毕。
\end{proof}

\leanverified{paper\_conclusion\_coordinate\_bundle\_kernel\_slab\_decomposition}
\begin{theorem}[内部坐标束屏幕核的板层分解]\label{thm:conclusion-coordinate-bundle-kernel-slab-decomposition}
对任意
\[
J\subseteq [n]:=\{1,\dots,n\},
\qquad
s:=|J|,
\]
令 \(S_J^{\mathrm{int}}\) 为全部法向平行于某个 \(e_j\)（\(j\in J\)）的内部共享面所组成之屏幕。对每个补坐标
\[
a_{J^c}\in\{0,\dots,2^m-1\}^{J^c},
\]
定义相应的 \(J\)-板层
\[
\Sigma_{a_{J^c}}
:=
\bigcup_{a_J\in \{0,\dots,2^m-1\}^{J}} C_{(a_J,a_{J^c})},
\]
并记其顶维示性链为 \([\Sigma_{a_{J^c}}]\)。则有自然直和分解
\[
\ker f_{S_J^{\mathrm{int}}}
\cong
\bigoplus_{a_{J^c}\in \{0,\dots,2^m-1\}^{J^c}}
\ZZ\,[\Sigma_{a_{J^c}}].
\]
特别地，
\[
\rank_{\ZZ}\ker f_{S_J^{\mathrm{int}}}=2^{m(n-s)}.
\]
\end{theorem}

\begin{proof}
由定理 \ref{thm:spg-internal-coordinate-bundle-screen-cost-closed-form} 的证明，
\(\Gamma_{S_J^{\mathrm{int}}}\) 的每一个内部连通分量都恰由固定补坐标 \(a_{J^c}\) 的全部顶维胞腔组成，而沿 \(J\) 中坐标变化的网格图在该分量内连通。故内部连通分量与所有 \(a_{J^c}\in\{0,\dots,2^m-1\}^{J^c}\) 一一对应。

另一方面，由定理 \ref{thm:spg-screen-kernel-connected-components}，
\[
\ker f_{S_J^{\mathrm{int}}}\cong \ZZ^{\,c(\Gamma_{S_J^{\mathrm{int}}})-1}.
\]
其构造性证明已表明：核中的一组自然基可取为“在某个内部连通分量上取常值、在其余分量上取零”的顶维链。将这一描述翻译回 dyadic 胞腔地址，便恰得到各个板层的示性链 \([\Sigma_{a_{J^c}}]\)。由于不同板层支撑互不相交，所得到的子群和为直和。最后由
\[
c(\Gamma_{S_J^{\mathrm{int}}})=2^{m(n-s)}+1
\]
即得秩公式
\[
\rank_{\ZZ}\ker f_{S_J^{\mathrm{int}}}=2^{m(n-s)}.
\]
证毕。
\end{proof}

\begin{theorem}[坐标束室中的对数模性与显式可见秩]\label{thm:conclusion-coordinate-bundle-logmodularity-visible-rank}
\leanverified{paper\_conclusion\_coordinate\_bundle\_logmodularity\_visible\_rank}
定义内部坐标束屏幕的亏秩函数
\[
\Delta_m(J):=\rank_{\ZZ}\ker f_{S_J^{\mathrm{int}}}
\qquad (J\subseteq [n]).
\]
则对任意 \(J,K\subseteq [n]\)，有精确乘法模性
\[
\Delta_m(J\cap K)\,\Delta_m(J\cup K)=\Delta_m(J)\,\Delta_m(K),
\]
等价地，
\[
\log_2\Delta_m(J)=m(n-|J|)
\]
是幂集格 \(\mathcal P([n])\) 上的模函数。并且相应可见秩满足
\[
\rank_{\QQ}\operatorname{im}\bigl(f_{S_J^{\mathrm{int}}}\otimes \QQ\bigr)
=
2^{mn}-2^{m(n-|J|)}.
\]
\end{theorem}

\begin{proof}
由定理 \ref{thm:conclusion-coordinate-bundle-kernel-slab-decomposition}，
\[
\Delta_m(J)=2^{m(n-|J|)}.
\]
于是
\[
\Delta_m(J\cap K)\,\Delta_m(J\cup K)
=
2^{m(n-|J\cap K|)}2^{m(n-|J\cup K|)}
=
2^{m(2n-|J|-|K|)}
=
\Delta_m(J)\,\Delta_m(K),
\]
其中用到了集合恒等式
\[
|J\cap K|+|J\cup K|=|J|+|K|.
\]
这便证明了乘法模性与对数模性。

另一方面，顶维链空间 \(C_n(\mathcal Q_m;\QQ)\) 的维数为 \(2^{mn}\)。由秩--零度公式，
\[
\rank_{\QQ}\operatorname{im}\bigl(f_{S_J^{\mathrm{int}}}\otimes \QQ\bigr)
=
2^{mn}-\rank_{\QQ}\ker\bigl(f_{S_J^{\mathrm{int}}}\otimes \QQ\bigr)
=
2^{mn}-2^{m(n-|J|)},
\]
其中最后一步利用了整数核是自由的，故张量到 \(\QQ\) 后秩保持不变。证毕。
\end{proof}

\begin{theorem}[全内部可见时的唯一残余全局模]\label{thm:conclusion-full-internal-visibility-unique-global-mode}
\leanverified{paper\_conclusion\_full\_internal\_visibility\_unique\_global\_mode}
令 \(J=[n]\)，并记
\[
\Omega_m:=\sum_{a\in \{0,\dots,2^m-1\}^n} C_a
\]
为全部顶维 dyadic 立方体之和。则
\[
\ker f_{S_{[n]}^{\mathrm{int}}}=\ZZ\,\Omega_m.
\]
换言之，在全部内部坐标方向都已可见时，部分边界观测只剩一个一维全局模。进一步，若 \(B\subseteq \overrightarrow{\mathcal F_m^{n-1}}\) 为任意非空外边界面族，则
\[
f_{S_{[n]}^{\mathrm{int}}\cup B}
\]
已经单射；亦即任意一个外边界参照都足以消灭该唯一残余模。
\end{theorem}

\begin{proof}
把定理 \ref{thm:conclusion-coordinate-bundle-kernel-slab-decomposition} 应用于 \(J=[n]\)。此时 \(J^c=\varnothing\)，故只剩下唯一板层，而其示性链恰为
\[
\Omega_m=\sum_{a} C_a.
\]
因此
\[
\ker f_{S_{[n]}^{\mathrm{int}}}=\ZZ\,\Omega_m.
\]

至于第二条，则正是推论 \ref{cor:spg-full-internal-screen-one-defect-boundary-closure} 的结论：全内部屏幕本身亏秩为 \(1\)，而任意非空外边界面族都足以把内部唯一连通分量与 \(\infty\) 接通，从而使扩张屏幕成为精确屏幕，即对应算子单射。证毕。
\end{proof}

\begin{corollary}[可见面的审计边际收益递减律]\label{cor:conclusion-screen-audit-marginal-diminishing-returns}
对任意屏幕 \(S\subseteq \overrightarrow{\mathcal F_m^{n-1}}\)，记其最小审计成本为
\[
C(S):=\rank_{\ZZ}\ker f_S.
\]
则 \(C\) 单调不增且满足边际收益递减：对任意
\[
A\subseteq B\subseteq \overrightarrow{\mathcal F_m^{n-1}},
\qquad
e\notin B,
\]
都有
\[
C(A)-C(A\cup\{e\})
\ge
C(B)-C(B\cup\{e\}).
\]
亦即：同一张新可见面在贫乏屏幕上所带来的审计成本下降，不小于它在富足屏幕上所带来的下降。
\leanverified{paper\_conclusion\_screen\_audit\_marginal\_diminishing\_returns}
\end{corollary}

\begin{proof}
由推论 \ref{cor:spg-partial-boundary-min-audit-cost-kernel-rank}，
\[
C(S)=\min_{(R,r_S)}\cdim(R).
\]
而定理 \ref{thm:conclusion-screen-rank-gain-closure-diminishing-returns} 已经表明最小审计成本的边际收益严格满足递减律。直接代入即可。证毕。
\end{proof}

\begin{theorem}[相对线性全息的精确补充维数]\label{thm:conclusion-relative-linear-holography-exact-supplement-dimension}
\leanverified{paper\_conclusion\_relative\_linear\_holography\_exact\_supplement\_dimension}
将 fixed-resolution 的分段常数空间
\[
V_m:=\operatorname{span}\{1_{C_a}:a\in \{0,\dots,2^m-1\}^n\}
\]
与顶维链空间 \(C_n(\mathcal Q_m;\QQ)\) 通过标准基识别。记 \(T_J\) 为由内部坐标束屏幕 \(S_J^{\mathrm{int}}\) 诱导的线性观测。若再附加 \(L\) 个标量线性泛函
\[
\ell_1,\dots,\ell_L\in V_m^\ast
\]
并要求联合映射
\[
x\longmapsto \bigl(T_Jx,\ell_1(x),\dots,\ell_L(x)\bigr)
\]
在整个 \(V_m\) 上单射，则必有
\[
L\ge 2^{m(n-|J|)}.
\]
并且此下界可达：取定理 \ref{thm:conclusion-coordinate-bundle-kernel-slab-decomposition} 中全部板层坐标作为附加泛函，即得恰用 \(2^{m(n-|J|)}\) 个标量的单射补全。
\end{theorem}

\begin{proof}
由定理 \ref{thm:conclusion-coordinate-bundle-logmodularity-visible-rank}，
\[
\rank T_J=2^{mn}-2^{m(n-|J|)}.
\]
因此其秩亏恰为
\[
\dim V_m-\rank T_J=2^{m(n-|J|)}.
\]
每附加一个标量线性泛函，联合映射的秩至多增加 \(1\)。若
\[
L<2^{m(n-|J|)},
\]
则联合秩仍严格小于 \(\dim V_m\)，故不可能单射。于是下界成立。

反过来，定理 \ref{thm:conclusion-coordinate-bundle-kernel-slab-decomposition} 已给出 \(\ker T_J\) 的一组显式板层基；取相应坐标泛函作为附加观测后，联合映射在核上成为单射，并与原可见部分互补，故整体单射。于是该下界可达。证毕。
\end{proof}

\begin{theorem}[相对同调亏秩对协议型 \texorpdfstring{$\mathsf{NULL}$}{NULL} 的实现]\label{thm:conclusion-relative-homology-corank-protocol-null}
\leanverified{paper\_conclusion\_relative\_homology\_corank\_protocol\_null}
沿用第 \ref{sec:logic-expansion-chain} 节把 \(\mathsf{NULL}\) 视为结构性失败标签的约定。固定部分屏幕 \(S\)，并在“显式精度索引且禁止隐式侧信息泄漏”的类型化制度下，定义确定性偏读出协议
\[
\mathrm{Read}_S
\]
为仅允许访问可见综合征
\[
y=f_S(x)
\]
而不允许访问任何额外审计坐标的读出。则以下三条件等价：
\begin{enumerate}
  \item \(\beta_n(\mathcal Q_m,A_{\neg S};\FF_2)>0\)；
  \item 存在 \(y\in \operatorname{im}(f_S)\) 使得
        \[
        \#f_S^{-1}(y)>1;
        \]
  \item \(\mathrm{Read}_S\) 在某个可见地址类上必返回协议型 \(\mathsf{NULL}\)。
\end{enumerate}
相反，若
\[
\beta_n(\mathcal Q_m,A_{\neg S};\FF_2)=0,
\]
则 \(f_S\) 单射，从而该读出协议中不存在由屏幕不足所导致的协议型 \(\mathsf{NULL}\)。
\end{theorem}

\begin{proof}
第一条与第二条的等价由定理 \ref{thm:conclusion-screen-entropy-audit-identity} 立即得到：\(\beta_n>0\) 当且仅当某个纤维大小大于 \(1\)。

现设存在 \(y\in \operatorname{im}(f_S)\) 使
\[
\#f_S^{-1}(y)>1.
\]
在所规定的类型纪律下，\(\mathrm{Read}_S\) 只能依赖于 \(y\) 本身，而不能额外访问任何 hidden audit 坐标。因此它不可能在同一可见综合征下，对两个不同的前像同时给出唯一且可核验的对象级回放。另一方面，该纤维非空，故失败并非来自不可实现的空纤维，也不是第 \ref{sec:logic-expansion-chain} 节意义下的 \(\mathsf{Undef}\)；它只可能被记录为由协议侧信息不足所导致的结构性失败，即协议型 \(\mathsf{NULL}\)。这便证明了第二条推出第三条。

反过来，若 \(\beta_n=0\)，则由定理 \ref{thm:conclusion-screen-entropy-audit-identity}，
\[
\#f_S^{-1}(y)=1
\qquad
\text{对一切 }y\in \operatorname{im}(f_S).
\]
故 \(f_S\) 为单射，存在仅依赖于可见综合征 \(y\) 的确定性逆译码，从而不再会出现由屏幕不足所导致的协议型 \(\mathsf{NULL}\)。于是第三条不能发生，得到其逆否命题。证毕。
\end{proof}

\begin{theorem}[线性屏幕实现中的 \texorpdfstring{$\mathsf{NULL}$}{NULL} 三分律塌缩]\label{thm:conclusion-screen-linear-readout-null-trichotomy-collapse}
\leanverified{paper\_conclusion\_screen\_linear\_readout\_null\_trichotomy\_collapse}
固定分辨率 \(m,n\)，并在 \(\FF_2\)-系数下考虑部分屏幕观测
\[
f_S:C_n(\mathcal Q_m;\FF_2)\longrightarrow \FF_2^{\,S},
\qquad
S\subseteq \overrightarrow{\mathcal F_m^{n-1}}.
\]
对任意可见数据 \(y\in \FF_2^{\,S}\)，定义三值读出
\[
\mathrm{Read}_S(y)\in \{\mathrm{Tot},\mathrm{ProtNULL},\mathrm{CollNULL}\}
\]
为
\[
\mathrm{Read}_S(y)=
\begin{cases}
\mathrm{Tot}, & f_S^{-1}(y)\ \text{为单点},\\
\mathrm{ProtNULL}, & f_S^{-1}(y)=\varnothing,\\
\mathrm{CollNULL}, & \#f_S^{-1}(y)\ge 2.
\end{cases}
\]
则有严格分类：
\[
\mathrm{Read}_S(y)=\mathrm{ProtNULL}
\iff
y\notin \operatorname{im}f_S,
\]
\[
\mathrm{Read}_S(y)=\mathrm{CollNULL}
\iff
y\in \operatorname{im}f_S
\ \text{且}\
\beta_n(\mathcal Q_m,A_{\neg S};\FF_2)>0,
\]
\[
\mathrm{Read}_S(y)=\mathrm{Tot}
\iff
y\in \operatorname{im}f_S
\ \text{且}\
\beta_n(\mathcal Q_m,A_{\neg S};\FF_2)=0.
\]
并且在第二种情形中，
\[
\#f_S^{-1}(y)=2^{\beta_n(\mathcal Q_m,A_{\neg S};\FF_2)}.
\]
\end{theorem}

\begin{proof}
由线性代数，\(f_S^{-1}(y)=\varnothing\) 当且仅当 \(y\notin \operatorname{im}f_S\)，故第一条成立。

若 \(y\in \operatorname{im}f_S\)，则纤维 \(f_S^{-1}(y)\) 是 \(\ker f_S\) 的一个仿射空间，因此
\[
\#f_S^{-1}(y)=2^{\dim_{\FF_2}\ker f_S}.
\]
而由本文前面已经反复使用的相对顶维同调刻画，
\[
\ker f_S\cong H_n(\mathcal Q_m,A_{\neg S};\FF_2),
\]
故
\[
\dim_{\FF_2}\ker f_S=\beta_n(\mathcal Q_m,A_{\neg S};\FF_2).
\]
于是，当该 Betti 数为零时，纤维为单点，读出为 \(\mathrm{Tot}\)；当其严格大于零时，纤维含至少两个点，读出为 \(\mathrm{CollNULL}\)。证毕。
\end{proof}

\begin{theorem}[延迟观测的几何决定时刻]\label{thm:conclusion-screen-delayed-observation-geometric-decision-time}
\leanverified{paper\_conclusion\_screen\_delayed\_observation\_geometric\_decision\_time}
在定理 \ref{thm:conclusion-screen-linear-readout-null-trichotomy-collapse} 的设定下，取一条单调屏幕链
\[
S_0\subseteq S_1\subseteq S_2\subseteq \cdots \subseteq \overrightarrow{\mathcal F_m^{n-1}}.
\]
对任意固定的隐藏顶维链 \(x\in C_n(\mathcal Q_m;\FF_2)\)，记
\[
y_t:=f_{S_t}(x).
\]
定义其决定时刻
\[
\tau:=\inf\{t\ge 0:\mathrm{Read}_{S_t}(y_t)=\mathrm{Tot}\}.
\]
则 \(\tau\) 与 \(x\) 无关，并且有显式公式
\[
\tau=\inf\{t\ge 0:\ c(\Gamma_{S_t})=1\}.
\]
更精确地，对一切 \(t\)：
\begin{enumerate}
  \item 若 \(t<\tau\)，则
        \[
        \mathrm{Read}_{S_t}(y_t)=\mathrm{CollNULL},
        \qquad
        \#f_{S_t}^{-1}(y_t)=2^{\,c(\Gamma_{S_t})-1};
        \]
  \item 若 \(t\ge \tau\)，则
        \[
        \mathrm{Read}_{S_t}(y_t)=\mathrm{Tot},
        \qquad
        f_{S_t}^{-1}(y_t)=\{x\};
        \]
  \item 一旦 \(t\ge \tau\)，则对所有 \(u\ge t\)，由 \(y_u\) 反演得到的唯一 bulk completion 仍为同一 \(x\)。
\end{enumerate}
\end{theorem}

\begin{proof}
由于 \(y_t=f_{S_t}(x)\)，故 \(y_t\in \operatorname{im}f_{S_t}\) 对所有 \(t\) 恒成立；因此 \(\mathrm{ProtNULL}\) 在这条观测轨道上不会出现。

由相对顶维同调核的闭式，
\[
\dim_{\FF_2}\ker f_{S_t}=c(\Gamma_{S_t})-1.
\]
再联用定理 \ref{thm:conclusion-screen-linear-readout-null-trichotomy-collapse}，可得
\[
\mathrm{Read}_{S_t}(y_t)=\mathrm{Tot}
\iff
\ker f_{S_t}=0
\iff
c(\Gamma_{S_t})=1.
\]
因此 \(\tau\) 只由屏幕几何决定，与具体的 \(x\) 无关；第一、第二条随即成立。

又因 \(S_t\subseteq S_u\) 蕴含 \(f_{S_u}\) 是在 \(f_{S_t}\) 之上增加若干坐标而得，故
\[
\ker f_{S_u}\subseteq \ker f_{S_t}.
\]
于是当某个时刻 \(t\) 已有 \(\ker f_{S_t}=0\) 时，对一切 \(u\ge t\) 仍有 \(\ker f_{S_u}=0\)。这就说明后续观测不会改写已经唯一决定的 bulk completion，第三条得证。证毕。
\end{proof}

\begin{theorem}[精确二元审计的同调比特下界]\label{thm:conclusion-screen-exact-binary-audit-homological-bit-bound}
沿用定理 \ref{thm:conclusion-screen-linear-readout-null-trichotomy-collapse} 的设定。记
\[
\beta:=\beta_n(\mathcal Q_m,A_{\neg S};\FF_2).
\]
称一个 \(b\)-比特审计补充为映射
\[
a:C_n(\mathcal Q_m;\FF_2)\longrightarrow \{0,1\}^b
\]
使得组合编码
\[
x\longmapsto \bigl(f_S(x),a(x)\bigr)
\]
为单射。则必有
\[
b\ge \beta.
\]
反之，存在恰好 \(b=\beta\) 的审计补充使上式单射。换言之，\(\beta\) 正是该屏幕上精确二元审计的最小比特成本。
\end{theorem}
\leanverified{paper\_conclusion\_screen\_exact\_binary\_audit\_homological\_bit\_bound}

\begin{proof}
由定理 \ref{thm:conclusion-screen-linear-readout-null-trichotomy-collapse}，对任意 \(y\in \operatorname{im}f_S\)，纤维 \(f_S^{-1}(y)\) 的基数恒为 \(2^\beta\)。若 \((f_S,a)\) 为单射，则对任意固定 \(y\)，其限制
\[
a|_{f_S^{-1}(y)}:f_S^{-1}(y)\to \{0,1\}^b
\]
也必须为单射。故必有
\[
2^\beta\le 2^b,
\]
即 \(b\ge \beta\)。

反过来，取 \(\ker f_S\) 的一组 \(\FF_2\)-基，并选择一补空间
\[
C_n(\mathcal Q_m;\FF_2)=\ker f_S\oplus W.
\]
任意 \(x\) 可唯一写成
\[
x=k+w,
\qquad
k\in \ker f_S,\ w\in W.
\]
定义 \(a(x)\) 为 \(k\) 在上述基下的坐标向量，则 \(a(x)\in \{0,1\}^{\beta}\)。若
\[
f_S(x)=f_S(x'),
\qquad
a(x)=a(x'),
\]
则 \(x-x'\in \ker f_S\)，且二者具有相同核坐标，故 \(x-x'=0\)，从而 \(x=x'\)。因此当 \(b=\beta\) 时该编码已经单射，下界可达。证毕。
\end{proof}

\endinput

\paragraph{屏幕亏秩的 R\'enyi 平坦化、零误差审计时间与拟阵信息律}
在部分屏幕的精确熵--审计恒等式、最优核坐标补全、内部坐标束板层分解与协议型 \(\mathsf{NULL}\) 的相对同调实现都已确立之后，还可把同一亏秩机制再推进一步：相对顶维 Betti 数不仅给出 Shannon 条件熵，而且在全阶 R\'enyi 塔上保持完全平坦；它同时等于零误差审计时间的二元基准值、决定条件碰撞矩的几何级数律，并把可见信息提升为一个由边界矩阵行空间诱导的可表示拟阵秩函数。于是，screen deficit 不再只是“缺多少补全位”的问题，而是统一组织了熵、时间、碰撞与可见秩的同一不变量。

\begin{theorem}[相对 Betti--R\'enyi 平坦律]\label{thm:conclusion-relative-betti-renyi-flatness}
\leanverified{paper\_conclusion\_relative\_betti\_renyi\_flatness}
沿用定理 \ref{thm:conclusion-screen-entropy-audit-identity} 的记号，并记
\[
N_{m,n}:=\dim_{\FF_2}C_n(\mathcal Q_m;\FF_2)=2^{mn},
\qquad
\beta_S:=\beta_n(\mathcal Q_m,A_{\neg S};\FF_2).
\]
则在均匀 bulk 先验下，部分屏幕信道 \(X\mapsto Y=f_S(X)\) 满足
\[
H(X\mid Y)=\beta_S,
\qquad
H(Y)=N_{m,n}-\beta_S,
\qquad
I(X;Y)=N_{m,n}-\beta_S.
\]
更强地，对任意阶 \(q\in [0,\infty]\) 的标准条件 R\'enyi 熵，只要它仅由各非空纤维上的条件分布决定，均有
\[
H_q(X\mid Y)=\beta_S.
\]
因此，部分屏幕所隐藏的信息量在全阶 R\'enyi 塔上完全平坦，并且该共同值恰由相对顶维 Betti 数给出。
\end{theorem}

\begin{proof}
由定理 \ref{thm:conclusion-screen-entropy-audit-identity} 的证明，对任意
\[
y\in \operatorname{im}(f_S)
\]
都有
\[
\#f_S^{-1}(y)=2^{\beta_S},
\]
且该数与 \(y\) 无关。由于 \(X\) 在 \(C_n(\mathcal Q_m;\FF_2)\) 上均匀，故条件分布 \(X\mid(Y=y)\) 在纤维 \(f_S^{-1}(y)\) 上均匀，从而
\[
H(X\mid Y=y)=\log_2\#f_S^{-1}(y)=\beta_S.
\]
对 \(Y\) 取平均即得
\[
H(X\mid Y)=\beta_S.
\]

又因为
\[
H(X)=N_{m,n},
\]
所以
\[
H(Y)=H(X)-H(X\mid Y)=N_{m,n}-\beta_S,
\]
并且
\[
I(X;Y)=H(X)-H(X\mid Y)=N_{m,n}-\beta_S.
\]

对任意固定 \(y\)，条件分布 \(X\mid(Y=y)\) 都是支撑在 \(2^{\beta_S}\) 个点上的均匀分布。故无论取 Shannon 熵、碰撞熵、最小熵还是任意阶 \(q\in(0,\infty)\setminus\{1\}\) 的标准 R\'enyi 熵，其值都等于
\[
\log_2 2^{\beta_S}=\beta_S.
\]
由于该值与 \(y\) 无关，故相应的条件 R\'enyi 熵也都等于 \(\beta_S\)。证毕。
\end{proof}

\begin{theorem}[条件碰撞矩的几何级数律与 Hankel 秩一刚性]\label{thm:conclusion-screen-collision-moment-hankel-rank-one}
\leanverified{paper\_conclusion\_screen\_collision\_moment\_hankel\_rank\_one}
沿用定理 \ref{thm:conclusion-relative-betti-renyi-flatness} 的设定。对每个 \(q\ge 1\)，定义条件碰撞矩
\[
M_q(S):=
\sum_{y}\PP(Y=y)\sum_x \PP(X=x\mid Y=y)^q.
\]
则有精确闭式
\[
M_q(S)=2^{-(q-1)\beta_S}.
\]
因此，序列 \(q\mapsto M_q(S)\) 是公比 \(2^{-\beta_S}\) 的纯几何级数；相应 Hankel 矩阵
\[
\mathcal H_S:=\bigl(M_{i+j}(S)\bigr)_{i,j\ge 1}
\]
满足
\[
\rank \mathcal H_S=1.
\]
从而，若某一隐藏机制的条件碰撞 Hankel 矩阵秩严格大于 \(1\)，则它不可能仅仅来自屏幕截断造成的纯相对同调盲区。
\end{theorem}

\begin{proof}
由定理 \ref{thm:conclusion-relative-betti-renyi-flatness}，给定任意
\[
y\in \operatorname{im}(f_S),
\]
条件分布 \(X\mid(Y=y)\) 在 \(2^{\beta_S}\) 个点上均匀。故
\[
\sum_x \PP(X=x\mid Y=y)^q
=
2^{\beta_S}\cdot 2^{-q\beta_S}
=
2^{-(q-1)\beta_S},
\]
且该值与 \(y\) 无关。对 \(y\) 再取平均，即得
\[
M_q(S)=2^{-(q-1)\beta_S}.
\]

于是对任意 \(i,j\ge 1\)，
\[
M_{i+j}(S)
=
2^{-((i+j)-1)\beta_S}
=
\bigl(2^{-(i-1)\beta_S}\bigr)\bigl(2^{-j\beta_S}\bigr).
\]
故 \(\mathcal H_S\) 为两个无限列向量的外积，因而
\[
\rank \mathcal H_S=1.
\]
证毕。
\end{proof}

\begin{theorem}[零误差审计时间的精确公式]\label{thm:conclusion-screen-zeroerror-audit-time}
\leanverified{paper\_conclusion\_screen\_zeroerror\_audit\_time}
沿用上述设定。设审计字母表大小为 \(A\ge 2\)。考虑如下零误差补全任务：先观察部分屏幕输出 \(Y=f_S(X)\)，再允许一个拥有 \(X\) 的编码器逐步发送 \(A\)-元审计符号，接收端须在最坏情形下精确恢复 \(X\)。则最小最坏情形延迟恰为
\[
T_\ast(S;A)=\left\lceil \frac{\beta_S}{\log_2 A}\right\rceil.
\]
特别地，在二元审计通道上，
\[
T_\ast(S;2)=\beta_S.
\]
因此，屏幕盲区所诱导的最小补全时间，不是额外的动态对象，而正是同一相对顶维亏秩按字母表底数重标后的时间化读出。
\end{theorem}

\begin{proof}
固定任意 \(y\in \operatorname{im}(f_S)\)。由定理 \ref{thm:conclusion-relative-betti-renyi-flatness}，
\[
\#f_S^{-1}(y)=2^{\beta_S}.
\]
若在 \(t\) 步之内发送 \(A\)-元审计符号，则至多产生 \(A^t\) 条不同转录。零误差恢复要求同一纤维中的不同点必须被不同转录区分，故必要条件为
\[
A^t\ge 2^{\beta_S},
\]
即
\[
t\ge \left\lceil \frac{\beta_S}{\log_2 A}\right\rceil.
\]

反之，由定理 \ref{thm:conclusion-screen-optimal-audit-kolmogorov-isometry} 的构造，固定 \(\ker f_S\) 的一组 \(\FF_2\)-基，并取直和分解
\[
C_n(\mathcal Q_m;\FF_2)=L_S\oplus \ker f_S.
\]
则任意 \(x\in f_S^{-1}(y)\) 都可唯一表示为
\[
x=s_S(y)+\sum_{j=1}^{\beta_S}\varepsilon_j k_j,
\qquad
\varepsilon_j\in \FF_2.
\]
因此，只需发送 \(\beta_S\) 个二进核坐标，便可在最坏情形下零误差恢复 \(x\)。把这 \(\beta_S\) 个比特按 \(A\)-元字母表分组，即得
\[
T_\ast(S;A)\le \left\lceil \frac{\beta_S}{\log_2 A}\right\rceil.
\]
与下界合并便得等式。证毕。
\end{proof}

\begin{theorem}[屏幕信息的拟阵秩结构与隐藏信息的超模律]\label{thm:conclusion-screen-information-matroid-supermodularity}
对任意屏幕
\[
S\subseteq \overrightarrow{\mathcal F_m^{n-1}},
\]
定义
\[
r_2(S):=\rank_{\FF_2}(f_S).
\]
则函数
\[
S\longmapsto I(X;Y_S)
\]
恰等于由顶维边界矩阵之行限制诱导的 \(\FF_2\)-可表示拟阵秩函数，即
\[
I(X;Y_S)=r_2(S).
\]
特别地，对任意屏幕 \(S,T\) 有
\[
I(X;Y_{S\cup T})+I(X;Y_{S\cap T})
\le
I(X;Y_S)+I(X;Y_T),
\]
且 \(I(X;Y_S)\) 随 \(S\) 单调不减。

对偶地，若定义隐藏信息
\[
J(S):=H(X\mid Y_S),
\]
则
\[
J(S)=N_{m,n}-r_2(S),
\]
从而 \(J\) 满足超模律
\[
J(S\cup T)+J(S\cap T)\ge J(S)+J(T),
\]
并随 \(S\) 单调不增。
\end{theorem}

\begin{proof}
固定 \(m,n\) 后，顶维边界算子在 \(\FF_2\) 上可由一个固定矩阵表示，而 \(f_S\) 恰是取其由 \(S\) 索引的那些行所得到的线性映射。故
\[
r_2(S)
=
\dim_{\FF_2}\operatorname{span}\{\text{boundary rows indexed by }S\},
\]
从而 \(r_2(\cdot)\) 正是一个 \(\FF_2\)-可表示拟阵的秩函数，因此自动满足单调性与次模性。

另一方面，\(X\) 在 \(C_n(\mathcal Q_m;\FF_2)\) 上均匀，而 \(f_S\) 为线性映射，因此 \(Y_S=f_S(X)\) 在像空间 \(\operatorname{im}(f_S)\) 上均匀。于是
\[
H(Y_S)=\log_2 |\operatorname{im}(f_S)|=r_2(S).
\]
又因 \(Y_S\) 是 \(X\) 的确定函数，
\[
I(X;Y_S)=H(Y_S)=r_2(S).
\]

再由定理 \ref{thm:conclusion-relative-betti-renyi-flatness}，
\[
J(S)=H(X\mid Y_S)=N_{m,n}-I(X;Y_S)=N_{m,n}-r_2(S).
\]
于是 \(J\) 继承了 \(r_2\) 的单调性对偶与超模性。证毕。
\end{proof}
\leanverified{paper\_conclusion\_screen\_information\_matroid\_supermodularity}

\begin{corollary}[内部坐标束屏幕上的精确面积律]\label{cor:conclusion-coordinate-bundle-exact-area-law}
对任意
\[
J\subseteq [n],
\qquad
s:=|J|,
\]
记 \(S_J^{\mathrm{int}}\) 为内部坐标束屏幕。则
\[
\beta_{S_J^{\mathrm{int}}}=2^{m(n-s)}.
\]
因此同时成立精确等式
\[
H(X\mid Y_{S_J^{\mathrm{int}}})=2^{m(n-s)}\ \text{bits},
\]
\[
T_\ast(S_J^{\mathrm{int}};2)=2^{m(n-s)},
\]
\[
\min_{(R,r_S)}\cdim(R)=2^{m(n-s)}.
\]
特别地，当 \(s=1\) 时，
\[
H(X\mid Y_{S_{\{j\}}^{\mathrm{int}}})
=
T_\ast(S_{\{j\}}^{\mathrm{int}};2)
=
\min_{(R,r_S)}\cdim(R)
=
2^{m(n-1)},
\]
即得到一条严格等式型面积律。
\end{corollary}

\begin{proof}
由定理 \ref{thm:conclusion-coordinate-bundle-kernel-slab-decomposition}，
\[
\rank_{\ZZ}\ker f_{S_J^{\mathrm{int}}}=2^{m(n-s)}.
\]
由于该核为自由阿贝尔群，其 \(\FF_2\)-约化维数也恰等于该秩，故
\[
\beta_{S_J^{\mathrm{int}}}=2^{m(n-s)}.
\]
再将此值分别代入定理 \ref{thm:conclusion-relative-betti-renyi-flatness} 与定理 \ref{thm:conclusion-screen-zeroerror-audit-time}，便得到前两条公式。

最后，由定理 \ref{thm:spg-internal-coordinate-bundle-screen-cost-closed-form}，
\[
\min_{(R,r_S)}\cdim(R)=2^{m(n-s)}.
\]
故三种量完全重合。证毕。
\end{proof}
\leanverified{paper\_conclusion\_coordinate\_bundle\_exact\_area\_law}

\begin{theorem}[内部坐标束屏幕的 dyadic block law]\label{thm:conclusion-coordinate-bundle-dyadic-block-law}
设
\[
J\subseteq [n],
\qquad
|J|=s.
\]
则内部坐标束屏幕满足
\[
a\!\left(S_J^{\mathrm{int}}\right)=2^{m(n-s)}.
\]
因此，当向 \(J\) 中再加入一个新坐标 \(j\notin J\) 时，审计缺口发生严格的 \(2^m\)-倍收缩：
\[
a\!\left(S_{J\cup\{j\}}^{\mathrm{int}}\right)
=
2^{-m}\,a\!\left(S_J^{\mathrm{int}}\right).
\]
在 window-\(6\) 原型 \((m,n)=(1,6)\) 中，这一律退化为精确二分链
\[
32\to 16\to 8\to 4\to 2\to 1.
\]
\end{theorem}
\begin{proof}
由推论 \ref{cor:conclusion-coordinate-bundle-exact-area-law}，
\[
a\!\left(S_J^{\mathrm{int}}\right)=2^{m(n-s)}.
\]
若 \(j\notin J\)，则
\[
a\!\left(S_{J\cup\{j\}}^{\mathrm{int}}\right)
=
2^{m(n-s-1)}
=
2^{-m}\,2^{m(n-s)}
=
2^{-m}\,a\!\left(S_J^{\mathrm{int}}\right).
\]
当 \((m,n)=(1,6)\) 且 \(|J|=1,2,3,4,5,6\) 时，依次得到
\[
2^5,2^4,2^3,2^2,2^1,2^0,
\]
也即
\[
32\to 16\to 8\to 4\to 2\to 1.
\]
证毕。
\end{proof}
\leanverified{paper\_conclusion\_coordinate\_bundle\_dyadic\_block\_law}

\begin{conjecture}[R\'enyi 平坦刚性]\label{conj:conclusion-screen-renyi-flatness-rigidity}
在有限状态 projection-word 语境中，设
\[
P_m:\Omega_m\to X_m
\]
为任意有限分辨率确定性投影，并在均匀先验下考虑条件纤维分布。若同时满足：
\begin{enumerate}
  \item 条件 R\'enyi 熵对全部阶 \(q\) 完全平坦；
  \item 条件碰撞矩的 Hankel 矩阵秩等于 \(1\)；
\end{enumerate}
则在每个连通可见分支上，\(P_m\) 的非空纤维应当形成某个有限阿贝尔群的仿射 torsor。等价地，其隐藏扇区应当是纯同调型的，而非共振型的。

dyadic 屏幕模型所对应的 torsor 正是
\[
H_n(\mathcal Q_m,A_{\neg S};\FF_2).
\]
因此，一旦条件碰撞 Hankel 矩阵秩严格大于 \(1\)，隐藏结构便不能再被纯 torsor 机制解释，而必须进入真正的类内模态或共振异质性世界。
\end{conjecture}

\endinput

\subsection{接口地图：四条可审计闭环与证明/语义分离}\label{subsec:conclusion-interface-map}
\begin{remark}[证明链与解释性语义的分离]\label{rem:conclusion-proof-vs-semantics}
下表将全文的主要“可审计接口”压缩为四条闭环：每条闭环仅记录协议内输入对象、可被引用的证书输出（定理/命题/推论/定义）以及对应的可复现脚本与产物路径。
其中凡以定理/命题/推论形式出现者构成正文的证明链条；表中标注为“语义”的短语仅用于对齐传统术语与参数读出，不作为任何证明的前提，也不在后续推导中作为假设使用。
\end{remark}

\begin{table}[H]
\centering
\scriptsize
\setlength{\tabcolsep}{4pt}
\renewcommand{\arraystretch}{1.15}
\caption{四条闭环的接口地图：输入对象、证书输出与可复现产物。}
\label{tab:conclusion-interface-map}
\begin{tabular}{p{0.13\linewidth} p{0.22\linewidth} p{0.33\linewidth} p{0.32\linewidth}}
\toprule
闭环（等价链） & 输入对象（协议内） & 输出证书（定理链锚点） & 可复现脚本/产物（实现锚点）\\
\midrule
\textbf{计算闭环}\\[-2pt]
ISA \(\Leftrightarrow\) 投影词 \(\Leftrightarrow\) 规范化重写
&
算术 ISA/微码程序（\S\ref{subsec:pom-isa}）；
投影词 \(W\in\mathbf{PW}\)；
FRACTRAN 程序（\S\ref{subsec:pom-fractran}）
&
\(\mathcal{P}\) 上的序门控平移刻画 FRACTRAN-step（定义 \ref{def:pom-fractran}，命题 \ref{prop:pom-fractran-prime-translation}）；
值层重写 \((\mathrm{RZ},\mathrm{RNF})\) 的终止与“末端至多一次 \texttt{PROJ\_NORM}”正规形（定义 \ref{def:pom-rewriting-val}，命题 \ref{prop:pom-rewriting-terminating}）；
并明确序投影 \(P_{\le}\) 为结构性边界（注 \ref{rem:pom-rewrite-boundary-order}）
&
\path{scripts/pom_isa_sim_v2.py}（注 \ref{rem:pom-isa-sim}）；
\path{scripts/exp_pom_rewriting_engine_demo.py}
\(\to\) 表 \ref{tab:pom_rewriting_engine_demo}（注 \ref{rem:pom-rewriting-engine-demo}）
\\
\midrule
\textbf{几何锚点}\\[-2pt]
window-\(6\) \(\Leftrightarrow\) 6D cut-and-project \(\Leftrightarrow\) \((\ZZ_2)^3\) 与 \(\delta=34\)
&
window-\(6\) boundary 三轴
\(
X_6^{\mathrm{bdry}}=\{\texttt{100001},\texttt{100101},\texttt{101001}\}
\)；
dyadic uplift 的二层覆盖差；
内部相位因子 \(T_{\mathrm{int}}\) 的模 \(2\) 影像（语义）
&
boundary 三条 dyadic 纤维独立生成中心嵌入 \((\ZZ_2)^3\hookrightarrow Z(\mathcal{G}_6^{\mathrm{bin}})\)（推论 \ref{cor:fold-bin-boundary-center-z2}）；
\(\delta=34\) 的三重同一性与二层覆盖/掩码/尾部分支的并置（推论 \ref{cor:bdry-delta34-triple-identity}）；
cut-and-project 的内部环面同调锚定与 phason \(2\)-torsion 读出（语义：注 \ref{rem:bdry-tint-h1-cartan-anchor}，注 \ref{rem:bdry-delta34-phason-2torsion}）
&
\path{scripts/exp_window6_c6_orbit_patisalam_seed.py}
\(\to\) 表 \ref{tab:fold6_boundary_sheet_pairs}（并见表 \ref{tab:fold6_bin_uplift_choice_collapse} 的 uplift-choice 坍缩）
\\
\midrule
\textbf{谱正性主干}\\[-2pt]
Carath\'eodory \(\Leftrightarrow\) Toeplitz--PSD \(\Leftrightarrow\) Schur/\allowbreak Verblunsky \(\Leftrightarrow\) CMV
&
视界 Carath\'eodory 对数导数接口 \(\mathscr{C}_\zeta\) 的矩（附录 \ref{subsubsec:app-horizon-weyl-carath-toeplitz}）；
或单点喷射数据 \(\{\ell_k\}\)（\S\ref{subsubsec:xi-verblunsky-cmv-from-local-jet}）
&
RH \(\Leftrightarrow\) Carath\'eodory 正性 \(\Leftrightarrow\) Toeplitz--PSD 层级（附录定理 \ref{thm:app-horizon-toeplitz-lmi}）；
Toeplitz 证书的标量化与 Verblunsky 系数生成（定理 \ref{thm:xi-verblunsky-from-local-jet}）；
Toeplitz 行列式乘积分解与最小失稳定位（定理 \ref{thm:xi-toeplitz-det-verblunsky}）；
有限窗 CMV（Hilbert--P{\'o}lya）酉谱实现（命题 \ref{prop:xi-hilbert-polya-cmv}）
&
\path{scripts/exp_sync_kernel_xi_cubic_cos_polynomial.py}
\(\to\)
\path{sections/generated/eq_sync_kernel_xi_cubic_cos_polynomial.tex},
\path{sections/generated/tab_sync_kernel_xi_cubic_roots.tex}
（注 \ref{rem:xi-hard-zero-certificate}）
\\
\midrule
\textbf{共振解释入口}\\[-2pt]
Poisson/Breit--Wigner（物理术语）\(\Leftrightarrow\) SU(1,1)--RHP \(\Leftrightarrow\) Maslov
&
内函数相位 \(\Theta(t)=\arg D(\tfrac12+\mathrm{i}t)\)；
离线零点参数 \((\gamma_k,\delta_k)\)；
正定跳跃数据与谱流参数（\S\ref{subsubsec:xi-debranges-canonical-rhp-adelic}）
&
相位导数的 Poisson 叠加律：离线零点贡献 Lorentz/Breit--Wigner 核（命题 \ref{prop:xi-phase-derivative-poisson-superposition}）；
SU(1,1)–RHP 正定跳跃可解性等价模板（定理 \ref{thm:xi-su11-rhp-equivalence}）；
零点计数差的 Maslov 指标公式（命题 \ref{prop:xi-maslov-spectralflow-entropy-quantization}）；
并在共振窗二步重标度口径下出现 Lorentz 预极限与宽度参数（注 \ref{rem:pom-resonance-poisson-lorentz-prelimit}）
&
\path{scripts/exp_xi_horizon_dtn_fourier_laplace_audit.py}
\(\to\)
\path{sections/generated/tab_xi_horizon_dtn_fourier_laplace_audit.tex}
（\S\ref{subsec:discussion-horizon-zk-toeplitz-area-law}；共振窗数据表 \ref{tab:fold_collision_shadow_spectral_packet_q9_17}）
\\
\bottomrule
\end{tabular}
\end{table}

\subsection{复现与审计指引}\label{subsec:repro-audit-pointer}
解释性综合、边界与可证伪点集中于第 \ref{sec:discussion} 节。附录汇集算子代数、相位塔、同步核与多尺度接口等支撑材料，以支撑正文的可审计推理链条与复核需求。

\paragraph{结语（可证伪的总论式定位）}
在本文固定的协议口径下，\(\Fold_m\) 的局部归一化规则在形式上极简；其“深度”来自它所诱导的有限状态对象在两条方向上的不可约结构：一是\textbf{谱}（包含 Perron 主根所锁定的代数尺度与由非半单 Jordan 缺陷所携带的残差模态），二是\textbf{对称}（由字典/轨道结构所强制的等变压缩与分块刚性）。
在同一谱世界中，核版 RH/de Bruijn--Newman 型阈值使临界相变可被整系数消元证书精确钉死（注 \ref{rem:pom-collision-rh-break-a4-threshold}；命题 \ref{prop:sync-kernel-weighted-newman-threshold}）；而 \(\mathbb{F}_p\) 上的特征式因子分解又给出一套与实谱正交的局部约化指纹（注 \ref{rem:pom-charpoly-modp-audit-method}），并在高阶端点处把 \(\varphi\) 固化为可由碰撞谱反演的自校准常数（命题 \ref{prop:pom-rq-qinfty-endpoint}）。
两类结构均可被压缩为有限证书（有理母函数与 Perron 主根：命题 \ref{prop:pom-multiplicity-composition-partition-rational}；Jordan 残差与秩/递推约束：定理 \ref{thm:fold-gauge-anomaly-A0-jordan}、定理 \ref{thm:fold-gauge-anomaly-hankel-rank3}；Weyl 轨道压缩与 uplift 分块：推论 \ref{cor:fold6-weyl-two-orbit-compression}、定理 \ref{thm:boundary-m10-b3c3-rigid-four-block-split}），
因此本文的统一叙述并非依赖类比的解释性拼接，而是由可审计输出所定义的可否证结构。
